\documentclass[11pt,a4paper]{article}
\usepackage[T1]{fontenc}
\usepackage[ngerman]{babel}
\usepackage{lmodern,microtype}
\usepackage[margin=23mm,headheight=15pt]{geometry}
\usepackage{amsmath,amssymb,mathtools}
\usepackage{booktabs,longtable,tabularx,array}
\usepackage{xcolor,tikz}
\usetikzlibrary{arrows.meta,positioning,calc,fit,backgrounds}
\usepackage{enumitem,xurl,hyperref,fancyhdr}
\definecolor{navy}{HTML}{16324F}
\definecolor{teal}{HTML}{087F8C}
\definecolor{amber}{HTML}{A85D08}
\definecolor{pale}{HTML}{F0F5F8}
\definecolor{red}{HTML}{9C3242}
\hypersetup{colorlinks=true,linkcolor=navy,urlcolor=teal,citecolor=teal,pdftitle={TFPT: Die gemeinsame Quelle und der Weg zur Gesamtlösung},pdfauthor={Sitzungssynthese für Stefan Hamann}}
\pagestyle{fancy}\fancyhf{}\fancyhead[L]{\small TFPT \textbar\ Gesamtsynthese der Sitzung}\fancyhead[R]{\small Stand: 29. September 2026}\fancyfoot[C]{\thepage}
\setlength{\parindent}{0pt}\setlength{\parskip}{5pt plus 1pt}
\setlength{\emergencystretch}{3em}
\setlist{itemsep=3pt,topsep=4pt,leftmargin=*}
\newcolumntype{Y}{>{\raggedright\arraybackslash}X}
\newcolumntype{L}[1]{>{\raggedright\arraybackslash}p{#1}}
\newcommand{\E}{\textcolor{teal}{\textbf{E}}}
\newcommand{\N}{\textcolor{teal}{\textbf{N}}}
\newcommand{\B}{\textcolor{amber}{\textbf{B}}}
\newcommand{\Opend}{\textcolor{amber}{\textbf{O}}}
\newcommand{\Xclaim}{\textcolor{red}{\textbf{X}}}
\newcommand{\C}{\mathbb C}\newcommand{\R}{\mathbb R}\newcommand{\Z}{\mathbb Z}
\newcommand{\tr}{\operatorname{tr}}\newcommand{\diag}{\operatorname{diag}}
\newcommand{\Tr}{\operatorname{Tr}}
\newcommand{\rank}{\operatorname{rank}}\newcommand{\Span}{\operatorname{span}}
\newcommand{\Id}{\mathrm{id}}\newcommand{\ket}[1]{|#1\rangle}\newcommand{\bra}[1]{\langle#1|}
\newcommand{\abs}[1]{\left|#1\right|}\newcommand{\norm}[1]{\left\|#1\right\|}
\newcommand{\file}[1]{\texttt{\nolinkurl{#1}}}
\newcommand{\key}[1]{\par\smallskip\noindent\colorbox{pale}{\parbox{\dimexpr\linewidth-2\fboxsep\relax}{{\color{navy}#1}}}\par\smallskip}
\newcommand{\picturetext}[1]{\textit{Bildlich:} #1}
\tikzset{box/.style={draw=navy,rounded corners=3pt,fill=pale,align=center,font=\small,inner sep=7pt},arr/.style={-{Latex[length=2mm]},thick,draw=teal},dasharr/.style={arr,dashed}}

\begin{document}
\begin{titlepage}
\vspace*{13mm}
{\Large\color{teal} TFPT \quad / \quad Universalraum \quad / \quad Prozess und Physik\par}
\vspace{16mm}
{\Huge\bfseries\color{navy} Die gemeinsame Quelle\par}
\vspace{5mm}
{\LARGE und der Weg zur Gesamtlösung\par}
\vspace{13mm}
{\large Gesamtsynthese der Sitzung vom 28. September 2026\\Fortgesetzt am 29. September 2026\par}
\vspace{8mm}
\key{\textbf{Korrigiertes Gesamtbild nach Prüfung der physikalischen Herleitungen:}\par
TFPT besitzt bereits konkrete Herleitungsketten für Quellenmassen, Yukawa-Hierarchien, Flavor-Mischungen und die Feinstrukturkonstante sowie bedingte gemeinsame Neutrino-, Skalen- und Kosmologieanschlüsse. Diese Ergebnisse sind keine erst künftig aus der neuen E8-Quelle zu erfindenden Zutaten. Sie legen fest, welche Antworten die gesuchte gemeinsame Quelle liefern muss. Die neue Quellenkonstruktion ist ein Kandidat für ihren gemeinsamen Ursprung; sie ersetzt die bestehenden Herleitungen nicht.\par
Die entscheidende Verbindung ist die ursprüngliche markierte Operator- und Determinantenfamilie: Dieselben Quellenzustände liefern das Flavorbudget und die elektromagnetische Antwort; derselbe Neutrinooperator muss Massen, Mischung und Zerfallsasymmetrie tragen. Abschnitt~\ref{sec:physicalconstraints} erhält diese Herleitungen. Die Fortsetzung in Abschnitt~\ref{sec:operatorbridge} beweist die ursprüngliche endliche Fuchs-Monodromie analytisch. Zwei wirkliche Halbwege dieses Systems realisieren nun exakt den vorhandenen Sechseroperator und seinen Flavor-Kernel. Die vollständige E8-Quelle enthält außerdem ein bestimmtes neutrales Feld vom Gewicht zwei und mit Normquadrat 48. Es realisiert den ursprünglichen Alpha-Rundlauf unter einer ausdrücklich angegebenen Kopplungs- und Zeitidentifikation. Die Fortsetzung in Abschnitt~\ref{sec:spinliftsynthesis} verbindet zusätzlich die primitive Windung mit der tatsächlichen Deckwirkung auf allen geladenen E8-Feldern. Derselbe Quellraum enthält jetzt auch die symmetrischen Neutrinopaaroperatoren, ein norm- und phasentreues Wörterbuch für die vollständigen Familientensoren sowie ihre gemeinsame Antwort mit den neutralen Quellenfeldern. Die ursprüngliche Eichgewichtung und der vollständige Flavor-Defektkern sind damit gemeinsam berechenbar. Die verbleibende Herkunftsfrage ist konkret: Welcher ursprüngliche markierte Prozess erzeugt gerade die vorhandenen Higgs-/Masseneinsetzungen und die erforderliche Nahtkopplung? Ihr gemeinsamer physischer Anschluss ist noch nicht bewiesen.}
\vspace{8mm}
Dieses Dokument führt die ursprünglichen TFPT-Paper, die Konsolidierungen vom 26.--28. September, die eingefügten Synthesevorschläge und die ausführbaren Ergebnisse dieser Sitzung zusammen. Es erklärt die Kette in einfacher Sprache und mit Skizzen, erhält die mathematischen Voraussetzungen und entwickelt daraus konkrete Forschungs- und Algorithmuswege.

\vfill
\textbf{Für Stefan Hamann}\par
Arbeitsfassung mit Quellen- und Befundregister. Die geometrischen Zeichnungen sind erklärende Schemata; sie behaupten keine bereits abgeleitete Gestalt der Raumzeit. Der Anwendungsteil interpretiert das gesprochene Wort ``hallieren'' vorläufig als Hylæan/Field Intelligence.
\end{titlepage}

\section*{Wie dieses Dokument zu lesen ist}
\addcontentsline{toc}{section}{Wie dieses Dokument zu lesen ist}
Der Auftrag lautet, die vollständige Verbindung zu finden. Deshalb ist der Maßstab dieses Textes eine gemeinsame Herleitung mit denselben Feldern, demselben Zustand und derselben Dynamik. Eine Sammlung passender Zahlen, ein funktionierendes Hilfsmodell oder eine große Zahl grüner Prüfungen reicht dafür nicht.

Die Lesereihenfolge ist bewusst gestuft: Zuerst kommt das Gesamtbild; anschließend die ursprüngliche Kette und ihre neuen Verbindungen; danach Universalraum, Raum und physische Realisierung; schließlich Anwendungen, Algorithmen und die konkrete Abschlussagenda. Die ausführbare Quellenkonstruktion steht in den Abschnitten \ref{sec:levelselection}--\ref{sec:physicalfollow}. Die anschließende Gesamtprüfung in Abschnitt~\ref{sec:jointclosure} verbindet Ursprung, K3, Quantisierung, lokale Zeit und Teilchenstatistik und korrigiert die bislang zu unspezifische Abschlussagenda. Wer zunächst nur die Bedeutung verstehen möchte, kann die eingerahmten Kernaussagen und Abschnitte \ref{sec:bigpicture}, \ref{sec:jointclosure} und \ref{sec:completion} lesen.

\begin{tabularx}{\linewidth}{lY}
\toprule
Zeichen & Bedeutung in diesem Dokument\\\midrule
\E & Exakte mathematische Aussage oder exakte endliche Rechnung innerhalb ausdrücklich genannter Voraussetzungen. Das Zeichen behauptet keine vollständige Formalisierung in Lean.\\
\N & Numerisch oder durch einen konkreten ausgeführten Test bestätigt; Reichweite durch Eingaben, Genauigkeit und Abdeckung begrenzt.\\
\B & Bedingter Anschluss: Die Folgerung trägt, wenn die benannte Identifikation oder Zusatzannahme erfüllt ist.\\
\Opend & Offene Herleitung, physische Identifikation, Auswahl oder empirische Behauptung.\\
\Xclaim & Eine konkrete Gleichsetzung, Konstruktion oder stärkere Behauptung ist durch einen Gegenfall ausgeschlossen.\\\bottomrule
\end{tabularx}

Die gesamte Sitzung ist hier inhaltlich konsolidiert, nicht als chronologisches Protokoll wiedergegeben. Überholte Zwischenstände bleiben dort sichtbar, wo ihre Korrektur die Interpretation verändert. Die Anhänge benennen die verwendeten Quellenfamilien, den ausgeführten Prüfstand und die weiterhin offenen Anschlüsse. Der dokumentierte Explorer-Schnappschuss vom 28.~September enthält 316 interne Prüfungen; frühere Zwischenstände vom 29.~September enthielten 343 beziehungsweise 367, der Massentransportstand enthielt 372 und die Wand-Prüfung 375. Der aktuelle Lauf mit dem wirklichen Majorana-Operatoranschluss besteht alle 380 internen Prüfungen. Diese Zahlen bezeichnen weder voneinander unabhängige Beweise noch eine vollständige Prüfung aller Paper.

\tableofcontents
\newpage

\section{Das Gesamtbild: Eine Quelle und mehrere Ansichten derselben Wirklichkeit}\label{sec:bigpicture}

\key{\textbf{Das verbindende Bild ist ein Instrument mit einer einzigen inneren Mechanik.} Topologie legt fest, welche Schleifen und Markierungen unterscheidbar sind. Algebra legt fest, wie Operationen zusammenwirken. Der Zustand bestimmt ihre gemeinsamen Antworten. Rekursion beschreibt, wie dieselben Antworten auf einer anderen Auflösung repräsentiert werden. Physik verlangt zusätzlich, dass dieses Instrument tatsächlich unsere räumlichen, zeitlichen und messbaren Prozesse beschreibt.}

Die erneute Prüfung korrigiert die bisherige Gewichtung dieses Dokuments. Die ursprünglichen Paper haben den Weg von den Markierungen zu physikalischen Größen bereits weit ausgearbeitet. Die jüngere Sitzung hat dagegen vor allem einen gemeinsamen Feld- und Prozesskandidaten untersucht. Das Gesamtbild muss beide Richtungen verbinden: vom Ursprung zu den bekannten Auslesungen und von diesen Auslesungen zurück zu den erlaubten Ursprüngen.

\begin{center}
\begin{tikzpicture}[node distance=7mm]
\node[box,text width=13.5cm] (p) {\textbf{Ursprung und Markierungen}\\P1, P2, Double Cover, $\mu_4$, Anker $(1,1,2)$};
\node[box,below=of p,text width=13.5cm] (l) {\textbf{Gemeinsamer Träger}\\markierte $D_5+A_3$-/E8-Struktur, drei Familien, Ladungen};
\node[box,below=9mm of l,xshift=-3.6cm,text width=6.0cm] (a) {\textbf{Vorhandener Flavorweg}\\Familienmonodromie, $R,Q,\Sigma$, $K,L$\\Quellenmassen, Yukawa, CKM/PMNS};
\node[box,below=9mm of l,xshift=3.6cm,text width=6.0cm] (b) {\textbf{Vorhandener Nahtweg}\\$c_3$, $48$ Zustände, gewichtetes Budget $41$\\$\phi_0$, $\alpha$, Skalenrelationen};
\node[box,below=29mm of a,xshift=3.6cm,text width=13.5cm] (o) {\textbf{Bereits berechnete gemeinsame Bedingungen}\\gleicher Seed, gleiche Markierung, gleiche Normierung; exakte und bedingte Ergebnisse getrennt};
\node[box,below=11mm of o,text width=13.5cm] (s) {\textbf{Gesuchte vollständige gemeinsame Quelle}\\derselbe Operator, Zustand und Prozess müssen beide Wege realisieren\\Die neue E8-Level-1-Quelle ist daran zu prüfen};
\node[box,below=of s,text width=13.5cm] (r) {\textbf{Universalraum und physische Realisierung}\\antworttreue Geschichten, lokale Dynamik, RG-/Messabbildung und Raumzeit};
\draw[arr] (p)--(l);\draw[arr] (l)--(a);\draw[arr] (l)--(b);
\draw[arr] (a)--(o);\draw[arr] (b)--(o);
\draw[dasharr] (o)--node[right,font=\scriptsize]{gemeinsame Auswahl}(s);
\draw[dasharr] (s)--(r);
\end{tikzpicture}
\end{center}

Durchgezogene Pfeile bezeichnen hier vorhandene Herleitungs- und Rechenwege mit den im Original ausgewiesenen Voraussetzungen. Sie erklären nicht jeden Endwert zu einem voraussetzungslosen Satz. Gestrichelte Pfeile bezeichnen die noch nachzuweisende gemeinsame Auswahl beziehungsweise physische Realisierung. Insbesondere liegen Massen und $\alpha$ nicht erst hinter einer vollständig neuen Raumzeittheorie auf einer leeren Aufgabenliste.

\key{\textbf{Die starke Verbindung lautet:} Derselbe Bauplan liefert unterschiedliche, miteinander gekoppelte Antworten. Das ungewichtete Zählen der Materiezustände ergibt $48$, die elektromagnetische Gewichtung zusammen mit dem Higgs ergibt $41$, und der Flavortransport verbraucht genau $40$ Einheiten plus den Higgsbeitrag. Der Nahtparameter bestimmt zugleich Massenverhältnisse und Mischungen. Ein gemeinsamer Ursprung muss diese vollständigen Operatorbeziehungen erhalten.}

\picturetext{Die ursprüngliche Theorie enthält bereits Zahnräder und berechnete Übersetzungsverhältnisse. Die neue Quellenkonstruktion soll zeigen, wie diese Zahnräder auf derselben Welle sitzen und tatsächlich zusammen laufen. Man darf weder die bekannten Übersetzungen vergessen noch aus mehreren einzeln laufenden Rädern schließen, dass sie schon miteinander verbunden sind.}

Der neue Quellkandidat beantwortet ganze Folgen von Feldeinsetzungen. Seine Algebra, sein Zustand und seine konforme Zeit sind unter der erklärten Quantisierung gemeinsam bestimmt. Das liefert einen konkreten Kandidaten für den Universalraum. Ob seine markierten Projektionen tatsächlich die ursprüngliche Flavorfiltration, den elektromagnetischen Determinantenkanal und die physische Dynamik liefern, ist eine stärkere Frage. Sie muss an den schon vorhandenen Antworten entschieden werden.

Eine eindeutige Darstellung einer gewählten Quelle ist noch keine eindeutige Wahl der Quelle. Die ursprünglichen physikalischen Ergebnisse sind dabei zusätzliche Auswahlbedingungen. Sie werden in Abschnitt~\ref{sec:physicalconstraints} ausdrücklich eingesetzt; eine Diskussion allein über Zentralzahl, Vakuum oder Raumzeitdimension würde den Auftrag verfehlen.

Die jüngste Fortsetzung verbindet diese Geometrie mit einer wirklichen Operation. Zwei Halbwege um eine ursprüngliche Fuchs-Punktur ergeben genau den Sechsertransport: Nach zwei Halbschritten ist eine Familienrunde vollendet, nach sechs Halbschritten ist der ursprüngliche Zustand wieder erreicht. Die sechs Einträge bezeichnen dabei zwei Zugänge zu jeweils drei Familienkomponenten. Sie sind keine sechs neu eingeführten Teilchensorten. Auf der neutralen Seite lässt sich die 48er-Auswahl in eine gemeinsame geometrische Antwort und eine zusätzliche, vollständig berechnete Quellenantwort zerlegen. Gerade diese zusätzliche Antwort liefert einen konkreten Kandidaten für den gesuchten Alpha-Port. Ihre Kopplung an den ursprünglichen Seam muss an demselben Quellenfunktional geprüft werden, das die Flavor-Defekte trägt.

Die zusätzliche Spin-$\Z_4$-Auswahl aus dem neuen Anhang ist nun an dieser gesamten Kette geprüft. Sie kann unter ihren Voraussetzungen den kleinsten Fünferträger auswählen. Sie muss aber auch die Felder zulassen, die die bereits vorhandene Neutrinomasse erzeugen. Der neue interne Spinlift verbindet die ursprüngliche Windung direkt mit den geladenen E8-Sektoren. In derselben Quelle existiert ein symmetrischer Neutrinopaaroperator: Sein Familienweg passt an den ursprünglichen Flavortransport, und der neutrale Alpha-Kandidat antwortet auf ihn mit einem berechneten Koeffizienten. Abschnitt~\ref{sec:spinliftsynthesis} führt diese Verbindungen und den Verträglichkeitstest mit dem Massensektor zusammen.

Die Wand-Lesart vom 29.~September ist in Abschnitt~\ref{sec:wallaudit} an derselben Quelle geprüft. Ihr behaupteter Dirac-Zwang entfällt: Die symmetrischen Quellpaare tragen $B-L=\pm2$ und erfüllen die bosonische Spin-Ladungsrelation. Die geometrische Verbindung über $V=S_X^+$ liefert dagegen eine richtige bedingte Identität für Familienindex und Anomalieklasse. Zur gesuchten Gesamtlösung gehört die Auswahl der tatsächlichen Verbindung und physischen Phase, die zugleich die bestehenden TFPT-Ausgaben reproduziert.

Abschnitt~\ref{sec:masszeromode} schließt nun die konkrete Operatorverbindung:
Die geladenen Paarfelder führen auf dem ursprünglichen Stromraum genau die
Operationen aus, aus denen der vorhandene schwere Neutrinomassenblock besteht.
Blendet man diesen schweren Teil aus, entstehen aus demselben Block die
leichte Seesaw-Masse, ein Determinantenbeitrag und eine berechnete
Gedächtnisantwort. Bildlich: Das schwere Zahnrad bleibt wirksam, auch wenn es
nicht mehr sichtbar ist. Damit gehören Massenanschluss und Vergröberung zu
derselben Rechnung. Aus den Ursprungsregeln herzuleiten bleibt die
Betriebsweise, die die vorhandenen Massenkoeffizienten und den neutralen
$\alpha$-Kanal gemeinsam auswählt.

\subsection{Was eine wirkliche Gesamtlösung leisten müsste}
Die Fortsetzung in Abschnitt~\ref{sec:nativeclosure} schließt weitere konkrete Stücke dieser Kette: den globalen Spin-Lift des ursprünglichen $10+6$-Trägers, seine charakteristische Antwort und die von den vier Quadratmarken bestimmte Mehrintervallstruktur. Deren Replica-Antwort $3/2$ entsteht aus einem tatsächlichen Ising-Twistfeld der verdoppelten E8-Quelle. Das fügt dem Gesamtbild einen Operatoranschluss hinzu; der historische Dreizustandstransfer ist damit noch nicht identifiziert.

Abschnitt~\ref{sec:pairingcompletion} führt diesen Anschluss weiter: Die gesamte Clockmatrix wird aus den normierten Replica-Paarungen und einem erklärten Instrument auf vorhandenen Reflexionen rekonstruiert. Die Unterscheidung zwischen Quellenzustand, Ereignisrecord, Erwartungswertanzeige und tatsächlichem Messausgang ist dabei entscheidend. Gerade sie macht deutlich, wie derselbe Prozess sowohl einen reichen Universalraum als auch eine einfache Clockanzeige besitzen kann.

Abschnitt~\ref{sec:anchorfibre} verbindet diese ganze Replica-Familie mit dem Anker und einer tatsächlichen markierten Quellenwirkung. Er berechnet ihre kohärente Rückkehr, den Gedächtniskern und die orthogonale Erweiterung der älteren Auslese. Die neue Dreier-Gitterbeschreibung bleibt dabei als Ladungsfaserung typisiert: Sie enthält die vollständige Quelle, besitzt aber noch keine ausgewählte räumliche Bewegung oder relativistische Lokalordnung.

Abschnitt~\ref{sec:sharedtime} entscheidet die gemeinsame kontinuierliche Zeit der erklärten Clockhierarchie vollständig. Er weist außerdem nach, dass geteilte Zufallsereignisse keine Eingriffe zwischen Registern transportieren und dass der vorgeschlagene unmittelbare C-Hop-Lift nicht unitär ist. Die vollständige Quelle besitzt dagegen bereits einen lokalen Eingriffs- und Transportprozess mit berechneter Antwort und Energie. Damit wird die weitere Suche an ihren tatsächlichen lokalen Wirkungen ausgerichtet.

Abschnitt~\ref{sec:inverseorigin} schließt nun den konstruktiven Rückweg von der markierten E8-Quelle zu ihrem ursprünglichen $10+6$-Fermionenträger. Die Feldregel, der Vakuumzustand und die Konformalzeit kommen dabei gemeinsam zurück. Der bewiesene lokale Rekonstruktionssatz reduziert ihren Ursprung auf eine zusammenhängende Frage: Liefert der rohe P1/P2-Prozess genau die markierten lokalen Felder vom Gewicht $1/2$, die diese Quelle erzeugen? Die Antwort wird dadurch konkret prüfbar, ohne sie vorwegzunehmen.

Eine Lösung muss mindestens vier Aussagen gemeinsam tragen:
\begin{enumerate}
\item \textbf{Herkunft:} Die Ausgangsprinzipien wählen die verwendete Struktur, statt nur mit ihr verträglich zu sein.
\item \textbf{Ausführbarkeit:} Die Struktur bestimmt Zustände, Operationen, Komposition und alle relevanten Antworten konsistent.
\item \textbf{Physische Identifikation:} Interne Ladungen, Raumzeit, Zustand, Zeit und beobachtbare Größen sind durch ausdrückliche Abbildungen verbunden.
\item \textbf{Prüfbarkeit:} Dieselbe Konstruktion liefert die behaupteten Vorhersagen mit dokumentierten Voraussetzungen und besteht unabhängige Beobachtungstests.
\end{enumerate}
Die Sitzung hat die zweite Aussage für einen gehaltvollen Quellkandidaten weitgehend konkret gemacht und mehrere zuvor getrennte Anschlüsse exakt verbunden. Sie hat die übrigen Aussagen dadurch enger formuliert. Sie hat sie nicht durch die Existenz einer grafischen Oberfläche ersetzt.

\section{Der Ursprung: P1, P2 und die rekursive Rückbestätigung}

\subsection{P1: Orientierung, Rückkehr und positive Antworten}
P1 beginnt mit einem orientierten Seam und einem positiven beziehungsweise reflexionspositiven Kern. Die primitive Windung fixiert im verwendeten TFPT-Vertrag die topologische Normierung
\begin{equation}
c_3=\frac{1}{8\pi}.
\end{equation}
Eine positive Antwort bedeutet bildlich: Die aus der Quelle berechneten Normen dürfen nicht negativ werden. Orientierung bedeutet, dass Hin- und Rückweg sowie eine umgekehrte Windung nicht einfach dieselbe markierte Operation sind. ``Primitiv'' bezeichnet hier die ursprüngliche Windungs- beziehungsweise Unzerlegbarkeitsbedingung. Daraus folgt nicht ohne weiteren Satz die Perron--Frobenius-Primitivität irgendeines später eingeführten Vergröberungskanals.

\picturetext{Auf einem verdrehten Band kann man nach einer Runde am gleichen Ort sein und dennoch eine andere Orientierung mitbringen. Die Ortsanzeige allein beschreibt den Zustand nicht vollständig. Genau diese Art verlorener Information motiviert Double Cover, Phasenregister und Carry.}

Double Cover, Möbiusbild und Spinstruktur sind verwandte Motive, aber keine austauschbaren Beweise. Eine zweifache Überlagerung ist zunächst eine topologische Aussage; eine Spinordarstellung, eine Eichladung und eine tatsächlich messbare Interferenzphase benötigen zusätzlich die jeweils passende Wirkung auf Feldern und Zuständen.

\subsection{P2: Die markierte Aufteilung $3+2$}
P2 stellt die markierte Aufteilung des fünfdimensionalen Trägers bereit. Auf der fundamentalen Darstellung wird der Generator als
\begin{equation}
Y=\diag\left(-\frac13,-\frac13,-\frac13,\frac12,\frac12\right)
\end{equation}
geschrieben. Die Spurfreiheit ist exakt: $3(-1/3)+2(1/2)=0$. Bei der entsprechenden Trägernormierung entstehen $k_Y=5/3$ und die bekannte Hochskalenrelation $\sin^2\theta_W=3/8$. Das sind Aussagen an der markierten Trägerstruktur. Ihre Identifikation mit laufenden, renormierten Messgrößen erfordert den zugehörigen physikalischen Anschluss.

Die Sitzung hat das Vorzeichenproblem der neuen $X$-Felder ausdrücklich aufgelöst: Sie tragen $\overline{\mathbf5}\otimes\mathbf4$, also bezüglich P2 die entgegengesetzten fundamentalen Ladungen. Der Name ``Fünferblock'' darf dieses Dualisieren nicht verdecken.

\subsection{Anker, Familiengeometrie und Bootstrap}
Die ursprüngliche Ankerstruktur verwendet
\begin{equation}
a=(1,1,2),\qquad e_1=4,\quad e_2=5,\quad e_3=2,
\qquad p_n=2+2^n.
\end{equation}
Damit ergeben sich
\begin{equation}
p_1p_2p_3=4\cdot6\cdot10=240,\qquad p_4-p_3=18-10=8,
\qquad 240+8=248.
\end{equation}
Diese Zahlen passen zur Wurzelzahl, zum Rang und zur Dimension von $E_8$. Die tragfähige Identifikation kommt jedoch aus der tatsächlich konstruierten Gitter- und Gruppenwirkung, nicht allein aus diesen Gleichheiten.

Die vierfach punktierte projektive Gerade hat $H_1\cong\Z^3$; die ursprüngliche Familienlesart benutzt diese Rang-drei-Struktur. Ein Divisor vom Grad vier auf $\mathbb P^1$ besitzt fünf globale Schnitte. Die Markierungen $1,i,-1,-i$ geben das $\mu_4$-Register. Dies erklärt, weshalb Vierermarken, ein Fünferträger und drei unabhängige Schleifen in derselben Geometrie auftreten können, ohne dieselbe Art von Dimension zu bezeichnen.

Die Origin-Story und ihre Lean-Aussagen enthalten echte Rückbestätigungen: Unter den jeweiligen primitiven Determinanten-, Rang- und Yukawa-Voraussetzungen kehren $3+2$, die Ladungsrelationen und die entsprechenden Ränge wieder. Diese Rekursion darf nicht unterschlagen werden. Ihr logischer Inhalt ist aber eine Äquivalenz beziehungsweise ein Schluss innerhalb benannter Voraussetzungen. Ein Zyklus ``A liefert B, B liefert A'' beweist die Selbstkonsistenz von A und B; er beweist nicht von selbst, dass jede andere Ausgangsstruktur unmöglich ist.

\key{Die Bootstrap-Rückbestätigung ist ein struktureller Gewinn: Die Teile stützen dieselben Markierungen. Der nächste Ursprungssatz muss zeigen, welche konkurrierenden Quantisierungen oder Realisierungen durch genau diese Markierungen ausgeschlossen werden.}

\subsection{Was die erneute Prüfung des behaupteten Ursprungsabschlusses ergibt}\label{sec:originproofaudit}
Die stärkste archivierte Fassung minimiert auf zulässigen einseitigen spektralen Bordismen den Defekt
\begin{equation}
\mathfrak D(B)=\bigl(|\operatorname{SF}|,\rank_{\rm ess},\deg^+_{\det},h_\Sigma^{\rm red}\bigr).
\end{equation}
Die Wohlordnung dieser diskreten Werte liefert bei nichtleerer Klasse einen minimalen \emph{Wert}. Sie beweist nicht, dass genau ein vollständiger Operator diesen Wert annimmt. Das Original benennt die dafür notwendige Klassifikation der minimalen Faser selbst, führt sie im anschließenden Eindeutigkeitsbeweis aber nicht aus: \file{01_boundary_kernel_source.tex}, Zeilen 543--593 und 650--682, unter \file{_archive/tfpt-45/source_extracts/}. Die spätere Master-Barriere ist null für $\mathcal B\cong\mathcal B_{\min}$ und sonst unendlich. Sie setzt die ausgewählte Klasse damit schon ein. Eine rückgekoppelte Auswahl wäre grundsätzlich zulässig; dafür müssten unabhängig definierte Bedingungen tatsächlich gemeinsam gelöst und ihre Eindeutigkeit bewiesen werden.

Der konkrete ausführbare Abschluss \file{verification/v165_premise_a_closure.py} behauptet, aus dem DtN-/Calder\'on-Symbol $|k|$ folge bereits der freie Fermionen-Zweipunktkern. Der dafür ausgeführte Test prüft Rang-, Dimensions- und Indexgleichheiten. Sein Vorgänger \file{v162_seam_transport_identification.py} setzt die Rang-acht-Polarisation als $P=(I+iA_{16})/2$ ein und prüft ihre Projektoreigenschaft. Beide Rechnungen liefern keinen Intertwiner vom ursprünglichen Kragenoperator zu dieser Polarisation. Das Hauptsymbol allein bestimmt weder die vollständige Operatorwirkung und ihre Randbedingungen noch die behauptete Feldquantisierung. Der entsprechende \emph{Quellenabschluss} ist durch diese Tests daher nicht bewiesen.

Auch die originale Definition $Z_{\rm rel}=Z_{\rm adm}/Z_{\rm ref}[0]$ liefert keine zusätzliche Auswahl des Feldgesetzes. Sie normiert ein bereits gegebenes Funktional; die Reflexionspositivität seines CAR-Kerns ist im folgenden Lemma eine ausdrückliche Voraussetzung. Die Ableitungen von $\log Z_{\rm rel}$ definieren anschließend die Korrelatoren, aber noch nicht das zuvor benötigte $Z_{\rm adm}$ selbst. Maßgebliche Stellen sind \file{02_carrier_source.tex}, Zeilen 2257--2311, und \file{04_qft_source.tex}, Zeilen 1499--1515.

Der noch mögliche geometrische Reparaturansatz ist ebenfalls präzise begrenzt. Die skalare Flachheitsbedingung für $|D_\theta|+M_f$ wählt keine Verbindung auf dem gesamten $V_{16}$ aus. Der vorhandene Master variiert $(\alpha,\chi_{\rm geo},\delta_{\rm ph},\rho_{\rm vac})$, nicht die blockmischende Verbindung dieses Trägers. Eine flache Gesamtverbindung wäre geometrisch möglich; ihre relative Lage zu den markierten Teilbündeln müsste zusätzlich aus den ursprünglichen Bedingungen folgen. Sie hier einzusetzen wäre eine neue Annahme.

\textbf{Reichweite dieses Befunds:} Der behauptete Quellenabschluss ist an einer bestimmten Beweisstelle nicht getragen. Es ist damit weder die Unmöglichkeit von TFPT noch die Existenz zweier vollständiger Modelle unter sämtlichen gemeinsamen Bedingungen bewiesen. Der fehlende Satz bleibt die Auswahl des vollständigen markierten lokalen Operators und seines Quellfunktionals aus den Ursprungsbedingungen. Weitere Rechnungen innerhalb der schon gewählten E8-Quelle schließen diesen Satz nicht nachträglich.

\section{Vom Double Cover zum $E_8$-Träger und zum Standardmodell}

\subsection{Die Verklebung ist eine konkrete Ladungsrechnung}
Die alte Architektur benutzt die $D_5$-Spinorstruktur
\begin{equation}
S^+\cong\Lambda^{\mathrm{even}}\C^5,\qquad \dim S^+=16,
\end{equation}
und den $A_3$-Registerpartner. Die Diskriminantengruppen der verwendeten Gitter tragen passende $\Z_4$-Klassen. Die Normbeiträge $5/4$ und $3/4$ addieren sich zu $2$. Die Index-vier-Verklebung liefert das gerade unimodulare Rang-acht-Gitter $E_8$; seine vier ursprünglichen Wurzelklassen haben Größen $52+64+60+64$.

Die zugehörige Zerlegung der adjungierten Darstellung lautet
\begin{equation}
\mathbf{248}=(\mathbf{45},\mathbf1)\oplus(\mathbf1,\mathbf{15})
\oplus(\mathbf{16},\mathbf4)\oplus(\overline{\mathbf{16}},\overline{\mathbf4})
\oplus(\mathbf{10},\mathbf6).
\end{equation}
Die Gleichung verbindet eine tatsächliche Gruppenwirkung mit Ladungssektoren. Die oft verführerische Gleichheit $16=16$ wäre allein unzureichend: Feldkomponenten eines geometrischen Operators sind nicht automatisch die sechzehn chiralen Materiezustände einer Generation.

\subsection{Was die Architektur liefert und was die Physik zusätzlich verlangt}
Die Architektur liefert einen gemeinsamen Organisationsraum für interne Ladungen, Register und mögliche Materieblöcke. Anomaliefreiheit ist eine notwendige Konsistenzbedingung. Sie wählt nicht automatisch eine Wechselwirkung, entkoppelt keine Spiegelmaterie und präpariert keinen physischen Zustand.

Für den Standardmodellanschluss müssen daher die tatsächlichen Ladungsdarstellungen, die chiral wirksamen Freiheitsgrade, die lokalen Wechselwirkungen, der Zustand und die Raumzeitdynamik zusammenkommen. Die Sitzung benutzt die $E_8$-Kaskade ausdrücklich nicht als vorausgesetzten universellen Entstehungsmechanismus. Sie bleibt eine mögliche Geschichte unter zusätzlichen Voraussetzungen; der direkte Gitterquellenweg benötigt sie nicht.

\section{Die verzahnten Clocks, Flavor und $\alpha$ im Gesamtgefüge}

\subsection{Mehrere Uhren können dieselbe Struktur ablesen}
Die ursprünglichen Clock-Strukturen enthalten
\begin{equation}
h(E_8)=30=2\cdot3\cdot5,\qquad \varphi(30)=8,
\end{equation}
mit den Exponenten $1,7,11,13,17,19,23,29$. Daneben steht die vierfache Markierung. $\mu_4$ ist keine Untergruppe eines zyklischen 30-Taktes. Der tatsächliche Anschluss verwendet unter anderem eine Wirkung der Form
\begin{equation}
G C_5G^{-1}=C_5^2,\qquad G^4=I.
\end{equation}
Die komplexe Reflexionsgruppe $G_{31}$ hat Ordnung $46080$, sechzig Reflexionen und Invariantengrade $8,12,20,24$. Diese Angaben erklären gemeinsame diskrete Symmetrien; sie stellen noch keine Uhr mit einer physisch festgelegten Sekundenskala dar.

In einer verwendeten Clock-Zerlegung gilt $c=J\sigma_F$ von Ordnung zwölf, mit $\sigma_F=c^4$ und $J=c^9$. Zwei Operatoren gleicher Ordnung zwölf müssen nicht konjugiert sein und nicht gleich auf den geladenen Feldern wirken. In der Sitzung wurde deshalb die tatsächliche Wirkung statt nur der Periodenlänge verglichen.

\subsection{Die bereits vorhandene physikalische Herleitungskette}\label{sec:physicalconstraints}

\key{Die Aussage ``Massen, Flavor, Yukawa und $\alpha$ müssen noch hergeleitet werden'' war zu pauschal. Die Originale enthalten konkrete Formeln, Operatoren und Begründungen. Die präzise Frage lautet: Realisiert ein und derselbe ursprüngliche Quellprozess diese vorhandenen Gesetze mit denselben Markierungen, Normierungen und Zuständen?}

Die erneute Prüfung benutzt die aktuellen PDF-Fassungen und ihre TeX-Quellen, die Originalverifier und die späteren Quellverträge. Im Standardmodell-PDF stehen die geladenen Leptonen auf Seite 8, die Quarkverhältnisse auf Seiten 10--11 und der CKM-Anschluss auf Seite 29. Die unterschiedlichen Beweisstärken sind in den Originalen selbst markiert; sie werden hier weder pauschal abgesenkt noch zusammen zu einem bereits bewiesenen physikalischen Abschluss erklärt.

\begin{longtable}{L{0.19\linewidth}L{0.45\linewidth}L{0.27\linewidth}}
\toprule
Bereich & Vorhandenes Ergebnis & Genaue Reichweite\\\midrule\endhead
Geladene Leptonen & Potenzen von $\phi_0$ und Koeffizienten $16/7,4/3,7/6$ aus der Familienmonodromie & Exakte Quellenverhältnisse; Legzuordnung, Maßstab und RG-Abbildung mitführen\\
Quarks & Exakte Verhältnisse aus Plücker-/Compiler-Auslesungen; Starrheit auf dem diskreten Selektorstratum & Kein weiterhin offenes vollständiges Hitchin-Problem für diese Verhältnisse; absolute Amplituden gesondert prüfen\\
Yukawa und CKM & Gemeinsame $K/L$-Hierarchie, konkrete Mischungswinkel und führende CP-Phase & Höhere CP-Korrektur und vollständiger gemeinsamer Operatoranschluss getrennt\\
Neutrinos & Endlicher Dirac-/Majoranaoperator, Seesaw, PMNS-Kanäle, Scalaron-Skalenkandidat & Nicht alle Zweige sind derselbe Ansatz; schwerer Sektor und Zerfallsquelle müssen gemeinsam passen\\
Feinstrukturkonstante & Konkrete Determinantenantwort, Potential und eindeutige positive Nullstelle & Exakte Lösung der angegebenen Gleichung; Herkunft als volle physische Determinantenvariation gesonderter Vertrag\\
Gravitation und Skalen & Stationäres Planckpotential, EW-/Vakuumenergie-Elimination und Scalaronroute & Bedingte gemeinsame Funktional- und Matchingannahmen; keine beliebig unabhängigen Skalen\\
\bottomrule
\end{longtable}

\subsection{Massen und Yukawa-Kopplungen aus derselben markierten Struktur}
Mit Zeilen $u,d,e$ lauten die ursprünglichen Operatoren beziehungsweise markierten Auslesetabellen
\begin{equation}
R=\begin{pmatrix}1&3&0\\1&5&2\\2&5&3\end{pmatrix},\quad
K=\begin{pmatrix}4&2&0\\4&3&2\\5&3&2\end{pmatrix},\quad
L=\begin{pmatrix}7&3&0\\7&5&2\\8&5&3\end{pmatrix}.
\end{equation}
Es gilt $K=R+Q\Sigma$, $L=R+Q(I+\Sigma)$ und $\sum_{f,j}L_{f,j}=40$. Der Seed und das gemeinsame Yukawa-Gesetz sind
\begin{align}
c_3&=\frac1{8\pi},\qquad \phi_0=\frac1{6\pi}+48c_3^4,
\qquad \lambda_Y=\sqrt{\phi_0(1-\phi_0)},\\
\hat y_{f,j}&=\lambda_Y^{L_{f,j}}\Lambda_{f,j}
=\pi c_{f,j}\phi_0^{K_{f,j}},\qquad
\hat m_{f,j}=\frac{v_{\rm geo}}{\sqrt2}\hat y_{f,j}.
\end{align}
Hüte kennzeichnen die Quellenwerte. Die Standard-RG-/Schwellenabbildung zu laufenden oder Polmassen ist ein eigener, benannter Rechenschritt. Sie macht die Quellenherleitung nicht wertlos; Quellenmassen dürfen aber nicht ungeprüft mit Messmassen gleichgesetzt werden.

Besonders konkret ist die Leptonenrechnung. Die ursprüngliche Familienmonodromie $M_F$ aus v117/v118 erfüllt
\begin{equation}
\det(I-tM_F)=1-t^3,\qquad d_-:=\det(I-M_F/2)=7/8,
\quad d_+:=\det(I+M_F/2)=9/8.
\end{equation}
Damit ergeben sich
\begin{equation}
c_e=\frac4{2d_-}=\frac{16}7,\quad
c_\mu=\frac3{2d_+}=\frac43,\quad
c_\tau=\frac{4d_-}3=\frac76.
\end{equation}
Die Vereinigung der Spektren von $M_F$ und $-M_F$ liefert die sechs Phasen: Familienrotation und Blattvorzeichen gehören zusammen. Die Monodromie wird hier $M_F$ genannt, damit sie nicht mit der anderswo ebenfalls $M_0$ genannten Gramform verwechselt wird. Der Originaltest setzt die etablierten Familienlegs ein; er leitet ihre physische Zuordnung nicht nochmals aus der Determinante her.

Folglich gilt exakt
\begin{equation}
(\hat m_e,\hat m_\mu,\hat m_\tau)
=\frac{v_{\rm geo}\pi}{\sqrt2}
\left(\frac{16}7\phi_0^5,\frac43\phi_0^3,\frac76\phi_0^2\right).
\end{equation}
Dies liefert zwei besonders klare gemeinsame Bedingungen:
\begin{equation}\label{eq:physicalleptonelimination}
\boxed{\frac78\frac{\hat m_\mu}{\hat m_\tau}=\phi_0,
\qquad
\frac{\hat m_e\hat m_\tau^2}{\hat m_\mu^3}=\frac{21}{16}.}
\end{equation}
Die zweite Gleichung eliminiert sowohl den gemeinsamen Maßstab als auch den Seed. Sie ist eine direkte algebraische Konsequenz der vorhandenen Formeln, kein neuer unabhängiger Messbefund und keine beanspruchte mathematische Neuheit.

Auch die Quarkverhältnisse sind konkret:
\begin{equation}
\frac{\hat m_u}{\hat m_d}=\frac{55}{117},\quad
\frac{\hat m_c}{\hat m_s}=\frac{34}{47}\phi_0^{-1},\quad
\frac{\hat m_t}{\hat m_b}=\frac3{26}\phi_0^{-2}.
\end{equation}
Die Originalkennungen \file{FLAV.RIGID.01/02} erhalten ihre Bedeutung: Diese Verhältnisse sind auf dem hergeleiteten diskreten Selektorstratum starr. Die absolute Quarknormierung folgt nicht schon daraus, dass eine Prüfroutine eingesetzte Amplitudenprodukte wieder zurückrechnet. Genau diese stärkere Behauptung muss von den tatsächlich geschlossenen Verhältnissen getrennt bleiben.

Für CKM stehen im Original
\begin{equation}
s_{12}=\lambda_Y,\qquad s_{23}=\frac{\phi_0}{1+\lambda_Y},\qquad
s_{13}=\frac{\lambda_Y^3}{3},\qquad
\delta_{\rm CKM}=\frac\pi3+3\lambda_Y^2+\Delta_\triangle.
\end{equation}
Die führende Phase ist berechnet; $\Delta_\triangle$ bezeichnet die ausdrücklich noch nicht geschlossene höhere Korrektur. Der Koeffizient drei darf nicht nach einem günstigeren Datenvergleich ausgetauscht werden.

\subsection{Warum $48$, $40$ und $41$ eine gemeinsame Einschränkung sind}
Die Kopplung zwischen den beiden Herleitungswegen ist genauer als eine Zahlenähnlichkeit. Auf einer ursprünglichen Familie gilt
\begin{equation}
\tr_{16}Y^2=\frac{10}3,\qquad
\Omega_{\rm adm}=3\cdot16=48,\qquad
\gamma=\tr_{5}Y^2=\frac56.
\end{equation}
Ungewichtet ergeben die drei Familien den Punkturnachtrag $48c_3^4$. Mit Fermion-/Skalargewicht und Eichnormierung folgt dagegen
\begin{equation}
b_1=\frac35\left(\frac23\,3\,\frac{10}3+\frac13\,\frac12\right)
=\frac{41}{10},\qquad
\sum L=48\gamma=40,\quad \sum L+N_\Phi=41.
\end{equation}
Der Higgs trägt den letzten Beitrag bei. Ein ungeladenes rechtshändiges Neutrino trägt nicht zur $Y^2$-Spur bei, gehört aber zum ungewichteten Zustandsbestand. Es darf daher nicht aus der 48-Zählung entfernt werden, nur weil die elektromagnetische Spur seine Entfernung nicht bemerkt. $5/3$ ist der Eich-Einbettungsindex; $5/4=(3/2)(5/6)$ ist die Nahtsuszeptibilität. Beide haben unterschiedliche Rollen.

Der elektromagnetische Ansatz enthält bereits einen tatsächlichen Operator:
\begin{align}
U_\Sigma^{\min}&=|\psi_\Sigma\rangle\langle\psi_\Sigma|,
\quad \norm{\psi_\Sigma}=1,\qquad q=48c_3^4e^{-2\alpha},\\
\mathcal W_\Sigma&=-\frac54\tr\log(I-qU_\Sigma^{\min})
=-\frac54\log(1-q),\\
\phi_s(\alpha)&=\frac1{6\pi}+q(1-q)^{-5/4}.
\end{align}
Rang eins bezeichnet den ausgewählten primitiven Transferkanal, nicht die gesamte Quelle. Der unresummierte Flavorseed $\phi_0$ und $\phi_s(\alpha)$ sind verwandte, aber unterschiedliche Antworten.

Auch das Potential ist ausdrücklich vorhanden:
\begin{align}
\Gamma_{U(1)}(\alpha)&=\frac{\alpha^4}{4}-\frac23c_3^3\alpha^3
-8b_1c_3^6\int_0^\alpha\log\bigl(1/\phi_s(a)\bigr)\,da,\\
\partial_\alpha\Gamma_{U(1)}&=
\alpha^3-2c_3^3\alpha^2-\frac45\,41c_3^6\log(1/\phi_s)=0.
\end{align}
Die positive Lösung ist $\alpha^{-1}=137.0359992168407\ldots$. Die neue Prüfung bestätigt diese interne Schließung. Der Originalvertrag \file{ALPHA.QUILLEN.EXACT.01} fragt genauer nach ihrer Herkunft als vollständige relative Determinantenvariation mit dem richtigen lokalen Gegenbeitrag. Ein vorhandenes Potential zu verneinen wäre falsch; es allein durch Integration der Zielgleichung physikalisch herzuleiten ebenso.

\picturetext{Einmal zählt man alle Plätze in einem Fahrzeug. Ein zweites Mal zählt man dieselben Plätze mit Gewichten nach ihrer Ladung. Die beiden Summen sind verschieden, aber man kann die Sitzordnung nicht in jeder Rechnung neu erfinden. Genau so müssen $48$, $41$, die Flavorzuordnung und die Nahtantwort aus demselben Bestand kommen.}

\subsection{Die bekannten Formeln schränken den Ursprung rückwärts ein}
Auf dem verwendeten kleinen Seedzweig folgt bereits
\begin{equation}
\phi_0=\frac78\frac{\hat m_\mu}{\hat m_\tau}
=\frac{1-\sqrt{1-4s_{12,\rm CKM}^2}}2.
\end{equation}
Die ursprünglichen weiteren, jeweils bedingt physisch gelesenen Kanäle lauten
\begin{equation}
\phi_0=e^{5/6}\sin^2\theta_{13,\rm PMNS}
=4\pi\beta_{\rm rad}
=\frac{4\pi}{4\pi-1}\Omega_b.
\end{equation}
Diese Gleichheiten verknüpfen dieselben Quellenformeln; sie sind kein heute neu ausgeführter globaler Fit an Beobachtungsdaten. Ihre Abhängigkeiten müssen mitgeführt werden, statt jede Auslesung als unabhängige Bestätigung zu zählen.

Die Rückrichtung ist dennoch streng: Für $c>0$ ist
\begin{equation}
\phi(c)=\frac43c+48c^4,\qquad \phi'(c)=\frac43+192c^3>0.
\end{equation}
Ein konsistenter positiver Flavorseed wählt somit innerhalb dieser bereits festgelegten Funktionsklasse genau ein $c$. Zusammen mit dem Budget $41$ ist dann auch die angegebene $\alpha$-Schließung festgelegt. Das eliminiert unabhängige Nachjustierungen zwischen den Sektoren. Es beweist keine Eindeutigkeit jedes vollständigen Quelloperators, der diese endlich vielen Antworten reproduziert.

\subsection{Auch die Skalen besitzen einen gemeinsamen Variationsansatz}
Im Metrologieoriginal ist die Planckfrage stärker formuliert als die bloße Wahl einer Einheit. Mit der ersten positiven Nahtfrequenz $\lambda_\Sigma$ und $\rho=\chi^2/\lambda_\Sigma^2$ lautet das Potential in Rand-Einheiten
\begin{equation}
\widehat V(\rho)=-6\widehat{\mathcal E}_D\rho
+\widehat\Gamma_D^{(4)}(\rho)
+\rho^2[-\log\det_{\rm adm}(I-\widehat U_\Sigma(\rho))].
\end{equation}
Unter $\widehat V'(0^+)<0$, $\widehat V''(\rho)>0$ für $\rho>0$ und $\widehat V'(\rho)\to+\infty$ gibt es genau ein positives stationäres $\rho_\star$. Dann gilt
\begin{equation}
\frac{\bar M_{\rm Pl}^2}{\lambda_\Sigma^2}=\frac{\rho_\star}{2\pi^2},\qquad
G_N\lambda_\Sigma^2=\frac{\pi}{4\rho_\star}.
\end{equation}
Das ist ein wirklicher bedingter Eindeutigkeitssatz. Sein Zahlenwert verlangt den konkreten Operator und die Prüfung der genannten Voraussetzungen. Eine Konvexitätsannahme allein berechnet $\rho_\star$ noch nicht.

Der ursprüngliche EW-/Vakuumenergieanschluss liefert, mit $A=\alpha^{-1}$,
\begin{equation}
\frac{v_{\rm EW}}{\bar M_{\rm Pl}}=5\beta_{\rm rad}^2e^{-(A+\delta_{\rm ph})/5},
\qquad \frac{\rho_\Lambda}{\bar M_{\rm Pl}^4}=\frac3{4\pi^2}e^{-2A}.
\end{equation}
Elimination derselben Größe $A$ ergibt exakt
\begin{equation}\label{eq:physicalscalebridge}
\boxed{\frac{\rho_\Lambda}{\bar M_{\rm Pl}^4}
=\frac3{4\pi^2}e^{2\delta_{\rm ph}}
\left(\frac{v_{\rm EW}}{5\beta_{\rm rad}^2\bar M_{\rm Pl}}\right)^{10}.}
\end{equation}
Die Branchwahl, Higgsnormalisierung und der Transportpol bleiben die ausdrücklich zu identifizierenden Eingänge. Innerhalb dieses Vertrags sind die beiden Skalen nicht unabhängig einstellbar.

Die Scalaronroute benutzt zusätzlich $M_{\rm scal}/\bar M_{\rm Pl}=c_3^{7/2}$ und liefert bedingt $A_s=N_*^2c_3^7/(24\pi^2)$, $n_s=1-2/N_*$, $r=12/N_*^2$. $N_*$ bleibt Szenarioeingang. Bei ihrer stärksten Normierungsbegründung ist eine konkrete Prüfung offen: Der in v36 isolierte $R^2$-Summand ist nicht der vollständige $a_4$ des dort verwendeten gewöhnlichen Spin-Diracoperators. Dessen vollständige gravitative Kombination ist proportional zu
\begin{equation}
5R^2-8R_{\mu\nu}^2-7R_{\mu\nu\rho\sigma}^2=11E_4-18C^2,
\end{equation}
modulo Divergenzen. Der relative, projizierte oder torsionshaltige TFPT-Masteroperator muss daher seinen vollständigen Koeffizienten liefern; die Standardrechnung allein liefert die behauptete Starobinsky-Normierung nicht \cite{chamseddineconnes}. Das grenzt diesen konkreten Begründungsweg ein und verwirft nicht die anderen Skalenrelationen.

\subsection{Der gemeinsame Operator ist im Original bereits präzise vorgegeben}
Der Flavoransatz schreibt tatsächlich
\begin{equation}
P_{\rm str}=P_{\Sigma,+}P_{\rm sing}P_\Theta,\quad
K_f^{\rm str}=(P_{\rm str}D_f^\dagger D_fP_{\rm str})^{-1/2},\quad
Y_f=\operatorname{Hol}(\nabla_F^\star)K_f^{\rm str},
\end{equation}
auf dem zulässigen, invertierbaren Projektionsraum. Der Transportpol wird über das ursprüngliche Determinantenfunktional von $D_{y,\delta}=yI-\delta U_6$ bestimmt. Für die neue Quelle ist deshalb die konkrete Gleichheit $Y_f^{\rm source}=Y_f$ mit denselben Familienlegs, $K/L$-Gradierungen und derselben Norm zu zeigen. Die neue E8-Quelle enthält nicht automatisch die Flavorhierarchie, nur weil sie die richtige Stromalgebra besitzt.

Zugleich muss ihre primitive Nahtauslesung das oben angegebene $U_\Sigma^{\min}$ und die $\alpha$-Variation ergeben; die Skalenvariation muss das genannte Planckpotential erzeugen. Die ursprüngliche Forderung nach einer gemeinsamen relativen Determinantenfamilie verbindet damit bereits vorhandene konkrete Gesetze. Sie ist wesentlich enger als ``finde irgendein Funktional''.

\key{\textbf{Veränderte Arbeitspriorität:} Die fundamentale Quelle muss am ursprünglichen markierten Flavor- und Nahtoperator identifiziert werden. $M_F$, $Q,\Sigma,K,L$, die $48/41$-Gewichte und die Skalenvariationen sind dabei gleichzeitig zu erhalten. Weitere Rechnungen nur an der neuen symmetrischen E8-Quelle ersetzen diesen Anschluss nicht.}

\subsection{Neutrinos: Die vorhandenen Herleitungen gemeinsam lösen}\label{sec:neutrinojoint}

Der Originalbestand liefert einen endlichen Diracoperator mit Dirac- und Majoranablock, den tatsächlichen Seesaw-Ausdruck und konkrete PMNS-Kanäle. Auf dem ursprünglichen Zweig stehen unter anderem $\sin^2\theta_{12}=1/3-\phi_0/2$, $\sin^2\theta_{13}=\phi_0e^{-5/6}$ und die hexagonale Phasenlesart $\delta_{\rm PMNS}=240^\circ$. Der spätere Pentagonzweig aus \file{FLAV.NUSCALE.06} benutzt dagegen eine andere Mischung und Phase; er ist ein ausdrücklich bedingter Alternativkandidat. Beide Zweige dürfen nicht zu einer einzigen voraussetzungslosen Herleitung zusammengezogen werden.

Die Scalaronverbindung ist eine echte sektorübergreifende Bedingung:
\begin{equation}
M_3=3M_{\rm scal},\qquad
M_{\rm scal}=c_3^{7/2}\bar M_{\rm Pl},\qquad y_{\nu,3}=y_t
\quad\text{im erklärten minimalen Matching.}
\end{equation}
Der Original-Ein-Schleifenlauf ergibt damit $m_{3,0}=0.05124031018\,\mathrm{eV}$. Die gemeinsame Einbettung dieses Matchings, der vollständigen Mischung und der Materieasymmetrie muss jedoch denselben Yukawaoperator verwenden.

Das wird am eingefrorenen späteren Kandidaten besonders konkret. Er setzt
\begin{equation}
M_R=M_{\rm scal}\diag(\epsilon,2\epsilon,3),\quad
\epsilon=\frac{\phi_0^2}{10},\qquad
Y_\nu=\diag(0,y_2,y_3)U^\dagger.
\end{equation}
In dieser Konvention tragen die Zeilen die schweren Neutrinos. Für jedes unitäre $U$ ist dann
\begin{equation}
Y_\nu Y_\nu^\dagger=\diag(0,|y_2|^2,|y_3|^2).
\end{equation}
Das leichteste schwere Neutrino ist entkoppelt. Sein Washout und seine Zerfallsquelle verschwinden. Der separat erfolgreiche Boltzmannlauf verwendet dagegen $\widetilde m_1=m_3/10$ sowie einen CP-Quellterm aus einer Davidson--Ibarra-Schranke mal $|\sin240^\circ|$. Diese beiden Teilrechnungen besitzen daher noch keinen gemeinsamen Yukawaoperator. Die Schranke ist nicht automatisch eine Gleichheit und die PMNS-Phase nicht automatisch die schwere Zerfallsphase \cite{davidsonibarra}.

\textbf{Die Fortsetzung benutzt einen bereits vorhandenen Originalansatz.} Die Verträge \file{NU.CI.01/02} beschränken die allgemeine Casas--Ibarra-Zuordnung \cite{casasibarra} durch Familienzyklus, Nahtreflexion, komplexe Orthogonalität und Determinante eins. Zur Unterscheidung vom Flavorresiduum $R$ heißen ihre Matrix und ihr Fouriertransformator hier $R_{\rm CI}$ und $F_3$:
\begin{align}
R_{\rm CI}(z)&=F_3\diag(1,e^z,e^{-z})F_3^\dagger,
\quad (F_3)_{jk}=\frac{\omega^{jk}}{\sqrt3},\quad \omega=e^{2\pi i/3},\\
R_{\rm CI}^{\mathsf T}R_{\rm CI}&=I,\qquad z\in\R.
\end{align}
Dies ist der identitätsverbundene reelle Zweig. Die zweite reelle Komponente mit beiden nichttrivialen Fourier-Eigenwerten negativ wird unten durch die Washoutbedingung ausgeschlossen. Die Symmetriebedingungen allein wählen noch keinen Wert von $z$.

Mit $v_D=v_{\rm EW}/\sqrt2$ lautet die gemeinsame Matrix am selben Matchingmaßstab
\begin{equation}
Y_\nu=\frac{i}{v_D}\sqrt{M_R}\,R_{\rm CI}(z)\sqrt{m}\,U^\dagger,
\qquad
-v_D^2Y_\nu^{\mathsf T}M_R^{-1}Y_\nu=U^*mU^\dagger.
\end{equation}
Das setzt ein gegebenes Niederenergiespektrum korrekt fort; es behauptet noch keine unabhängige Auswahl dieses Spektrums aus dem Ursprung. Im summierten schweren Kanal fällt $U$ aus $Y_\nu Y_\nu^\dagger$ heraus. Bei einzelnen Leptonflavors gilt diese Vereinfachung nicht.

Schreibe $c=\cosh z$, $h=\sinh z$. Die erste Zeile ist $(d,a-ib,a+ib)$ mit
\begin{equation}
d=\frac{1+2c}{3},\qquad a=\frac{1-c}{3},\qquad b=\frac{h}{\sqrt3}.
\end{equation}
Für $m_1=0$ erhält man exakt
\begin{equation}
\widetilde m_1=(m_2+m_3)f(z),\qquad
f(z)=\frac{(c-1)^2+3h^2}{9}.
\end{equation}
Der bestehende bedingte Anker $\widetilde m_1=m_3/10$ legt daher den Betrag fest:
\begin{equation}\label{eq:nuactualselection}
\boxed{\cosh z_\star=
\frac{1+3\sqrt{1+4m_3/[10(m_2+m_3)]}}4.}
\end{equation}
Es bleiben genau die beiden Orientierungen $z=\pm z_\star$. Die zweite reelle Komponente hat mindestens $\widetilde m_1=4(m_2+m_3)/9$ und kann den Anker nicht erfüllen. Der zuvor hypothetisch eingesetzte Selektor $z=6\log(3/2)$ ist damit tatsächlich prüfbar: Für den eingefrorenen Massensatz liefert er $0.7871\,\mathrm{eV}$ Washout statt $0.005124\,\mathrm{eV}$. Diese zwei bestehenden Bedingungen passen in dieser Klasse nicht zusammen.

\textbf{Washout und CP sind jetzt gemeinsam gekoppelt.} Der tatsächliche CP-Zähler lautet
\begin{equation}
\operatorname{Im}\sum_jm_j^2(R_{\rm CI})_{1j}^2
=2ab(m_3^2-m_2^2).
\end{equation}
Daraus folgt in der hierarchischen ungeflavorten SM-Näherung
\begin{equation}
\varepsilon_1^{\rm hier}
=\frac{3M_1(m_3-m_2)}{16\pi v_D^2}\,
\frac{\sqrt3\sinh z}{2\cosh z+1},
\end{equation}
mit einheitlichen Masseneinheiten und der angegebenen CP-Vorzeichenkonvention. Die gleiche Norm, die den Washout festlegt, bestimmt damit den CP-Betrag. Besonders kompakt lässt sich $z$ vollständig eliminieren. Setze
\begin{equation}
w=\frac{\widetilde m_1}{m_2+m_3},\qquad s=\sqrt{1+4w},\qquad
\mathcal C=\frac{16\pi v_D^2\varepsilon_1^{\rm hier}}{3M_1(m_3-m_2)}.
\end{equation}
Dann gilt innerhalb dieser ursprünglichen Klasse exakt
\begin{equation}\label{eq:nuwashoutcp}
\boxed{\mathcal C^2=\frac{w(3s+5)}{(s+1)^3}.}
\end{equation}
Für endliches $w$ bleibt $|\mathcal C|<\sqrt3/2$. Der früher eingesetzte Betrag $|\sin240^\circ|=\sqrt3/2$ wird in dieser Familie erst bei $|z|\to\infty$ erreicht, während zugleich der Washout divergiert. Er ist hier eine Grenzamplitude, keine unabhängig einsetzbare endliche Phase. Diese Schlussfolgerung benutzt die vorhandene Familie und den vorhandenen Anker; sie führt keine neue angepasste Textur ein.

\textbf{Die Rückwirkung auf die Masse wird mitgelöst.} Würde man den alten Wert $m_{3,0}$ unverändert behalten, ergäbe die Washoutgleichung $|z|=0.4829422092$. Dann stiege aber die dritte schwere Yukawa-Zeilennorm um den Faktor $1.08594$; insbesondere wäre auch der größte Singularwert größer als der alte Top-Matchingwert. Deshalb darf der alte Massennahetreffer nicht zusammen mit dieser neuen Matrix als unverändert bestätigt ausgegeben werden.

Die ursprüngliche Einfamilienrechnung definiert keine vollständige $3\times3$-Basiszuordnung. Als notwendige basisinvariante Fortsetzung ihres dominanten Yukawa-Matchings verlangen wir
\begin{equation}
\sigma_{\max}(Y_\nu)=y_t(M_3).
\end{equation}
Dies ist ausdrücklich schwächer als eine vollständig begründete Matrixgleichheit $Y_\nu=Y_u$ im minimalen $10_H$-Modell. Mit der im Original verwendeten gemeinsamen Rückskalierung $R_\downarrow=0.7944014077$ kürzt sich deren gemeinsamer Faktor aus der relativen Normgleichung heraus. Setze
\begin{equation}
S=\diag\bigl(\sqrt{M_1/M_3},\sqrt{M_2/M_3},1\bigr),
\qquad z=z(m)\text{ aus Gleichung~\eqref{eq:nuactualselection}}.
\end{equation}
Dann lautet das gleichzeitig zu lösende System kompakt
\begin{equation}\label{eq:nujointmassnorm}
\boxed{\lambda_{\max}\!\left[
S R_{\rm CI}(z(m))\diag(0,m_2,m)R_{\rm CI}(z(m))^\dagger S
\right]=m_{3,0}.}
\end{equation}
Die nichtverschwindenden Eigenwerte stammen aus einer $2\times2$-Gramform. Auf dem positiven Zweig wachsen deren Diagonaleinträge und der Betrag ihres Nebendiagonaleintrags streng mit $m$. Ihr größter Eigenwert wächst daher streng, beginnt bei $(M_2/M_3)m_2<m_{3,0}$ und geht gegen unendlich. Es gibt genau eine positive Massenlösung; die beiden Orientierungen bleiben konjugiert.

Mit den unveränderten schweren Massen aus der eingefrorenen v3-Datei und deren $m_2=0.008654478609\,\mathrm{eV}$ folgt
\begin{equation}
\boxed{m_3=0.0435180861582\,\mathrm{eV},\quad
z=\pm0.477382171162,\quad
\widetilde m_1=0.00435180861582\,\mathrm{eV}.}
\end{equation}
Dabei ist $m_2$ im Original ausdrücklich daten-kalibriert. Washout-Anker, CI-Symmetrieklasse und Scalaron-Matching bleiben bedingt. Dies ist eine eindeutige gemeinsame Konsistenzlösung der angegebenen notwendigen Bedingungen, keine neue unabhängige Neutrinovorhersage und kein voller $10_H$-Abschluss. Die vollständige $3\times3$-Rechnung und die reduzierte Gramrechnung wurden unabhängig gegengeprüft.

Im vorhandenen Boltzmannpfad wurde danach die tatsächliche CP-Quelle dieser Matrix verwendet. Unter der bisherigen ungeflavorten $N_1$-Näherung ergibt sich:
\begin{center}
\begin{tabular}{lrr}
\toprule
 & Hierarchische Schleife & Endliche Massenschleifen\\\midrule
$|\varepsilon_1|$ & $1.57943\cdot10^{-7}$ & $1.36236\cdot10^{-7}$\\
$\kappa_f$ & $0.1129051$ & $0.1129051$\\
$|\eta_B|$ & $1.71193\cdot10^{-10}$ & $1.47665\cdot10^{-10}$\\\bottomrule
\end{tabular}
\end{center}
Die endlichen, nichtresonanten Ein-Schleifenfunktionen \cite{covirouletvissani} sind hier wichtig, weil $M_2=2M_1$ keine starke Hierarchie ist. Der in dem alten Programm gespeicherte Vergleichswert beträgt $6.1\cdot10^{-10}$; der gemeinsame Anschluss erreicht in dessen Näherung rund $24\,\%$ davon. Dieser Vergleich benutzt den eingefrorenen Datenstand, keine heute neu recherchierte Messung.

Die Rechnung erfolgte direkt mit \file{kappa_f_ode} im bestehenden \file{experiments/ftransfer/leptogenesis_boltzmann/fboltzmann_solve.py}. Die ursprünglichen schweren Massen wurden aus \file{experiments/nu-scalaron-falsification/hypotheses/nu_scalaron_v3.yaml} übernommen. Die Algebra benutzt \file{nu01_casas_ibarra_seam.py} und prüft die Hypothese aus \file{nu02_transport_selector.py} im Verzeichnis \file{experiments/theory-contracts/}. Eine vollständige Flavor-Dichtematrixrechnung mit mehreren schweren Spezies und Schwellen-RG wurde nicht ausgeführt. Die ODE übernimmt bewusst die bisherige Niedrigenergie-/Matching-Näherung des Originals. Insbesondere beweist das verlorene Nahetreffen bei dessen großer RG-Schemenunsicherheit noch keinen Ausschluss der gesamten Leptogeneseroute. Auch die physische Zuordnung des P1-Orientierungszeichens zu positivem Baryonenüberschuss ist dadurch nicht hergeleitet.


\picturetext{Für dasselbe schwere Neutrino sind Kopplungsstärke, erzeugte Asymmetrie und ihr späteres Auswaschen miteinander verbunden. Man kann diese drei Regler nicht einzeln aus drei erfolgreichen Rechnungen übernehmen. Die gemeinsame Gleichung zeigt nun genau, welche Einstellungen zugleich möglich sind.}

\textbf{Bedeutung für die fundamentale Frage:} Die vorhandene Neutrinogeometrie liefert bereits den nötigen nichttrivialen Kanal. Ihre gemeinsamen Bedingungen lösen einen Teil der bisher freien Zuordnung und schließen eine konkrete ältere Gap-Identifikation aus. Die ursprünglichen Massen-, Mischungs- und Kosmologieansätze sind dadurch engere Anforderungen an dieselbe Quelle geworden. Der Befund betrifft diesen gekoppelten Ansatz und diese Transfernäherung; er widerlegt weder den gesamten Neutrinozweig noch den unabhängigen kosmologischen $\Omega_b$-Readout.



\subsection{Nachweise der erneuten Prüfung}
Die Originalkennungen sind insbesondere \file{FLAV.LEPTONC.01}, \file{FLAV.RIGID.01/02}, \file{FLAV.CKM.01}, \file{FLAV.PMNS.02}, \file{FLAV.NUSCALE.05/06}, \file{EM.WARD.01} und \file{ALPHA.QUILLEN.EXACT.01}. Neu ausgeführt wurden v20 (12/12), v71 (7/7), v118 (6/6), v88 (11/11), v3 (14/14) und v485 (5/5); diese Läufe prüfen ihre angegebenen Formeln und Operatorrechnungen. Eine neue Ausführung der gesamten Lean- oder Repository-Suite wird damit nicht behauptet.

Maßgebliche Vertiefungen: \file{tfpt_2_standard_model.tex}, Zeilen 131, 310, 450, 573, 1751 und 1862; \file{tfpt_1_architecture_e8.tex}, Zeilen 3968--4050; \file{_archive/tfpt-45/source_extracts/03_em_flavor_source.tex}, Zeilen 94--230, 811--846 und 1257--1273; \file{05_metrology_source.tex} im selben Ordner, Zeilen 622 und 1370--1525. Das spätere \file{source-flavor-joint-dictionary-20260921/PROOF.md} erhält die vollständigen markierten Operatorprodukte. Sein endliches Wörterbuch ersetzt noch nicht die Auswahl ihrer physischen Quellenmetrik. Der ältere Archiv-Monodromiesatz wird hier nicht für die vollständige Auswahl verwendet: Seine drei angegebenen konjugiert-diagonalen Monodromien multiplizieren sich zu $I$, wodurch die vierte die gleichzeitig verlangte nichttriviale Punkturspektralbedingung verletzt. Die ausführbare v117/v118-Familienmonodromie bleibt die geprüfte Grundlage der hier erhaltenen Leptonformeln.

\section{Fortsetzung: Der gemeinsame Flavor- und Nahtanschluss}\label{sec:operatorbridge}

\key{\textbf{Die gemeinsame Verbindung ist konkreter geworden.}
Das ursprüngliche Fuchssystem liefert jetzt einen wirklichen Wegoperator, der exakt zum bereits verwendeten Sechsertransport äquivalent ist. Auf der anderen Seite liefert die ursprüngliche E8-Quelle eine konkrete neutrale Anregung mit Gewicht zwei und Normquadrat 48. Die gemeinsame Aufgabe ist dadurch keine Suche nach irgendeiner passenden Matrix mehr: Es ist die Herleitung der angegebenen Nahtkopplung und ihrer gemeinsamen Antwort mit den ursprünglichen Flavor-Defekten.}

\subsection{Der ganze endliche Operator besitzt bereits ein gemeinsames Wörterbuch}
Der Originalvertrag \file{source-native-clock-flavor-bridge-20260922/PROOF.md} liefert eine feste invertierbare Matrix $S$ mit
\begin{equation}
S^\dagger S=N,\qquad SPS^{-1}=M,\qquad
S(-G)S^{-1}=U,\qquad Q_c=SQS^{-1}.
\end{equation}
$N$ ist die dort ausgewählte positive Familienform, $P$ der ganzzahlige Dreierzyklus, $M$ die Familienmonodromie und $U$ die Vierteldrehung. Die Orientierungsgerade ist in $-G$ ausdrücklich mitgeführt. Für den verdoppelten Raum folgt deshalb
\begin{equation}\label{eq:fullflavorintertwiner}
(S\oplus S)\bigl[yI-\delta(P\oplus(-P))\bigr](S\oplus S)^{-1}
=yI-\delta(M\oplus(-M)).
\end{equation}
Mit dem zu $N$ gehörenden Adjungieren werden gleichzeitig $D^\dagger D$, Resolventen, inverse Quadratwurzeln, Polarfaktoren und Determinanten transportiert. Das ist stärker als die Gleichheit von Eigenwertlisten. Es ist keine neue freie Wahl eines Flavoroperators.

Die ursprüngliche positive Transportantwort und der inverse Operator bleiben verschieden:
\begin{equation}
D=V|D|,\qquad K=|D|^{-1},\qquad
(D^\dagger D)^{-1}=K^2,\qquad D^{-1}=KV^\dagger .
\end{equation}
Die zusätzlich im Yukawagesetz benutzte Familienholonomie ist dadurch noch nicht als $V^\dagger$ identifiziert. Ebenso ist die ursprüngliche Wortgewichtung $\lambda_Y^L$ kein automatischer Ersatz der geometrischen Reihe in $\delta/y$. Beide tragen eigene, bereits markierte Daten.

\subsection{Die Vierer- und Sechsergeometrie gehören zu demselben Überlagerungsproblem}\label{sec:flavorcover}
\textbf{Zunächst lässt sich die bisher numerische Fuchs-Monodromie analytisch schließen.}
Das ursprüngliche System aus v117 ist $\Psi'=A(z)\Psi$, mit
\begin{equation}
A(z)=\sum_{k=0}^3\frac{U^kA_0U^{-k}}{z-i^k},\qquad
A_0=\begin{pmatrix}
1/2&\sqrt2/6&0\\
\sqrt2/6&1/4&\sqrt5/12\\
0&\sqrt5/12&1/4
\end{pmatrix},\qquad U=\diag(1,i,-i).
\end{equation}
Die Substitution $x=z^4$, $\Psi(z)=\diag(1,z,z^3)v(x)$ liefert exakt
\begin{equation}
\frac{dv}{dx}=\left(\frac{B_0}{x}+\frac{B_1}{x-1}\right)v,\qquad
B_0=\begin{pmatrix}
0&0&0\\-\sqrt2/6&-1/4&0\\0&-\sqrt5/12&-3/4
\end{pmatrix},\quad B_1=A_0,
\end{equation}
und
\begin{equation}
B_\infty=-B_0-B_1
=\begin{pmatrix}
-1/2&-\sqrt2/6&0\\0&0&-\sqrt5/12\\0&0&1/2
\end{pmatrix}.
\end{equation}
Die lokalen Exponenten bei $0,1,\infty$ sind somit
$(0,-1/4,-3/4)$, $(0,1/3,2/3)$ und $(-1/2,0,1/2)$.

Bei $\infty$ muss die ganzzahlige Exponentendifferenz geprüft werden; die Eigenwerte allein genügen dort nicht. Mit
\begin{equation}
S_\infty=\begin{pmatrix}
1&-\sqrt2/3&\sqrt{10}/36\\
0&1&-\sqrt5/6\\0&0&1
\end{pmatrix},\qquad
D_\infty=S_\infty^{-1}B_\infty S_\infty
=\diag(-1/2,0,1/2)
\end{equation}
lautet das System in $t=1/x$:
\begin{equation}
\frac{D_\infty}{t}+\frac{C_\infty}{1-t},
\qquad C_\infty=-S_\infty^{-1}A_0S_\infty,\qquad
(C_\infty)_{31}=0.
\end{equation}
Für eine lokale Lösung $H(t)t^{D_\infty}$, $H_0=I$, ist die Rekursion
\begin{equation}
(n-\operatorname{ad}D_\infty)H_n
=C_\infty\sum_{j=0}^{n-1}H_j.
\end{equation}
Nur $n=1$, Eintrag $(3,1)$, ist resonant. Genau dort verschwindet die rechte Seite; bei $n\ge2$ ist der Rekursionsoperator invertierbar. Es entsteht kein logarithmischer Jordananteil. Die drei lokalen Monodromien haben deshalb tatsächlich die Ordnungen $4,3,2$.

Ihr Bild ist ein Quotient der endlichen sphärischen Dreiecksgruppe $\Delta(4,3,2)\cong S_4$. Ein echter Quotient von $S_4$ besitzt kein Element der Ordnung vier. Das Bild ist daher treu; seine dreidimensionale Darstellung ist irreduzibel. Die Determinante eins wählt die Rotationsdarstellung, also die Standarddarstellung mal Vorzeichencharakter.

Die Markierung stimmt ebenfalls: Aus der Quadratäquivarianz folgt bei $\Psi(0)=I$ die Identität $\Psi(iz)=U\Psi(z)U^{-1}$. Die positive $x=0$-Schleife wirkt deshalb mit $U^{-1}$. Die exakten Zielmatrizen erfüllen $(MU^{-1})^2=I$. Das markierte Paar ist somit analytisch bis auf gemeinsamen Basiswechsel bestimmt. Weil $U$ drei verschiedene Eigenwerte hat, ist ein noch verbleibender Basiswechsel diagonal; insbesondere folgen
\begin{equation}
\boxed{\diag M=(0,-i/2,+i/2),\qquad |M_{22}|=\tfrac12}
\end{equation}
ohne ODE-Numerik. Die genaue gedruckte Offdiagonal-Vorzeichenfassung legt eine Basiszeichenkonvention fest. Der Betrag $1/2$ ist ein Monodromiedatum und nicht der kontinuierliche Transportpol $\delta_{\rm ph}$.

\textbf{Auf der ursprünglichen $z$-Ebene erzeugen die vier Punkturmonodromien $A_4$ mit zwölf Elementen.} Erst die zusätzliche Quadratrotation erweitert das Paket zu $S_4$ mit 24 Elementen. Diese Unterscheidung ist für die Deckfläche wichtig. Die Reduktion, der Resonanzkoeffizient und die Gruppenrelation wurden unabhängig exakt gegengeprüft; der Beweis ist nicht als neuer Lean-Beweis ausgegeben.

Für die exakten Originalmatrizen aus v117 gibt es den unitären Basiswechsel
\begin{equation}
F=\begin{pmatrix}
1&0&0\\
0&-(1+i)/2&(-1+i)/2\\
0&(1-i)/2&(1+i)/2
\end{pmatrix},\qquad
F^\dagger MF=P_c=\begin{pmatrix}0&0&1\\1&0&0\\0&1&0\end{pmatrix},
\qquad
F^\dagger UF=R_c=\begin{pmatrix}1&0&0\\0&0&1\\0&-1&0\end{pmatrix}.
\end{equation}
Die vier lokalen Monodromien $U^kMU^{-k}$ sind damit signierte Permutationsmatrizen. Ihre Permutationsanteile sind abwechselnd $P_c,P_c^{-1},P_c,P_c^{-1}$. Genau diese Verzweigung an $z=1,i,-1,-i$ realisiert die dreiblättrige Kurve
\begin{equation}\label{eq:flavorg2}
w^3=\frac{z^2-1}{z^2+1},
\qquad
\eta=z(1-w^3)
\quad\Longrightarrow\quad
\boxed{\eta^2=1-w^6.}
\end{equation}
Die Riemann--Hurwitz-Rechnung ist $2g-2=3(-2)+4(3-1)=2$, also $g=2$. Die sechs hyperelliptischen Verzweigungspunkte sind genau $\mu_6$. Die verbleibenden Vorzeichen der Monodromiematrizen sind Daten eines Linienlokalsystems auf dieser Kurve. Die Kurve allein enthält sie noch nicht.

Die ursprüngliche Quadratrotation hebt ausdrücklich zu
\begin{equation}
r:(w,\eta)\longmapsto
\left(\frac1w,-\frac{i\eta}{w^3}\right),\qquad
r^2(w,\eta)=(w,-\eta),\quad r^4=I.
\end{equation}
Auch der Sechserzyklus ist auf derselben Kurve vorhanden. Für $\omega=e^{2\pi i/3}$ setze
\begin{equation}
c(w,\eta)=(\omega w,\eta),\qquad
s(w,\eta)=(-w,\eta).
\end{equation}
Die sechs Punkte lassen sich als $p_a=\omega^a$, $q_a=-\omega^a$ ordnen. In diesem Punktlabelraum wirken
\begin{equation}
c=\diag(P_c,P_c),\qquad
s=\begin{pmatrix}0&I\\I&0\end{pmatrix},\qquad T=sc.
\end{equation}
Der explizite unitäre Intertwiner
\begin{equation}
J_{\rm pts}=\frac1{\sqrt2}
\begin{pmatrix}F^\dagger&F^\dagger\\F^\dagger&-F^\dagger\end{pmatrix}
\end{equation}
erfüllt
\begin{equation}\label{eq:pointhexagon}
\boxed{J_{\rm pts}^\dagger TJ_{\rm pts}=M\oplus(-M),\qquad
J_{\rm pts}^\dagger sJ_{\rm pts}=\diag(I,-I).}
\end{equation}
Damit ist der endliche Sechseroperator ein tatsächlicher markierter Punkttransport. Er wurde nicht aus der Zahl sechs erraten.

\begin{center}
\begin{tikzpicture}[node distance=7mm]
\node[box,text width=12.8cm] (q) {\textbf{Originales Quadrat}\quad $1,i,-1,-i$\\drei Familien; signierte lokale Monodromie};
\node[box,below=of q,text width=12.8cm] (g) {\textbf{Dieselbe Deckkurve}\quad $w^3=(z^2-1)/(z^2+1)$\\zwei Henkel; Vorzeichen-Linienlokalsystem mitführen};
\node[box,below=of g,text width=12.8cm] (h) {\textbf{Zweite Ansicht}\quad $\eta^2=1-w^6$\\sechs Verzweigungspunkte; $M\oplus(-M)$ als Punkttransport};
\node[box,below=of h,text width=12.8cm] (f) {\textbf{Noch zu identifizierende Feldwirkung}\\derselbe Transport auf geladenen Defekten, ihrem Zustand und den vollständigen Antworten};
\draw[arr] (q)--(g);\draw[arr] (g)--(h);\draw[dasharr] (h)--(f);
\end{tikzpicture}
\end{center}

\picturetext{Vier Marken, drei Blätter und sechs ausgezeichnete Punkte können Ansichten derselben geometrischen Konstruktion sein. Die mitgeführten Vorzeichen sind dabei wie zusätzliche Anweisungen auf den Blättern: Das Drehen der Zeichnung genügt noch nicht, um die Feldoperation richtig zu drehen.}

Insbesondere ist $s$ nicht die hyperelliptische Involution. Letztere fixiert alle sechs Verzweigungspunkte; $s$ vertauscht die beiden Dreierbahnen und induziert $z\mapsto1/z$. Der hier bewiesene Punktintertwiner darf deshalb weder als Feldintertwiner noch als Ableitung des ursprünglichen Nahtprojektors ausgegeben werden.

\subsection{Der Sechsertransport entsteht aus zwei ursprünglichen Halbwegen}\label{sec:pathclock}
\textbf{Ein globaler Inversionsautomorphismus des einzelnen Flachbündels ist hierfür nicht nötig.}
Wähle einen regulären Ausgangspunkt $x\in(0,1)$ und setze $x'=1/x$.
Der Kreis mit Mittelpunkt $a=(x+x')/2$ und Radius $r=(x'-x)/2$ erfüllt
$a^2-r^2=1$ und ist unter $\iota(z)=1/z$ invariant. Sein unterer Halbkreis $p$ führt von $x$ nach $x'$. Der Bildweg $\iota p$ führt auf dem oberen Halbkreis zurück. Zusammen umlaufen sie $+1$ genau einmal positiv. Der ganze Kreis liegt im rechten Halbplan und umschließt keine weitere Fuchs-Punktur.

Seien $P=\operatorname{PT}(p):E_x\to E_{x'}$ und
$Q=\operatorname{PT}(\iota p):E_{x'}\to E_x$ die tatsächlichen Paralleltransporte der ursprünglichen Verbindung. Dann gilt
\begin{equation}
QP=M_x,\qquad PQ=PM_xP^{-1}.
\end{equation}
Die endliche Monodromie besitzt eine positive invariante hermitesche Form, etwa durch Mittelung einer positiven Form über ihre endliche Gruppe. Paralleltransport setzt sie zur flachen Fasermetrik fort. Wählt man den Rahmen bei $x'$ durch Transport des orthonormalen Rahmens bei $x$, so erhält man
\begin{equation}\label{eq:pathT}
P=I,\quad Q=M,\qquad
T=\begin{pmatrix}0&M\\I&0\end{pmatrix},\qquad
T^2=\diag(M,M),\qquad T^6=I.
\end{equation}
Die rohen ODE-Matrizen in einem beliebigen holomorphen Rahmen sind deshalb nicht automatisch bezüglich der gewöhnlichen euklidischen Matrixmetrik unitär. Der Satz verwendet die flache hermitesche Metrik.

Mit dem expliziten unitären Wörterbuch
\begin{equation}
J=\frac1{\sqrt2}\begin{pmatrix}I&I\\M&-M\end{pmatrix}
\end{equation}
folgt exakt
\begin{equation}\label{eq:pathclockdictionary}
\boxed{J^\dagger T^{-1}J=M\oplus(-M)=U_6,\qquad
J^\dagger T^3J=\diag(I,-I).}
\end{equation}
Die ursprüngliche Clock entspricht bei dieser positiven Umlaufkonvention dem inversen Halbschritt. Diese Orientierung wird ausdrücklich mitgeführt. $T^3$ ist eine komplexlineare Involution mit zwei komplex dreidimensionalen Eigenräumen; sie ist kein antiunitärer Reflexionsswap.

Für den bereits gegebenen Flavoroperator folgt damit
\begin{equation}
D^{\rm path}_{y,\delta}=yI-\delta T^{-1},\qquad
\boxed{J^\dagger D^{\rm path}_{y,\delta}J=yI-\delta U_6.}
\end{equation}
Dasselbe $J$ transportiert die positive Gramform, ihre regulierte Resolvente, den inversen Quadratwurzelkernel und die Determinante
$\det D=y^6-\delta^6$. Es werden weder neue Flavorparameter eingesetzt noch bestehende Massen- und Yukawa-Herleitungen ersetzt. Veränderungen von $x$ innerhalb $(0,1)$ ändern nur den parallel transportierten Rahmen und damit den Operator durch unitäre Konjugation.

\begin{center}
\begin{tikzpicture}
\node[box,text width=4.1cm] (left) at (0,0) {Zugang $x$\\drei Familienkomponenten};
\node[box,text width=4.1cm] (right) at (8,0) {Zugang $1/x$\\dieselben Feldtypen};
\draw[arr,bend right=28] (left.south) to node[below,font=\small] {unterer Halbweg: $P$} (right.south);
\draw[arr,bend right=28] (right.north) to node[above,font=\small] {oberer Halbweg: $Q$} (left.north);
\node[align=center,font=\small,text=navy] at (4,-2) {zwei Halbschritte: $M$\qquad sechs Halbschritte: $I$\\Original-Flavoroperator: $yI-\delta T^{-1}$};
\end{tikzpicture}
\end{center}

\picturetext{Ein Paket mit drei Familienkomponenten wird über zwei Zugänge weitergereicht. Nach der Hin- und Rückrunde trägt es die ursprüngliche Familienänderung $M$. Drei solcher Runden schließen sich. Die sechs Schritte sind damit eine Operation am vorhandenen Paket, keine neue Zahlenspielerei.}

\textbf{Der genaue Feldumfang bleibt erhalten.}
$E_x\oplus E_{x'}$ ist ein Raum von Randdaten an zwei Zugängen. Er ist nicht der sechsdimensionale Raum horizontaler Lösungen: Auf einer eingeschnittenen verbundenen Domäne erfüllen deren Werte bereits $s(x')=Ps(x)$ und bilden nur einen dreidimensionalen Graphen. Die Quantisierung der gesamten Feldquelle muss daher den ursprünglichen Testfunktionenraum und seinen Riemann--Hilbert-Kernel mitführen. Die endliche Gleichung~\eqref{eq:pathclockdictionary} allein behauptet weder eine globale Zustandsinvarianz noch die Gleichheit $T^3=\mathsf iC$ mit dem ursprünglichen Calder\'onoperator.

\subsection{Die ursprüngliche Determinantenlinie trägt bereits primitive Windung}\label{sec:fuchswinding}
Hier liefert das Fuchssystem selbst eine weitere Herkunftsverbindung. Aus den ursprünglichen Residuen folgt
\begin{equation}
\tr A(z)=\frac{4z^3}{z^4-1},\qquad
\frac{d}{dz}\log\det\Psi(z)=\tr A(z),\qquad \Psi(0)=I.
\end{equation}
Deshalb gilt exakt
\begin{equation}
\boxed{\det\Psi(z)=1-z^4.}
\end{equation}
Die Phase $u_F=(1-z^4)/|1-z^4|$ besitzt auf dem oben konstruierten positiven Kreis um $+1$ Windung $+1$. Die tatsächliche Fuchs-Determinantenlinie stellt somit dieselbe primitive $K^1(S^1)$-Klasse dar, die P1 verlangt. Der ganzzahlige Anteil stammt aus der logarithmischen Erweiterung; aus $\det M=1$ allein könnte man ihn nicht rekonstruieren.

Die Kontur gehört zum Satz: Ein gemeinsamer Umlauf um alle vier Fuchs-Punkte hat Windung vier. Dass der ursprüngliche normale P1-Seamkreis gerade auf den einzelnen markierten Punkturkreis abgebildet wird, ist durch die bloße Klassengleichheit noch nicht bewiesen. Der nächste geometrische Herkunftsschritt ist somit eine konkrete Kontur- und Randoperatorabbildung, keine neue freie Wahl der Indexzahl.

\subsection{Die Familienoperationen wirken in der bereits vorhandenen Quelle}
Die rekonstruierte Fermionenquelle enthält $(10,1)\oplus(1,6)$ unter $\mathrm{Spin}(10)\times SU(4)$. Für $g\in SU(3)$ gilt ausdrücklich
\begin{equation}
\widehat g=\diag(1,g)\in SU(4),\qquad
\Lambda^2\widehat g\simeq g\oplus\bar g.
\end{equation}
Somit sind $M$ und $U$ wirkliche innere Symmetrien ihres Familienzweigs. Die vorhandenen sechs reellen Familienfermionen entsprechen drei komplexen Fermionen. Ihre Wirkung $M\oplus\bar M$ ist jedoch nicht $M\oplus(-M)$. Das Original-Sechserregister umfasst eine zusätzliche Sheet-Markierung. Die beiden Spin-Lifts $\widehat M$ und $-\widehat M$ haben dieselbe Vektorwirkung, aber entgegengesetzte Wirkung auf dem Familien-Spinor. Eine zusätzliche Identifikation mit dieser internen Spin-Liftwahl wäre gesondert zu beweisen. Im oben konstruierten Wegoperator entsteht das Vorzeichen bereits aus den beiden Zugängen; eine Gleichsetzung mit diesem zentralen Spin-Liftzeichen wird dafür nicht vorausgesetzt.

Der passende Anschluss an die Fuchsverbindung besteht aus Monodromiedefekten derselben Fermionenquelle. Die freie Fermionenkonstruktion bestimmt deren Antworten durch den Riemann--Hilbert-Kernel und Wick-Determinanten~\cite{gavrylenkomarshakov}. Bei der Originalkonvention $\Psi'=A\Psi$ lautet der entsprechend transponierte Kernel
\begin{equation}
\mathcal K(z,w)=\frac{[\Psi(w)^{-1}\Psi(z)]^{\mathsf T}}{z-w}.
\end{equation}
Die Defekte, ihre Sektoren und ihre Normierung sind Teile dieser Einsetzung. Eine gewöhnliche unverdrehte E8-Vakuumkorrelation wird damit nicht stillschweigend zur ursprünglichen Fuchsantwort.

Die Originalresiduen haben die Exponenten $0,1/3,2/3$. In der freien Fermionennormierung entsprechen ihnen $Q=1$ und $h=5/18$. Die gleichphasigen Vertreter $0,1/3,-1/3$ hätten dagegen $Q=0$, $h=1/9$. Gleiche Umlaufphasen erlauben keinen unbemerkten Austausch der geladenen Defektzustände. Auch der Infinity-Sektor mit Residuum $-\diag(2,1,1)$ muss erhalten bleiben. Diese genaue Markierung ist wichtiger als eine bloße Übereinstimmung der Gruppenordnungen.

Für die vollständige Randfeldwirkung ist zudem der Spinanteil getrennt zu führen. Auf dem rein internen getwisteten Schleifenraum $f(\theta+2\pi)=Mf(\theta)$ erfüllt die Halbtranslation $f(\theta)\mapsto f(\theta+\pi)$ dieselbe Sechserrelation und kommutiert mit $-i\partial_\theta$ und dessen Spektralprojektoren. Enthält die Gesamtsektion dagegen zusätzlich das NS-Spinvorzeichen, lautet ihre Randbedingung $f(\theta+2\pi)=-Mf(\theta)$ und die sechste Potenz ist $-I$. Dann wäre der rohe Halbwegdeterminant $y^6+\delta^6$. Der Original-Flavordeterminant benötigt daher die erklärte interne Transportkonvention mit getrenntem Spin-/Halbdichtefaktor oder die entsprechende andere Spinstruktur. Ein künstlich gewähltes stationäres Schleifenvakuum ersetzt nicht die tatsächlichen getwisteten Quellenmatrixelemente.

\subsection{Eine explizite gemeinsame Klasse von Randoperatoren}
Die endliche Flavorgram-Matrix besitzt außerdem eine genaue Darstellung als Dirichlet-zu-Neumann-Antwort. Für $y,\delta>0$, $y\ne\delta$, einen zyklischen Shift $S_N$ und
\begin{equation}
t=\frac{\min(y,\delta)}{\max(y,\delta)}\in(0,1),\qquad
ma=-\log t,\qquad a=\ell/N
\end{equation}
ist der Sprungoperator von $-\partial_x^2+m^2$ auf den $N$ Kreisintervallen
\begin{equation}
R_N(m)=\frac{m}{\sinh(ma)}
\bigl[2\cosh(ma)I-S_N-S_N^\dagger\bigr]
=\boxed{\frac{2m}{|y^2-\delta^2|}D_{y,\delta}^\dagger D_{y,\delta}.}
\end{equation}
Die Gleichheit folgt durch Einsetzen von $e^{-ma}=t$. Sie identifiziert eine Operatorfamilie; die Interpretation des internen $C_6$-Registers als sechs physische Kreisorte folgt daraus nicht.

Die Intervall-Determinanten und die BFK-Verklebung ergeben in dieser eindimensionalen Realisierung
\begin{equation}
2^{-N}\left(\frac{2\sinh(ma)}m\right)^N\det R_N(m)
=e^{m\ell}(1-e^{-m\ell})^2.
\end{equation}
Die lokalen Faktoren dürfen bei Variationen nicht fortgelassen werden. Die allgemeine Verklebungsformel enthält entsprechende lokale Normierungen~\cite{bfkoriginal,leebfk}. Durch Eliminieren zusätzlicher Schnittstellen reduziert dieselbe Zwölfer-Sprungmatrix exakt zur Vierer- beziehungsweise Sechsermatrix. Das ist eine gemeinsame Rechenverfeinerung desselben Kreises, keine neue Physik.

Auf dem algebraischen Originalzweig
\begin{equation}
\delta_\star=\left(\frac{794-7\sqrt{9961}}{2187}\right)^{1/6}
=0.5932772801\ldots
\end{equation}
lauten die positiven Rundläufe für $y=1,1/3,2/3$:
\begin{equation}
q_+=0.0436060052\ldots,\quad0.0314576422\ldots,\quad0.4966996535\ldots .
\end{equation}
Der elektromagnetische Rundlauf ist dagegen $q_{\rm EM}=48c_3^4e^{-2\alpha}\simeq0.0001185614$. Damit ist eine unmittelbare Identifikation aller Antworten mit derselben skalaren Massenspezialisierung ausgeschlossen. Verschiedene markierte Kanäle einer gemeinsamen Familie bleiben möglich.

Für $y<\delta$ bleibt der holomorphe Rundlauf $(\delta/y)^6>1$; die positive Darstellung benutzt seinen Kehrwert. Insbesondere gilt
\begin{equation}
\det(D^\dagger D)=\delta^{12}\bigl[1-(y/\delta)^6\bigr]^2.
\end{equation}
Der Faktor $\delta^{12}$ trägt zur Flavorvariation bei. Auch die Phase von $\det D=y^6-\delta^6$ wird durch den positiven Gramoperator allein nicht dargestellt.

\subsection{Der gemeinsame neutrale Port ist jetzt genauer bestimmt}\label{sec:neutralport}
Im Original-Mark-Greenoperator mit renormierter Diagonale null steht die dimensionslose Matrix
\begin{equation}
C_{\rm marks}=\operatorname{circ}(0,1,2,1),\qquad
G_{\rm marks}=-4c_3\log2\,C_{\rm marks}.
\end{equation}
Für den geometrischen Viererschritt $T_4$ liefert der Halbturn-gerade Projektor $E_+=(I+T_4^2)/2$ exakt
\begin{equation}
\frac14E_+C_{\rm marks}E_+=P_0,\qquad
P_0=\frac14\begin{pmatrix}1\\1\\1\\1\end{pmatrix}
\begin{pmatrix}1&1&1&1\end{pmatrix},
\qquad \det(I-qP_0)=1-q.
\end{equation}
Das benennt einen konkreten geometrischen Port. $C_{\rm marks}$ ist nicht die ursprüngliche Nahtinvolution; $E_+=P_{\Sigma,+}$ bleibt eine gesonderte Feldidentifikation.

Die Materieseite besitzt einen ladungstreuen Anschluss: Für $\mathcal R=S^+\otimes H_F$ ist
\begin{equation}
|\mathbb I_{\mathcal R}\rangle\!\rangle
=\sum_{a=1}^{48}|a\rangle\otimes|\bar a\rangle,\qquad
\|\,|\mathbb I_{\mathcal R}\rangle\!\rangle\,\|^2
=\tr_{\mathcal R}I=48
\end{equation}
der kanonische neutrale Paarungstensor. Er ist unter $V\otimes\bar V$ invariant. Dadurch können geschlossene Materieantworten in einen neutralen Port eingehen, ohne verschieden geladene Einzelteilchen zu identifizieren.

\textbf{Die fehlende gemeinsame Kopplungsidentität ist damit ausdrücklich}
\begin{equation}\label{eq:neutralselfenergy}
\boxed{\Sigma_{\rm matter}(g)
=\mathcal V^\dagger G_{\rm matter}(g)\mathcal V
=48g\,P_0,\qquad g=c_3^4e^{-2\alpha}.}
\end{equation}
Hier müssen $\mathcal V$ und $G_{\rm matter}$ aus der ursprünglichen Quelle stammen. Die neutrale Norm $48$ allein berechnet diese Selbstenergie nicht. Die Originalfassung bezeichnet $\delta_{\rm top}=48c_3^4$ selbst als Besetzungsinterpretation und definiert anschließend den Rang-eins-Transfer mit diesem Wert: \file{02_carrier_source.tex}, Zeilen 1729--1746, und \file{03_em_flavor_source.tex}, Zeilen 94--120.

Ein präziser höherer Antworttest trennt die Möglichkeiten:
\begin{align}
-\log(1-48g)&=48g+1152g^2+O(g^3),\\
-48\log(1-g)&=48g+24g^2+O(g^3).
\end{align}
Ein gemeinsamer Port und 48 unabhängige Determinanten haben dieselbe führende Spur, aber verschiedene Rückkopplungen. Die ursprüngliche zweite Defektkorrektur allein unterscheidet sie noch nicht: Mit $\beta=5/4$ beginnen beide Ausdrücke
\begin{align}
48g(1-48g)^{-\beta}&=48g+2880g^2+155520g^3+O(g^4),\\
48g(1-g)^{-48\beta}&=48g+2880g^2+87840g^3+O(g^4).
\end{align}
Der zweite Logdet-Jet beziehungsweise dritte Defekt-Jet ist deshalb der erste hier trennende gemeinsame Quellenkoeffizient. Diese Feststellung erhält die bereits berechnete Alpha-Gleichung; sie präzisiert die noch benötigte Herkunft ihrer Resummierung.

\subsection{Die tatsächliche Quelle liefert einen bestimmten neutralen Kandidaten}\label{sec:neutralfield}
Die ursprüngliche Level-eins-E8-Quelle enthält die Spinorströme $(16,4)$ mit 64 Komponenten und ihre Gegenfelder. Unter der erklärten Familienauswahl $SU(4)\to SU(3)\times U(1)$ zerfällt dieser Block in $(16,3)\oplus(16,1)$. Die erste Komponente ist der hier untersuchte 48er-Materiereadout. Seine Identifikation mit dem ursprünglichen topologischen Familienraum bleibt Teil des markierten Wörterbuchs; die Rechnung benutzt keine vierte physische Generation.

Für jeden normierten Wurzelstrom sei der symmetrische neutrale Paarvertex
\begin{equation}
B_\beta=\tfrac12\bigl(:E_\beta E_\beta^\dagger:
 +:E_\beta^\dagger E_\beta:\bigr)
=\tfrac12:(\beta\cdot h)^2:.
\end{equation}
Die zweite Gleichung folgt aus der ursprünglichen Level-eins-Gitter-OPE; $h$ sind die acht normierten Heisenbergströme. Für eine Menge $\mathcal R$ dieser Wurzeln setze
\begin{equation}
A_{\mathcal R}=\sum_{\beta\in\mathcal R}\beta\beta^{\mathsf T},\qquad
B_{\mathcal R}=\tfrac12:h^{\mathsf T}A_{\mathcal R}h:.
\end{equation}
Die exakte Rechnung an den ursprünglichen Wurzeln ergibt
\begin{align}
A_{64}&=16I_8,& B_{64}&=16T_{E_8},\\
\operatorname{spec}A_{48}&=\{12^{[5]},16^{[2]},4^{[1]}\},&
B_{48}&=12T_{D_5}+16T_{A_2}+4T_{U(1)},\\
B_{16}&=4T_{D_5}+12T_{U(1)},&
B_{48}+B_{16}&=16T_{E_8}.
\end{align}
Die Zentralzahlen der drei kommutierenden Stressanteile sind $(5,2,1)$. Das hier verwendete $U(1)$ gehört zur Familienzerlegung; es ist nicht ohne weitere Ladungsabbildung die Standardmodell-Hyperladung.

Diese Gleichungen erklären positiv, wie vollständige Quelle und 48er-Auswahl zusammengehören. Bereits die 48 geladenen Wurzelströme und ihre Gegenfelder erzeugen unter den ursprünglichen Chevalley-Klammern wieder alle E8-Wurzeln: $96\to202\to240$. Dabei kehren auch die 16 zunächst ausgelassenen Spinorwurzeln zurück. Eine Materieauslese darf diese 16 unangezeigt lassen; eine unter den ursprünglichen Feldoperationen abgeschlossene Quelle darf sie nicht löschen.

Setze $T=T_{E_8}$. Die Projektion von $B_{48}$ auf den gemeinsamen Stresszustand liefert
\begin{equation}\label{eq:nativeresidual}
\boxed{B_{48}=12T+R,\qquad R=4T_{A_2}-8T_{U(1)}.}
\end{equation}
Die Vakuumpaarungen sind
\begin{equation}
\langle T,T\rangle=4,\quad
\langle B_{48},T\rangle=48,\quad
\langle R,T\rangle=0,\quad
\boxed{\langle R,R\rangle=48.}
\end{equation}
Daher ist $\|B_{48}|0\rangle\|^2=624=576+48$.
$R$ ist ein tatsächliches neutrales Primärfeld vom konformen Gewicht zwei: Der vierte Pol seiner OPE mit $T$ verschwindet wegen $4\cdot2-8\cdot1=0$. Es ist kein Stresstensor einer neuen Theorie mit Zentralzahl null. Neutralität und Gewicht zwei allein wählen es nicht eindeutig; erst die erklärte 48er-Auswahl und die orthogonale Subtraktion des gemeinsamen Stressanteils bestimmen diese Richtung.

\picturetext{Die Antwort der 48 Materiefelder enthält einen Anteil, der das ganze Instrument gemeinsam bewegt, und eine zusätzliche Bewegung im Familienbereich. Nach Abziehen der gemeinsamen Bewegung bleibt ein konkret bestimmter Zustand. Seine Stärke ist genau 48 und seine konforme Energie genau zwei. Die beiden Zahlen stammen jetzt aus demselben Quellenoperator.}

\subsection{Ein exakter Quellenport realisiert den Alpha-Rundlauf bedingt}\label{sec:heatport}
Dieser Zustand liefert eine konkrete Realisierung der zuvor nur als Ziel formulierten Gleichung~\eqref{eq:neutralselfenergy}. Sei
\begin{equation}
|r\rangle=R_{-2}|0\rangle,\qquad
L_0|r\rangle=2|r\rangle,\qquad
\langle r|r\rangle=48,\qquad
|p_0\rangle=\tfrac12(1,1,1,1)^{\mathsf T}.
\end{equation}
Definiert man den Quellenzugang und die euklidische Quellenfortsetzung durch
\begin{equation}\label{eq:conditionalvertex}
\mathcal V=c_3^2|r\rangle\langle p_0|,\qquad
G(\alpha)=e^{-\alpha L_0},
\end{equation}
so folgt auf der vollständigen Quelle exakt
\begin{equation}\label{eq:actualheatport}
\boxed{\mathcal V^\dagger G(\alpha)\mathcal V
=48c_3^4e^{-2\alpha}P_0.}
\end{equation}
Damit liefert derselbe erklärte Port auch
$\det(I-\mathcal V^\dagger G\mathcal V)=1-48c_3^4e^{-2\alpha}$, einschließlich seiner höheren Logdet-Koeffizienten. Der Exponent zwei ist nun das Gewicht eines vorhandenen Feldzustands; die 48 ist dessen berechnete Norm. Es wird kein neuer kontinuierlicher Parameter angepasst.
Der bereits im Original verwendete Suszeptibilitätsfaktor $5/4$ wird durch diese Portidentität nicht zusätzlich hergeleitet; seine bestehende Herleitung und seine gemeinsame Normierung sind weiterhin mitzuführen.

\textbf{Status: exakte Konstruktion unter drei benannten Identifikationen, noch kein Ursprungssatz.}
Zu beweisen ist, dass die ursprüngliche Naht tatsächlich (i) den stressorthogonalen Zustand $|r\rangle$ koppelt, (ii) die Amplitude $c_3^2$ dieser Kopplung festlegt und (iii) die elektromagnetische Koordinate $\alpha$ als genau diesen konformen Fortsetzungsparameter verwendet. Die bisher geprüften Originaldefinitionen von $P_{\rm prim}$, $P_{\rm sing}$ und $P_\Theta$ geben diese Abbildung nicht an. Insbesondere bedeutet die ursprüngliche primitive Windungsbedingung nicht automatisch die Projektion auf ein Virasoro-Primärfeld.

Der Port ist eine Abbildung in einen bestimmten Quellenzustand. Er ist nicht automatisch dasselbe wie eine beliebige lokale Kopplung an das Feld $R(z)$. Dessen Dreipunktantwort lautet
\begin{equation}
\langle R(z_1)R(z_2)R(z_3)\rangle
=\frac{-384}{z_{12}^2z_{23}^2z_{31}^2}.
\end{equation}
Diese nichtverschwindende Antwort muss bei lokalen Mehrfacheinsetzungen mitgeführt werden. Sie schließt außerdem eine vakuumerhaltende Quellenautomorphie aus, die $R$ schlicht auf $-R$ abbildet. Eine solche Nahtparität kann diesen Anteil nicht ohne weiteres löschen. Der vollständige Quellprozess bleibt reicher als seine eine geometrische Portanzeige.

\subsection{Flavor und neutrale Quelle besitzen bereits eine gemeinsame berechnete Antwort}\label{sec:jointneutralward}
Der neue Kandidat lässt sich unmittelbar am ursprünglichen Flavor-Hintergrund prüfen. Dafür wird die normierte Riemann--Hilbert-Defektantwort derselben drei komplexen Familienfermionen verwendet, mit den ursprünglichen Residuen und dem vorgeschriebenen ausgehenden Sektor. In dieser Darstellung liefern die Fermionen-Wardidentitäten
\begin{equation}
\langle T_{\rm Familie}(z)\rangle_F=\tfrac12\tr A(z)^2,
\qquad
\langle T_{U(1)}(z)\rangle_F=\tfrac16\bigl(\tr A(z)\bigr)^2.
\end{equation}
Die symmetrische Stressdefinition entfernt den Ableitungsterm; der verbundene $U(1)$-Stromzweipunktteil hat weiterhin den Koeffizienten drei. Da $T_{\rm Familie}=T_{A_2}+T_{U(1)}$, erhält man durch Einsetzen der ursprünglichen Verbindung
\begin{equation}\label{eq:jointRresponse}
\boxed{\langle R(z)\rangle_F
=2\bigl[\tr A(z)^2-(\tr A(z))^2\bigr]
=\frac{-20z^6+(52/9)z^2}{(z^4-1)^2}.}
\end{equation}
Das ist eine konkrete Antwort des neutralen Kandidaten auf die bereits vorhandenen Flavor-Defekte. Sie wurde unabhängig symbolisch gegengeprüft. Da $R$ ein Primärfeld vom Gewicht zwei ist, ist diese Antwort als quadratisches Differential zu transformieren; eine zusätzliche Schwarzian-Korrektur entsteht für $R$ nicht.

Die lokalen Doppelpolkoeffizienten an jeder der vier Punkturen betragen $-8/9$. Am unendlichen Ende gilt $z^2\langle R(z)\rangle_F\to-20$. Die negativen Werte sind kein Positivitätsverstoß: $R$ ist eine Differenz von Stressantworten und kein positiver Energieoperator. Die Originalresiduen $(0,1/3,2/3)$ liefern lokal
\begin{equation}
h_{A_2}=1/9,\qquad h_{U(1)}=1/6,\qquad
4h_{A_2}-8h_{U(1)}=-8/9.
\end{equation}
Der lokale Rest bei Unendlich besitzt die Exponenten $(-2,-1,-1)$; im ausgehenden Bra-Sektor entsprechen sie der Ladung $+4$. Er kompensiert die vier Defekte mit Ladung $+1$. Triviale Monodromie am unendlichen Ende bedeutet somit nicht Vakuumladung. Genau diese ganzzahlige Quelleninformation ginge bei einer Beschreibung ausschließlich durch die Monodromiematrizen verloren.

Auch der erste lokale Normtest ist jetzt berechenbar. Für einen normierten Highest-Weight-Defektzustand $|\sigma\rangle$ mit den obigen Gewichten folgt aus der Virasoro-Relation
\begin{equation}\label{eq:defectRnorm}
\|R_{-2}|\sigma\rangle\|^2
=16(4h_{A_2}+1)+64(4h_{U(1)}+1/2)
=\boxed{880/9},
\end{equation}
während die Vakuumnorm 48 beträgt. Eine bloße Mittelwertsubtraktion ändert den Zustand $R_{-2}|\sigma\rangle$ nicht. Das ist ein lokaler Sektortest, keine Behauptung, dass die gesamte Vierdefektgeometrie aus einem einzigen Highest-Weight-Zustand besteht.

\key{\textbf{Was der gemeinsame Befund entscheidet:}
Der Flavortransport und ein exakter neutraler Alpha-Port sind nun innerhalb der vorhandenen Quellenstruktur konkret darstellbar. Ihre gemeinsame Antwort ist jedoch zustands- und markierungsabhängig und bereits ausdrücklich berechnet. Die gesuchte Nahtabbildung muss deshalb erklären, welcher Quellenkanal ausgewertet wird und wie seine Normierung mit den Flavor-Einsetzungen zusammenhängt. Verschiedene Einsetzungen desselben Quellenfunktionals sind zulässig; ihre gemeinsamen Antworten dürfen nicht durch getrennte Vakuumannahmen ersetzt werden.}

\subsection{Die vollständige neutrale Antwort im ursprünglichen Flavorprozess}\label{sec:fullflavorRR}
Die gemeinsame Einpunktantwort bestimmt bereits einen Teil der Verträglichkeit.
Der vollständige verbundene Zweipunktkern liefert nun die tatsächliche Antwort
zwischen zwei verschiedenen Einsetzungsorten. Entscheidend ist die unveränderte
Feldkonvention des ursprünglichen Fuchs- und Riemann--Hilbert-Systems:
\[
 \Psi'(z)=A(z)\Psi(z),\qquad j(z)=\tr A(z),\qquad
 \mathcal K(z,w)=[\Psi(w)^{-1}\Psi(z)]^{\mathsf T}/(z-w).
\]
Im normierten quasifreien Riemann--Hilbert-Defektmatrixelement folgt aus
symmetrischer Fermion-Punktaufspaltung~\cite[Abschnitte 3.3 und 4.1]{gavrylenkomarshakov}
\begin{align}
 \langle T_{\rm Familie}(z)T_{\rm Familie}(w)\rangle_{F,c}
 &=\frac{3}{2(z-w)^4}
   +\frac{\tr\!\left(A(z)A(w)\right)}{(z-w)^2},\\
 \langle T_{U(1)}(z)T_{U(1)}(w)\rangle_{F,c}
 &=\frac{1}{2(z-w)^4}+\frac{j(z)j(w)}{3(z-w)^2}.
\end{align}
Die $A_2$- und $U(1)$-Faktorisierung ergibt für $R=4T_{\rm Familie}-12T_{U(1)}$
\begin{equation}\label{eq:fullflavorRR}
\boxed{\begin{aligned}
 \langle R(z)R(w)\rangle_{F,c}
 &=\frac{48}{(z-w)^4}\\[-2pt]
 &\quad+\frac{16}{(z-w)^2}
 \left[\tr\!\left(A(z)A(w)\right)+j(z)j(w)\right].
\end{aligned}}
\end{equation}
Der zusätzliche Term ist für die ursprüngliche Vierdefektverbindung vollständig rational:
\begin{equation}
 \langle R(z)R(w)\rangle_{F,c}-\frac{48}{(z-w)^4}
 =\frac{32(99z^3w^3+4z^2+5zw+4w^2)}
 {9(z-w)^2(z^4-1)(w^4-1)}.
\end{equation}
Das Indexzeichen $c$ bedeutet, dass das Produkt der Einpunktwerte bereits
abgezogen wurde. Der Zusatzterm ist also keine bloße Verschiebung des Mittelwerts.

Der lokale Doppelpolkoeffizient lautet exakt
\begin{equation}
 C_2(z)=16\left[\tr A(z)^2+j(z)^2\right]
 =\frac{32z^2(99z^4+13)}{9(z^4-1)^2}.
\end{equation}
Er stimmt mit der unabhängigen OPE-Identität
$C_2=64\langle T_{\rm Familie}\rangle_F-8\langle R\rangle_F$ überein.
Jede der vier Punkturen liefert $224/9$; dadurch wird auch die lokale Norm
$48+2(224/9)=880/9$ aus~\eqref{eq:defectRnorm} wiedergewonnen.
Bei $w=0{,}43$ und $z=0{,}17$ ist der verbundene Kern numerisch
$10572{,}47509$, gegenüber $10503{,}83390$ im Vakuum.

Die Reihenfolge im Kernel ist wesentlich. Die Primärquelle verwendet
$\varphi'=\varphi A_{\rm rechts}$. Für unsere linke Fuchs-Gleichung gilt
$\varphi=\Psi^{\mathsf T}$ und $A_{\rm rechts}=A^{\mathsf T}$.
In der neutralen Wick-Antwort kürzen sich die Fundamentallösungen deshalb
vollständig; es bleibt $\tr(A(z)A(w))$.
Die zunächst untersuchte vertauschte Matrixreihenfolge hätte stattdessen
$\tr(A(z)P(z,w)A(w)P(w,z))$ geliefert. Sie gehört zu einer anderen lokalen
Feldrahmung und wird hier ausdrücklich nicht als ursprüngliche RHP-Antwort
übernommen. Die korrigierte Antwort ist rational und einwertig, wie es die
Invarianz von $R$ unter den Familienmonodromien verlangt.

Das Programm \file{tfpt_explorer/source_flavor_correlator.py} berechnet
jetzt diese eindeutige Antwort aus der Originalverbindung. Die
Reihenfolgenprüfung gehört zur Regression; lokale Einpunkt- und
OPE-Ergebnisse bleiben unverändert. Die ausgewiesene Reichweite ist das
ursprüngliche Defektmatrixelement. Eine positive globale Dichtematrix oder
physische Photonpolarisation wird daraus nicht stillschweigend abgeleitet.

\key{\textbf{Die neue gemeinsame Bedingung:}
Eine Nahtabbildung im Flavor-Hintergrund muss den gesamten Kern
\eqref{eq:fullflavorRR} in der richtigen ursprünglichen Feldkonvention abbilden.
Das bloße Ersetzen durch die Vakuumzahl 48 ist unzureichend.
Ein vollständiger $E_+$-Port mit Referenzdeterminante kann trotzdem zulässig
sein: Die gerade lokale Korrektur allein entscheidet diese globale,
normalisierte Abbildung noch nicht.}

\picturetext{Der Familienweg bewegt sich innerhalb desselben Instruments.
Deshalb verändert er auch, wie dieses Instrument auf zwei neutrale Anfragen
gemeinsam reagiert. Diese Veränderung ist jetzt ausgerechnet. Sie gehört zur
gemeinsamen Mechanik, die eine vollständige Lösung erhalten muss.}

\paragraph{Der ursprüngliche Randport ist bereits berechenbar.}
Die Originalverifikationen v484 und v485 enthalten den Viermarken-Kern
und seinen logarithmischen Referenzabzug. Ihre Matrix lautet
\[
 C_{\rm mark}=\begin{pmatrix}
 0&1&2&1\\1&0&1&2\\2&1&0&1\\1&2&1&0
 \end{pmatrix},\qquad
 K_{\rm mark}=-c_3\log2\,C_{\rm mark}.
\]
Hier bezeichnet $K_{\rm mark}$ bereits das Orbitmittel des ursprünglichen
Mark-Greenkerns. Mit der zyklischen Markverschiebung $Q_4$,
$E_{\rm mark,+}=(I+Q_4^2)/2$ und $p_0=(1,1,1,1)^{\mathsf T}/2$ folgt exakt
\[
 E_{\rm mark,+}C_{\rm mark}E_{\rm mark,+}=4P_0,\quad
 P_0=p_0p_0^\dagger,\quad
 p_0^\dagger K_{\rm mark}p_0=-4c_3\log2.
\]
Der gerade Markraum hat Dimension zwei; die komprimierte Antwort hat Rang eins.
In der Originalkonvention $u=\varepsilon c_3\log2$ sind
\begin{equation}
 \det(I-uC_{\rm mark})=(1-4u)(1+2u)^2,\qquad
 \det(I-uE_{\rm mark,+}C_{\rm mark}E_{\rm mark,+})=1-4u.
\end{equation}
Die schon vorhandene Alpha-Vorgabe $q=48c_3^4e^{-2\alpha}$ würde dadurch
exakt getroffen, wenn der ursprüngliche Kontaktparameter die Identifikation
\begin{equation}\label{eq:originalportmatching}
 \boxed{\varepsilon(\alpha)=\frac{12c_3^3}{\log2}e^{-2\alpha}}
\end{equation}
besitzt. Dies ist eine aus beiden vorhandenen Antworten berechnete
\emph{Matchingbedingung}; sie wird hier nicht als unabhängige Herleitung der
Kontaktkopplung ausgegeben.

Die jetzt präzisierte Anschlussfrage lautet daher: Welche ursprüngliche
Variation $\partial_\alpha D_\xi$ beziehungsweise Rand-zu-$R$-Einsetzung
erzeugt~\eqref{eq:originalportmatching} und denselben Referenzabzug im
Flavor-Hintergrund? v485 subtrahiert die logarithmische Singularität von
$|D|^{-1}$. Das definiert noch keine endliche Diagonale des neuen
$R$-$R$-Kerns vierter Polordnung an den Fuchsdefekten.
Ebenso ist der Mark-Halbdrehungsprojektor nicht allein durch die Bezeichnung
mit dem ursprünglichen Calderón- oder Nahtprojektor identifiziert.
Der fehlende Gegenstand ist diese konkrete Einsetzungs- und
Referenzabbildung, nicht ein generell unbekannter ursprünglicher Randkernel.
Die Formeln stammen aus \file{verification/v484_seam_contact_unit.py},
\file{verification/v485_contact_diagonal_closed.py} und dem
Alpha-Transfer in \file{_archive/tfpt-45/source_extracts/03_em_flavor_source.tex}.

\subsection{Die 48 und die 41 sind zwei Antworten derselben geladenen Quelle}\label{sec:chargeweightedsource}
Der neue neutrale Zustand $R$ trägt die ungewichtete Familienantwort.
Die ursprüngliche Hyperladung ist dagegen durch die $3+2$-Markierung bereits festgelegt:
\[
y=(-1/3,-1/3,-1/3,1/2,1/2,0,0,0),\qquad
Y(\beta)=y\cdot\beta,\qquad J_Y=y\cdot h.
\]
Es wird also keine neue Ladung gewählt. Auf den ursprünglichen 48 Wurzeln ergeben sich
\[
\sum_{\beta\in\mathcal R_{48}}Y(\beta)^2=10,
\qquad \|y\|^2=5/6,\qquad k_Y=2\|y\|^2=5/3.
\]
Die tatsächlichen Multiplizitäten sind $18$ Felder mit $Y=1/6$, jeweils
$9$ mit $Y=-2/3,1/3$, $6$ mit $Y=-1/2$ sowie jeweils $3$ mit $Y=1,0$.
Das sind die bereits bekannten drei Familien einschließlich der neutralen Komponente.

Mit denselben neutralen Paarvertices $B_\beta$ wie in Abschnitt~\ref{sec:neutralfield} definiert die Ladungsmarkierung
\begin{align}
 B_Y&=\sum_{\beta\in\mathcal R_{48}}Y(\beta)^2B_\beta
     =\frac52 T+R_Y,\\
 R_Y&=\frac5{24}R+S_Y,\qquad
 S_Y=\frac12:h^{\mathsf T}A_{S_Y}h:,
 \quad (A_{S_Y})_{ij}=6(y_i y_j-\delta_{ij}y_i^2)\quad(i,j\leq5).
\end{align}
Außerhalb des $D_5$-Blocks ist $A_{S_Y}$ null. Alle drei Zustände sind in derselben
ursprünglichen Quelle berechnet. Wegen $\tr A_{R_Y}=0$ ist auch $R_Y$ ein neutrales
Primärfeld vom Gewicht zwei. Die positive Quellenform liefert exakt
\begin{equation}\label{eq:neutralchargegram}
\operatorname{Gram}(R,R_Y)=
\begin{pmatrix}48&10\\10&35/3\end{pmatrix},\qquad
\det\operatorname{Gram}(R,R_Y)=460,\qquad \|S_Y\|^2=115/12.
\end{equation}
Eine gemeinsame isometrische Abbildung dieser beiden Zustände muss deshalb beide
Richtungen erhalten. Dies verlangt nicht, den ursprünglichen einzelnen topologischen
Rang-eins-Port durch einen Rang-zwei-Port zu ersetzen: Die eine Anzeige und der sie
tragende vollständige Quellenraum sind unterschiedliche Objekte.

Die konkrete Eichstromantwort macht den Unterschied sichtbar:
\begin{align}
 \langle R(z)J_Y(w)J_Y(u)\rangle_\Omega&=0,\\
 \langle R_Y(z)J_Y(w)J_Y(u)\rangle_\Omega
 &=\frac{115/36}{(z-w)^2(z-u)^2}.
\end{align}
Das folgt direkt aus der Wick-Kontraktion $y^{\mathsf T}A y$. $R$ liegt vollständig
im Familienanteil und hat mit diesem $D_5$-Strom keinen entsprechenden Vakuum-Dreipunktterm.
$R_Y$ enthält zusätzlich die durch die tatsächlichen Ladungen bestimmte Richtung $S_Y$.
Beide Gleichungen sind chirale Quellenantworten; keine von beiden wird hier als
vierdimensionale Photonpolarisation ausgegeben. Die ursprüngliche Trennung zwischen
topologischer Besetzung und gewichteter Eichvariation bleibt erhalten.

Unter der bereits im physikalischen Vertrag erklärten Identifikation mit
vierdimensionalen Weylfeldern und einem komplexen Higgsdoublett folgt aus denselben Ladungen
\[
 b_Y=\frac23\underbrace{\tr_{48}Y^2}_{10}
       +\frac13\underbrace{\tr_HY^2}_{1/2}=\frac{41}{6},
 \qquad b_1=b_Y/k_Y=\frac{41}{10}.
\]
Das ergänzt den vorhandenen Compilerbefund aus v2 um den expliziten Quellenanschluss.
Die vierdimensionale Schleifeninterpretation bleibt eine eigene Voraussetzung.

Schließlich antwortet derselbe ladungsgewichtete Operator auf alle sechs behaltenen
Neutrinopaare aus Abschnitt~\ref{sec:sourcepairs}:
\begin{equation}\label{eq:twojointpairresponses}
 R_0B_{ab}=4B_{ab},\qquad
 R_{Y,0}B_{ab}=-\frac53 B_{ab},\qquad a,b\in\{1,2,3\}.
\end{equation}
Alle Paare haben den $D_5$-Impuls $(-1)^5$ und keine $D_5$-Oszillatoren.
Daher trägt $S_{Y,0}$ genau $-5/2$ bei, während $(5/24)R_0=5/6$ liefert.
Die Paare bleiben hyperladungsneutral: $Y(B_{ab})=0$. Ein neutraler zusammengesetzter
Operator ist nicht der Ladungsoperator selbst. Da beide Antworten auf dem ganzen
Sechserraum skalar sind, wählen sie keine relative Majoranaphase aus.

\picturetext{Die Quelle enthält dieselben Materiefelder in zwei verschiedenen Rechnungen:
Einmal zählt jedes Feld gleich viel, einmal zählt seine tatsächliche Ladung mit.
So hängen die 48 und die 41 zusammen. Die Neutrinopaare reagieren auf beide
Rechnungen in fest bestimmter Weise. Diese gemeinsame Antwort verbindet die Teile,
entscheidet aber noch nicht, welcher Massenzustand tatsächlich entsteht.}

Die exakte Rechnung steht in \file{tfpt_explorer/source_charge_response.py}.
Sie verwendet die ursprünglichen Wurzeln und die positive Heisenbergform; der
vorhandene Vierdimensionalitäts- und Higgsvertrag wird ausdrücklich mitgeführt.

\subsection{Welche gemeinsame Reduktion die vorhandenen Antworten wirklich erhält}
Für einen positiven Operator $A=D^\dagger D$ auf erhaltenen und eliminierten Zuständen gilt
\begin{equation}\label{eq:fullsourceschur}
P(A+z)^{-1}P
=\bigl[A_{11}+z-A_{12}(A_{22}+z)^{-1}A_{21}\bigr]^{-1}.
\end{equation}
Der Komplementterm ist das genaue Gedächtnis der nicht angezeigten Zustände. Bei positiver invertierbarer Komplementresolvente würde die Gleichheit mit $(A_{11}+z)^{-1}$ bereits $A_{12}=0$ erzwingen. Ein exakt abgeschlossener endlicher Block ist deshalb eine stärkere Annahme als eine antworttreue Reduktion mit Gedächtnis.

Am neuen ursprünglichen Wegoperator lässt sich diese Regel vollständig ausführen. Aus
\begin{equation}
D^{\rm path}=\begin{pmatrix}yI&-\delta I\\-\delta M^2&yI\end{pmatrix}
\end{equation}
ergeben sich mit $a=y^2+\delta^2$ die Schurantworten
\begin{equation}
S_D=yI-\frac{\delta^2}{y}M^2,\qquad
S_{D^\dagger D}=\frac{y^2}{a}S_D^\dagger S_D.
\end{equation}
Die genaue Determinantenverklebung lautet
\begin{equation}
\boxed{\det(D^\dagger D)=a^3\det S_{D^\dagger D}
=y^6|\det S_D|^2=(y^6-\delta^6)^2.}
\end{equation}
Eine Vergröberung auf einen der beiden Zugänge erhält also die Flavorantwort, wenn die induzierte Metrik und der Beitrag des eliminierten Zugangs erhalten bleiben. Eine einfache Kompression $P D^\dagger DP$ ist ein anderer Operator. Ebenso würde die bloße Auswahl eines $T^3$-Eigenraums einen der beiden Faktoren $y^3\mp\delta^3$ verlieren. Die ursprüngliche faktorisierte Hilbertstruktur erlaubt dagegen eine Nahtprojektion auf einem eigenen Faktor, ohne den gesamten Flavortransport zu halbieren.

Die ursprüngliche Masterwirkung enthält ausdrücklich eine gemeinsame Operatorfamilie. Aus dem Produkt kommutierender Selektoren $P_{\rm adm}=P_{\rm prim}P_{\rm sing}P_\Theta$ folgt jedoch keine orthogonale Aufteilung in vier unabhängige physische Wirkungen. Selbst eine tatsächliche Blockzerlegung beweist noch nicht, dass jeder Block von nur einer Masterkoordinate abhängt. Das entsprechende Argument in \file{02_carrier_source.tex}, Zeilen 3117--3123, benötigt diese zusätzlichen Eigenschaften.

Die vorhandenen Formeln erzeugen dabei keinen unmittelbaren Alpha--Flavor-Widerspruch: Der kanonische Seed $\phi_0$ im Flavorweg ist von $\phi_{\rm seam}(\alpha)$ unterschieden. Auch ein Test allein der gemischten Ableitungen genügt nicht. Für $D_{\rm fl}=s(\alpha)D_0(\delta)$ gilt
\begin{equation}
\log\det(D_{\rm fl}^\dagger D_{\rm fl})
=2n\log s(\alpha)+\log\det(D_0^\dagger D_0).
\end{equation}
Die gemischte Ableitung verschwindet, während die Alpha-Erstableitung einen zusätzlichen Beitrag enthält. Entscheidend bleiben daher derselbe vollständige Operator, dieselben lokalen Subtraktionen und die gemeinsam berechneten Erstableitungen.

\key{\textbf{Zusammenhängendes Ergebnis:}
Der ursprüngliche Fuchsprozess realisiert den Flavoroperator durch wirkliche Wege und trägt lokal bereits primitive Determinantenwindung. Die vorhandene E8-Quelle realisiert den Alpha-Rundlauf durch den konkreten bedingten Port~\eqref{eq:actualheatport}; seine Antwort im Flavor-Hintergrund ist durch~\eqref{eq:jointRresponse} festgelegt. Damit ist die tragende gemeinsame Auswahlfrage auf eine bestimmte Naht-, Zustands- und Normierungsabbildung zugespitzt. Sie muss den topologischen P1-Kreis, den internen Flavortransport, den Spinanteil und diesen neutralen Quellenzugang gleichzeitig erhalten. Die bestehenden Massen-, Yukawa- und Alphaformeln bleiben die Bedingungen an diese Abbildung.}

\subsection{Ausführbare Nachweise dieser Fortsetzung}
\file{tfpt_explorer/flavor_path_transport.py} integriert die ursprüngliche Verbindung entlang beider Halbwege und unabhängig entlang des geschlossenen Kreises. Es berechnet die Determinantenwindung aus den integrierten Lösungen und vergleicht die vollständigen endlichen Greenmatrizen. \file{tfpt_explorer/neutral_source_response.py} benutzt die ursprünglichen E8-Wurzeln, ihre Stressmomentmatrizen und den tatsächlichen Heisenbergzustand vom Grad zwei. Es hält den bedingten Alpha-Port und den abweichenden lokalen Defektzustand getrennt.

Die Quellenstellen sind insbesondere \file{verification/v117_monodromy_weyl_a3.py}, \file{v118_hexagon_family_dictionary.py}, \file{v498_celestial_wp5b_singular_vector.py}, \file{matter-geometry-round15/PROOF.md} unter \file{experiments/theory-contracts/}, sowie \file{_archive/tfpt-45/source_extracts/01_boundary_kernel_source.tex}, Zeilen 192--262, und \file{04_qft_source.tex}, Zeilen 219--307. Der Riemann--Hilbert-Feldanschluss verwendet die normierten Defektmatrixelemente und Wick-Determinanten aus \cite[Abschnitt 4.1]{gavrylenkomarshakov}; ihre Interpretation als physische Nahtkorrelation ist zusätzlich zu prüfen.

\section{Die Windung wirkt auf der ganzen Quelle: neuer Anschluss und gemeinsamer Massentest}\label{sec:spinliftsynthesis}

\key{\textbf{Der neue positive Herkunftssatz:}
Die primitive Determinantenwindung des ursprünglichen Flavorprozesses bestimmt beim internen Spin-Lift ein zusätzliches Minuszeichen. In der vorhandenen E8-Verklebung ist dieses Zeichen genau die Involution, die den 120er-Tensorsektor vom 128er-Spinorsektor unterscheidet. Damit ist ein Sechserumlauf auf tatsächlichen geladenen Quellenfeldern konstruiert. Die Identifikation dieses inneren Zeichens mit vierdimensionaler Fermionstatistik bleibt eine weitere, an den vorhandenen Massenmechanismen zu prüfende Aussage.}

Diese Fortsetzung prüft den Anhang \file{30ea245e-2aae-4de0-8aad-8e3655f3e12b} und das vollständig gelesene \href{https://claude.ai/artifact/QKsCJRb79D8Z3pbxqV54n2}{Claude-Artefakt zur TFPT-Prozesskette}. Dessen drei neue Contractordner und der angegebene Commit lagen im lokalen Repository nicht vor. Die dort genannten 109 und 78 Prüfungen sowie die Enumeration von 7560 Einbettungen werden nicht als hier ausgeführte Prüfungen ausgegeben. Die folgenden Aussagen sind am ursprünglichen Fuchssystem, den tatsächlichen E8-Wurzeln, der bestehenden Majoranaroute und den zitierten Primärarbeiten geprüft.

\subsection{Eine weggeworfene Windung wäre eine weggeworfene Quellenwirkung}
Die Eigenwerte der lokalen ursprünglichen Residuen sind $(0,1/3,2/3)$. Die Spurlosfassung $(0,1/3,-1/3)$ hat dieselben exponentierten Eigenwerte. Ihre Verbindungspfade sind dennoch verschieden:
\begin{equation}
u(t)=\diag(1,e^{2\pi it/3},e^{4\pi it/3}),\qquad
u_0(t)=\diag(1,e^{2\pi it/3},e^{-2\pi it/3}).
\end{equation}
Beide enden am gleichen $SU(3)$-Element. Die Determinante des ersten Pfads windet sich einmal, die des zweiten keinmal. Genau die erste Klasse trägt der integrierte Originalkreis, weil $\det\Psi(z)=1-z^4$ lokal eine einfache Nullstelle hat.

Auf dem reellen sechsdimensionalen Familienvektorraum hebt eine $U(3)$-Wirkung nach $\operatorname{Spin}(6)$ mit
\begin{equation}
\mathcal S(u)=\det(u)^{-1/2}\Lambda^\bullet u .
\end{equation}
Die Quadratwurzel muss entlang des Pfades fortgesetzt werden. Auf dem chiralen Viererraum $\Lambda^{\rm even}\C^3$ erhält man deshalb am ursprünglichen Ende
\begin{equation}\label{eq:nativespinlift}
\boxed{\widehat M=-\Lambda^{\rm even}M\in SU(4),\qquad
\widehat M^3=-I_4,\qquad \widehat M^6=I_4.}
\end{equation}
Der rohe Fuchs-Transport ist in seinen gedruckten Koordinaten nicht euklidisch unitär. Die bereits bestimmte positive flache Metrik beziehungsweise eine polare Retraktion mit festgehaltenem Endpunkt liefert die hier verwendete unitäre Pfadklasse. Die Windungszahl bleibt erhalten. Die Gleichung ersetzt den ursprünglichen Transport somit nicht durch eine frei passende Matrix.

\picturetext{Der gewöhnliche Dreieranzeiger speichert nur die Endphasen. Der mitgeführte Spinrahmen erinnert sich zusätzlich daran, wie die Runde durchlaufen wurde. Nach drei Runden steht an bestimmten Quellenfeldern noch ein Minuszeichen; erst nach sechs Runden ist auch dieses verschwunden.}

\subsection{Der Spinrahmen trifft genau die ursprüngliche E8-Verklebung}
Im Original gilt als wirkliche Untergruppe von E8
\begin{align}
G&=\frac{\operatorname{Spin}(10)\times SU(4)}
 {\langle(\omega_{10},-iI_4)\rangle},\\
248&=(45,1)\oplus(1,15)\oplus(10,6)
\oplus(16,4)\oplus(\overline{16},\overline4).
\end{align}
Die diagonale Vierergruppe ist der \emph{Kern} dieser Darstellung: Sie wirkt auf E8 trivial. Die einzeln wirksamen Faktorzentren erfüllen dagegen
\begin{equation}\label{eq:deckspinidentity}
[(1,-I_4)]=[(-1_{\operatorname{Spin}(10)},1)]
=\eta_{D_8},\qquad
\eta_{D_8}=+I_{120}\oplus(-I_{128}).
\end{equation}
Alle 240 Originalwurzeln bestätigen dieselbe Charaktergleichheit; die acht Cartanrichtungen gehören zum geraden 120er-Sektor. Die Vorzeichen respektieren die ursprünglichen Wurzelklammern und damit die ganze von ihnen erzeugte Quellenalgebra.

Für $g=[(1,\widehat M)]\in G\subset E_8$ folgt auf der vollständigen Quelle
\begin{equation}\label{eq:fullsourcesix}
\boxed{g^3=\eta_{D_8},\qquad g^6=I.}
\end{equation}
Dies ist ein Anschluss an geladene Felder. Er kommutiert mit dem $\operatorname{Spin}(10)$-Faktor; die inneren Materieladungen werden nicht miteinander vertauscht. Auf dem Familienvektor $6=\Lambda^2 4$ verschwindet das zusätzliche zentrale Minuszeichen. Die Spinorfelder $(16,4)$ behalten es. Das unterscheidet diese Operation vom Zweifaseroperator $U_6=M\oplus(-M)$: Beide besitzen Ordnung sechs, wirken aber zunächst auf verschiedenen Darstellungsräumen.

\begin{center}
\begin{tikzpicture}[node distance=6mm]
\node[box,text width=12.8cm] (w) {\textbf{Ursprünglicher Weg}\\$\det\Psi=1-z^4$; ein Punkturumlauf hat Windung $1$};
\node[box,below=of w,text width=12.8cm] (l) {\textbf{Familien-Spinlift}\\$\widehat M=-\Lambda^{\rm even}M$; nach drei Umläufen $-I_4$};
\node[box,below=of l,text width=12.8cm] (e) {\textbf{Dieselbe geladene E8-Quelle}\\$g^3=+I_{120}\oplus(-I_{128})$, $g^6=I$};
\node[box,below=of e,text width=12.8cm] (p) {\textbf{Gemeinsam zu identifizierende physische Felder}\\Statistik, drei Familien und zulässiger Neutrino-Massenoperator};
\draw[arr] (w)--(l);\draw[arr] (l)--(e);\draw[dasharr] (e)--(p);
\end{tikzpicture}
\end{center}
Die drei ersten Kästen sind an denselben Originaldaten verbunden. Der gestrichelte Übergang verlangt ein physisches Feldwörterbuch.

\subsection{Was der neue Anomalieselektor tatsächlich auswählt}
Der interessante Gedanke des Anhangs ist, zusätzlich eine vierdimensionale Spin-$\Z_4$-Struktur mit $\omega^2=(-1)^F$ zu verlangen. Für gleichorientierte Weylfermionen zählt deren Anomalie den Unterschied der $\omega=+i$- und $\omega=-i$-Zustände modulo 16~\cite[Abschnitt 3, Gleichung (7)]{razamattong}. Auf dem ursprünglichen Halbspinor gilt feldweise
\begin{equation}
S^+=\Lambda^{\rm even}\C^5=10_1\oplus\overline5_{-3}\oplus1_5,\qquad
\boxed{X=5-2N=5(B-L)-4Y.}
\end{equation}
Alle 16 Komponenten tragen denselben $\Z_4$-Charakter. Das sterile linkschirale $\nu^c$ ist unter diesen Voraussetzungen nötig: Ohne dieses Feld verbleibt der Index $15$ statt $0$ modulo 16. Diese konditionale Verbindung ist in der Anomalieliteratur bereits bekannt~\cite[Abschnitt 4.3]{garciamontero}.

Für $m$ gleichartige Familien eines chiralen Halbspinors von $\operatorname{Spin}(2g)$ mit ungeradem $g$ lautet der Index
\begin{equation}
\nu=\pm m\,2^{g-1}\pmod{16}.
\end{equation}
Ist $m$ ungerade, insbesondere beim vorhandenen Dreifamilienreadout, erzwingt Anomaliefreiheit $g\ge5$. \textbf{Minimalität innerhalb dieser Klasse wählt $g=5$.} Vorausgesetzt sind vierdimensionale Weylfelder, dieselbe Spin-$\Z_4$-Struktur auf allen Feldern, dieser Halbspinor als Materieträger und keine weitere Anomaliekompensation. Der Satz leitet diese Voraussetzungen nicht selbst her. Mit unverändertem Volumengenerator $X_g=g-2N$ sind die Indizes für $g=1,3,5$ gleich $1,12,0$; die $4$ des Anhangs bei $g=3$ entspricht dem invertierten Generator.

Die Anomalie der noch unaufgeteilten 64 Quellenkomponenten $(16,4)$ verschwindet unter beiden Faktorzentren. Sie entscheidet deshalb nicht allein, welcher Faktor Materie und welcher Familienmultiplikität bedeutet. Auch die vier Marken und das $A_3$-Komplement bleiben eigenständige, bereits geometrisch begründete Daten. Der bedingte Selektor kann eine bisherige Dimensionsannahme ersetzen; er beweist nicht aus sich heraus P1, P2, die $3+2$-Markierung und Raumzeit.

\subsection{Der ganze Massensektor entscheidet über die neue Zusatzannahme}\label{sec:spinmajorana}
Eine gemeinsame Lösung muss dieselbe Symmetrie auch auf den bereits verwendeten Higgs- und Neutrinofeldern realisieren. Das ursprüngliche Standardmodell-Paper benutzt, weil ein elementares $\overline{126}_H$ im erklärten Hull fehlt, den Einsatz
\begin{equation}
\frac{(16_F\,\overline{16}_H)^2}{\Lambda}.
\end{equation}
Quelle: \file{tfpt_2_standard_model.tex}, Zeilen 2003--2006, und \file{verification/v488_majorana_clebsch_door.py}. Unter einer gewöhnlichen internen $\Z_4$ ist der Einsatz neutral und zulässig.

Für die \emph{zusätzliche} Spin-$\Z_4$-Identifikation ist die Feldgruppe dagegen
\begin{equation}
G_f=\frac{\operatorname{Spin}(1,3)\times\Z_4}
 {\langle(-1,\omega^2)\rangle}.
\end{equation}
Jedes einzelne Feld muss auf dem Quotientenkern trivial wirken:
\begin{equation}
r_{\rm Lorentz}(-1)\,r_X(\omega^2)=I.
\end{equation}
Der ursprüngliche Weyl-$16_F$ besteht diesen Test mit $(-I)(-I)=I$. Ein skalares $\overline{16}_H$ besteht ihn mit $(+I)(-I)=-I$ nicht. Der skalare $10_H$ besteht ihn. Die Neutralität des gesamten dimension-fünf-Einsatzes repariert die fehlende Darstellung seines einzelnen skalaren Spinorfelds nicht.

\key{\textbf{Das gemeinsame Ergebnis schränkt die erlaubte Lösung ein:}
Der neue Spin-$\Z_4$-Selektor und die alte skalare $\overline{16}_H$-Realisierung können nicht gleichzeitig unverändert vorausgesetzt werden. Das widerlegt weder die bereits berechneten Neutrinowerte noch TFPT insgesamt. Die physische Herkunft des Massenoperators und die Herkunft der Feldstatistik müssen gemeinsam geschlossen werden.}

Das nackte Paar $\nu^c\nu^c$ hat $X=10\equiv2\pmod4$. Ein neutraler Scalaron und der oben berechnete neutrale Zustand $R$ tragen die fehlende Gegenladung nicht. Ein bosonischer Paarsektor mit Gegenladung $2$ modulo vier könnte sie tragen; seine Existenz als Quelloperator, seine physische Kopplung und eine mögliche Kondensation sind getrennte Aussagen. Insbesondere bricht ein Majoranamassenterm diese $\Z_4$, wenn kein transformierender Einsatz mitgeführt wird~\cite{garciamontero}.

\subsection{Der benötigte geladene Paarkanal ist in der tatsächlichen Quelle vorhanden}\label{sec:sourcepairs}
Der Konflikt wurde unmittelbar in der bestehenden Quelle weiterverfolgt. Es wurde kein neues elementares Higgsfeld eingeführt. Die ursprünglichen vier Singletströme $C_a$ im markierten Spinorblock besitzen Wurzeln $\alpha_a$ mit
\begin{equation}
(\alpha_a,\alpha_b)=1+\delta_{ab},\qquad
C_a(z)C_b(w)=p_ap_b\,\epsilon(\alpha_a,\alpha_b)
(z-w)^{1+\delta_{ab}}\bigl[V_{\alpha_a+\alpha_b}(w)+\cdots\bigr].
\end{equation}
$p_a$ und $\epsilon$ sind die tatsächlichen Originalphasen und der ursprüngliche Gitterkokzyklus. Bei der Konvention $X=5-2N$ trägt dieser $C$-Zweig $X=+5$; der adjungierte Zweig trägt $-5$. Damit gilt im einfachen unitären Level-eins-Quotienten
\begin{equation}
\boxed{:C_aC_b:=0\quad\text{für alle }a,b.}
\end{equation}
Dies wurde an allen 16 Paaren mit der ursprünglichen affinen Vakuumform nachgerechnet: Die Normen der Grad-zwei-Vektoren sind exakt null. Das ist eine wirkliche Nullrelation, keine lediglich fehlende Implementierung.

Die höhere Operatorentwicklung liefert dagegen konkrete nichtverschwindende Zustände:
\begin{align}
A_{ab}&=C_{a,-2}C_{b,-1}|0\rangle,\quad a<b,
&h&=3,&\norm{A_{ab}}^2&=1,\\
B_{aa}&=C_{a,-3}C_{a,-1}|0\rangle,
&h&=4,&\norm{B_{aa}}^2&=1,\\
B_{ab}&=(C_{a,-3}C_{b,-1}+C_{b,-3}C_{a,-1})|0\rangle,\quad a<b,
&h&=4,&\norm{B_{ab}}^2&=2.
\end{align}
Die sechs $A_{ab}$ sind antisymmetrisch in den Familienindizes. Die zehn $B_{ab}$ sind symmetrisch. Ihre Gramformen sind diagonal mit den angegebenen positiven Normen. Die symmetrischen Grad-vier-Vektoren sind außerdem orthogonal zu den Ableitungen der antisymmetrischen Grad-drei-Vektoren.

Der Test aller acht positiven einfachen Wurzel-Nullmoden weist tatsächliche höchste Gewichte nach:
\begin{equation}
A_{01}:\ (00002;010)\ \longleftrightarrow\ (126,6),\qquad
B_{00}:\ (00002;200)\ \longleftrightarrow\ (126,10).
\end{equation}
Die Darstellung beziehungsweise ihre Konjugierte richtet sich nach der gewählten Spinorchiralität; der adjungierte Quellenoperator liefert den Gegenkanal. Die höchsten Gitterimpulse besitzen Normquadrat sechs beziehungsweise acht. Ihre kleinsten konformen Gewichte sind daher drei beziehungsweise vier; die angegebenen Zustände erreichen diese Schranken. Eine beliebige Tensorproduktzerlegung hätte weder diese erste Auftretensstufe noch die Nullrelation im Grad zwei gezeigt.

Auch die Quellkopplung ist berechnet. Mit dem normierten Operator $B_{00}=C_{0,-3}C_{0,-1}|0\rangle$ liefert der ursprüngliche Kokzyklus
\begin{equation}\label{eq:nativemajoranacoupling}
\boxed{\langle B_{00}^{\dagger}(z_3)\,C_0(z_1)\,C_0(z_2)\rangle
=\frac{z_{12}^{\,2}}{z_{13}^{\,4}z_{23}^{\,4}}.}
\end{equation}
Die normierte OPE-Zahl eins wurde damit nicht zusätzlich angepasst. Die unten beschriebenen Grad-zwei-Zeugen besitzen entsprechend den Abstandsfaktor $1/(z_{13}^2z_{23}^2)$ mit Koeffizient eins. Eine vierdimensionale Yukawa-Kopplung verlangt weiterhin die physische Feldabbildung.

\key{\textbf{Der positive Neutrinoanschluss:}
Die vollständige Quelle enthält einen bosonischen symmetrischen Spinor-Paaroperator mit $X=\pm10$, also Ladung zwei modulo vier. Er ist mit der zusätzlichen Spin-$\Z_4$-Statistikregel vereinbar und besitzt den für eine Majoranakopplung benötigten Darstellungstyp. Dass ein elementares $126_H$ im Grad-eins-Hull fehlt, bedeutet somit nicht, dass dieser zusammengesetzte Kanal in der vollständigen Quelle fehlt.}

Die Rechnung erhält auch die anderen Paarstrukturen der vorhandenen Yukawaherleitungen. Als endlichdimensionale Darstellung gilt
\begin{equation}
\operatorname{Sym}^2(16\otimes4)
=(10,10)\oplus(126,10)\oplus(120,6).
\end{equation}
Im Grad-zwei-Produkt existieren höchste Quellenvektoren vom Typ $(10,10)$ und $(120,6)$ mit Norm eins und Gitterlänge zum Quadrat vier. Der $(126,10)$-Anteil ist dort null und erscheint erst im Grad vier. Die symmetrische Paarabbildung hat folglich in dieser 2080-dimensionalen Eingabe den Bildrang $100+720=820$ und den Kern der Dimension $1260$. Diese Zahlenaussage folgt aus der multiplikitätsfreien Zerlegung und den geprüften höchsten Vektoren; sie behauptet keinen zusätzlich durchgerechneten 2080-mal-2080-Gramtest.

\picturetext{Die Quelle erlaubt nicht jedes Zusammenschieben zweier Materiezeichen. Sie bestimmt, welche Verbindung sofort entsteht, welche nur mit einer zusätzlichen Bewegung entsteht und welche erst auf der nächsten Anregungsstufe auftaucht. Die Yukawa- und Neutrinokanäle lassen sich damit in derselben Mechanik verorten.}

\subsection{Der symmetrische Paartransport schließt an den ursprünglichen Flavorweg an}
Der benutzte Spinlift wirkt auf
\begin{equation}
\Lambda^{\rm even}\C^3=1\oplus\Lambda^2 3=1\oplus\overline3.
\end{equation}
Sein symmetrischer Familienpaarraum zerfällt daher als $10\to1\oplus\overline3\oplus\overline6$; der konjugierte Kanal enthält $1\oplus3\oplus6$. Der Sechserraum ist der Raum symmetrischer Matrizen für drei Familien. Seine Komponenten sind noch keine bestimmten Erwartungswerte.

Die Konjugation lässt sich für die \emph{tatsächliche diskrete Originalmonodromie} exakt mitführen. Setze
\begin{equation}
B_{\rm real}=\begin{pmatrix}1&0&0\\0&0&-1\\0&-1&0\end{pmatrix}.
\end{equation}
Für die bereits verwendeten Originalmatrizen $M,U$ gilt
$M B_{\rm real}=B_{\rm real}\overline M$ und
$U B_{\rm real}=B_{\rm real}\overline U$.
Die Hodgeabbildung $H$ mit
$(12)\mapsto e_3^*$, $(13)\mapsto-e_2^*$, $(23)\mapsto e_1^*$
erfüllt $H\Lambda^2(T)=\overline T H$ für $T=M,U$.
Also verbindet $J=B_{\rm real}H$ die Außenquadratwirkung mit der Originalwirkung. Einfeldseitig bleibt das echte Windungszeichen erhalten:
\begin{equation}
J(-\Lambda^2 M)=-M J.
\end{equation}
Im symmetrischen Paar hebt es sich dagegen auf:
\begin{equation}\label{eq:majoranapairtransport}
\boxed{\operatorname{Sym}^2(J)\,
\operatorname{Sym}^2(-\Lambda^2 M)
=\operatorname{Sym}^2(M)\,\operatorname{Sym}^2(J).}
\end{equation}
Der Anschluss gilt jetzt auch in der gedruckten Quellenbasis: $(C_1,C_2,C_3)$ entspricht der Reihenfolge $(23),(13),(12)$. Die vorhandenen Originalphasen wurden durch die vier einfachen anhebenden und absenkenden $SU(3)$-Wurzeloperationen geprüft. Mit der umkehrenden Permutation $P$ ist der vollständige Komponentenintertwiner
\begin{equation}
J_C=B_{\rm real}HP
=\begin{pmatrix}1&0&0\\0&0&-1\\0&1&0\end{pmatrix}.
\end{equation}
Er transportiert die tatsächlich behaltenen Quellenkomponenten nach $-M$ und ihre symmetrischen Paare nach $\operatorname{Sym}^2M$; die analoge markierte Relation gilt für $U$. Es wird keine Gleichheit von $3$ und $\overline3$ für die ganze kontinuierliche Gruppe $SU(3)$ behauptet. Ebenso folgt aus~\eqref{eq:majoranapairtransport} keine bereits ausgewählte Neutrinomatrix.

Die Operatoren und ihre inneren Ladungen sind somit vorhanden; ihr Paartransport passt an den ursprünglichen Flavorweg. Zur vorhandenen physischen Massenskala führt erst der weitere Herkunftssatz: Welcher dieser Quellenoperatoren wird zu einem lokalen vierdimensionalen Feld, welcher Zustand erzeugt seinen Einsatz und wie ergeben sich daraus die bereits abgeleitete Skala und die Familienmatrix? Das chirale Gewicht vier ist keine vierdimensionale Ingenieursdimension. Die Konstruktion macht das elementare skalare $16_H$ nicht quotientzulässig; sie liefert einen konkreten zusammengesetzten Alternativkanal innerhalb derselben Quelle.

\subsection{Der neutrale Alpha-Kandidat antwortet auf diesen Neutrinokanal}\label{sec:neutralpairward}
Die neue Paarstruktur kann mit dem bereits bestimmten neutralen Feld $R=4T_{A_2}-8T_{U(1)}$ in derselben Quelle geprüft werden. Verwende dieselbe Auswahl, in der $C_0$ der ausgelassene Familien-Singulettzugang und $C_1,C_2,C_3$ der Dreieranteil sind. Für einen reinen normierten Gitterzustand $|\lambda\rangle$ gilt mit der ursprünglichen Matrix $A_R$
\begin{equation}
R_0|\lambda\rangle=\tfrac12\lambda^{\mathsf T}A_R\lambda\,|\lambda\rangle.
\end{equation}
Der symmetrische Paarzustand $B_{11}$ hat $\lambda=2\alpha_1$. Seine Familienanteile besitzen
\begin{equation}
h_{A_2}=4/3,\qquad h_{U(1)}=1/6,\qquad
\boxed{R_0 B_{11}=4B_{11}.}
\end{equation}
Da $R$ mit den Nullmoden des gewählten $SU(3)$ kommutiert, gilt derselbe Eigenwert auf seinem gesamten symmetrischen Sechsermodul. Der antisymmetrische Paarzustand $A_{12}$ trägt dagegen $h_{A_2}=1/3$, $h_{U(1)}=1/6$ und $R_0A_{12}=0$. Der reine Singulettpaarzustand $B_{00}$ antwortet mit $R_0B_{00}=-12B_{00}$. Diese Werte wurden aus denselben Originalwurzeln und derselben zuvor berechneten Matrix $A_R$ unabhängig nachgerechnet.

Für den normierten Paarzeugen im ausgewählten Dreieranteil ergibt sich damit die konkrete Quellenkorrelation
\begin{equation}\label{eq:neutralmajoranaward}
\boxed{\langle B_{11}^{\dagger}(\infty)\,R(z)\,B_{11}(0)\rangle
=\frac{4}{z^2}.}
\end{equation}
Der Ausdruck ist eine Matrixantwort im geladenen Paarzustand, kein nichtverschwindender geladener Vakuumerwartungswert. Die Quelle besitzt nun eine berechnete gemeinsame Antwort zwischen dem Alpha-Kandidaten und dem symmetrischen Neutrinopaar. Es wird keine zusätzliche freie Kopplungszahl eingesetzt.

\key{\textbf{Wie die Kette jetzt zusammenklickt:}
Die ursprüngliche Windung liefert einen geladenen Quellenlift. Derselbe Lift transportiert den vorhandenen symmetrischen Neutrinopaarraum wie der ursprüngliche Flavorprozess. Derselbe neutrale Zustand, dessen Vakuumnorm die 48 des bedingten Alpha-Ports liefert, antwortet auf diesen Paarraum mit dem exakt berechneten Koeffizienten vier. Die Herkunft einer gemeinsamen physischen Realisierung ist dadurch an konkrete Mehrpunktantworten gebunden. Eine kondensierte Phase, eine vierdimensionale Kopplung und die bereits bekannte Massenskala folgen aus dieser Quellenidentität allein noch nicht.}

\subsection{Warum die vorhandene Fermionisierung den Statistikschritt nicht überspringt}
Für die tatsächliche chirale E8-Fermionisierung gelten
\begin{equation}
(E_8)_1=O\oplus S,\qquad
\mathcal F_{\rm NS}=O\oplus V,\qquad
\mathcal F_{\rm R}=S\oplus C.
\end{equation}
Die E8-Decksymmetrie hat Zeichen $(+O,-S)$; Fermionparität in den neuen Sektoren hat Zeichen $(+O,-V)$ beziehungsweise $(+S,-C)$~\cite[Abschnitt 5.1]{boylesmith}. Das sind verschiedene Operatoren auf verschiedenen Sektoren.

Die 16 lokalen CAR-Fermionen bilden den $\operatorname{Spin}(16)$-Vektor $(10,1)\oplus(1,6)$. Dort wirkt $\omega_{10}$ als $\diag(-I_{10},+I_6)$, also $\omega_{10}^2=I$. Auf jedem dieser Fermionen ist $(-1)^F=-I$. Die konkreten 16 CAR-Felder sind daher nicht bereits die 16 Weylkomponenten einer Standardmodellfamilie. Die ursprünglichen $(16,4)$-Felder besitzen dagegen konformes Gewicht $5/8+3/8=1$ und sind bosonische Quellenströme.

Auch der ursprüngliche geometrische Vierteldrehungslift ist zu erhalten: v522 erfüllt $\alpha_{\rm quarter}^4=(-1)^F$, also eine NS-$\Z_8$-Relation. Durch $\omega=\alpha_{\rm quarter}^2$ erhält man formal eine Halbdrehung mit $\omega^2=(-1)^F$; das identifiziert nicht die ursprüngliche Vierteldrehung mit dem inneren Zentrum. Der Herkunftssatz~\eqref{eq:fullsourcesix} bewahrt diese Typunterschiede.

\subsection{Die Majoranapaare bleiben als lokale CAR-Komposite erhalten}\label{sec:carpairdictionary}
Die Sektorunterscheidung lässt einen konstruktiven gemeinsamen Anschluss zu.
Die geraden Paaroperatoren liegen im gemeinsamen $D_8$-Vakuummodul $O$.
Für die hier auftretenden ganzzahligen Impulse $\lambda_i\in\{-1,0,1\}$ setze
\[
\chi_i^\pm=V_{\pm e_i},\qquad
P_\lambda=:\!\prod_{i:\lambda_i\ne0}\chi_i^{\lambda_i}\!:
\]
in wachsender Koordinatenfolge. Die ursprünglichen E8- und CAR-Kokzyklen auf $D_8$
sind durch die im Programm angegebene quadratische Eins-Kokette $\eta$ verbunden.
Mit den ursprünglichen Stromphasen $p_a$ lautet der erklärte Abgleich
\[
c_{ab}=p_ap_b\epsilon_{E_8}(\alpha_a,\alpha_b)\eta(\alpha_a+\alpha_b).
\]
Er gilt auf dem gesamten gemeinsamen $D_8$-Gitter; eine Ausdehnung auf lokale
halbzahlige CAR-Felder wird daraus nicht behauptet. Die Operatorabbildung ist
\begin{align}
 \Phi(B_{aa})&=c_{aa}P_{2\alpha_a},\\
 \Phi(B_{ab})&=c_{ab}:H_{\alpha_a-\alpha_b}P_{\alpha_a+\alpha_b}:\quad(a\ne b).
\end{align}
Die Diagonalzustände enthalten acht CAR-Komponenten und haben Normquadrat eins.
Die außerdiagonalen Zustände enthalten sechs CAR-Komponenten sowie einen
Stromoszillator auf den beiden übrigen Richtungen und haben Normquadrat zwei.
Alle besitzen chirales Gewicht vier, $X=10$, $Y=0$ und gerade CAR-Parität.

Auf dem behaltenen Sechserraum werden die beiden gemeinsamen Antworten
\eqref{eq:twojointpairresponses} erhalten. Der volle Zehnerraum besitzt dagegen
keine unzulässig angenommene Invarianz: $B_{00}$ hat $R_0=-12$;
bei $B_{0i}$ betragen der normierte Erwartungswert $-4$ und die Varianz $32$.
Hier entsteht unter $R_0$ eine zusätzliche Oszillatorrichtung außerhalb des
ursprünglichen Zehnerraums. Das Programm bewahrt diese Restwirkung.

Dieser Anschluss gilt für die Paare, nicht für die einzelnen ursprünglichen
Materiefelder. Jedes lokale NS-CAR-Wort beliebigen Grades besitzt ganzzahligen
Gitterimpuls $p\in\Z^8$ und deshalb
\[
 X=-2\sum_{i=1}^5p_i\in2\Z.
\]
Die ursprüngliche Halbspinorfamilie trägt dagegen die ungeraden $X$-Werte
$5,1,-3$. Weitere NS-Nachfahren können diesen Unterschied nicht beseitigen.
Das gemeinsame Wörterbuch enthält daher lokale CAR-Felder, C/X-Ramond-Intertwiner
und ihre wieder lokalen geraden Paaroperatoren. Diese Sektoren müssen beim
Raumzeitanschluss gemeinsam erhalten bleiben~\cite{boylesmith,barronfermion}.

\picturetext{Zwei passende Materie-Defekte lassen sich zu einem lokalen Paar
zusammensetzen. Für dieses Paar gibt es jetzt eine ausdrückliche Darstellung
mit den vorhandenen Fermionen der Quelle. Das ersetzt den Einzelschritt zum
physikalischen Teilchen nicht, bewahrt aber bereits dessen benötigten Massensektor.}

\subsection{Die vollständige Neutrinomatrix besitzt ein norm- und phasentreues Quellenwörterbuch}\label{sec:neutrinodictionary}
Die vorhandenen Neutrinoherleitungen liefern ganze komplexe Matrizen.
Ihr Quellenanschluss darf sich deshalb nicht auf Massen, Spur oder Norm
beschränken. Der bereits konstruierte symmetrische Sechserraum hat genau
die dafür benötigte bilineare Struktur. Für eine komplex-symmetrische
Koeffizientenmatrix $A$ sei
\begin{equation}
 \Gamma(A)=\sum_a A_{aa}B_{aa}+\sum_{a<b}A_{ab}B_{ab}.
\end{equation}
Die 36 ursprünglichen affinen Paarungen ergeben
\begin{equation}\label{eq:pairisometry}
 \langle\Gamma(A),\Gamma(B)\rangle=\tr(A^\dagger B).
\end{equation}
Dabei haben diagonale Paare Normquadrat eins und außerdiagonale Paare
Normquadrat zwei. Die rechte Seite ist deshalb genau das volle
Frobenius-Skalarprodukt, einschließlich der komplexen Kreuzterme.

Auch das tatsächliche Auslesen ist berechnet. Mit $C(v)=\sum_a v_a C_a$ gilt
\begin{equation}\label{eq:pairtomography}
 \langle\Gamma(A)^\dagger(\infty)C(v)(z_1)C(v)(z_2)\rangle
 =(z_1-z_2)^2\,v^{\mathsf T}A^\dagger v.
\end{equation}
Die sechs kohärenten Einsetzungen
$e_1,e_2,e_3,e_1+e_2,e_1+e_3,e_2+e_3$ rekonstruieren daher jedes $A$.
Die gemessene Größe ist hierbei eine \emph{komplexe Amplitude mit Phasenreferenz}.
Sechs Betragsquadrate reichen nicht: Sie unterscheiden im Allgemeinen eine
Matrix und ihre komplex konjugierte Matrix nicht.

Der Familienbasiswechsel wird aus den vorhandenen Generatoren konstruiert.
Setze $F=(e_1,Me_1,M^2e_1)$ und $S=F^\dagger J_C$, wobei $J_C$ der zuvor
an den Originalströmen geprüfte Intertwiner ist. $S$ ist unitär und verbindet
den Quellenvektor $v$ mit der ursprünglichen zyklischen Darstellung $w=Sv$.
Die einzelne C-Wirkung behält ihr Minuszeichen gegenüber dem Familienumlauf;
im symmetrischen Paar hebt es sich auf. Ein kovarianter physischer
Massentensor $K$ zieht sich dagegen als $K_C=S^{\mathsf T}KS$ zurück.
Der richtige Zustandsanschluss ist deshalb
\begin{equation}\label{eq:massformencoding}
 \mathcal E(K)=\Gamma\!\left(\overline{S^{\mathsf T}KS}\right),\qquad
 \frac{\langle\mathcal E(K)^\dagger(\infty)C(v)(z_1)C(v)(z_2)\rangle}
 {(z_1-z_2)^2}=(Sv)^{\mathsf T}K(Sv).
\end{equation}
Die Konjugation ist wesentlich: $\Gamma$ ist komplexlinear auf
Zustandskoeffizienten, $\mathcal E$ dagegen antilinear auf Massentensoren.
Es gilt $\langle\mathcal E(K),\mathcal E(L)\rangle=\tr(L^\dagger K)$.
Ein untypisierter Ähnlichkeitstransform würde diese physische bilineare
Antwort nicht erhalten.

\paragraph{Schwere und leichte Neutrinos bleiben verschiedene physische Felder.}
Der konkrete D5-Singulettpaarkanal besitzt $X=+10$ und gehört zum
$\nu^c\nu^c$- beziehungsweise konjugierten schweren Majorana-Anschluss.
Sein direkter Familientensor ist die ursprüngliche schwere Form
\[
 M_R=M_{\rm scal}\,\diag(\varepsilon,2\varepsilon,3),
 \qquad \varepsilon=\phi_0^2/(2g_{\rm car}),
\]
in der erklärten Q$_+$-Eigenbasis. Auch sie wird durch
$\mathcal E(M_R)$ samt allen sechs Rückleseamplituden dargestellt.
Für den dimensionslosen Tensor $D_M=M_R/M_3$ ergibt sich ausdrücklich
\[
 D_M=\diag(\varepsilon/3,2\varepsilon/3,1),\qquad
 \norm{\mathcal E(D_M)}^2=1+5\varepsilon^2/9.
\]
Das ist die direkte schwere Quellenübersetzung. Die anschließende bereits
verwendete normierte Seesaw-Rechnung lautet
\[
 K_\nu=-\widehat Y^{\mathsf T}D_M^{-1}\widehat Y,\qquad
 \widehat Y=i\sqrt{D_M}\,R_{\rm CI}\sqrt{m_{\rm niedrig}}\,U^\dagger.
\]
Hier hat $\widehat Y$ in der erklärten gemeinsamen RG-Normierung die Einheit
$\sqrt{\rm eV}$, während $K_\nu$ in eV steht. $\widehat Y$ ist daher nicht
mit dem dimensionslosen physischen Yukawa-Operator gleichzusetzen.
Die vorhandene Wahl der gemischten $\varepsilon$-Einsetzung bleibt dabei
ein ausgewiesener Ansatz; ihre Übereinstimmung mit dem physisch markierten
Quellenrahmen wird nicht allein durch einen Basiswechsel bewiesen.
Der leichte Massentensor entsteht anschließend über die vorhandene
Seesaw-Regel aus $M_R$ und $Y_\nu$.
Ein leichter $3\times3$-Tensor kann mathematisch in denselben Sechserraum
kodiert werden. Daraus folgt aber keine direkte Identifikation eines
linken Standardmodell-Neutrinopaares mit dem geladenen $X=10$-Quellenoperator.
Diese Feldtypgrenze wird in der folgenden Phasenprobe erhalten.

\paragraph{Die bereits vorhandenen gemeinsamen Neutrinobedingungen bleiben erhalten.}
Das Programm verwendet den ursprünglichen Operator $R_{\rm CI}(z)$ und
die eingefrorenen Eingaben von \file{nu_scalaron_v3.yaml}.
Die bestehende gleichzeitige Bedingung aus größtem Dirac-Singularwert und
$\widetilde m_1=m_3/10$ aus Abschnitt~\ref{sec:neutrinojoint} ergibt erneut
\[
 m_3=0{,}0435180861582\ {\rm eV},\qquad |z|=0{,}477382171162.
\]
Dies ist die bereits dokumentierte bedingte Konsistenzlösung, keine neue
Massenprognose. Insbesondere bleiben die Kalibrierung von $m_2$, der
Washout-Ansatz, die gewählte CI-Symmetrieklasse und der erklärte gemeinsame
Skalenfaktor erhalten; vollständige Schwellen-RG- und Flavor-Boltzmann-Dynamik
werden dadurch nicht zusätzlich geprüft.

\paragraph{Die verbleibende Phase wird nun zu einer konkreten Quellenantwort.}
Im ausdrücklich phasenoffenen Zweig mit $m_1=0$ verbleibt eine physische
Majorana-Phase. Mit dem ursprünglichen $U_0$ lautet diese Familie
\begin{equation}
 K_\nu(\gamma)=\overline{U_0}\,\diag(0,m_2e^{-i\gamma},m_3)\,U_0^\dagger.
\end{equation}
Die unveränderten leichten Massen, PMNS-Winkel und Dirac-CP-Invariante,
$YY^\dagger$, die ungeflavorten schweren CP-Invarianten sowie
$\norm{\mathcal E(K_\nu)}^2=m_2^2+m_3^2$ wählen $\gamma$ nicht aus.
Mit $v_j=S^\dagger u_j$, wobei $u_j$ die Spalten von $U_0$ sind,
liefern die ursprünglichen Paaramplituden dagegen
\begin{equation}\label{eq:majoranaphaseresponse}
 \frac{\mathcal A_2}{\mathcal A_3}
 =\frac{m_2}{m_3}e^{-i\gamma}.
\end{equation}
Die volle Quelle unterscheidet somit genau die relative Phase, die ihre
neutralen Normantworten ausblenden. Dieser Befund gilt für die im YAML
ausdrücklich als unbeschränkte Majorana-Phasen geführte Klasse.
Wird die gedruckte vollständige komplexe Yukawa-Matrix als unveränderlicher
Ansatz übernommen, ist $\gamma=0$ bereits durch diesen Ansatz festgelegt.
Eine Variation von $\gamma$ erhält dann nicht alle gedruckten Matrixeinträge.
Da $YY^\dagger$ unverändert bleibt, behebt die Phase auch nicht das zuvor
gefundene Defizit der ungeflavorten gemeinsamen Baryogeneserechnung.

\paragraph{Welche Symmetrie die Matrix tatsächlich auswählen darf.}
Die beiden ursprünglichen markierten Generatoren liefern einen zusätzlichen
exakten Test. In der zyklischen Darstellung seien sie $P$ und $Q$.
Sie erfüllen $P^3=Q^4=(PQ)^4=I$.
Für eine gemeinsame Charakterlinie des symmetrischen Paarraums folgen daraus
$\lambda^3=\mu^4=(\lambda\mu)^4=1$, also $\lambda=1$.
Eine $P$-invariante symmetrische Matrix besitzt die Form
\[
 A=\begin{pmatrix}a&b&b\\b&a&b\\b&b&a\end{pmatrix}.
\]
Die tatsächliche $Q$-Wirkung ersetzt sie durch
$\left(\begin{smallmatrix}a&b&-b\\b&a&-b\\-b&-b&a\end{smallmatrix}\right)$.
Proportionalität zu $A$ erzwingt $b=0$ und bei $A\ne0$ den Charakter $\mu=1$.
Die einzige gemeinsame Linie ist somit $A=aI$; in der ursprünglichen
C-Basis lautet sie
\[
 A_C=a\begin{pmatrix}1&0&0\\0&0&1\\0&1&0\end{pmatrix}.
\]
Ihr Rang ist drei, ihre drei Takagiwerte sind gleich.
Der vorhandene leichte Massentensor besitzt hingegen Rang zwei und
Takagiwerte $(0,m_2,m_3)$ mit $m_2\ne m_3$, für jede Majorana-Phase.
Eine zusätzlich geforderte vollständige Invarianz würde diesen vorhandenen
Tensor daher ausschließen.

Diese zusätzliche Invarianz steht nicht im ursprünglichen Neutrinovertrag:
NU.CI.01 verlangt die Zyklus- und Realitätsbedingungen an $R_{\rm CI}$;
die gedruckte schwere Matrix ist eine Generationenraum-/Q$_+$-Einsetzung.
Die passende weiterzuprüfende Quellenregel ist deshalb eine
\emph{kovariante} Abhängigkeit von den markierten Defekten $\mathcal D$,
\[
 K(g\mathcal D)=\rho(g)^{-\mathsf T}K(\mathcal D)\rho(g)^{-1},
\]
sofern die Quellenwirkung diese Kovarianz realisiert. Sie verlangt keinen
konstanten Fixvektor. Die Markierung kann die Familienrichtung unterscheiden;
die physische Higgs-/Masseneinsetzung muss ihre geladene Kopplung liefern.

\paragraph{Der jetzt genau bestimmte Auswahlgegenstand.}
Alle sechs Paare tragen $X=10$, während die neutralen Felder und die
Familien-Defekte diese D5-Ladung nicht kompensieren. Im ursprünglichen
$X$-invarianten Vakuum gilt folglich
\[
 \langle\mathcal E(K)\rangle=0.
\]
Eine Matrix als geladene Einsetzung zu berechnen ist damit etwas anderes
als ihre Auswahl durch einen physisch realisierten Hintergrund.
Der verschwindende Einpunktwert schließt einen eichinvarianten
Higgs-/Majorana-Mechanismus nicht aus. Der physische Anschluss kann einen
zugehörigen Higgsinsertionskanal oder einen insgesamt neutralen höheren
Korrelator verwenden; dann muss genau dessen ursprüngliche Abbildung
angegeben und ausgewertet werden.
Gesucht ist der tatsächlich erzeugte Koeffizient in \emph{diesem}
symmetrischen Sechserkanal, mit seiner Skala und Orientierung sowie dem
zugehörigen ladungskompensierenden Anschluss. Die vorhandene vollständige
Neutrinoherleitung legt fest, welchen Tensor dieser Anschluss reproduzieren muss.
\file{tfpt_explorer/source_neutrino_dictionary.py} berechnet nun seine
Isometrie, die ursprüngliche Generatorabbildung und die vollständige
phasensensitive Rücklesung.

\picturetext{Die Massentabelle lässt sich jetzt verlustfrei in die gemeinsame
Quelle übersetzen und wieder auslesen. Auch die Informationen, die beim
bloßen Wiegen der Zustände verloren gehen, bleiben erhalten. Der nächste
Herkunftsschritt muss erklären, warum der ursprüngliche Prozess genau
diese Tabelle in genau diesen Quellenzustand schreibt.}

\subsection{Ein gemeinsamer Transport für primitive Windung, Massentensor und neutrale Antwort}\label{sec:masssourceconnection}
Die Fortsetzung vom 29.~September arbeitet unmittelbar am vorher benannten
Anschluss. Die ursprüngliche Fuchsverbindung bestimmt nicht nur die
Familienendphasen. Ihr tatsächlicher Spin-Lift induziert auf den sechs
symmetrischen Paaren eine Verbindung, aus der sich die \emph{ganze}
ursprüngliche Verbindung wieder zurückgewinnen lässt. Auch die primitive
Determinantenwindung und die bereits berechnete neutrale RHP-Antwort lassen
sich an diesem selben Transport ablesen.

\paragraph{Die ursprüngliche Dreierwirkung muss mit ihrem Spin-Lift eingehen.}
Setze $d=\det\Psi$, $\Psi'=A\Psi$ und $t=\tr A$.
Die bereits an den tatsächlichen C-Strömen geprüfte Basisabbildung ist
$J=HP=\operatorname{diag}(1,-1,1)$, wobei $P$ die native Reihenfolge
$(23,13,12)$ in die Außenpotenzreihenfolge $(12,13,23)$ übersetzt und
$H$ die dreidimensionale Hodge-Abbildung ist. Die Identität
\[
 H(\Lambda^2\Psi)H^{-1}=d\,\Psi^{-\mathsf T}
\]
ergibt auf dem behaltenen C-Dreier den komplexifizierten internen Spintransport
\begin{equation}\label{eq:massspinconnection}
 S_C=d^{-1/2}P^{\mathsf T}(\Lambda^2\Psi)P
     =d^{1/2}J^{\mathsf T}\Psi^{-\mathsf T}J,
 \qquad
 B:=S_C'S_C^{-1}=\frac t2 I-J^{\mathsf T}A^{\mathsf T}J.
\end{equation}
Die Quadratwurzel wird entlang des ursprünglichen Wegs fortgesetzt.
Am unitären Ende ist das genau der schon nachgewiesene gewundene Spin-Lift;
die rohen holomorphen Fuchskoordinaten werden nicht als orthonormal erklärt.
Die drei Cartangewichte der tatsächlichen C-Ströme,
$(-1,1,1)/2$, $(1,-1,1)/2$, $(1,1,-1)/2$, prüfen insbesondere den
Faktor $t/2$ unabhängig von einer Endpunktmatrix.

\paragraph{Zustand und Massenkoeffizient besitzen entgegengesetzte Orientierung.}
Für einen symmetrischen Quellenkoeffizienten $X_0$ und eine duale symmetrische
Massenform $K_0$ lautet der faserweise Transport
\begin{equation}\label{eq:massdeterminantcarry}
 \begin{aligned}
 X(z)&=S_C(z)X_0S_C(z)^{\mathsf T},
 &\det X(z)&=d(z)\det X_0,\\
 K(z)&=S_C(z)^{-\mathsf T}K_0S_C(z)^{-1},
 &\det K(z)&=d(z)^{-1}\det K_0.
 \end{aligned}
\end{equation}
Denn $\det S_C=d^{1/2}$. Die beiden Vorzeichen sind daher keine
wählbaren Konventionen innerhalb derselben Zuordnung: Die Paarzustände und
ihre kontrahierenden Koeffizienten sind dual. Das skalare Matrixelement
$v^{\mathsf T}K v$ bleibt erhalten, wenn gleichzeitig $v\mapsto S_Cv$.

Für $\Psi(0)=I$ ist $d=1-z^4$. Auf dem ursprünglichen positiv orientierten
Kreis um $+1$ trägt die Determinantenlinie des invertierbaren Paarzustands
somit Windung $+1$, diejenige der dualen Massenform Windung $-1$.
Bei umgekehrter Kontur kehren sich beide Vorzeichen um. Eine konstante
Phasenänderung des ursprünglichen Basisrahmens ändert diese Windungen nicht.
Eine über der Punktur selbst singuläre, windende Rahmentransformation würde
dagegen die erklärte logarithmische Erweiterung ändern und ist hier kein
unbemerkter Basiswechsel.

Der ursprüngliche schwere Tensor ist für $\epsilon>0$ invertierbar:
\[
 K_0=M_R/M_3=\operatorname{diag}(\epsilon/3,2\epsilon/3,1),
 \qquad \det K_0=2\epsilon^2/9.
\]
Seine bereits berechnete native Einsetzung hat dieselbe nichtverschwindende
Determinante bis auf eine konstante Basisphase. Die Identität gilt für jede
korrekt basierte invertierbare Einsetzung und ersetzt nicht die noch zu
begründende physische Basiszuordnung. Der leichte Seesaw-Ausgang hat dagegen
Rang zwei und Determinante null; für ihn wird keine solche
Determinantenwindung behauptet.

\textbf{Zwei trennende Kontrollen.} Würde man den Faktor $d^{-1/2}$ im
tatsächlichen Außenpotenz-Lift weglassen, entstünde für $\det X$ Windung
vier. Würde man die C-Felder direkt mit dem ursprünglichen Vektordreier
$\Psi$ transportieren, entstünde Windung zwei. Der korrekte Spintransport
liefert Windung eins. Damit trifft die primitive Klasse die vollständige
Tensorwirkung und nicht nur eine Zahl am Endpunkt.

Die \emph{Determinantenlinie} schließt nach einem Punkturumlauf. Ein
allgemeiner voller Tensor wird dabei durch die Familienmonodromie verändert;
der unabhängige direkte Wegtest erhält ihn erst nach drei Umläufen zurück.
Die Gleichung setzt also keinen unveränderten Massentensor nach jeder
einzelnen Runde voraus. Ebenso wenig beweist sie schon einen
vierdimensionalen Vortex-Indexsatz oder die Gleichheit des vollständigen
P1-Nahtoperators mit diesem Familienoperator.

\paragraph{Die sechs Paarkanäle behalten die vollständige Verbindung.}
Die induzierte Verbindung auf $\operatorname{Sym}^2\mathbb C^3$ ist
\begin{equation}\label{eq:sixconnection}
 \mathcal L(X)=BX+XB^{\mathsf T}.
\end{equation}
Aus ihrer Wirkung auf $E_{ii}$ liest man $2B_{ii}$ auf dem diagonalen
Eintrag und $B_{ki}$ auf dem Eintrag $(k,i)$ ab. Damit sind sämtliche neun
komplexen Einträge von $B$ bekannt. Es folgt
\[
 \boxed{A=(\tr B)I-JB^{\mathsf T}J^{\mathsf T}.}
\]
Die Abbildung der Verbindung ist daher injektiv. Eine beliebige
$6\times6$-Matrix liegt nicht notwendig in ihrem neundimensionalen Bild;
die Implementierung prüft die vollständige induzierte Form beim Rücklesen.

Für beliebige, auch nichtkommutierende $3\times3$-Matrizen $A,D$ gilt
\begin{equation}\label{eq:pairtraces}
 \Tr_6\mathcal L_A=2\tr A,\qquad
 \Tr_6(\mathcal L_A\mathcal L_D)=5\tr(AD)-(\tr A)(\tr D).
\end{equation}
Dies folgt aus
$\Tr_{\mathrm{Sym}^2}(B\otimes I+I\otimes B)=4\tr B$
und der entsprechenden bilinearen Spuridentität
$5\tr(BC)+\tr B\tr C$. Die Rechnung gilt für ganze Operatoren;
ein Vergleich ihrer Eigenwertmengen wäre dafür nicht ausreichend.

\paragraph{Die neutrale Antwort ist eine Kontraktion genau dieser Verbindung.}
Im festgehaltenen ursprünglichen Fuchsrahmen und mit derselben symmetrischen
RHP-Normalordnung wie in Abschnitt~\ref{sec:fullflavorRR} folgt exakt
\begin{align}
 \langle R(z)\rangle_F
 &=\frac25\left[\Tr_6\mathcal L(z)^2-(\Tr_6\mathcal L(z))^2\right],\label{eq:massRconnection}\\
 \langle R(z)R(w)\rangle_{F,c}
 &=\frac{48}{(z-w)^4}
 +\frac{16}{5(z-w)^2}\left[\Tr_6\bigl(\mathcal L(z)\mathcal L(w)\bigr)
 +\frac32\Tr_6\mathcal L(z)\Tr_6\mathcal L(w)\right].\label{eq:massRRconnection}
\end{align}
Die rechte Seite stimmt mit dem bereits berechneten vollständigen
rationalen Flavor-Defektkern überein. Es wird keine zweite neutrale Quelle
eingesetzt. Die Spuren einer Verbindungsform sind bei beliebigen lokalen
Rahmenwechseln nicht für sich invariant; deren Ward-/Normalordnungsterme
müssen dann ebenfalls transformiert werden. Die Gleichungen gelten im
erklärten gemeinsamen Originalrahmen, nicht als neue eichinvariante
vierdimensionale Wirkung.

\paragraph{Auch der nötige lokale Ableitungsterm lässt sich aus derselben Quelle ergänzen.}
Eine faserweise Darstellung allein berechtigt noch nicht dazu, gewöhnliche
Ableitungen von ortsabhängig gedrehten Feldern einzusetzen. Hier ist die
Korrektur konkret. In einer Spaltenkonvention für die Feldkomponenten sei
\[
 C(z+s)C(z)^{\mathsf T}=s\mathcal V+s^2\mathcal W+\cdots,
 \quad \mathcal V^{\mathsf T}=-\mathcal V,
 \quad \mathcal B_{\rm raw}=\tfrac12(\mathcal W+\mathcal W^{\mathsf T}).
\]
Der konstante Term fehlt wegen der tatsächlichen Nullrelation
$:C_iC_j:=0$. $\mathcal V$ ist der bereits vorhandene antisymmetrische
Grad-drei-Kanal. Für eine Feldverbindung $b$ liefert das Zurücktransportieren
des ersten Felds mit $I-sb+O(s^2)$ vor dem Zusammenführen
\begin{equation}\label{eq:covariantpairjet}
 \boxed{\mathcal B_{\rm cov}
 =\mathcal B_{\rm raw}-\tfrac12(b\mathcal V-\mathcal V b^{\mathsf T}).}
\end{equation}
Bei $C'=GC$ und $b'=G'G^{-1}+GbG^{-1}$ entsteht im rohen symmetrischen
Koeffizienten der Zusatz
$\tfrac12(G'\mathcal V G^{\mathsf T}-G\mathcal V G'^{\mathsf T})$.
Gleichung~\eqref{eq:covariantpairjet} hebt ihn exakt auf:
$\mathcal B'_{\rm cov}=G\mathcal B_{\rm cov}G^{\mathsf T}$.
Die Spaltenkonvention der Feldkomponenten ist dabei von den dualen
Koeffizienten zur Kombination von Zuständen zu unterscheiden.
Wegen des verschwindenden konstanten OPE-Terms wird an dieser Stelle kein
zweiter Verbindungsjet benötigt. Im festen gewöhnlichen Ableitungsrahmen
bleibt der Grad-drei-Beitrag aber real vorhanden. Eine reine Kürzung auf
sechs Grad-vier-PBW-Zustände wäre somit nicht unter lokalen Rahmenwechseln
abgeschlossen.

\begin{center}
\begin{tikzpicture}[node distance=5mm]
\node[box,text width=12.6cm] (a) {\textbf{Eine ursprüngliche Verbindung $A$}\\
Vier Marken, vollständige Residuen, festgehaltene Erweiterung};
\node[box,below=of a,text width=12.6cm] (b) {\textbf{Ihr Spintransport und die sechs Paarkanäle}\\
$B=\tfrac12(\tr A)I-J^{\mathsf T}A^{\mathsf T}J$,
\quad $\mathcal L(X)=BX+XB^{\mathsf T}$};
\node[box,below=of b,text width=12.6cm] (c) {\textbf{Drei verbundene Antworten}\\
Primitive Windung \quad vollständiger Tensortransport \quad neutrale $R$-Antwort};
\node[box,below=of c,text width=12.6cm] (d) {\textbf{Physische Auswahl noch herzuleiten}\\
Ursprüngliche Rand-/Higgsbedingung für $X_0$ und ursprüngliche Nahtkopplung};
\draw[arr] (a)--(b);\draw[arr] (b)--(c);\draw[dasharr] (c)--(d);
\end{tikzpicture}
\end{center}

\picturetext{Ein gemeinsames Getriebe bewegt jetzt nachweislich die
Orientierung der Massentabelle und die neutrale Anzeige. Sein mitgeführtes
Vorzeichen sorgt dafür, dass aus einer ursprünglichen Runde auch genau
eine Runde der Paar-Determinante wird. Welche Massentabelle in dieses
Getriebe eingesetzt wird, ist eine andere Frage: Die Transportregel
bewegt jede Anfangstabelle korrekt. Sie wählt deren Werte noch nicht aus.}

\textbf{Bedeutung für die fundamentale Auswahl.}
Die Verbindung zwischen P1-artiger Windung, ursprünglichem Familienweg,
symmetrischem Massensektor und neutraler Antwort ist damit als gemeinsames
internes Transportgesetz geschlossen. Offen bleibt die physische
Ursprungsabbildung: Die homogene Gleichung
$X'=BX+XB^{\mathsf T}$ besitzt für \emph{jede} Anfangsmatrix eine Lösung.
Sie kann daher weder die schwere Textur noch deren Skala auswählen.
Auch rohe Änderungen der Matrixsingulärwerte durch einen nichtunitären
Rahmen sind keine erzeugte Massenhierarchie: Die gleichzeitig transportierte
positive Metrik erhält die normierten Größen.
Gesucht ist nun die ursprüngliche Randbedingung beziehungsweise die
tatsächlich erzeugte geladene Higgsantwort in diesem Kanal. Im letzteren
Fall müsste ihre Quelle in
$X'-BX-XB^{\mathsf T}=\mathcal J_H$
aus der ursprünglichen Wirkung berechnet werden; $\mathcal J_H$ wird hier
nicht zur gewünschten Antwort passend definiert. Dieselbe Wirkung muss
die weiterhin bedingte Nahtkopplung $c_3^2|R\rangle\langle p_0|$ und die
physische Zeitzuordnung tragen.

\file{tfpt_explorer/source_mass_transport.py} implementiert die induzierte
Verbindung, ihre vollständige Umkehrung, den endlichen Transport, die
kovariante lokale Ergänzung und die gemeinsamen Antwortformeln.
Die direkte ursprüngliche $3\times3$-Wegintegration und eine separat
integrierte $6\times6$-Gleichung stimmen auf beiden Umlauforientierungen
mit einem relativen Fehler unter $4\cdot10^{-12}$ überein.
Diese numerische Kontrolle ergänzt die ausgeschriebenen algebraischen
Identitäten; sie ist kein physischer Herkunfts- oder Kondensationsbeweis.

\subsection{Die Wand-Lesart: gemeinsamer geometrischer Anschluss und korrigierte Neutrino-Weiche}\label{sec:wallaudit}
\key{\textbf{Die Wand-Idee liefert einen bedingten geometrischen Anschluss, aber keinen Zwang zu Dirac-Neutrinos.} Die volle ursprüngliche Quelle enthält bereits Standardmodell-Singuletts mit $B-L=\pm2$. Ein spezielles Spinorbündel verbindet außerdem Familienindex und Anomalieklasse. Die entscheidende gemeinsame Auswahl betrifft die Verbindung, den physischen Zustand und die Kopplung; sie folgt nicht aus dem Fehlen eines Feldes in der ersten 248-dimensionalen Stufe.}

Das am 29.~September erneut gelesene Claude-Artefakt und der beigefügte Wand-Bericht werden hier gegen die Primärliteratur und den nativen TFPT-Code geprüft. Die dort genannten 27 externen Prüfungen und die nicht vorliegenden Commitinhalte werden nicht als eigene Ausführung verbucht. Die Aussage über die vollständige Quelle lässt sich dagegen unmittelbar im vorhandenen Quellprogramm entscheiden.

\paragraph{Was Hořava--Witten tatsächlich verbindet.}
Die Konstruktion beginnt mit elfdimensionaler Supergravitation und zwei zehndimensionalen Randkomponenten. Deren chirale Anomalie verlangt unter diesen Voraussetzungen 248 Vektormultipletts pro Wand; die vollständige Faktorisierung und das heterotische Wörterbuch führen zu $E_8$~\cite{horavawittenorigin,horavawittenwall}. Der relevante irreduzible Gravitationsanteil ist $\Tr R^6$ in einem Zwölfform-Anomaliepolynom. Eine Quartikidentität allein ersetzt diese Rechnung nicht. Die zweidimensionale Quellenanomalie $c_-=8$ und dieses zehn-dimensionale Inflow sind verschiedene Aussagen, solange ihre geometrische Abbildung fehlt.

Die bereits bekannte Verzweigung
\begin{equation}
248=(45,1)\oplus(1,15)\oplus(16,4)
\oplus(\overline{16},\overline4)\oplus(10,6)
\end{equation}
passt zur heterotischen Wahl eines $SU(4)$-Eichbündels mit Kommutant $Spin(10)$. Das Eich-$SU(4)$ wirkt auf Bündelfasern. Das isomorphe Tangential-$Spin(6)$ wirkt auf innere Raumrichtungen und deren Spinoren. Eine Gleichsetzung dieser Wirkungen ist eine zusätzliche, prüfbare Bündelwahl. Auch
\begin{equation}
d=10-2(8-g)=2g-6
\end{equation}
ist nur innerhalb der hinzugefügten zehn-dimensionalen Wand- und Tangentialeinbettung gültig. Sie leitet weder diese Wand noch die Einbettung aus P1/P2 her. Die vierdimensionale $\Z_{16}$-Anomalie aus Abschnitt~\ref{sec:spinliftsynthesis} darf dabei nicht unverändert für Kandidaten anderer Raumzeitdimension verwendet werden: Ihr Eingang ist bereits ein vierdimensionales fermionisches Symmetrieproblem.

Auch die Familienplätze müssen wirklich übersetzt werden. In einer heterotischen Entwicklung
\begin{equation}
\lambda^{aA}(x,y)=\sum_i\psi_i^a(x)\,\xi_i^A(y)
\end{equation}
bezeichnet $A=1,\ldots,4$ den $SU(4)$-Bündelindex, während $i$ die unabhängigen inneren Nullmoden und damit die vierdimensionalen Familien zählt. Die drei beibehaltenen TFPT-Plätze $C_1,C_2,C_3$ sind deshalb nicht bereits durch das gemeinsame Zeichen $(16,4)$ mit drei solchen Nullmoden identifiziert. Dafür braucht es einen Intertwiner, der insbesondere die ursprüngliche Familienmonodromie erhält.

Ebenso wenig folgt die Spin-Ladungsrelation automatisch für das volle masselose heterotische Spektrum. Ein ungebrochenes $\mathcal N=1$-chirales Multiplett enthält Skalar und Fermion in derselben Eichdarstellung. Ist $\omega$ nur das entsprechende Eichzentrum, wirkt $\omega^2$ auf beide gleich, $(-1)^F$ dagegen verschieden. Eine zusätzliche diagonale Symmetrie, Projektion oder SUSY-Brechung muss ausdrücklich definiert werden, bevor die Relation als physische Spektralbedingung verwendet wird.

\paragraph{Die stärkste richtige geometrische Identität.}
Sei $X$ eine geschlossene orientierte Spin-Sechsmannigfaltigkeit, und werde ausdrücklich $V=S_X^+$ gewählt. Für die formalen Krümmungswurzeln $x_1,x_2,x_3$ von $TX$ sind die vier Spinorgewichte
\begin{equation}
\frac12(x_1+x_2+x_3),\quad\frac12(x_1-x_2-x_3),\quad
\frac12(-x_1+x_2-x_3),\quad\frac12(-x_1-x_2+x_3).
\end{equation}
Die elementarsymmetrischen Polynome ergeben unmittelbar
\begin{equation}\label{eq:wallchern}
c_1(S_X^+)=0,\qquad c_2(S_X^+)=-\frac12p_1(TX),\qquad
c_3(S_X^+)=e(TX),
\end{equation}
mit umgekehrtem Vorzeichen des letzten Terms bei vertauschter Chiralität. Da $\widehat A$ keinen Term vom Grad sechs und $c_1(V)=0$ hat, folgt
\begin{equation}\label{eq:wallindex}
\operatorname{ind}D_{S_X^+}
=\int_X\widehat A(TX)\operatorname{ch}(S_X^+)
=\frac12\int_X c_3(S_X^+)=\frac{\chi(X)}2.
\end{equation}
Diese spezielle Eulerformel ist also korrekt. Für ein allgemeines heterotisches $SU(4)$-Bündel gilt dagegen $\operatorname{ind}D_V=\tfrac12\int_Xc_3(V)$; es ist keine Eulerformel für $TX$. Ein Index zählt Nettochiralität, nicht automatisch die Abwesenheit zusätzlicher vektorartiger Paare. Ein konkreter Gegenvergleich ist das heterotische Modell von Braun et al.: Dort hat $X$ die Hodgezahlen $(3,3)$, also $\chi(X)=0$, während $c_3(\widetilde V)=\pm54$ auf der neunfachen Überlagerung drei Nettofamilien auf $X$ liefert~\cite{braunheterotic}.

Auf einer Calabi--Yau-Dreifaltigkeit wird auch $c_2(S_X^+)=c_2(TX)$: Aus $c_1(TX)=0$ folgt $p_1(TX)=-2c_2(TX)$. Damit passen Familienindex und die topologische Bianchi-Bedingung
\begin{equation}
c_2(TX)-c_2(V)-c_2(V_{\rm hidden})=[W]
\end{equation}
mit trivialem verborgenem Bündel und $[W]=0$ zusammen. Dies löst die \emph{kohomologische} Bedingung; eine lokale Lösung der gesamten Feldgleichungen folgt daraus noch nicht.

Die Verbindung bleibt entscheidend: Bei Calabi--Yau-Holonomie $SU(3)$ zerfällt $S_X^+$ als $\mathbf1\oplus\mathbf3$, bis auf Konjugation. Die Levi-Civita-Verbindung in diesem Bündel lässt $E_6$ statt $Spin(10)$ übrig. Das wurde zusätzlich an den ursprünglichen 240 E8-Wurzeln nachgerechnet: Orthogonalität zum tatsächlichen $A_2$ liefert 72 Wurzeln plus sechs Cartanrichtungen, also Dimension 78; für das volle $A_3$ sind es 40 plus fünf, also Dimension 45. Eine gleiche abstrakte Strukturgruppe legt ihre Verbindung und Holonomie nicht fest.

\paragraph{Ein konkreter Anschluss, ohne eine neue Innengeometrie zu erfinden.}
Braun--Candelas--Davies--Donagi verwenden einen bekannten Raum mit
$(h^{1,1},h^{2,1})=(1,4)$, $\chi=-6$ und $\pi_1=\Z_{12}$. Ihr Abschnitt~3 untersucht Deformationen von $TX\oplus\mathcal O^{\oplus r}$; Anhang~B skizziert die mögliche Rang-vier-Konstruktion als Hilfsargument; das ausgearbeitete Modell verwendet Rang fünf. Die Topologie bleibt dabei erhalten~\cite{braundeformation}. Für die Wand-Lesart ist deshalb der präzise Kandidat eine stabile irreduzible, mit dem Quotienten verträgliche $SU(4)$-Deformation von $TX\oplus\mathcal O$. Sie könnte $c_2$ und den Drei-Familien-Index erhalten und zugleich den Kommutanten verkleinern. Ihre tatsächliche Verbindung, Nullmoden, Wilsonlinien und Auswahl aus TFPT müssen gemeinsam bestimmt werden. Volle Tangentialholonomie $Spin(6)$ darf dabei nicht zugleich als Calabi--Yau-Eigenschaft gefordert werden.

Die Deformationsrichtungen liegen unter den üblichen Calabi--Yau-Voraussetzungen in
\begin{equation}
H^1(\operatorname{End}(TX\oplus\mathcal O))=
H^1(\operatorname{End}TX)\oplus H^1(TX)\oplus H^1(T^*X).
\end{equation}
Die beiden letzten Summanden mischen den Dreier- und den Singulettanteil. Integrabilität, Stabilität und Äquivarianz sind zusätzliche Bedingungen, keine Folgen bloß nichtverschwindender Kohomologie. Eine solche Bündeldeformation wird physikalisch als Higgs-Hintergrund gelesen; sie ist noch nicht mit dem nativen symmetrischen $B_{ab}$-Kanal identifiziert. Genau hier wäre eine konkrete Gleichsetzung von Wirkungen statt einer Dimensionszählung erforderlich.

Dafür existiert ein konkreter mathematischer Fortsetzungssatz: Li--Yau zeigen für eine gegebene kleine holomorphe Deformationsfamilie einer Summe zweier stabiler Grad-null-Bündel, dass bei beidseitig nichtverschwindenden gemischten Kodaira--Spencer-Komponenten eine passende Hermitian--Yang--Mills-Metrik fortgesetzt werden kann. Ihr Satz~4.3 behandelt unter weiteren Voraussetzungen das gekoppelte Strominger-System~\cite{liyautorsion}. Das ist ein vorhandenes Konstruktionswerkzeug für die bedingte Wand-Realisierung; die Deformationsfamilie selbst und ihr TFPT-Ursprung sind weiterhin Eingang.

\paragraph{Die Herleitung eines Ausschlusses aus der 248 ist widerlegt.}
Wir verwenden die bereits originale Normierung
\begin{equation}
Q_{\rm em}=T_3+Y,\qquad X=5(B-L)-4Y,\qquad \omega=i^X.
\end{equation}
Für die linkshändigen $\nu^c$ gilt $(X,Y,B-L)=(5,0,1)$. Die zwei exemplarischen nativen Zustände
\begin{align}
B_{11}&=C_1(-3)C_1(-1)\Omega,& \norm{B_{11}}^2&=1,\\
B_{12}&=\bigl(C_1(-3)C_2(-1)+C_2(-3)C_1(-1)\bigr)\Omega,
&\norm{B_{12}}^2&=2
\end{align}
haben $X=10$, $Y=0$, $B-L=2$. Alle acht positiven und negativen Wurzelgeneratoren von $SU(3)_c\times SU(2)_L$ annihilieren sie im einfachen unitären Quotienten; auch ihre Cartan-Gewichte verschwinden. Dieselbe Rechnung wurde für alle zehn symmetrischen Paare und alle zehn adjungierten Gegenladungen ausgeführt. Sie sind damit wirkliche SM-Singuletts der Quelle, keine bloß passend benannten Vektoren.

Die Spin-Ladungsrelation lässt sie zu:
\begin{equation}
X=\pm10\equiv2\pmod4,\qquad \omega=-1,\qquad
\omega^2=+1=(-1)^F\quad\text{für einen Skalar}.
\end{equation}
Sie verbietet einen \emph{ungerade} geladenen Skalar wie den zuvor geprüften $\overline{16}_H$ unter dem zusätzlichen Spin-$\Z_4$-Postulat. Sie verbietet keinen bosonischen Ladung-zwei-Paarkanal. Die chiral bosonische Quelle allein beweist weiterhin keine vierdimensionale Skalarrealisierung.

Auch die niedrigste mögliche Stufe lässt sich ohne endliche Abschneideannahme bestimmen. Ein SM-Singulettmomentum mit $Y=0$ hat in den ersten fünf E8-Koordinaten denselben Wert $t$. Aus $X=-10t=10$ folgt $t=-1$. Die übrigen drei Koordinaten sind dann ganzzahlig und haben wegen der geraden E8-Gesamtsumme eine ungerade Summe. Also gilt
\begin{equation}
\norm{p}^2\ge5+1=6,\qquad h=\frac{\norm p^2}{2}+N\ge3.
\end{equation}
Die antisymmetrischen Gewicht-drei-Zustände erreichen diese Schranke; die symmetrische Familien-$10$ tritt auf Gewicht vier auf, wie in Abschnitt~\ref{sec:sourcepairs} bewiesen. Das Fehlen in Gewicht eins ist richtig, seine Übertragung auf die vollständige Quelle ist falsch.

\paragraph{Was über Dirac oder Majorana wirklich entscheidet.}
Bei einer \emph{exakt erhaltenen} schützenden Rest-$\Z_4$ lauten die Auswahlregeln für hyperladungsneutrale Operatoren:
\begin{center}
\begin{tabular}{lrl}
\toprule Operator & $X$ & Rest-$\Z_4$\\\midrule
$L H_u\nu^c$ & $0$ & erlaubt\\
$\nu^c\nu^c$ & $10$ & verboten\\
$(L H_u)^2$ & $-10$ & verboten\\
$(\nu^c)^4$ & $20$ & erlaubt\\
$(u^cd^cd^c)^2$ & $-10$ & verboten\\
$B^\dagger_{ab}\nu^c_a\nu^c_b$ & $0$ & erlaubt\\\bottomrule
\end{tabular}
\end{center}
Hier bedeutet erlaubt ausschließlich: von dieser Symmetrie nicht ausgeschlossen. Das bestimmt weder eine vierdimensionale Kopplung noch eine Zerfallsrate. Die elektroschwache Higgsphase erfordert die Eichkompensation
\begin{equation}
\widehat\omega=\omega e^{2\pi iY}
=e^{5\pi i(B-L)/2},\qquad \widehat\omega H_u=H_u.
\end{equation}
Auf hyperladungsneutralen Operatoren bleibt die gezeigte Auswahlregel gleich. Die globale Gruppenordnung auf quarkischen Bruchladungen hängt zusätzlich von den Eichzentrumsquotienten ab; aus dieser Formel allein wird keine neue globale Gruppe behauptet.

Heeck--Rodejohann enthalten ausdrücklich einen Skalar $\chi$ mit $B-L=-2$ und einen Skalar $\phi$ mit Ladung vier. Ihr Dirac-Zweig verlangt $\langle\chi\rangle=0$ bei nichtverschwindendem $\langle\phi\rangle$~\cite{heeckquadruple}. Gerade dieses zitierte Beispiel widerlegt die Gleichsetzung von Feldabwesenheit und Dirac-Schutz. Ein nachgewiesener Neutron--Antineutron-Übergang würde den hier beschriebenen exakten Restgruppen-Zweig ausschließen, nicht jede denkbare Theorie mit Dirac-Neutrinos.

\paragraph{Warum der vorhandene Familienprozess die Phase noch nicht auswählt.}
Alle zwölf ursprünglichen $A_3$-Wurzeln besitzen $X=0$; seine Cartanoperatoren sind ebenfalls neutral. Der bisherige Familientransport und das neutrale $R$ erhalten diese Ladung. Im invarianten Quellvakuum folgt deshalb exakt
\begin{equation}
\langle B_{ab}\rangle_\Omega
=\langle\omega B_{ab}\omega^{-1}\rangle_\Omega
=-\langle B_{ab}\rangle_\Omega=0,
\end{equation}
während $\langle B^\dagger B\rangle$ positiv ist. Die volle Transportgleichung aus Abschnitt~\ref{sec:masssourceconnection} erhält entsprechend die Null-Lösung. Sie kann einen eingesetzten Massentensor transportieren, aber seine Erzeugung nicht aus einer homogenen Gleichung ableiten. Die Wahl einer anderen Phase oder einer geladenen Randpräparation benötigt die zugehörige ursprüngliche Wirkung. In einer Eichtheorie ist ein geladener Einpunktwert zudem nur nach Eichfixierung ein direktes Ordnungsparameterbild; eine eichinvariante Higgsdiagnose verwendet passende Korrelationen. Aus dem Vakuum-Einpunktwert wird hier kein Verbot einer Higgsphase gefolgert.

Ein neutraler Skalaron ist vor Symmetriebrechung ebenfalls nicht dasselbe Feld wie ein komplexer Skalar mit $B-L=2$. Eine mögliche spätere Verbindung wäre dessen neutrale Radialmode in einer Darstellung $\sigma=(v+s)e^{ia/v}/\sqrt2$ samt hergeleiteter Mischung mit dem Skalaron. Der Ledger bezeichnet $\sigma=\text{Skalaron}$ als bedingt. Die entsprechende Zeile in \file{verification/v254_inner_fluctuations.py} prüft wörtlich den Wahrheitswert \texttt{True}; der frische Lauf mit fünf grünen Prüfungen beweist diese Identifikation daher nicht.

\paragraph{Die gemeinsame Bedingung für den weiteren Lösungsweg.}
Die Wand-Lesart darf die vorhandenen $\alpha$-, Massen- und Flavorherleitungen nicht durch freie Kompaktifizierungsparameter ersetzen. Schon dimensionsanalytisch ist $g_{10}^2$ von der Dimension Länge hoch sechs, während $c_3=1/(8\pi)$ dimensionslos ist. In einer unverwobenen Produktkompaktifizierung mit konstantem Hintergrund gilt schematisch $g_4^2=g_{10}^2/\operatorname{Vol}(X)$; Normierung, Schwellen und Warp-/Dilatonkorrekturen sind zusätzlich zu bestimmen. Der gleiche Gruppenname liefert diesen Zahlenanschluss nicht.

Der konsistente positive Auftrag ist deshalb enger und zusammenhängend: Aus dem ursprünglichen markierten Prozess muss ein gemeinsames physisches Wirkungs- und Zustandsfunktional hervorgehen. Seine Lösung muss den bereits vorhandenen Masseneinsatz im erlaubten geladenen Kanal tragen, den ursprünglichen Familientransport erhalten und zugleich den neutralen Nahtport reproduzieren. Eine heterotische Bündelrealisierung wäre daran zu messen: an der tatsächlichen Verbindung und ihrem Spektrum, an ihrem physischen Masseneinsatz und an den gleichen Kopplungen. Zuse-/Hypergraph-Regeln und der Universalraum bleiben mögliche Prozess- beziehungsweise Vergröberungsdarstellungen dieses Funktionals; sie wählen es nicht allein durch ihre Darstellung aus.

\picturetext{Die Wand ist eine mögliche Bauform für dasselbe Instrument. Im vollständigen Instrument ist das für Majorana-Massen benötigte Bauteil bereits vorhanden. Der Familienweg beschreibt, wie es mitbewegt wird. Gesucht ist jetzt die aus den Ursprungsregeln folgende Betriebsweise, die dieses Bauteil tatsächlich aktiviert und dabei gleichzeitig die schon berechneten Massen, Mischungen und die Feinstrukturkonstante liefert.}

\subsection{Vom Quellenfeld zum wirklichen Massenblock und zur gemeinsamen Vergröberung}\label{sec:masszeromode}
\key{\textbf{Der innere Operatoranschluss ist jetzt explizit geschlossen.}
Die zuvor konstruierten Paarfelder erzeugen durch ihre tatsächliche Modenwirkung
genau die symmetrischen Matrixeinträge des ursprünglichen schweren
Neutrinoblocks. Dieselbe Einsetzung liefert nach Eliminierung der schweren
Zustände die vorhandene Seesaw-Matrix, einen Determinantenbeitrag und eine
vollständige Gedächtnisantwort. Die Wahl der Einsetzung und ihr physischer
Ursprung bleiben dabei die benannten Voraussetzungen.}

\textbf{Die entscheidende Prüfung betrifft die Wirkung.}
Ein Zustand mit passenden Ladungen genügt noch nicht als Massenoperator.
Für ein homogenes Quellenfeld $v$ vom chiralen Gewicht $h$ bezeichnet
$o(v)=v_{(h-1)}$ seine gradbewahrende Mode~\cite{dongradical}.
Aus der ursprünglichen Stromalgebra folgt für
$u_{ab}=C_a(-3)C_b(-1)\Omega$
\begin{align}
Y(u_{ab},z)&=\frac12:\partial^2 C_a(z)C_b(z):,\\
o(u_{ab})&=\frac12\sum_{m\in\Z}(m+1)(m+2)
 :C_a(m)C_b(-m):.
\end{align}
Auf dem gesamten Stromraum $V_1$ bleiben genau die Terme $m=0,1$ übrig:
\begin{equation}\label{eq:massmodeexact}
\boxed{o(u_{ab})\big|_{V_1}=C_b(0)C_a(0)+3C_b(-1)C_a(1).}
\end{equation}
Der Term $m=-1$ besitzt Koeffizient null; die übrigen Terme verschwinden
durch Grad und Vakuumbedingung. Die Formel ist damit keine frei entworfene
Matrix, sondern eine Folge des vorhandenen Feldoperators.

Setze $|a\rangle=C_a(-1)\Omega$ und
$|\bar a\rangle=C_a^\dagger(-1)\Omega$. Mit den ursprünglichen
Gitterphasen ergibt die Rechnung für alle vier Familienplätze
\begin{align}
o(B_{aa})\big|_{V_1}&=|a\rangle\langle\bar a|,\\
o(B_{ab})\big|_{V_1}&=|a\rangle\langle\bar b|
                    +|b\rangle\langle\bar a|\quad(a<b),\\
o(B_{ab}^\dagger)\big|_{V_1}&=
  \bigl(o(B_{ab})\big|_{V_1}\bigr)^\dagger.
\end{align}
Diese Gleichheiten wurden als vollständige PBW-Ausgaben auf allen
248 Basisrichtungen geprüft, einschließlich der Cartanrichtungen.
Zusätzlich erklärt die Ladung das gesamte Stützgebiet: $X(B)=10$ kann im
$V_1$-Spektrum nur $X=-5$ mit $X=+5$ verbinden. Diese beiden Räume sind
genau die vier $C^\dagger$- beziehungsweise $C$-Zustände.
Für die sechs Paarfelder mit Indizes $1,2,3$ verbleibt deshalb ein
sechsdimensionaler Wirkungsraum; die übrigen 242 Richtungen werden annihiliert.
Auf höheren Quellengraden darf diese endliche Aussage nicht fortgeschrieben
werden. Die unteren Feldmoden erzeugen dort weitere Zustände.

Die Gewicht-zwei-Paarfelder haben auf $V_1$ die Wirkung null.
Die Gewicht-drei-Felder liefern stattdessen
$|a\rangle\langle\bar b|-|b\rangle\langle\bar a|$.
Damit ist die Symmetrie der Familienmatrix tatsächlich entscheidend:
Ein lokaler Majorana-Massenterm für gleichartige Weylfelder benötigt eine
symmetrische Familienform. Die Gewicht-vier-Felder liefern genau diese
Form. Die triviale Wirkung auf dem Vakuum-Top-Level $V_0$ widerspricht
dem nicht; die hier berechnete Wirkung liegt auf $V_1$, also jenseits der
gewöhnlichen Zhu-Algebra $A_0$~\cite{barronhigherzhu}.

\textbf{Der vollständige Block und die richtige Konjugation.}
Sei $K=K^{\mathsf T}$ eine schwere Massenform in der nativen Dreierbasis.
Der bereits benutzte Zustandskodierer lautet $\Gamma(\overline K)$.
In der Reihenfolge $(C^\dagger,C)$ folgt nun als Operatorgleichheit
\begin{equation}\label{eq:sourcemajoranaoperator}
D_K=o(\Gamma(\overline K))+o(\Gamma(\overline K))^\dagger
=\begin{pmatrix}0&K\\K^\dagger&0\end{pmatrix},\qquad
\Tr D_K^2=2\Tr K^\dagger K=2\norm{\Gamma(\overline K)}^2.
\end{equation}
Das Spektrum besteht aus den positiven und negativen Singularwerten von
$K$. Der Block ist selbstadjungiert, ungerade bezüglich der vorhandenen
Teilchen-/Antiteilchen-Chiralitätsgraduierung und verträglich mit der
reellen Struktur aus Tausch und komplexer Konjugation. Diese endlichen
Eigenschaften beweisen noch keine Raumzeitstatistik.

Die ursprüngliche Familienabbildung $S$ wird dabei beibehalten:
$K_C=S^{\mathsf T}K_{\rm original}S$ und
$W=\diag(\overline S,S)$ geben
\begin{equation}
W D_{K_C}W^\dagger=D_{K_{\rm original}}.
\end{equation}
Ein direkter Vergleich mit \file{verification/v252_full_finite_triple.py}
setzt diesen Quellenblock in den vorhandenen 96-dimensionalen Operator
ein. Er ersetzt exakt dessen Majoranablock; die gleichzeitig vorhandenen
komplexen Dirac-/Yukawablöcke bleiben gleich. Für den bisherigen bedingten
schweren Ansatz entsteht somit genau
$K_{\rm original}=M_R/M_3=\diag(\varepsilon/3,2\varepsilon/3,1)$.
Die Werte dieser Einsetzung werden hierbei aus dem Original übernommen.

Zwei unterschiedliche Wirkungen des neutralen Feldes müssen auseinanderbleiben:
\begin{equation}
R_0\Gamma(K)=4\Gamma(K),\qquad
R_0\big|_{\Span(C_{1,2,3},C_{1,2,3}^\dagger)}=I_6,
\qquad [R_0,D_K]\big|_{V_1}=0.
\end{equation}
Ein Gewicht-zwei-Feld wirkt nicht wie ein gewöhnlicher
Ladungsableitungsoperator. Das vertraute Vergleichsbeispiel ist
$L_0B=4B$, aber $[L_0,B_{(3)}]=0$.
Für gemischte Paare $B_{0i}$ gilt der letzte Kommutator gerade nicht,
weil der ausgelassene Platz $C_0$ den $R_0$-Eigenwert $-3$ trägt.
Auf dem behaltenen Block gilt außerdem
$\Tr(R_0D_K^2)=\Tr D_K^2$.
Diese neutrale Antwort misst die Gesamtnorm, wählt aber keine
Hierarchie oder relative Majoranaphase.

\textbf{Derselbe Block erzeugt Seesaw, Determinante und Gedächtnis.}
Die bestehende bedingte Neutrinorechnung liefert die schwere Skala
$m_*=M_3$ in eV und die normierte Matrix $\widehat Y$ mit Einheit
$\sqrt{\mathrm{eV}}$. Setze
\begin{equation}
H=D_{M_R/M_3},\qquad
\widehat V=\diag(\widehat Y^{\mathsf T},\widehat Y^\dagger),\qquad
z=E/m_*.
\end{equation}
Der ursprüngliche endliche Neutrinooperator besitzt nach der entsprechenden
Umsortierung die dimensionshomogene Form
\begin{equation}
\mathcal D_{\rm phys}=
\begin{pmatrix}
0&\sqrt{m_*}\,\widehat V\\
\sqrt{m_*}\,\widehat V^\dagger&m_*H
\end{pmatrix}.
\end{equation}
Alle Einträge tragen hier die Einheit eV. Die exakte Elimination des
schweren Blocks liefert
\begin{align}
\Sigma(z)&=\widehat V(zI-H)^{-1}\widehat V^\dagger,\\
\Sigma(0)&=\begin{pmatrix}0&K_\nu\\K_\nu^\dagger&0\end{pmatrix},
&K_\nu&=-\widehat Y^{\mathsf T}(M_R/M_3)^{-1}\widehat Y,\\
\det(EI-\mathcal D_{\rm phys})
&=\det(EI-m_*H)\,
\det\bigl(EI-\Sigma(E/m_*)\bigr).
\end{align}
Die Determinantenidentität gilt außerhalb der schweren Pole und setzt sich
als entsprechende algebraische Identität fort. Sie ist eine Aussage über
diesen endlichen quadratischen Block. Eine vollständige Majorana-
Pfaffianwirkung oder die ursprüngliche elektromagnetische relative
Determinante ist damit noch nicht identifiziert.

Ist $Hu_j=\lambda_j u_j$, so gilt mit positiven semidefiniten Residuen
\begin{equation}
\Sigma(z)=\sum_j\frac{\mathcal R_j}{z-\lambda_j},\qquad
\mathcal R_j=(\widehat V u_j)(\widehat V u_j)^\dagger\ge0,\qquad
\mathcal M(u)=\sum_j e^{-iu\lambda_j}\mathcal R_j.
\end{equation}
Die formale retardierte Fourierantwort bezüglich $z$ enthält
$-i\Theta(u)\mathcal M(u)$. Bei einer physischen Zeitdeutung wäre
$u=m_*t$ und die Zeitkern-Normierung entsprechend $m_*\mathcal M(m_*t)$;
eine solche physische Clockidentifikation wird hier nicht hergeleitet.
Außerdem erzeugen anfängliche schwere Zustände eine eigene inhomogene
Antwort, die beim Ausblenden mitgeführt werden muss.

Für den vorhandenen bedingten Neutrinozweig besitzt die Matrixantwort sechs
verschiedene Pole
\begin{equation}
\lambda_j=\pm1,\quad
\pm1.8848376655\cdot10^{-4},\quad
\pm9.4241883277\cdot10^{-5}.
\end{equation}
Alle sechs Residuen sind von null verschieden und von Rang eins. Somit sind
alle sechs Moden in dieser vollständigen Matrixantwort sichtbar; ihre
exakte endliche lineare Realisierung kann diese Pole nicht weglassen.
Das ist kein allgemeiner unterer Speicherbedarf eines beliebigen
Quanten- oder Universalraummodells. Die direkte Resolventenrechnung und
die getrennte Spektralsumme stimmen relativ besser als $5\cdot10^{-16}$
überein; die statische Quellenantwort stimmt mit der bestehenden leichten
Seesaw-Matrix bis auf $4.1\cdot10^{-17}$ eV überein.
Diese Zahlen prüfen die gemeinsame Implementierung, nicht eine neue
Messvorhersage.

\picturetext{Die schweren Zustände sind wie eine verborgene Mechanik hinter
einem sichtbaren Hebel. Bewegt man den Hebel sehr langsam, spürt man nur
ihren wirksamen Widerstand: Das ist hier die leichte Massenmatrix. Ändert
man die Anregung, antwortet die verborgene Mechanik mit ihren eigenen
Moden weiter. Genau diese Rückwirkung muss eine korrekte Vergröberung
speichern. Masse, Determinantenantwort und Gedächtnis entstehen in dieser
Rechnung aus derselben gekoppelten Struktur.}

\textbf{Was dadurch über die gesuchte gemeinsame Lösung feststeht.}
Der Majoranablock benötigt auf der Ebene dieser inneren Operatorwirkung
keine hinzugefügte fremde Quelle. Die Wahl einer festen geladenen
Einsetzung bleibt jedoch eine echte Wahl:
$\omega D_K\omega^{-1}=-D_K$. Ein physischer Massenterm verlangt daher
die passende Higgs-/Phasenrealisierung und dieselbe ursprüngliche
Wirkung, welche auch die Nahtkopplung auswählt. Die hier konstruierte
lineare Abbildung stellt \emph{jede} symmetrische Massenform dar; ihre
Existenz allein bevorzugt die gewünschte Form nicht.
Ebenso wenig bestimmt die allgemeine Schuridentität den besonderen
$\alpha$-Kanal mit Gewicht $5/4$ und Kopplung $c_3^2$.
Der nächste Herkunftsschritt muss deshalb am ursprünglichen gemeinsamen
Variationsproblem ansetzen: Es muss die bereits eingeschränkten
geladenen Einsetzungen, ihre Normierung und die primitive Nahtantwort
als zusammengehörigen stationären Prozess auswählen.

\subsection{Was die Raum- und Literaturbefunde für dieselbe Kette beitragen}
Am tatsächlichen 30-dimensionalen Blochoperator des markierten $A_3$-Netzes wurden die Bandbefunde unabhängig bestätigt. Auf dem zehn-dimensionalen Eigenraum zum Bandniveau $3$ hat die erste Richtungsableitung die charakteristischen Polynome
\begin{align}
\chi_{(1,0,0)}(\ell)&=\tfrac14\ell^4(\ell^2-6)(2\ell^2-9)^2,\\
\chi_{(1,1,1)}(\ell)&=\tfrac14\ell^4(\ell^2-36)(2\ell^2-9)^2.
\end{align}
Nach Normierung der zweiten Richtung folgen die positiven Steigungsquadrate $(6,9/2,9/2)$ und $(12,3/2,3/2)$. Diese Bandsteigungen sind ohne Feld- und Zeitabbildung keine gemessenen Materiegeschwindigkeiten. Der kompakte Würfelmodus mit alternierenden Amplituden erfüllt für alle $k$ exakt $L(k)\psi=3\psi$; nicht jeder Zustand auf seinen acht Knoten bleibt lokalisiert.

Der QCA-Anschluss von D'Ariano--Perinotti ist ein gezieltes Vorbild: Unter bereits vorausgesetzten fermionischen Feldern, lokaler unitärer linearer Entwicklung und diskreter Isotropie entstehen in der untersuchten minimalen Klasse BCC-Weyl-Automaten~\cite{darianojo}. Das reziproke Verhältnis von BCC und FCC/$A_3$ ersetzt die Quellenabbildung nicht. Für aus der Quelle herzuleitende Kantenoperatoren $A_h$ lautet der nächste notwendige Test
\begin{equation}
U(k)=\sum_h e^{ik\cdot h}A_h,\qquad
\sum_h A_h^\dagger A_h=I,\qquad
\sum_{h'-h=d}A_h^\dagger A_{h'}=0\quad(d\ne0).
\end{equation}
Der Graph liefert Adressen und Nachbarschaft. Die geladenen Operatoren samt ihren Phasen müssen die Aktualisierungsregel bestimmen. Hier können Zuses und Wolframs Prozessideen ansetzen, ohne die bestimmte TFPT-Quelle durch einen beliebig gewählten Automaten zu ersetzen.

Die weiteren Literaturhinweise liefern folgende Reichweiten:
\begin{itemize}
\item Beems Schur-Anschluss verlangt $k_{2d}=-k_{4d}/2$ und $c_{2d}=-12c_{4d}$. Die positive E8-Quelle mit $(k,c)=(1,8)$ ist daher nicht unmittelbar dieselbe Schur-Algebra einer unitären $4D$-$\mathcal N=2$-Theorie. Das ist kein allgemeines Verbot eines vierdimensionalen Anschlusses~\cite{beemchiral}.
\item Costello--Paquette enthält die Ausnahmegruppe E8. Die dortigen celestial Ströme $J[0,0]$ besitzen jedoch keinen zentralen Level-eins-Doppelpol. Ihr OPE ist nicht bereits das OPE unserer affinen Quelle~\cite{costellopaquette}.
\item Die Arbeit von Araki et al.\ vom August 2026 untersucht acht Majorana-Flavors auf einem zweidimensionalen euklidischen Gitter. Die Arbeit beweist keine $3+1$-dimensionale Spin(10)-Spiegelentkopplung. Hasenfratz--Xu unterscheidet Anomaliefreiheit von einer tatsächlich ausgewählten Wechselwirkung mit symmetrischem Gap~\cite{arakidomain,hasenfratzguide}.
\end{itemize}
Die 16-Leiter des Artefakts ist ein Netz bedingter Abbildungen. Halbspinordimension, Zahl der CAR-Felder und Ordnung der Anomaliegruppe sind keine austauschbaren Gegenstände. Auch Rochlins Signaturbedingung betrifft geschlossene glatte Spin-Viermannigfaltigkeiten; sie verbietet weder das E8-Gitter noch eine E8-Plumbing mit Rand.

\subsection{Das gemeinsame Feldwörterbuch und seine noch herzuleitende Auswahl}
Die Rechnung ersetzt eine Vermutung durch die exakten Gleichungen~\eqref{eq:nativespinlift}--\eqref{eq:fullsourcesix}. Die vorhandenen Flavor-, Alpha- und Massenherleitungen bleiben die Bedingungen an eine gemeinsame physische Realisierung. Diese muss gleichzeitig den ursprünglichen Windungstransport erhalten, die drei physischen Familien mit richtigen Ladungen und richtiger lokaler Statistik ausweisen, denselben zulässigen Neutrinopaar- und Higgsmechanismus reproduzieren und den neutralen Alpha-Port samt seiner Antwort im Flavor-Hintergrund koppeln.

Die Fortsetzung schließt dabei konkrete Übersetzungen: Die 48er-Zählung und
die Hyperladungsantwort werden aus denselben ursprünglichen Wurzeln gewonnen;
die vollständige komplexe Neutrinomatrix besitzt ein norm- und phasentreues
Quellenwörterbuch; ihre geraden Paarfelder bleiben unter Fermionisierung erhalten;
und der ursprüngliche Flavortransport bestimmt den ganzen neutralen
Zweipunktkern. Diese Verbindungen stammen aus derselben Quellenalgebra.
Eine Übereinstimmung bloßer Dimensionen oder Zahlen wird dafür nicht verwendet.

Abschnitt~\ref{sec:masssourceconnection} führt diesen Zusammenhang bis zum
gemeinsamen Transportgesetz weiter: Die Majoranapaar-Verbindung bewahrt die
vollständige ursprüngliche Verbindung, trägt die primitive
Determinantenwindung und liefert durch feste Spurkontraktionen dieselbe
neutrale Antwort. Die lokale Feldfortsetzung enthält den ausdrücklich
berechneten Grad-drei-Korrekturterm. Der noch fehlende physische Einsatz
ist dadurch genauer bestimmt; die homogene Transportgleichung selbst
wählt seine Anfangswerte nicht aus.

\key{\textbf{Der gemeinsame noch zu beweisende Schritt ist jetzt eine bestimmte Abbildung.}
Der ursprüngliche markierte Prozess muss seine Naht-, Higgs- und
Masseneinsetzungen in genau diese Quellenkanäle übersetzen.
Die Abbildung muss gleichzeitig den normierten Alpha-Rundlauf,
die vollständigen Massentensoren einschließlich Phasen und den
Flavor-Defektkern reproduzieren. Die vorhandenen physikalischen
Herleitungen geben ihre Zielwerte vor; das neue Wörterbuch legt fest,
welche Operatoren und gemeinsamen Antworten dabei erhalten werden müssen.}

Im Universalraum bedeutet dies eine antworttreue Übersetzung ganzer
Einsetzungsgeschichten. Zwei Geschichten gelten erst dann als gleich,
wenn auch ihre späteren gemeinsamen Antworten übereinstimmen.
Das neue Phasenbeispiel ist ein konkreter Fall: Gleiche Massen und gleiche
neutrale Normantwort bedeuten noch nicht denselben Zustand, weil eine
kohärente Paarantwort die relative Phase unterscheidet. Ein Zuse-Automat
oder Hypergraph kann diesen Prozess repräsentieren, sofern seine Regel
diese Operatoren, Ladungen und Antworten transportiert. Eine frei gewählte
Graphregel liefert ihre ursprüngliche Auswahl dadurch noch nicht.

Eine Antwort auf nur die Gleichung $\omega^2=(-1)^F$ genügt nicht, falls sie den verwendeten Massenmechanismus entfernt. Umgekehrt muss ein erfolgreicher Massenanschluss nicht die hier zusätzlich geprüfte Spin-$\Z_4$-Annahme verwenden. Die beiden Wege werden an denselben tatsächlichen Quellenfeldern unterschieden.

\file{tfpt_explorer/source_spin_lift.py} führt den internen Spinlift, die Originalwurzelcharaktere, die 16 Feldladungen und den gemeinsamen Quotiententest aus. \file{tfpt_explorer/source_majorana_pairs.py} prüft die tatsächlichen affinen Paarzustände, ihre Normen, Nullrelationen, höchsten Gewichte und den Anschluss an den ursprünglichen Flavortransport. Abschnitt~\ref{sec:neutralpairward} verbindet diese Paare mit der neutralen Quelle. Die Rechnung ist an den Explorer und die bestehende Flavor-Tour angeschlossen. Algebraische Aussagen, numerisch integrierter Weg und physische Zusatzannahmen bleiben getrennt.

\section{Geprüfter Anhang: Orientierung und mehrteilige Antworten erhalten}\label{sec:orientationaudit}

Der eingefügte Vorschlag vom 29.~September mit der Kennung \file{9df45383-826d-4f53-8687-bf615454c26c} enthält relevante Ergänzungen. Die folgenden algebraischen Aussagen wurden unabhängig neu berechnet; die angegebenen wissenschaftlichen Arbeiten wurden am Original geprüft. Die im Text verlinkten externen Rechnerpakete standen hier nicht zur Verfügung. Ihre QWZ-Tabellen und Ausführungsprotokolle werden deshalb nicht als eigene verifizierte Ergebnisse übernommen.

\subsection{Der C-Ladungsanschluss der Vierzustandsalgebra stimmt}
Die ursprünglichen vier C-Wurzeln erfüllen $(\alpha_a,\alpha_b)=1+\delta_{ab}$. Ihre isolierten Ladungsshifts antikommutieren bei verschiedenen Indizes. Die entsprechende Form über $\mathbb F_2$ ist $B=I+\mathbf1\mathbf1^{\mathsf T}$ und erfüllt $B^2=I$, also $\rank B=4$.

Nach Festlegung der zentralen Quadrate repräsentieren
\begin{equation}
\Gamma_0=X\otimes I,\quad \Gamma_1=Y\otimes I,\quad
\Gamma_2=Z\otimes X,\quad \Gamma_3=Z\otimes Y
\end{equation}
diese Relationen als $\mathrm{Cl}_4(\C)\cong M_4(\C)$. Die 16 geordneten Monome sind unabhängig. Die vollständige Enumeration ergibt 15 maximale kommutierende Pauli-Kontexte und 60 verschiedene Rang-eins-Projektoren. Die vier $\Gamma_a$ allein erzeugen 32 Gruppenelemente; das zusätzlich zugelassene $iI$ vervollständigt die 64-elementige Pauli-Gruppe.

Das ist ein echter algebraischer Anschluss an die ursprünglichen Ladungsvorzeichen. Es ist noch keine Identifikation der ganzen G31-, P2- oder Clockwirkung. Vor allem gilt im ursprünglichen Kokzyklus $U_a^2=-U_{2\alpha_a}$: Das Quadrat verschiebt weiterhin Ladung. Ein fester zentraler Charakter und sein vierdimensionaler Darstellungsraum ersetzen diese volle Ladungsquelle nicht.

\subsection{Das vorgeschlagene Carry ändert die kohärente Kompositionsregel}
Die algebraische Rechnung
\begin{equation}
\widetilde U_a=\Gamma_a\otimes U_a,\qquad
[\widetilde U_a,\widetilde U_b]=0\quad(a\ne b)
\end{equation}
ist richtig. Sie liefert jedoch \textbf{keine getreue Einbettung derselben ursprünglichen Operationen}.

\textbf{Beweis.} Angenommen, ein nichtnulliger Intertwiner $V$ erfülle $VU_a=\widetilde U_aV$ für alle $a$. Aus $U_aU_b=-U_bU_a$ folgt
\begin{equation}
\widetilde U_a\widetilde U_bV
=-\widetilde U_b\widetilde U_aV.
\end{equation}
Die rechten Operatoren kommutieren und sind invertierbar. Also ist $2\widetilde U_a\widetilde U_bV=0$ und damit $V=0$, ein Widerspruch.

Eine zusätzliche Präparation oder Auslese könnte bestimmte Antworten reproduzieren; sie müsste ausdrücklich konstruiert werden. Für bloße Konjugationskanäle verschwindet ein skalares Minus ohnehin, während kohärente Pfadsummen es unterscheiden können. Auch bleibt $\sum_a\alpha_a=(-2,-2,-2,-2,-2;0,0,0)\ne0$ trotz geschlossener Dreierprojektion. Der Anhang repariert damit ein ausgewähltes Vorzeichenproblem in einer veränderten Darstellung; er liefert noch keinen räumlichen Lift der vollständigen Quelle.

\picturetext{Zwei Minuszeichen ergeben ein Plus. Das kann eine neue Rechenregel vereinfachen. Wenn das ursprüngliche Minus aber eine beobachtbare Interferenz bestimmt, hat man damit die Regel verändert. Ein zusätzliches Register muss diese Antwort nachweislich wiedergeben können.}

\subsection{Ein exakter Zeuge gegen Kompression auf Paaransichten}
Setze auf drei Qubits
\begin{equation}
P=X\otimes X\otimes I,\quad Q=I\otimes Y\otimes Y,\quad
R=X\otimes Z\otimes Y,\qquad PQ=iR,
\end{equation}
und
\begin{equation}
\rho_\pm=\frac18\left(I+\frac P2+\frac Q2\pm\frac R2\right).
\end{equation}
Beide Zustände besitzen die Eigenwerte $(2\pm\sqrt3)/16$, jeweils vierfach. Ihre sämtlichen Ein- und Zweiteilreduktionen sind gleich:
\begin{equation}
\rho_{AB}=\tfrac14(I+XX/2),\qquad
\rho_{BC}=\tfrac14(I+YY/2),\qquad \rho_{AC}=I/4.
\end{equation}
Trotzdem liefert dieselbe zugelassene Paaroperation $e^{-itQ}$
\begin{equation}
\langle P\rangle_\pm(t)=\tfrac12\cos(2t)\pm\tfrac12\sin(2t).
\end{equation}
Die anfänglichen Steigungen sind $+1$ und $-1$. Gleiche Paaransichten und gleiche globale Spektren bestimmen somit nicht alle zukünftigen Paarantworten.

Auch die geordneten Zustandsantworten unterscheiden sich:
\begin{align}
J_\pm&=i\tr\rho_\pm[K_{AB},K_{BC}]
=\mp\frac{(\log3)^2}{4},
&K_A&=-\log\rho_A,\\
z_\pm&=\tr\rho_\pm\,\rho_{AB}\rho_{BC}
=\frac{12\pm i}{128},
&\frac{z_\pm}{|z_\pm|}&=\frac{12\pm i}{\sqrt{145}}.
\end{align}
Das ist ein vollständig endlicher Beweis der benötigten zusätzlichen Information. Er ist kein E8-Vakuum und keine Messung einer topologischen Zentralzahl.

\textbf{Anschluss an die Hauptfrage:} Der Universalraum muss zwei Geschichten getrennt halten, sobald ein gemeinsamer zulässiger Folgeeingriff sie unterscheiden kann. Dasselbe gilt in Abschnitt~\ref{sec:flavorcover}: Gleiche Punktpermutationen erlauben nicht das Weglassen des orientierten Linienlokalsystems. Die operative Antworttreue ist der gemeinsame Maßstab.

\subsection{Chirale Diagnostik: bestätigte Literatur, ausdrücklich bedingte Übertragung}
Für einen vollrangigen exakten Quanten-Markovzustand gilt
\begin{equation}
I(A:C\mid B)=0
\ \Longrightarrow\
K_{AB}+K_{BC}=K_B+K_{ABC}
\ \Longrightarrow\ J=0.
\end{equation}
Der letzte Schritt folgt durch zyklische Spur und $[\rho_R,K_R]=0$. Eine nichtverschwindende solche Antwort darf daher nicht durch vollständige Entfernung der betreffenden räumlichen bedingten Korrelationen ausgelöscht werden. Dies verbietet keine zeitliche Markov- oder GKLS-Dynamik.

Yangs Arbeit vom 8.~September 2026 untersucht die Robustheit unter gleichmäßig kontrollierten angenäherten Markovbedingungen und passenden Deformationen~\cite{yangmodular}. Eine Energielücke allein liefert diese Voraussetzungen nicht. Für stabile Replica-Phasen braucht man zusätzlich relative Fehlerkontrolle oder eine passende untere Schranke des komplexen Überlappungsbetrags.

Die geordnete Replica-Auslese
\begin{equation}
z_{a,b}=\tr\rho\,\rho_{AB}^a\rho_{BC}^b,\qquad
\omega_{a,b}=z_{a,b}/|z_{a,b}|
\end{equation}
besitzt in den jeweils begründeten Klassen den Grenzwert
\begin{equation}
\omega_{a,b}\longrightarrow
\exp\!\left[-\frac{\pi i c_-}{12}
\left(\frac a{a+1}+\frac b{b+1}-\frac{a+b}{a+b+1}\right)\right].
\end{equation}
Gass und Levin beweisen ihre Auswertung insbesondere für geeignete gapped freie Fermion-Grundzustände und String-net-Grundzustände; sie nennen Gegenbeispiele gegen eine uneingeschränkte Universalität~\cite{gasslevinreplica}. Die Replica-Verklebung von Sheffer et al.\ setzt geeignete große Regionen einer passenden $2+1$-dimensionalen Bulkphase voraus~\cite{shefferdefects}. Beide Resultate werden hier nicht unmittelbar auf vier Punkte des isolierten TFPT-Randnetzes übertragen.

Für eine zulässige Realisierung mit $c_-=8$ folgen korrekt
\begin{equation}
\boxed{\omega_{1,1}=e^{-2\pi i/9},\qquad
\omega_{1,2}=e^{-5\pi i/18}}
\qquad(-40^\circ,\ -50^\circ).
\end{equation}
Das sind bedingte Zielsignaturen. Die erste bestimmt $c_-$ nur modulo $72$, beide gemeinsam nur modulo $288$. Die beiden Winkel allein wählen E8 daher nicht eindeutig aus. Die Zweiintervall-Amplitude $3/2$ und diese orientierten Replica-Phasen haben unterschiedliche Gebietsgeometrien; ihre gemeinsame Realisierung muss nachgewiesen werden.

\subsection{Der stärkste neue Selektor verbindet Orientierung und gesamten Randinhalt}\label{sec:chiralselector}
Für dieselbe geeignete Bulk-/Randquelle sei der chirale Zusammenhang $J=\pi(c_L-c_R)/3$ begründet. Für die Vakuum-Intervallentropie ihrer Rand-CFT auf der Linie gilt bei festem Regulator- und Lorentzrahmen~\cite{castroentropy}
\begin{equation}
s=\frac{dS_{\rm Rand}}{d\log\ell}=\frac{c_L+c_R}{6}.
\end{equation}
Dann ergibt sich
\begin{equation}
c_L=3s+\frac{3J}{2\pi},\qquad
c_R=3s-\frac{3J}{2\pi},\qquad
\boxed{s\ge\frac{|J|}{2\pi}.}
\end{equation}
Die Schranke folgt aus Unitarität, $c_L,c_R\ge0$. Ihre Sättigung schließt gegenläufige gaplose CFT-Beiträge aus. Insbesondere
\begin{equation}
J=\frac{8\pi}{3},\qquad s=\frac43
\quad\Longrightarrow\quad (c_L,c_R)=(8,0).
\end{equation}
Mit tatsächlich vorhandenen unitären affinen E8-Strömen folgt dann aus $248k/(k+30)\le8$ und $k\in\Z_{>0}$ das bereits bekannte $k=1$.

Dieser gemeinsame Test ist eine nützliche Verschärfung der Herkunftsprüfung. Seine beiden Eingangswerte wurden hier jedoch nicht aus dem ursprünglichen P1/P2-Rohprozess berechnet. Einsetzen der bereits gewählten E8-Level-1-Zielquelle wäre eine Konsistenzprüfung. Außerdem schließt der Test keine beliebigen zusätzlichen gegappten Sektoren und keine sämtlichen anderen Bulk-Realisierungen aus.

Die ebenfalls genannte lokale Eichkonstruktion von Thorngren, Preskill und Fidkowski ist ein passender konstruktiver Literaturanschluss, aber keine Herleitung dreier Raumdimensionen oder des vollständigen TFPT-Standardmodells~\cite{thorngrendisentanglers}. Sie behält konkrete Voraussetzungen an lokale Schaltungen und eine geeignete gegappte Grenzfläche. Massen, Yukawa und Alpha bleiben die bereits vorhandenen gemeinsamen Anforderungen aus Abschnitt~\ref{sec:physicalconstraints}; sie werden durch diese neuen Orientierungstests nicht ersetzt.

\key{\textbf{Was aus dem Anhang übernommen wird:}
Der Rückanschluss der abstrakten Vierzustandsalgebra an die C-Ladungszeichen ist korrekt. Der Drei-Qubit-Zeuge beweist den Informationsverlust durch reine Paaransichten. Die orientierten Replica-Antworten und der gemeinsame $J/s$-Selektor liefern zusätzliche, bedingte Anforderungen an dieselbe Quelle. Die Carry-Vereinfachung wird ausdrücklich nicht als bereits gelöster Raumlift übernommen.}

\section{Die wichtigste neue Verbindung: Die ursprünglichen 24 Felder erzeugen eine Quelle}\label{sec:source}

\subsection{Die direkte Konstruktion}
Aus dem positiven geraden unimodularen $E_8$-Ladungsgitter und einem kompatiblen Kokzyklus entsteht in der erklärten Level-1-Gitterquantisierung
\begin{equation}
\mathcal H_{E_8}=\widehat{\bigoplus_{\lambda\in E_8}}\mathcal F_\lambda,
\qquad \Omega=\ket0,
\qquad L_0=\frac{\norm{\lambda}^2}{2}+\sum_{j=1}^8\sum_{n>0}nN_{j,n}.
\end{equation}
Die grundlegenden geladenen Vertexantworten haben die Form
\begin{equation}
\left\langle\prod_i V_{\alpha_i}(z_i)\right\rangle
=\delta_{\sum_i\alpha_i,0}\,\varepsilon(\alpha_1,\ldots,\alpha_n)
\prod_{i<j}(z_i-z_j)^{\langle\alpha_i,\alpha_j\rangle}.
\end{equation}
Ladungsneutralität, Kokzyklus und Gitterpaarung bestimmen hier gemeinsam die Antwort. Die Unimodularität begründet die Holomorphie der Gitter-VOA; sie muss nicht als zusätzliche freie Auswahl erneut eingeführt werden. Der Übergang einer geeigneten unitären stromerzeugten VOA zu einem stark lokalen konformen Netz ist durch entsprechende Voraussetzungen und Sätze abgesichert \cite{cklw}.

\key{Die zusätzliche Ursprungshypothese ist klar lokalisierbar: Die physische TFPT-Quelle ist diese minimale positive lokale Level-1-Vakuumquantisierung des ursprünglichen markierten Gitters. Innerhalb dieser Hypothese sind Vakuum, Ladungsfelder und konforme Zeit gemeinsam festgelegt.}

Die direkte Gitterroute ist bereits in den Originalen angelegt, insbesondere in \file{v459_seam_lattice_voa_route.py} und der Matter-Geometry-Route. Ein noch fehlender spezieller CAR-/QWZ-Skalierungsnachweis ist deshalb kein universelles Hindernis für jede Konstruktion der Quelle. Er bleibt erforderlich, wenn genau dieser mikroskopische Ursprung behauptet wird.

\subsection{$C_4+X_{20}$: Der kleine Erzeugersatz}
Die ursprünglichen Felder sind
\begin{equation}
C_a,\quad X_{ia},\qquad a=0,1,2,3,\quad i=0,1,2,3,4,
\end{equation}
zusammen mit ihren Adjungierten. Dabei transformiert $X$ als $\overline{\mathbf5}\otimes\mathbf4$, $C$ als $\mathbf1\otimes\mathbf4$. In der ursprünglichen Wurzelbasis ist P2 diagonal; in der nativen Prozessbasis muss derselbe Generator passend transformiert werden. Seine Wirkung ist
\begin{equation}
Q(Y)=-\sum_{ij}Y_{ij}S_{ji},\qquad Q(Y)|_X=-Y^T\otimes I_4,
\qquad Q(Y)|_C=0.
\end{equation}
Die Sitzung hat diese Zuordnung auf allen 24 ursprünglichen Feldern geprüft.

Die vierzig Wurzeln aus $X$ und $X^\dagger$ schließen auf 72 Wurzeln und acht Cartanrichtungen: die $80$-dimensionale $A_8$-Stromalgebra. Mit den acht Wurzeln aus $C$ und $C^\dagger$ wächst der Abschluss dagegen
\begin{equation}
48\longrightarrow140\longrightarrow240\quad\text{Wurzeln},
\qquad240+8=248.
\end{equation}
Es werden somit keine fremden Erzeuger nachträglich angefügt. Die ursprünglichen vier $C$-Felder liefern genau die Vervollständigung, die dem $X$-Block allein fehlt.

\begin{center}
\begin{tikzpicture}[node distance=9mm]
\node[box,text width=4cm] (x) {$X_{ia}$ und $X_{ia}^\dagger$\\20 Felder plus Adjungierte};
\node[box,right=of x,text width=3cm] (a8) {$A_8$\\72 Wurzeln + 8 Cartan};
\node[box,right=of a8,text width=3cm] (e8) {$E_8$\\240 Wurzeln + 8 Cartan};
\node[box,below=of a8,text width=3cm] (c) {$C_a,C_a^\dagger$\\ursprünglicher Viererblock};
\draw[arr] (x)--(a8);\draw[arr] (a8)--(e8);\draw[arr] (c)--(e8);
\end{tikzpicture}
\end{center}

\picturetext{Die zwanzig Felder erzeugen mit ihren Adjungierten bereits eine große geschlossene Instrumentengruppe. Vier weitere, schon ursprünglich vorhandene Stimmen verbinden sie mit den fehlenden Klangrichtungen. Danach ist das gesamte Instrument erreichbar.}

\subsection{Die $A_8$-Erweiterung und ihre Grenzen}
Mit ursprünglichen $X$-Wurzeln lässt sich eine einfache $A_8$-Basis konkret schreiben:
\begin{align}
\beta_i&=X_{i+1,0}-X_{i,0} &&(i=0,\ldots,3),\\
\beta_4&=-X_{4,0},\\
\beta_{5+a}&=X_{0,a}-X_{0,a+1} &&(a=0,1,2).
\end{align}
Hier bezeichnet $X$ in der Gleichung den Ladungsvektor des entsprechenden Feldes. Die Gramdeterminante ist $9$; die Smith-Normalform der Inklusion besitzt sieben Einträge $1$ und einen Eintrag $3$. Der \emph{Gitterquotient} ist deshalb $E_8/A_8\cong\Z_3$. Die Wurzeln zerfallen in $72+84+84$ und die Ströme in
\begin{equation}
\mathbf{248}=\mathbf{80}\oplus\mathbf{84}\oplus\overline{\mathbf{84}}.
\end{equation}
Für die verwendete Orientierung liegen die $C$-Wurzeln im $\Lambda^6$-Sektor. Die $SU(9)/\Z_3$-Stromalgebraerweiterung besitzt bekannte Vorläufer \cite{distler}; neu in dieser Sitzung ist ihr expliziter Anschluss an die ursprünglichen markierten TFPT-Felder und die ausführbare Antwortkette.

Der Gitterquotient $\Z_3$ ist nicht ohne Weiteres die Familienclock $\sigma_F$. Gleiche Gruppenordnung ersetzt keine identische Wirkung. Die noch zu prüfende Gleichsetzung muss auf denselben Feldern, Phasen und Zuständen formuliert werden.

\section{Code, Quartik und Rekursion: Was die endlichen Modelle wirklich zeigen}\label{sec:finite}

\subsection{Vom Binärcode zu den 60 Ereignissen}
Der verwendete binäre Code ist
\begin{equation}
C=RM(1,3)=[8,4,4],\qquad W_C(y)=1+14y^4+y^8.
\end{equation}
Unter der Gaussian-$\mu_4$-Wirkung werden die 240 reellen $E_8$-Wurzeln zu sechzig komplexen Strahlen. Ihre Projektoren lassen sich als gemeinsame Pauli-Eigenprojektoren darstellen. Die fünfzehn primitiven Kontexte enthalten je vier Strahlen. Andere orthogonale Gruppierungen erzeugen dadurch nicht automatisch neue primitive Kontexte.

Zu jedem Strahl gehört die Reflexion
\begin{equation}
r_\ell=I_4-2\ket{\psi_\ell}\!\bra{\psi_\ell},
\qquad \bar U_4=\frac1{60}\sum_\ell r_\ell^{\otimes4},
\qquad G_4=\frac35I-\bar U_4.
\end{equation}
Dabei ist $\bra{\psi}$ die adjungierte Zeile von $\ket\psi$. Das Spektrum von $\bar U_4$ lautet, mit Multiplizitäten als Hochzahlen,
\begin{equation}
(-1)^1,\quad(-1/5)^{45},\quad(1/15)^{135},\quad(1/5)^{70},\quad(3/5)^5.
\end{equation}
Der obere Fünferraum ist durch eine Lücke $2/5$ getrennt. Auf diesem Code gilt ein echter Intertwiner
\begin{equation}
r_\ell^{\otimes4}V=VT_{ab}.
\end{equation}
Sechzig Ereignisse werden auf fünfzehn logische Transpositionen abgebildet, mit vier Ereignissen pro Faser. Die Transpositionen, perfekten Matchings, Pauli-Kontexte und Hecke-Marken haben zwar jeweils fünfzehn Elemente, sind jedoch unterschiedliche Objekte. Nur die tatsächlich konstruierten Abbildungen verbinden sie.

\subsection{Warum die Quartik einen Fünfer erzeugt}
Für die konkrete Pauli-Gruppe der Ordnung 64 wurde der Molien-Ausdruck
\begin{equation}
M(t)=\frac{1-t^{16}}{(1-t^4)^5}
\end{equation}
ausgewertet. Die Dimensionen in den Graden $0,4,8,12,16$ sind $1,5,15,35,69$. Somit ist der Fünfer der erste nichtkonstante Pauli-invariante Polynomraum, \emph{wenn} der kleinste positive Grad als Auswahlregel verlangt wird. Die Auswahlregel selbst ist eine zusätzliche Annahme, solange ihre Herkunft nicht gezeigt ist.

Die projizierte Vierfachpräparation $V^\dagger x^{\otimes4}$ liefert sechs Nullsummenkoordinaten auf der Igusa-Quartik:
\begin{equation}
\sum_i z_i=0,\qquad
\left(\sum_i z_i^2\right)^2=4\sum_i z_i^4.
\end{equation}
Dies ist eine spezielle Präparationsfläche im logischen Fünfer. Ein beliebiger Zustand dieses Fünfers liegt nicht notwendig auf ihr. Die fünf Quartikpolynome der Abbildung sind Pauli-invariant und bilden zusammen einen $G_{31}$-Kovarianten. Die Igusa-Quartik ist deren quartische Relation in den Zielkoordinaten. Beides ist von den skalaren $G_{31}$-Invariantengraden $8,12,20,24$ zu unterscheiden.

Die Projektoren erfüllen $P=SQ_P$ und
\begin{equation}
\rank(S,Q_P,P,S-P)=(35,16,5,30).
\end{equation}
Paarablesungen erfassen zunächst zehn Operatorrichtungen. Mit den ausdrücklich zugelassenen aktiven Paarkontrollen wächst der Beobachtungsraum $10\to20\to25$. Vollständige Beobachtbarkeit ist hier ein Ergebnis der gesamten Kontrollfamilie; sie folgt nicht schon aus einer passiven Paarablesung.

\subsection{Der Vierertensor und die 60/40-Zerlegung}
Der symmetrische Delta-Tensor $A$ und der ursprüngliche Simplextensor $C$ liefern
\begin{equation}
E=A-2C,\qquad F_I=4C-A/3,\qquad G=C+A/6.
\end{equation}
Die normierten Richtungen $\widehat E$ und $\widehat F_I$ sind orthogonal, und
\begin{equation}
\widehat G=\sqrt{3/5}\,\widehat E+\sqrt{2/5}\,\widehat F_I.
\end{equation}
Die 60/40-Aussage bezeichnet also die quadrierten Normanteile dieser beiden Richtungen. Sie ist keine universelle Zerlegung physischer Energien. Als $125\times5$-Matrix gelesen liefert $G_{abcd}/\sqrt2$ den ursprünglichen Dreierencoder.

\subsection{Zwei unterschiedliche, jeweils sinnvolle Rekursionen}
Die ältere gemeinsame Ereignisbindung ist
\begin{equation}
k=I-\frac1{15}\sum_{ab}T_{ab}\otimes T_{ab}.
\end{equation}
Sie besitzt das Spektrum $0^1,(2/5)^5,(2/3)^9,(4/5)^{10}$ und den eindeutigen neutralen Zustand $\Omega_5$. Werden die beiden Ereignislabels unabhängig gezogen, entsteht stattdessen die skalare Energie $16/25$. Die gemeinsame Ereignisidentität ist deshalb eine reale Voraussetzung.

Das Dreieck $k_{12}+k_{23}+k_{31}$ besitzt einen Fünfergrundraum bei Energie $4/5$ und eine Lücke $2/5$. Das ist ein exakter rekursiver Codezweig. Das räumliche Inzidenznetz ist jedoch dreieckfrei. Deshalb darf derselbe Dreiecksencoder nicht ohne neue Identifikation zur lokalen Raumregel erklärt werden.

Der spätere ladungskovariante Pfad verwendet dagegen
\begin{equation}
h_{\rm cov}=\frac56(I-P_\Omega),\qquad
(Wx)_{abc}=\frac{\delta_{ab}x_c+x_a\delta_{bc}}{\sqrt{12}}.
\end{equation}
Er wirkt auf $U\otimes\overline U\otimes U$ und erfüllt
\begin{equation}
(X_1-X_2^T+X_3)W=WX.
\end{equation}
Sein Grundraum ist fünfdimensional, seine Grundenergie $2/3$ und seine Lücke $1/3$. Der Ladungsgenerator wird wirklich durch die Rekursion transportiert.

Vollständige kollektive native $S_6$-Symmetrie zusammen mit Erhaltung des tatsächlichen P2-Markers verbindet die $5$-, $9$- und $10$-Sektoren und erzwingt die Form
\begin{equation}
H=cI+\kappa(I-P_\Omega).
\end{equation}
Positivität und ein eindeutiger neutraler Nullzustand verlangen $\kappa>0$ nach entsprechender Offsetwahl. Der konkrete Wert $5/6$ benötigt zusätzlich die erhaltene ursprüngliche Spur. Eine mögliche Instrumentenoperation ist dabei nicht automatisch eine Symmetrie des Hamiltonoperators.

\subsection{Was bei der Blockkomposition erhalten bleibt}
Die positive Neun-Port-Vervollständigung erfüllt in beiden Orientierungen
\begin{equation}
H_9V=V(4I+h).
\end{equation}
Auf einem vorgegebenen endlichen Blockgraphen bleibt das Produkt der Codes daher unter einer Summe solcher Links invariant: $PHQ=0$. Die ursprünglichen einzelnen schwachen Restlinks haben dagegen Leakage. Die Neun-Port-Regel ist somit eine zusätzliche, ausdrücklich benannte Kompositionsregel. Wenn die Graphstruktur selbst dynamisch variiert, darf der Offset $4I$ nicht pauschal als bedeutungslos entfernt werden.

Der ältere Quartikencoder zerlegt sich gegenüber dem Pfadencoder als
\begin{equation}
V_G=\sqrt{2/3}\,W+\sqrt{1/3}\,R,
\qquad P_G(t)=\frac{5+4\cos t}{9}.
\end{equation}
Unter beliebigen kollektiven Ladungsoperationen wächst der gewählte Restfünfer auf den gesamten 70er-Sektor. $5+70=75$ ist deshalb der nächste notwendige Träger dieser Operationen. Weitere ursprüngliche Operationen erreichen auch den 45er-Sektor. Die vollständige Quelle muss diese Fortsetzungen bewahren; ein kleiner Restträger kann nicht allein wegen seiner Bequemlichkeit geschlossen erklärt werden.

\section{Die gemeinsame Zeit: Der alte und der neue Quellblock passen wirklich zusammen}

Die beiden Beschreibungen $D_5+A_3$ und $A_4+A_3+\mathfrak u(1)$ sind jetzt am gleichen achtfachen Heisenbergraum verbunden. Mit $J_5$ als Fünfermatrix aus lauter Einsen lauten die Projektoren
\begin{align}
P_{A_4}&=\diag(I_5-J_5/5,0_3),\\
P_{A_3}&=\diag(0_5,I_3),\\
P_{\mathfrak u(1)}&=\diag(J_5/5,0_3).
\end{align}
Sie summieren sich zu $I_8$, und $P_{A_4}+P_{\mathfrak u(1)}=P_{D_5}$. Deshalb ist der tatsächliche Virasorovektor
\begin{equation}
\omega=\frac12\sum_{j=1}^8h_j(-1)^2\Omega
\end{equation}
in beiden Zerlegungen gleich. Damit gilt auf dieser Quelle
\begin{equation}
T_{E_8}=T_{A_8}=T_{D_5}+T_{A_3}
=T_{A_4}+T_{A_3}+T_{\mathfrak u(1)}.
\end{equation}
Die Zentralzahlen erfüllen $8=5+3=4+3+1$, aber diese Zahlen sind nur eine Konsequenz des stärkeren Projektorbeweises. Gleich ist der Stress-Tensor samt Moden, nicht lediglich seine Zentralzahl.

\begin{center}
\begin{tikzpicture}
\node[box,text width=4.8cm] (left) {$D_5$\quad +\quad $A_3$\\$5$\quad +\quad $3$};
\node[box,right=18mm of left,text width=6.2cm] (right) {$A_4$\quad +\quad $\mathfrak u(1)$\quad +\quad $A_3$\\$4$\quad +\quad $1$\quad +\quad $3$};
\draw[arr] (left)--node[above,font=\small]{derselbe $T$}(right);
\node[below=7mm of right,align=center,font=\small] {Die acht Richtungen werden anders gruppiert.\\Die gemeinsame konforme Zeit bleibt gleich.};
\end{tikzpicture}
\end{center}

Auf den tatsächlichen ursprünglichen Feldern ergibt sich
\begin{equation}
h(X)=\frac25+\frac38+\frac9{40}=1,
\qquad h(C)=0+\frac38+\frac58=1.
\end{equation}
Die ergänzende normierte $\mathfrak u(1)$-Richtung ist
$(1/\sqrt5,\ldots,1/\sqrt5,0,0,0)$; ihre Ladungen sind $3/(2\sqrt5)$ auf $X$ und $-\sqrt5/2$ auf $C$. Sie steht orthogonal zur P2-Hyperladung. Ihre Benennung als $U(1)$ erlaubt daher keine Gleichsetzung mit Elektromagnetismus oder Hyperladung.

Die geometrische modulare beziehungsweise konforme Zeit ist innerhalb des Vakuumnetzes bestimmt. Ein interner Clock-Zeiger, eine diskrete Survivalmatrix, $L_0$, eine räumliche Translation und die Laborzeit sind trotzdem unterschiedliche Objekte. Die physische Realisierung muss ihre passende Verflechtung liefern.

\section{Das ausführbare Quellgesetz und die echte Hyperkante}\label{sec:ward}

\subsection{Eine Regel für alle endlichen Stromwörter}
Für $F_n=\omega(J_{A_1}(z_1)\cdots J_{A_n}(z_n))$ bestimmt die affine Ward-Rekursion
\begin{equation}\label{eq:ward}
F_n=\sum_{j=2}^n\frac{\kappa(A_1,A_j)}{(z_1-z_j)^2}
F_{n-2}^{\widehat1,\widehat j}
+\sum_{j=2}^n\frac{1}{z_1-z_j}F_{n-1}^{[1,j]},
\quad F_0=1,\quad F_1=0.
\end{equation}
Im zweiten Term wird $A_j$ durch $[A_1,A_j]$ ersetzt. Alle Zwischenströme bleiben in der wirklichen 248-dimensionalen Wurzel-/Cartanalgebra. Die Implementierung nutzt exakte rationale Positionen und rationale Arithmetik; sie ist keine Näherung durch Gaußpaarungen.

Die öffentlichen Funktionen \file{evaluate_source_word} und \file{connected_source_word} nehmen $C$-/ $X$-Felder, ihre Adjungierung und paarweise verschiedene endliche Positionen entgegen. Der Kern ist für beliebige endliche Wörter formuliert; die Webschnittstelle begrenzt die Wortlänge aus Laufzeitgründen auf acht. Diese Grenze ist keine Aussage über die Physik. Exakte Ergebnisse werden als Bruch erhalten; ein nicht darstellbarer Gleitkommawert wird nicht als gültige Zahl ausgegeben.

\subsection{Eine vollständig verbundene Sechserantwort}
Für die Folge
\begin{equation}
X_{00},\ X_{10}^{\dagger},\ X_{11},\ X_{21}^{\dagger},\ X_{22},\ X_{02}^{\dagger}
\end{equation}
ist die Gesamtladung null. Jede der $2^6-2=62$ echten nichtleeren Teilmengen besitzt dagegen eine nichtverschwindende Gesamtladung. Alle ihre Vakuumkorrelatoren verschwinden, ebenso alle fünfzehn vollständigen Zweierpaarungsprodukte. Die ganze Antwort ist trotzdem
\begin{equation}\label{eq:hex}
F_6=\frac{1}{(z_1-z_2)(z_1-z_6)(z_2-z_3)
(z_3-z_4)(z_4-z_5)(z_5-z_6)}.
\end{equation}
Bei $(-5,-2,-1,1,3,7)$ ist der Wert exakt $1/576$. Er wurde unter allen 720 gleichzeitigen Permutationen von Feldern und zugehörigen Positionen verglichen. Wird nur die erste Position auf $-4$ verschoben, ergibt sich $1/352$; bei $(-4,-1,0,2,5,9)$ ergibt sich $1/936$.

\begin{center}
\begin{tikzpicture}[scale=1.0]
\fill[teal!9] (90:1.9)--(30:1.9)--(-30:1.9)--(-90:1.9)--(-150:1.9)--(150:1.9)--cycle;
\foreach \ang/\lab in {90/1,30/2,-30/3,-90/4,-150/5,150/6}{\node[circle,draw=teal,fill=white,minimum size=7mm] (v\lab) at (\ang:1.9) {\lab};}
\node[align=center,font=\small] at (0,0) {ganze Antwort\\$1/576$};
\node[align=left,text width=7.2cm,right=20mm of v2] {\textbf{Alle sechs gemeinsam zählen.}\\Jede echte Teilgruppe ist ladungsungleichgewichtig und hat Antwort null. Die ganze Gruppe ist neutral und antwortet.};
\end{tikzpicture}
\end{center}

Mit
\begin{equation}
Z(t)=\omega\!\left(\prod_i(1+t_iJ_i(z_i))\right)
\end{equation}
ist der verbundene Beitrag der Koeffizient von $t_1\cdots t_n$ in $\log Z$. Weil hier alle echten Teilantworten verschwinden, ist die volle Sechserantwort bereits vollständig verbunden. Ein unabhängiger Vergleich über sämtliche 203 Mengenpartitionen eines anderen gemischten Sechserworts ergab volle Antwort $1/720$ und verbundenen Anteil $-1/8640$. Eine Cartan-Viererantwort liefert den passenden Gauß-Gegenfall: volle Antwort ungleich null, verbundener Anteil null.

\key{Diese Hyperkante ist ein tatsächlich berechneter Koeffizient derselben Quelle. Sie ist kein zusätzlich postuliertes Sechskörper-Hamiltontermchen und beweist keine nichtlokale Signalübertragung.}

\subsection{Die $W$-Rekursion ist eine besondere Antwort dieses Gesetzes}
An den harmonischen Marken $(-1,0,1,\infty)$ ergibt der vollständige $X$-Vierertensor
\begin{equation}
F_{ia,jb,kc,ld}=(\delta_{ij}\delta_{kl}+\delta_{il}\delta_{jk})
(\delta_{ab}\delta_{cd}+\delta_{ad}\delta_{bc}).
\end{equation}
Definiert man
\begin{equation}
(N_d)_{abc,x}=\delta_{ab}\delta_{cx}+\delta_{ax}\delta_{bc},
\qquad W_d=\frac{N_d}{\sqrt{2(d+1)}},
\end{equation}
so ist $W_d^\dagger W_d=I_d$. Nach der passenden Anordnung der Tensorfaktoren folgt
\begin{equation}
W_{\rm gemeinsam}=W_5\otimes W_4,
\qquad N_5^\dagger N_5=12I_5,\quad N_4^\dagger N_4=10I_4,
\quad N_{\rm gemeinsam}^\dagger N_{\rm gemeinsam}=120I_{20}.
\end{equation}
Der mittlere Faktor ist jeweils konjugiert. Alle 8000 Ausgabetripel der dünn besetzten Abbildung wurden geprüft. Bei festem Familienindex wird derselbe Vierertensor zu $2(\delta_{ij}\delta_{kl}+\delta_{il}\delta_{jk})$. Mit der tatsächlichen reflexionspositiven Paarung ist er $K=2I+10P_\Omega$, sodass $I-K/12=h_{\rm cov}$ folgt.

Damit verbindet sich die früher separat motivierte Bindung direkt mit der tatsächlichen Quellenantwort. Der Geltungsbereich bleibt jedoch konkret: Ein gemischtes Wort $C_0,X_{00}^\dagger,X_{01},C_1^\dagger$ liefert $1/2$ statt des naiven Wertes $1$ einer Erweiterung zu $W_6\otimes W_4$. Auch ein naiver einziger $W_{20}$ ersetzt das Produkt nicht.

Im Koaleszenzgrenzwert eines Sechserworts schließt die Antwort auf $(W\otimes\overline W)\Omega$. Bei endlichem Verhältnis $t=\epsilon/L$ ist die Leakage positiv und beginnt mit $t^2/3$. Die Sektoren $5_W+5_Z+70+45$ bleiben deshalb erforderlich. Die affine Gramrechnung zeigt beim Niveau $J_{-1}X$ die Eigenwerte $6^5,2^{45},0^{70}$; beim Niveau $J_{-2}X$ sind diese um eins erhöht. Die Quelle besitzt also einen Nachfahrenturm. Der kleine $W$-Code beschreibt einen kontrollierten Grenzsektor dieses Turms.

\subsection{KZ, Fusionsbäume und echte Komposition}
Die betrachtete $SU(5)$-Level-1-KZ-Verbindung ist flach, obwohl überlappende Residuen nicht miteinander kommutieren. Zwei klassische Tensorkoordinaten bedeuten dabei nicht zwei unabhängige physische Level-1-Fusionsblöcke. Die Partnersektoren des vollständigen $E_8$-Feldes kompensieren die gebrochenen Phasen des isolierten Faktors.

Die rationalen Hecke-Charakterphasen erhalten Wurzelklammer, Paarung, Ward-Regel und diagonalen Casimir. Damit ist ihr Anschluss an die vollständige Quelle real. Eine Hecke-Inklusion wird dadurch nicht zu einem KZ-Transport. Ebenso sind die $E_8$-Ströme Vakuumnachfahren und keine zusätzlich eingeführten adjungierten $E_8$-Level-1-Primärfelder. Analytische Aussagen zu Sewing und Faktorisierung benötigen die entsprechenden Konvergenzvoraussetzungen \cite{gui,guizhang}.

\section{Die Quellen- und Clockprüfungen: Welche Alternativen wirklich ausgeschlossen sind}

\subsection{Die frühe Kanalfamilie und die vollständige Normierung}
Eine frühe Phase der Sitzung betrachtete die endliche Familie $\Phi_b$, $0\le b\le2/9$. Ihr Populationsschatten bleibt gleich, während eine Kohärenz den Eigenwert
\begin{equation}
\lambda_{12}=2b-\frac19
\end{equation}
besitzt. Damit ist die reine Populationsauslese in dieser Klasse unzureichend, um den Kanal zu bestimmen. Die Familie ist kein Gegenmodell zu sämtlichen TFPT-Prinzipien, weil nicht alle ursprünglichen Markierungen und höheren Antworten in ihre Definition eingingen.

Der kohärente Lift erfüllte auf dem Code die Normierung, aber nicht auf dem ganzen 256-dimensionalen Raum:
\begin{equation}
q=\sum_\mu L_\mu^\dagger L_\mu,\qquad
q|_{\rm code}=I,\qquad q|_{\Lambda^4\C^4}=\frac{10}{9}I.
\end{equation}
Die funktionale Reparatur $B_\mu=L_\mu q^{-1/2}$ normiert den gesamten Lift und erhält die Code-Krauszweige. Sie wählt den Parameter $b$ nicht. Dass $q^{-1/2}$ zur endlichen erzeugten Algebra gehört, beweist auch keine einfache lokale physische Realisierung. Der frühere HTML-Gesamtbericht dokumentiert diesen Zwischenstand und darf nicht als letzte Aussage der Sitzung gelesen werden.

\subsection{Die ursprüngliche Clock ist ein konkreter markierter Operator}
Der v487-Survivaloperator lautet
\begin{align}
D&=\begin{pmatrix}1&0&0\\0&0&1\\0&1&0\end{pmatrix},\qquad P_0=\diag(1,0,0),\\
M&=\frac12I+\frac16D+\frac13P_0,\qquad \operatorname{spec}M=(1,2/3,1/3).
\end{align}
Hier ist $D$ der Decktausch von Ordnung zwei, während $\sigma_F$ die Familienmarkierung von Ordnung drei bezeichnet. Seine Eindeutigkeit gilt in der ursprünglichen treuen markierten Dreierklasse. Die Fortsetzung
$C=\frac12I+\frac16S_{13}+\frac13WW^\dagger$ auf den Trimer setzt zusätzliche Identifikationen voraus und bestimmt außerhalb des Seedsektors nicht automatisch alles.

Auch die beiden früher gleichgesetzten Fünferclocks sind verschieden. Der Riemann--Roch-Charakter ist $(5,1,1,1)$, der zyklische Codecharakter $(5,3,5,3)$. Der Hom-Raum ist neundimensional, sein maximaler Intertwinerrang aber nur drei. Es gibt also Intertwiner, jedoch keinen invertiblen markierten Intertwiner zwischen diesen beiden Fünferwirkungen. Ein gemeinsamer Siebenerträger repariert nur diese isolierten $C_4$-Wirkungen.

Der volle unmarkierte Cartan-$M_4$-Anschluss scheiterte an der ursprünglichen $\sigma_F$-Markierung: Die Fixalgebren $M_2\oplus\C\oplus\C$ und $M_2\oplus M_2$ besitzen Dimension sechs beziehungsweise acht. Die korrekte native Fixebene ist $\C^2$, mit zwölf Wurzeln und drei $\mu_4$-Strahlen. Dort ergibt $E_i=(2/3)P_i$ die ursprüngliche Trine.

Ihre Erwartungswertmatrix
\begin{equation}
B=\frac1{18}\begin{pmatrix}13&1&4\\1&13&4\\4&4&10\end{pmatrix}
\end{equation}
ist keine Übergangsmatrix aufeinanderfolgender gleicher Trinemessungen: Jede einzelne Ausgangswahrscheinlichkeit einer solchen Messung ist höchstens $2/3$, während $B_{11}=13/18$ größer ist. Dies ist ein kleiner, entscheidender Gegenfall gegen die falsche Interpretation.

\subsection{Die korrekt markierte Ereignisfaser}
Native Reflexionswörter realisieren die registrierte v976-Faser
\begin{equation}
w_t=(1/2+t,t,t,1/18-t,2/9-t,2/9-t),
\quad q=6t,\quad0\le t\le1/18.
\end{equation}
Die GNS-Antwort mit linken Einsetzungen ist $FAAF=q/9$. Die ältere Vorgabe $1/27$ wählt deshalb $q=1/3$, sofern ihre unabhängige Quellenidentifikation gilt. Dies ist eine konkrete Möglichkeit, wie eine zusätzliche ursprüngliche Antwort eine zuvor offene Faser auswählt.

Für jede Faser ist der stationäre reversible $S_3$-Random-Walk mit Gleichverteilung $1/6$ und Faltungstransfer $T_q$ bei allen endlichen Historientests reflexionspositiv. Diese Historienfunktion ist nicht automatisch die zuvor genannte linksinsertierte $M_2$-GNS-Antwort. Eine reversible kontinuierliche Einbettung mit nichtnegativen $S_3$-Ereignisraten existiert dagegen genau im geprüften Bereich $2/9\le q\le1/4$. Bei $q=1/3$ ist eine der Raten $\log(3/4)/6<0$. Ausgeschlossen ist somit diese spezielle klassische kontinuierliche Ereignisratenklasse. Eine allgemeine Quanten-Markov-Einbettung ist damit nicht widerlegt.

Die frühere $b$-Familie auf $M_4$ und diese spätere markierte $q$-Faser auf $\C^2$ sind nicht dieselbe Parametrisierung derselben vollständigen Quelle. Ihr jeweiliger Prüfbereich muss erhalten bleiben.

\subsection{Höhere Antworten unterscheiden sogar gleiche Wurzelzerfallsraten}
Die Phasenkanalprüfung betrachtete eine ausdrücklich beschränkte Klasse zufälliger unitärer Generatoren
\begin{equation}
\mathcal L_O(A)=\frac1{|O|}\sum_{m\in O}(U_m^\dagger A U_m-A).
\end{equation}
Die nichttrivialen Charakterorbits haben Größen $15,120,120$. Ein Deckorbit und ein Braidorbit liefern bei derselben gesamten Sprungrate auf allen 240 Wurzelströmen dieselbe Zerfallsrate $16/15$. Auf einem Grad-zwei-Vertex mit $\beta=(1+i)\alpha$ und $\norm\beta^2=4$ unterscheiden sie sich dagegen: Rate null beziehungsweise $16/15$.

Das ist ein exakter Unterschied in einer höheren Zukunftsantwort. Er bewahrt unter anderem die untersuchte native Symmetrie, Ladungsaddition und Vakuumkovarianz. Er erzwingt aber nicht schon die ursprüngliche primitive Ereignisauswahl, alle Clockantworten, sämtliche registrierten Mehrzeitwörter oder die physische Raumzeitabbildung. Die korrekte Aussage lautet deshalb: \emph{Nicht-Eindeutigkeit innerhalb dieser Phasenkanalklasse}. Daraus folgt weder eine Mehrdeutigkeit des festgelegten Gittervakuums noch eine Nicht-Eindeutigkeit nach sämtlichen TFPT-Constraints.

\section{Hecke, Double Cover und geladener Carry: Was bei einer Vergröberung mitwandern muss}\label{sec:carry}

\subsection{Die fünfzehn Marken sind an die ursprünglichen Ereignisse angeschlossen}
Mit $\pi=1+i$ gilt im verwendeten Gaussian-Gitter
\begin{equation}
L/\pi L\cong\mathbb F_2^4,\qquad
M_y=\{x:\overline h(x,y)=0\},\qquad
\chi_y(x)=(-1)^{\overline h(x,y)}.
\end{equation}
Alle sechzig nativen Ereignisse wurden tatsächlich auf die fünfzehn Marken abgebildet, einschließlich der Äquivarianzprüfung. Vier Ereignisse pro Marke bleiben dennoch verschiedene Operationen.

Der konkrete Deckcharakter mit Kennung 152 liefert ein $D_8$-Untergitter von Index zwei: 112 gerade Wurzeln und 128 ungerade Spinorwurzeln. $V_{D_8}$ zusammen mit seinem tatsächlichen Spinorkoset stellt dieselbe $E_8$-Quelle wieder her. Es wird kein unabhängiges neues Materiefeld erfunden.

Die Projektion auf den geraden Teil ist jedoch kein Liehomomorphismus. Für eine ungerade Wurzel gilt
\begin{equation}
P[e_\alpha,e_{-\alpha}]=[e_\alpha,e_{-\alpha}]
\ne0=[Pe_\alpha,Pe_{-\alpha}].
\end{equation}
Die zunächst entfernten Felder können gemeinsam eine erhaltene Cartanantwort erzeugen. Das ist genau der Mechanismus, wegen dessen eine heute unsichtbare Information für spätere Prozesse erhalten bleiben muss.

Beim ersten echten Untergitter sinkt der Rang der geerbten reduzierten Paarung von vier auf zwei. Das kanonische Dual stellt weitere Tests bereit; es erzeugt nicht automatisch zusätzliche Vertexfelder. Über die ersten Kinder ergeben sich die vorhandenen 256 Vorzeichencharaktere.

\subsection{Bei größerer Tiefe reichen Vorzeichen nicht aus}
Für eine ganzzahlige Inklusion $M=K\Z^8\subseteq L$ und einen Kindcharakter $c$ ergibt
\begin{equation}
\theta=K^{-T}c,\qquad
\chi_\theta(x)=e^{\pi i\theta\cdot x},\qquad
\chi_\theta(Kz)=(-1)^{c\cdot z}
\end{equation}
eine Fortsetzung auf derselben Quelle. Dabei können Phasenordnungen $2,2,4,4,8$ auftreten. Für Ordnung vier ist $(I+U)/2$ kein Projektor. Der richtige Sektorprojektor lautet bei Ordnung $N$
\begin{equation}
\Pi_r=\frac1N\sum_{k=0}^{N-1}e^{-2\pi i rk/N}U_\theta^k.
\end{equation}
Die kohärente Speicherung
\begin{equation}
V\psi=\sum_r\ket r\otimes\Pi_r\psi
\end{equation}
ist isometrisch. Für $\theta/2=\ell/N$ mit $\ell\in\Z^8$ ist das Sektorlabel $r(x)=\ell\cdot x\pmod N$. Wenn ein Feld die Ladung um $\alpha$ ändert, muss sein Verbraucher den Speicher gleichzeitig verschieben:
\begin{equation}
\widetilde E_\alpha=S_{r(\alpha)}\otimes E_\alpha,
\qquad\widetilde E_\alpha V=VE_\alpha.
\end{equation}
Ein gespeichertes Label ohne diesen Verbraucher wäre keine vollständige Rekursion. Wird der Speicher verworfen, entsteht die Dephasierung $\rho\mapsto\sum_r\Pi_r\rho\Pi_r$; die Information ist dann tatsächlich verloren.

\begin{center}
\begin{tikzpicture}[node distance=9mm]
\node[box,text width=3.1cm] (s) {voller Zustand\\Ladung + Phase};
\node[box,right=of s,text width=4.1cm] (c) {reduzierte Darstellung\\plus kohärenter Carry};
\node[box,right=of c,text width=3.1cm] (e) {geladene Operation\\ändert beide Teile};
\draw[arr] (s)--(c);\draw[arr] (c)--(e);
\node[below=7mm of c,align=center,font=\small,text width=13cm] {Carry ist die mitgeführte Information, auf die ein späterer Test reagieren kann.\\Er ist kein kostenloser Ersatz für Gedächtnis.};
\end{tikzpicture}
\end{center}

Nach $n$ echten Index-zwei-Inklusionen mit der geerbten Gitterform ist die arithmetische Bilanz
\begin{equation}
[L:M_n]=2^n,\qquad |M_n^\#/M_n|=4^n,\qquad |L/M_n|=2^n.
\end{equation}
Die ursprüngliche Quelle belegt nur einen bestimmten Teil der gesamten dualen Klassen. Ihre geladene Vervollständigung rekonstruiert die Quelle. Der Index $2^n$ ist keine abgeleitete Anzahl räumlicher Zellen.

\subsection{Der Lift muss die wirkliche Operatorwirkung erhalten}
Auf dem unverdrehten Cartanraum wird die native Reflexion durch $s_\alpha s_{J\alpha}$ realisiert. Die gleiche Pauli-/Heisenberg-Paarung eines verdrehten Moduls reicht dagegen nicht: Die native Reflexion hat eine $3+1$-Aufteilung mit $|\tr|^2=4$, während die untersuchten Pauli-ambigen Cliffordimplementer $2+2$ oder $1+1+1+1$ und $|\tr|^2=8$ oder $0$ besitzen. Das schließt die behauptete Identifikation dieser Wirkungen aus.

Geladene Cartanlifts benötigen Charakterphasen. Die Wörter $0,1,0$ und $2,1,2$ haben dieselbe nackte Gitterwirkung, unterscheiden sich aber um Deckcharakter 152. Eine fortsetzbare Darstellung der geprüften endlichen Cartan-Liftgruppe speichert deshalb Gitterwirkung plus acht Charakterbits. Die allgemeinen Hecke-Fortsetzungen benötigen dagegen rationale Phasenkoordinaten $\theta\bmod2$. Sie braucht nicht die gesamte alte Wortliste. Ein ganzzahliges Relationszeugnis schließt in der geprüften Liftaufgabe eine globale Sektion auch unter kontinuierlichen Toruskorrekturen aus.

Der ältere Quartiklift ist ein anderer tatsächlicher $E_8$-Automorphismus:
\begin{equation}
g_5=-T_e,\qquad g_4=e^{-i\pi/4}r_e,
\qquad U_X=\overline g_5\otimes g_4.
\end{equation}
Auf $X$ gilt $U_X^2=-iI$; der Lift hat Ordnung acht. Seine adjungierte Spur $16$ unterscheidet ihn von den involutiven Cartanlifts mit Spuren $-8$ oder $24$. Nur 16 der 60 nativen Transporte erhalten den fest gewählten P2-Generator; die übrigen transformieren den Ladungsrahmen kovariant. Sie sind deshalb nicht automatisch physische Ereignisse, die dieselbe feste Hyperladung erhalten.

\section{Universalraum: Der Raum dessen, was eine Zukunft noch unterscheiden kann}\label{sec:universal}

\key{Zwei Vorgeschichten gehören genau dann zum selben wirksamen Zustand, wenn keine zugelassene spätere Handlung oder Messung sie unterscheiden kann. Welcher Universalraum gemeint ist, hängt deshalb vom tatsächlichen Operations- und Auslesesystem ab.}

\subsection{Operationale Äquivalenz und lineare Vorhersage}
Für eine festgelegte Klasse zulässiger Zukunftstests $T$ ist
\begin{equation}
h\sim h'\quad\Longleftrightarrow\quad
\mathcal P(T\mid h)=\mathcal P(T\mid h')\quad\text{für alle }T.
\end{equation}
Wenn das Gewicht der vorbereitenden Geschichte relevant ist, müssen unnormierte gemeinsame Antworten verglichen werden. Wenn kohärente Referenzen und Interferenz zulässig sind, gehören auch diese zum Testsystem. Eine Aussage über alle Zukunftstests darf nicht auf die gerade sichtbare Bildschirmablesung verkürzt werden.

Für eine skalare Wortantwort $s$ bildet
\begin{equation}
H_{u,v}=s(uv)
\end{equation}
die Hankelmatrix. Ihre Zeilen sind Zukunftsantworten von Vorgeschichten. Bei endlichem Rang entspricht $\rank H$ der Dimension einer minimalen linearen Wortrealisierung unter den üblichen algebraischen Voraussetzungen \cite{carlyle,thon}. Daraus folgen weder Positivität noch automatisch eine minimale Quantensystemdimension.

\subsection{Die positive kohärente Rekonstruktion}
Für Operatorwörter $K_u$ ist der kohärente Kern
\begin{equation}
D(u,v)=\omega(K_u^\dagger K_v).
\end{equation}
Er enthält Kreuzterme zwischen Geschichten. Auf einer unitalen $C^*$-Algebra mit normiertem positivem Zustand liefert die GNS-Konstruktion
\begin{equation}
\mathcal N_\omega=\{A:\omega(A^\dagger A)=0\},\qquad
\mathcal H_\omega=\overline{\mathfrak A/\mathcal N_\omega}.
\end{equation}
Das Nullideal ist unter Linksmultiplikation stabil; die zyklische Darstellung ist bis auf unitäre Äquivalenz eindeutig \cite{segal}. Bei einer allgemeinen Algebra unbeschränkter Felder sind gemeinsame Domänen und analytische Voraussetzungen zusätzlich zu behandeln. Die pauschale Behauptung einer Darstellung in beschränkten Operatoren für jede positive Feld-$*$-Algebra wäre zu stark.

Die GNS-Konstruktion rekonstruiert eine Darstellung aus einer gegebenen positiven Quelle. Sie wählt nicht voraussetzungslos die Quelle, ihre Positivität, die Born-Auslese oder die physisch zulässigen Eingriffe. Ein beliebiger positiver Gramkern auf Geschichtslabels ist außerdem noch kein kausal normalisierter Mehrzeitprozess. Prozess-Tensoren beziehungsweise Quantennetzwerke verlangen zusätzliche Kompositions- und Normierungsbedingungen \cite{combs}.

\subsection{Ein Beispiel, das die drei Ebenen sofort trennt}
Das implementierte Instrument $K_e=U_e/\sqrt{15}$ ergibt bei alleiniger Beobachtung seiner Ereignislabels
\begin{equation}
p(e_n\cdots e_1)=15^{-n}.
\end{equation}
Die zugehörige klassische Wort-Hankelmatrix hat Rang eins. Für denselben Operatorprozess mit $\omega(A)=\tr(A)/5$ besitzt der kohärente Kern dagegen den berechneten GNS-Rang 25, weil die Operatorwörter $M_5(\C)$ aufspannen.

Das ist kein Widerspruch. Die Labelauslese sieht eine einfache Folge gleichgewichteter Ereignisse. Ein kohärenter Zukunftstest kann zusätzliche Operatorunterschiede erkennen. Ebenso erzeugen $I$ und $-I$ dieselbe isolierte Quantenkanalkarte, aber verschiedene GNS-Vektoren; erst eine zugelassene kohärente Referenz macht eine relative Phase sichtbar.

Damit wird die stärkste eingefügte Universalprozess-Synthese präzisiert: Operationale Zukunftsklassen, lineare Hankelrealisierung und positiver GNS-Raum sind drei verbundene Konstruktionen, deren Gleichsetzung passende Tests und Voraussetzungen benötigt.

\subsection{Die exakte Bedingung für eine verlustfreie Rekodierung}
Eine Reduktion $R$ ist unter einer Fortsetzung $F$ geschlossen, wenn
\begin{equation}\label{eq:quotient}
RF=\overline F R.
\end{equation}
Im linearen Fall ist insbesondere $F(\ker R)\subseteq\ker R$ erforderlich. Zusätzlich müssen alle relevanten Auslesungen durch $R$ faktorisieren. Dann folgt durch Induktion die Gleichheit der Antworten für beliebig lange zulässige Operationswörter.

Die Hecke-Carry-Intertwiner liefern eine konkrete Instanz dieser Regel. Der Universalraumanschluss ist deshalb nicht bloß eine Metapher. Er sagt genau, welche Information auf einer gröberen Darstellung bleiben und wie sie bei einer Operation weiterwirken muss.

\subsection{Warum ein Fixpunkt des Gesetzes keine universelle Vergessensmaschine ist}
Ein primitiver positiver Operator kann unter passenden endlichen Voraussetzungen einen eindeutigen normierten Fixpunkt besitzen. Doch P1-Primitivität ist nicht automatisch diese Operatorprimitivität; auch ein endlicher Satz darf nicht ohne Grenzwertkontrolle auf einen unendlichen Raum übertragen werden.

Für einen CPTP-Kanal $\mathcal R$ mit CPTP-Linksinverser $\mathcal L$ auf der gesamten betrachteten Zustandsklasse gilt
\begin{equation}
\mathcal L\mathcal R=\Id
\quad\Longrightarrow\quad
\norm{\mathcal R(\rho)-\mathcal R(\sigma)}_1
=\norm{\rho-\sigma}_1.
\end{equation}
Beide Ungleichungen folgen aus der Kontraktivität der Spurweite. Eine derart verlustfreie Abbildung auf einem nichttrivialen festen Zustandsraum kann deshalb nicht zugleich alle unterscheidbaren Zustände in denselben primitiven Attraktor treiben.

Ein eindeutiges \emph{Gesetz} kann viele Zustände haben. Ein reduzierter Kanal kann nach bewusstem Vergessen relaxieren, während die volle Entwicklung einschließlich Umgebung Information bewahrt. Zukunftsvollständigkeit des Prozessraums allein behauptet keine Linksinvertierbarkeit seiner Vergröberung. Diese Unterscheidungen verhindern einen falschen Kurzschluss von ``rekursiv konsistent'' zu ``einziger physischer Weltzustand''.

\section{Raum, Topologie und Ausbreitung: Was aus dem Graphen folgt}

\subsection{Die räumliche Konstruktion ist markiert}
Die Tutte--Coxeter-Zelle besitzt dreißig Knotentypen, fünfundvierzig Kanten, Grad drei und Zyklusrang sechzehn. Der verwendete markierte periodische $A_3$-Cover wählt daraus einen Rang-drei-Periodenquotienten. In der benutzten Normierung ergibt die kleine Blochwelle
\begin{equation}
\lambda_{\min}(k)=\frac85\norm{k}^2+O(\norm{k}^4).
\end{equation}
Dies ist ein konkreter dreidimensionaler Ausbreitungszweig. Die Dimension stammt aus der gewählten Periodenmarkierung, weder aus den sechzehn unabhängigen Graphschleifen noch aus dem Rang acht des internen Ladungsgitters. Die universelle Überlagerung des endlichen Graphen wäre dagegen ein Baum.

Zehn disjunkte markierte Pfade decken die dreißig Typen. Fünf Blocktypen tragen $U$, fünf $\overline U$. Damit besitzt die ladungskovariante Pfadrekursion einen wirklichen Ort in diesem markierten Netz. Die Auswahl der Perioden und ihre physische Bedeutung bleiben als Voraussetzungen sichtbar.

\subsection{Die effektive Metrikantwort braucht die gesamte Zelle}
Ein naiver Ausdruck aus lokalen zweiten Sprungmomenten,
$\sum_e r_e\,\Delta x_e\Delta x_e^T$, muss noch keine globale effektive Metrik ergeben. An einzelnen Knoten ist er sogar singulär. Der interne harmonische Korrektor beschreibt, wie sich die Knotentypen einer ganzen Periode gemeinsam einstellen.

Die überprüfbare effektive Antwort wird deshalb aus der Blochkrümmung gewonnen:
\begin{equation}
D_{ij}=\frac12\left.\frac{\partial^2\lambda_{\min}(k)}{\partial k_i\partial k_j}\right|_{k=0}.
\end{equation}
In der isotropen Ausgangsnormierung ist $D=(8/5)I_3$. Die Sitzung berechnet auch die Antwort auf geänderte Kantenraten mit dem vollständigen Zellkorrektor. Das ist mehr als eine gezeichnete Raumvermutung, aber noch keine abgeleitete Lorentzmetrik. Diffusive beziehungsweise spektrale Kleinimpulsantwort, unitäre Wellenausbreitung und relativistische Kausalstruktur müssen passend unterschieden werden.

\picturetext{Die scheinbare Weglänge in einem Straßennetz hängt nicht nur von einer Kreuzung ab. Manche Wege zwingen zu Umwegen innerhalb des Viertels. Erst wenn das ganze wiederkehrende Viertel berücksichtigt ist, erhält man die effektive Ausbreitung über große Distanzen.}

\subsection{Existenzsätze erhalten, ihre Reichweite nicht ausweiten}
Die ältere thermodynamische Dynamikarbeit enthält bereits Lieb--Robinson-Schranken, Zeitentwicklung, Gauss-Kompatibilität sowie Grundzustands-/KMS-Existenz für die dort spezifizierte lokale Hamiltonklasse mit endlichen Faktoren auf $\Z^2$. Diese Ergebnisse werden hier nicht erneut als ungelöste reine Existenzfragen behandelt. Sie begründen aber auch nicht von selbst die ausgewählte räumliche $3+1$-Realisierung der vollständigen geladenen Quelle.

Die richtige Anschlussfrage lautet: Welcher dieser bereits verstandenen lokalen Mechanismen trägt genau die jetzt festgelegte Quelle, ihre ursprünglichen Markierungen und ihre physischen Auslesungen? Dabei müssen dieselbe Zeit und derselbe Zustand durch die Konstruktion laufen.

\section{Zuse und Wolfram: Der produktive gemeinsame Kern und die nötigen Unterschiede}

Zuses \emph{Rechnender Raum} motiviert den Gedanken, physische Entwicklung als lokale regelhafte Verarbeitung zu beschreiben. Für diese Sitzung ist die nützliche Konsequenz: Eine Theorie muss sagen, was einen Zustand ausmacht, welcher lokale Schritt erlaubt ist und wie große Strukturen aus wiederholter Komposition entstehen \cite{zuse}. Die konkrete TFPT-Quelle liefert dafür eine reichere Algebra als ein bloßes Bitgitter; daraus folgt keine historische Prioritätsbehauptung und kein Beweis, dass das Universum ein bestimmter Zellularautomat ist.

Wolframs Multiway- und Hypergraphbilder helfen, Zustände, mögliche Aktualisierungen und Abhängigkeiten getrennt darzustellen. Seine eigene Behandlung gewichteter Multiwaygraphen zeigt bereits, dass eine gleichmäßige Verteilung auf Nachfolger bei anderer Foliation keine robuste Normierung garantieren muss \cite{wolfram}. Deshalb darf eine graphische Verzweigung nicht stillschweigend ihre physikalischen Wahrscheinlichkeiten festlegen.

Die konkrete Verbindung zu den Ergebnissen dieser Sitzung ist die folgende Typentrennung:
\begin{longtable}{L{0.21\linewidth}L{0.34\linewidth}L{0.37\linewidth}}
\toprule
Graphart & Knoten und Kanten & Was zusätzlich zu prüfen ist\\\midrule\endhead
Ladungsgraph & Wurzeln, Cartanrichtungen, nichtverschwindende Klammern & Wirkliche Darstellung, Kokzyklus, Markierungen\\
Rechengraph & Teilwörter und Ward-Reduktionsschritte & Äquivalenz verschiedener Auswertungsbäume\\
Antwort-Hypergraph & Verbundene Mehrfeldkoeffizienten & Quelle, Zustand, Positionen; keine automatische Kausalinterpretation\\
Geschichtsgraph & Physische Operationen und ihre Verkettung & Instrumente, Phasen, kausale Normierung\\
Raumgraph & Lokale Zellen und räumliche Nachbarschaft & Periodenwahl, Ausbreitung, Kontinuums- und Zeitanschluss\\
Theoriegraph & Aussagen und ihre Begründungsabhängigkeiten & Originalbeleg und Voraussetzungen; ein Pfad ist kein Beweis\\\bottomrule
\end{longtable}

Die besonders wichtige positive Verbindung ist: Das ausgeführte Quellgesetz besitzt unterschiedliche exakte Auswertungsbäume, während sein verbundener Anteil echte höhere Antwort-Hyperkanten liefert. Beides stammt aus derselben Algebra. Die Bäume dürfen optimiert werden; die Quellantwort muss erhalten bleiben. Werden diese Auswertungsbäume stattdessen als verschiedene Welthistorien gezählt und erneut gewichtet, ändert man das bereits festgelegte Gesetz.

Eine gute Gesamtkonstruktion würde die Typen durch explizite Abbildungen verbinden: Ein räumliches lokales Ereignis wirkt durch einen Operator; dieser verändert einen Zustand und seinen Carry; das erzeugt eine überprüfbare Mehrzeitantwort; eine gröbere Darstellung berechnet dieselbe Antwort. Der Graph ist dann eine Darstellung der vollständigen Regel und kein Ersatz für sie.

\section{Fortsetzung: Die Auswahl von Level eins ist auf einen Ursprungssatz reduziert}\label{sec:levelselection}

\key{\textbf{Ein positiver Auswahlsatz ist jetzt verfügbar.} Innerhalb einer unitären positiven affinen $E_8$-Quelle genügt entweder eine unabhängig begründete Zentralzahl höchstens acht oder die Nullrelation eines einzigen echten langen Wurzelstroms, um Level eins zu erzwingen. Das verbleibende Problem ist die Herkunft dieser konkreten Eigenschaft aus der ursprünglichen Quelle.}

\subsection{Was Zentralzahl acht tatsächlich erzwingt}
Die Vergleichsklasse ist ausdrücklich festgelegt: eine einfache unitäre Vertexoperatoralgebra vom CFT-Typ, mit positivem Energiegrad, eindeutigem Vakuum im Grad null und einer tatsächlichen kompakten $E_8$-Stromalgebra im Grad eins. Lange Wurzeln haben Normquadrat zwei. Der affine Level ist dann eine positive ganze Zahl $k$; die von den Strömen erzeugte Algebra ist $L_k(\mathfrak e_8)$ \cite{donglin}.

Ihr Sugawara-Stress-Tensor hat
\begin{equation}
c_{\rm aff}(k)=\frac{248k}{k+30},\qquad
c_{\rm ges}\ge c_{\rm aff}(k)\ge8.
\end{equation}
Die erste Ungleichung folgt aus der Positivität der verbleibenden konformen Komponente. Wenn der Ursprung unabhängig $c_{\rm ges}\le8$ liefert, folgt zwingend $k=1$ und $c_{\rm ges}=8$. Noch stärker:
\begin{equation}
\norm{\omega_{\rm ges}-\omega_{\rm Sug}}^2
=\frac{c_{\rm ges}-c_{\rm aff}}2=0.
\end{equation}
Damit stimmen die Stress-Tensoren selbst überein. Bei Erzeugung durch die Ströme ist die Quelle $L_1(\mathfrak e_8)\cong V_{E_8}$; der Übergang zum lokalen Netz benötigt die einschlägigen starken Lokalitätsvoraussetzungen \cite{cklw}. Die Eindeutigkeit des affinen Vakuumsektors wählt noch keinen angeregten physischen oder kosmologischen Zustand aus.

\textbf{Die kritische Richtung:} Rang acht gibt zunächst eine Untergrenze für die Zentralzahl, keine Obergrenze. Auch ``nur diese Ströme, keine entkoppelten Zusatzsysteme'' genügt nicht: $L_2(\mathfrak e_8)$ ist eine unitäre, von denselben 248 Stromrichtungen erzeugte Alternative mit $c=31/2$. Sie ist ein Gegenfall zu genau diesem schwächeren Auswahlargument, keine Behauptung, sämtliche übrigen TFPT-Ursprungsbedingungen zu erfüllen.

Die kanonische Gitterquantisierung aus Abschnitt~\ref{sec:source} liefert ihrerseits bereits Level eins: Der Wurzelvertex besitzt bei Normquadrat zwei einen Doppelpol. Sobald diese Quantisierung als ursprüngliche Quelle identifiziert ist, ist ihr Level keine weitere freie Wahl. Genau diese Identifikation darf aber nicht durch den schon daraus berechneten Doppelpol bewiesen werden.

\subsection{Ein einziger tatsächlicher Wurzelstrom genügt}
Für einen normierten langen Wurzelstrom und sein wirkliches Adjungiertes gilt
\begin{equation}
[e_m,f_n]=h_{m+n}+mk\delta_{m+n,0},\quad
[h_0,e_{-1}]=2e_{-1},\quad e_{-1}^{\dagger}=f_1.
\end{equation}
Anwendung auf das Vakuum ergibt
\begin{equation}\label{eq:rootnull}
\norm{e_{-1}^2\Omega}^2=2k(k-1).
\end{equation}
Somit erzwingen $e_{-1}\Omega\ne0$ und $e_{-1}^2\Omega=0$ den Level eins. Die zweite Gleichung ist eine Aussage über einen vollständigen affinen Modus. Die Nilpotenz irgendeiner endlichen Pauli-Matrix würde dafür nicht genügen.

Diese Nullrelation besitzt eine konkrete ursprüngliche Realisierung, sobald zwei unabhängige gleichchirale CAR-Arten vorhanden sind. Mit der tatsächlichen Viertelholonomie-Polarisation
\begin{equation}
a_{i,j}\Omega=0\ (j\le0),\qquad
a^\dagger_{i,j}\Omega=0\ (j\ge1)
\end{equation}
liefert der neutrale Wurzelstrom
\begin{equation}
E_{-1}=\sum_j a^\dagger_{1,j}a_{2,j+1},\qquad
E_{-1}\Omega=a^\dagger_{1,0}a_{2,1}\Omega,\qquad E_{-1}^2\Omega=0.
\end{equation}
Seine erste Norm ist eins; seine Energie ist $1/4+3/4=1$. Die gemeinsame chemische Viertelverschiebung fällt im neutralen Strom weg.

Der bereits hergeleitete einzelne Rohkanal liefert dagegen einen Dichtestrom $J_{-1}=\sum_j a_j^\dagger a_{j+1}$ mit
\begin{equation}
J_{-1}^2\Omega=a_{-1}^\dagger a_1\Omega-a_0^\dagger a_2\Omega,
\qquad\norm{J_{-1}^2\Omega}^2=2.
\end{equation}
Beide Rechnungen wurden mit den ursprünglichen CAR-Matrizen geprüft; das Fenster $j=-1,0,1,2$ enthält für diese zwei Anwendungen alle nichtverschwindenden Beiträge. Ein einzelner Kanal darf deshalb nicht als das benötigte unabhängige Artenpaar ausgegeben werden. Der vorhandene Kanalvertrag findet einen komplexen chiralen Kanal je Rand, mit entgegengesetzter Orientierung der Ränder; die acht Reihen sind keine acht unabhängigen gleichchiralen Ströme.

\picturetext{Die ganze Anlage muss nicht nochmals aus 248 Einzelteilen erraten werden. Ein einziges richtig angeschlossenes Zahnrad kann die gemeinsame Übersetzung festlegen. Es muss jedoch wirklich ein Zahnrad derselben Anlage sein. Ein ähnlich aussehender Schalter erfüllt die benötigte Relation nicht.}

Der kleinste hinreichende Herkunftssatz lautet damit: Die Rohquelle liefert einen vakuum-, moden-, adjungierten- und ladungsverträglichen langen Wurzelstrom mit Einbettungsindex eins und der Nullrelation~\eqref{eq:rootnull}. Die alte freie R2-Route mit 16 reellen Majoranas liefert eine hinreichende Mehrkanalrealisierung, sobald ihre Quellidentifikation gilt. Ihre Herkunft wird durch die bloße Zahl 16 nicht nachgeholt.

\subsection{Warum die bisherige Fünferrekursion den Level nicht auswählt}
Die tatsächliche Chevalley-Ward-Rechnung lässt sich ohne Level-eins-Vorentscheidung ausführen. Für die Felder $X_{ia},X_{jb}^\dagger,X_{kc},X_{ld}^\dagger$ an $(-1,0,1,\infty)$ folgt exakt
\begin{equation}
F_k=k^2(A+D)+k(B+C),
\end{equation}
mit
\begin{align}
A&=\delta_{ij}\delta_{kl}\delta_{ab}\delta_{cd},&
D&=\delta_{il}\delta_{jk}\delta_{ad}\delta_{bc},\\
B&=\delta_{ij}\delta_{kl}\delta_{ad}\delta_{bc},&
C&=\delta_{il}\delta_{jk}\delta_{ab}\delta_{cd}.
\end{align}
Bei festem Familienindex bleibt $k(k+1)(\delta_{ij}\delta_{kl}+\delta_{il}\delta_{jk})$. Nach Gesamtnormierung ist der Fünfertensor daher für alle positiven $k$ derselbe. Die normierte Fünferrekursion bestätigt ihre Form, sieht aber diese Levelalternative nicht.

Zwei gemischte Antworten unterscheiden sie unmittelbar: Für Carrierindizes $(0,0,1,1)$ ergeben Familienindizes $(0,0,1,1)$ den Wert $k^2$ und $(0,1,1,0)$ den Wert $k$. Das Verhältnis ist $1/k$. Es bleibt unter einer gemeinsamen Feldnormierung erhalten.

\begin{center}
\begin{tabular}{ccccc}
\toprule
$k$ & $c_{\rm aff}$ & gemischtes Verhältnis & $\norm{e_{-1}^2\Omega}^2$ & neutrale Restnorm$^2$\\\midrule
1 & 8 & 1 & 0 & 0\\
2 & $31/2$ & $1/2$ & 4 & $15/4$\\
3 & $248/11$ & $1/3$ & 12 & $80/11$\\\bottomrule
\end{tabular}
\end{center}
Die letzte Spalte ist $\norm{\omega_{\rm Sug}-\omega_{\rm Heis}}^2=120(k-1)/(k+30)$. Im Grad zwei besitzt die Level-eins-Gitterquelle 4124 Zustände, die affinen Quellen mit $k\ge2$ dagegen 31124. Die Differenz 27000 entsteht durch die Level-eins-Nullrelationen. Diese exakten Kennzeichen ersetzen keine unabhängige Herleitung ihrer Erfüllung durch den Rohprozess.

Die neue Rechnung prüft alle 225 Paare von Gleichheitsmustern der vier Carrier- und Familienindizes; die Erweiterung auf konkrete Indexbelegungen verwendet die bereits geprüfte Tensor-Kovarianz. Das volle $W_5\otimes W_4$ fordert im affinen Vergleich tatsächlich $k=1$. Wurde dieser Tensor zuvor aus $E_{8,1}$ hergeleitet, ist er aber kein unabhängiger Ursprungsbeweis für dieselbe Quelle.

\subsection{Primitive Windung und Determinanten haben unterschiedliche Normierungen}
Die alte Rechnung \file{v470_alpha_inflow_level.py} identifiziert $|C|=1$ mit dem primitiven $U(1)$-Inflow und dem betreffenden Koeffizienten der $\alpha$-Gleichung. Sie identifiziert nicht ohne zusätzlichen Abbildungssatz den affinen $E_8$-Level.

Bei einer Determinante in einer endlichen Darstellung $\rho$ multipliziert der Dynkinindex $d_\rho$ den grundlegenden affinen Kokzyklus. In der Normierung langer Wurzeln auf zwei hat die adjungierte 248er-Darstellung von $E_8$ Index 60; der kleinste nichtverschwindende Darstellungsindex ist 60 \cite{laszlosorger}. Deshalb ist ``eine adjungierte Fermiondeterminante'' kein automatischer Level-eins-Beweis. Die ursprüngliche Determinantenlinie muss auf die grundlegende Linie samt Verbindung mit der tatsächlich passenden Normierung abgebildet werden. Die bereits formulierte kontinuierliche Quillen-Brücke bleibt genau dafür relevant.

\section{Konstruktiver Abschluss innerhalb der Quelle: Felder, Antworten und Gedächtnis}\label{sec:constructive}

\key{\textbf{Die Quelle hat jetzt eine zweite vollständige, ausführbare Darstellung.} Die $X$-Felder sind explizite Fermionbilineare; die $C$-Felder sind die erforderlichen geladenen Erweiterungsfelder. Mit dem ursprünglichen Kokzyklus ergibt sich eine geschlossene Formel für sämtliche endlichen C/X-Wörter. Dieselbe Formel liefert einen positiven Geschichtskern und eine exakte Darstellung ihres gesamten Nachfahrengedächtnisses.}

\subsection{Die neun Hilfsfermionen und die vier Erweiterungsfelder}
Die aus den ursprünglichen Wurzeln gebaute $A_8$-Isometrie bildet alle 20 X-Wurzeln exakt auf
\begin{equation}
\alpha(X_{ia})=e_{5+a}-e_i\quad(0\le i<5,\ 0\le a<4)
\end{equation}
im neundimensionalen Nullsummenraum ab. Neun komplexe chirale Hilfsfermionen realisieren deshalb die $A_{8,1}$-Ströme durch
\begin{equation}
X_{ia}(z)=:\!\psi_{5+a}^\dagger(z)\psi_i(z)\!:.
\end{equation}
Der gemeinsame $U(1)$-Anteil wird abgespalten: $c=9-1=8$. Die tatsächlichen 8000 Tripelklammern erfüllen
\begin{equation}
[[X_{ia},X_{jb}^\dagger],X_{kc}]
=\delta_{ij}\delta_{bc}X_{ka}+\delta_{ab}\delta_{jk}X_{ic}.
\end{equation}
Die ursprünglichen Phasen passen also zu den Matrixeinheiten; es genügt nicht nur eine Dimensionsgleichheit. Die P2-Wirkung wird mitgeführt: Die diagonalen Werte
\[
(-1/3,-1/3,-1/3,1/2,1/2,0,0,0,0)
\]
geben auf $X_{ia}$ genau $-y_i$.

Die vier C-Felder besitzen in denselben Koordinaten die Ladungen
\begin{equation}\label{eq:Csu9}
q(C_a)=\sum_{i=0}^4e_i+e_{5+a}-\frac23\sum_{r=0}^8e_r.
\end{equation}
Sechs Einträge sind $1/3$, drei sind $-2/3$; die Norm ist zwei und die P2-Ladung null. Das sind Gewichte des $\Lambda_6$-Sektors. Das zugehörige Sechsfermionprodukt hätte zunächst Gewicht drei. Nach Abspaltung seiner gemeinsamen U1-Ladung sechs bleibt
\begin{equation}
h(C)=\frac62-\frac{6^2}{2\cdot9}=1.
\end{equation}
Die C-Felder sind somit geladene Verzweigungsvertexfelder, keine gewöhnlichen neutralen Sechsfermionpolynome. Zusammen mit dem konjugierten $\Lambda_3$-Sektor ergeben sie genau die ursprüngliche lokale Erweiterung
\begin{equation}
V_{E_8}=V_{A_8}\oplus V_{A_8+\Lambda_3}\oplus V_{A_8+\Lambda_6},
\qquad248=80+84+84.
\end{equation}
Dieser bekannte mathematische Zusammenhang \cite{distler} ist hier an die tatsächlichen 24 Originalfelder angeschlossen. Die neun Hilfsfermionen sind eine Darstellungshilfe; ihre Anzahl behauptet weder neun Standardmodellteilchen noch eine neue räumliche Dimension.

\subsection{Die Sechser-Hyperkante ist ein elementarer Fermionring}
Für den bekannten Sechserzeugen gibt es in der Liste der zwölf elementaren Fermionfelder genau eine zulässige Wick-Paarung. Ihr Gesamtvorzeichen ist positiv. Sie liefert unabhängig von der Strom-Ward-Rekursion
\begin{equation}
F_6=\frac1{z_{12}z_{16}z_{23}z_{34}z_{45}z_{56}}=\frac1{576}
\end{equation}
an den bisherigen sechs Positionen. Alle 15 Paarungen ganzer Ströme bleiben null. Damit ist eine wichtige falsche Schlussfolgerung behoben: Nichtgaußische Stromantworten schließen eine zugrunde liegende Gaußdarstellung durch \emph{zusammengesetzte} Felder nicht aus. Sie schließen den Ersatz der Ströme durch lineare quasifreie Felder mit denselben Paarungen aus.

\picturetext{Sechs ganze Bauteile passen paarweise nicht zusammen. Jedes enthält aber zwei Anschlüsse. Über diese inneren Anschlüsse schließen sich alle sechs zu einem Ring. Wer jedes Bauteil durch einen einzigen Anschluss ersetzt, verliert den Ring und damit die berechnete gemeinsame Antwort.}

Das entspricht einem präzisen Hypergraphen: Die Sechserkante ist in den Stromvariablen irreduzibel; in einer reicheren, tatsächlich hergeleiteten Felddarstellung besitzt sie innere Struktur. Zuse- und Wolfram-Bilder können solche Darstellungen beschreiben. Der Wechsel der Darstellung muss die Antwort erhalten; er erzeugt keine neue Weltregel.

\subsection{Die volle C/X-Regel als geschlossenes Produkt}
Jedes Originalfeld ist ein phasenmarkierter Wurzelvertex $s_iV_{\alpha_i}$. Mit dem ursprünglichen bimultiplikativen E8-Kokzyklus $\varepsilon$ gilt bei verschiedenen endlichen Positionen
\begin{equation}\label{eq:latticeoracle}
F_n=\delta_{\sum_i\alpha_i,0}
\left(\prod_i s_i\right)
\left(\prod_{i<j}\varepsilon(\alpha_i,\alpha_j)\right)
\prod_{i<j}(z_i-z_j)^{(\alpha_i,\alpha_j)}.
\end{equation}
Die Potenzen sind ganzzahlig. Die Formel folgt direkt aus der Gittervertex-Konstruktion \cite{frenkelkac}; sie ist hier mit den Originalwurzeln, ihrer kompakten Adjungierung und ihren tatsächlichen Kokzyklusphasen implementiert. Bei neutraler Ladung enthält sie auch sämtliche geladenen Zwischenwege, die der Ward-Auswerter rekursiv durchläuft.

Ein neuer entscheidender Zeuge verlangt die Erweiterung wirklich:
\begin{equation}
C_0,C_1,C_2,X_{00},X_{11},X_{22},X_{33},X_{43}.
\end{equation}
Seine E8-Gesamtladung ist null. In der Hilfsfermiondarstellung beträgt die gemeinsame U1-Ladung jedoch 18. Ein gewöhnlicher Fermion-Vakuum-Wickwert wäre deshalb null; die volle ursprüngliche Quelle ergibt an $z=0,\ldots,7$
\begin{equation}
F_8=-\frac1{4032000}.
\end{equation}
Gitterprodukt und ursprüngliche Ward-Rekursion liefern denselben Wert. Die $\Z_3$-Erweiterung ist somit eine wirksame Komponente des Quellgesetzes, kein bloßes Etikett.

Für reine C/X-Wörter benötigt die Produktformel quadratisch viele exakte Paaroperationen. Sie vermeidet die wachsende Zahl rekursiver Ward-Verzweigungen. Die Bitkomplexität großer rationaler Zahlen kommt hinzu. Dies ist eine konkrete Vereinfachung des TFPT-Rechners, keine Behauptung einer neuen Erfindung der Gittervertexformel. Cartaneinsetzungen und beliebige Nachfahren sind nicht Eingaben dieser speziellen API; die ursprüngliche Ward-Regel bleibt für ihre entsprechenden Aufgaben vorhanden.

\subsection{Ein expliziter Universalraum für radiale Quellgeschichten}
Es sei
\begin{equation}
\ket{u}=\prod_{i=1}^m s_iV_{\alpha_i}(z_i)\Omega,
\qquad1>|z_1|>\cdots>|z_m|\ge0.
\end{equation}
Das ist eine radiale Feldgeschichte, keine bereits normierte Folge von Detektormessungen. Definiere
\begin{equation}
Q_u=\sum_i\alpha_i,\qquad p_{u,n}=\sum_i\alpha_i z_i^n\quad(n\ge1)
\end{equation}
und den Vorfaktor $A_u$ als das Produkt rechts in~\eqref{eq:latticeoracle} ohne Neutralitätsdelta. Die Heisenberg-Konstruktion gibt dann den \emph{vollständigen} Zustandsvektor
\begin{equation}\label{eq:fullhistory}
\ket{u}=A_u\exp\!\left(\sum_{n\ge1}\frac{p_{u,n}\cdot a_{-n}}{n}\right)\ket{Q_u}.
\end{equation}
Die Überlappung zweier solcher Geschichten ist daher
\begin{align}\label{eq:historyexact}
K(u,v)&=\delta_{Q_u,Q_v}\overline{A_u}A_v
\exp\!\left(\sum_{n\ge1}\frac{\overline{p_{u,n}}\cdot p_{v,n}}{n}\right)\nonumber\\
&=\delta_{Q_u,Q_v}\overline{A_u}A_v
\prod_{i,j}(1-\overline{z_{u,i}}z_{v,j})^{-(\alpha_{u,i},\alpha_{v,j})}.
\end{align}
Die Positivität folgt aus der tatsächlichen Hilbertraumdarstellung. Sie wurde zusätzlich an einer exakten Gram-Matrix und unabhängig über radial adjungierte Original-Ward-Wörter geprüft. Die implementierte API verwendet reelle rationale Positionen mit $|z|<1$, einschließlich negativer Positionen, in der angegebenen radialen Ordnung; die Formel erklärt den allgemeineren komplexen Radialbereich.

\key{Damit ist der Universalraum für diese Quellgeschichten konkret: Eine Geschichte hinterlässt ihre Gesamtladung, ihre Amplitude und ihre gesamte Modenfolge. Diese Daten bestimmen den unnormierten Zustandsvektor und damit alle späteren wohldefinierten linearen Quellantworten auf diesem Zustand.}

Die unendliche Modenfolge muss dabei nicht als unendliche Liste gespeichert werden. Ihre erzeugende Funktion ist rational:
\begin{equation}\label{eq:momentmemory}
P_u(t)=\sum_{n\ge1}p_{u,n}t^n
=\sum_i\alpha_i\frac{z_i t}{1-z_i t}.
\end{equation}
Eine endliche Geschichte besitzt also eine endliche exakte Beschreibung \emph{aller} dieser Nachfahrenmomente. Bei wachsender Geschichte kann die Zahl ihrer Pole wachsen; das ist kein Beweis eines universellen festen endlichen Gedächtnisrangs. Operationale Äquivalenz bei einer eingeschränkten Menge von Zukunftsmessungen kann außerdem gröber sein als die Gleichheit dieser unnormierten Vektoren.

Wird links ein neues Feld $s_\beta V_\beta(w)$ mit $1>|w|>|z_1|$ eingesetzt, lautet die genaue Fortschreibung
\begin{align}
Q'&=Q+\beta,&P'(t)&=P(t)+\beta\frac{wt}{1-wt},\\
A'&=A\,s_\beta\,\varepsilon(\beta,Q)
\prod_i(w-z_i)^{(\beta,\alpha_i)}.
\end{align}
Beim verwendeten bimultiplikativen Kokzyklus ist $\varepsilon(\beta,Q)=\prod_i\varepsilon(\beta,\alpha_i)$. Wenn eine Rekodierung diese Daten oder exakt denselben Zustandsvektor erhält, bleiben die nachfolgenden Quellantworten erhalten. Das schließt die Verbindung von geladenem Carry, positivem Kern und Nachfahrengedächtnis für diese konkrete Geschichtenklasse. Die Wahl physischer Messinstrumente und ihre Kausalnormierung bleibt eine zusätzliche Aufgabe.

\begin{center}
\begin{tikzpicture}[node distance=7mm]
\node[box,text width=13.2cm] (f) {\textbf{Dieselben ursprünglichen Felder}\quad $C_4+X_{20}$ mit Ladungen und Kokzyklus};
\node[box,below=of f,xshift=-3.45cm,text width=5.9cm] (w) {\textbf{Antworten}\quad Ward-Rekursion\quad $=$ Gitterprodukt für Wurzelwörter};
\node[box,below=of f,xshift=3.45cm,text width=5.9cm] (g) {\textbf{Geschichtszustand}\quad $Q$, $A$, rationale Funktion $P(t)$\quad alle Nachfahrenmomente};
\node[box,below=31mm of f,text width=13.2cm] (k) {\textbf{Positiver Universalraumkern}\quad $K(u,v)=\langle u|v\rangle$\quad exakte Fortschreibung und konforme Bewegung};
\node[box,below=of k,text width=13.2cm] (p) {\textbf{Noch zu identifizieren}\quad ursprüngliche physische Quelle und Raumzeitfelder\quad gleicher Zustand, gleiche Ladungen, gleiche Zeit};
\draw[arr](f)--(w);\draw[arr](f)--(g);\draw[arr](w)--(k);\draw[arr](g)--(k);\draw[dasharr](k)--(p);
\end{tikzpicture}
\end{center}

\section{Die Quelle bewegt sich wirklich: Eine exakte Zeitentwicklung mit ihrem Rest}\label{sec:mobius}

\subsection{Die höheren Moden gehören zur Bewegung}
Der bisherige Einwand $L_0=1$ auf allen Stromzuständen vom Grad eins ist richtig, aber er bedeutet nicht, dass die Quelle keine nichttriviale geometrische Zeit besitzt. Mit der positiven Hermiteschen Strommetrik und tatsächlichem Adjungieren bilden
\begin{equation}
\ket{A,n}=\frac{J_{-(n+1)}^A\Omega}{\sqrt{k(n+1)}}\quad(n\ge0)
\end{equation}
einen normierten Nachfahrenturm. Aus $[L_m,J_n]=-nJ_{m+n}$ folgen die Leiterkoeffizienten $\sqrt{(n+1)(n+2)}$ und $\sqrt{n(n+1)}$. Der antihermitesche Generator
\begin{equation}
G=\frac{L_{-1}-L_1}{2},\qquad U(s)=e^{sG},\qquad r=\tanh(s/2)
\end{equation}
liegt bereits in derselben konformen Quelle. Seine exakte Wirkung lautet
\begin{equation}\label{eq:mobiustower}
U(s)\ket{A,0}=(1-r^2)\sum_{n\ge0}\sqrt{n+1}\,r^n\ket{A,n}.
\end{equation}
Dies ist die entsprechende bekannte $SU(1,1)$-Kohärenzstruktur \cite{perelomovstates}, hier aus den wirklichen Strommoden gewonnen. Es wurde kein zusätzlicher Oszillator-Hamiltonoperator gewählt.

Die Modengewichte und der exakte Rest sind
\begin{equation}
p_n=(n+1)(1-r^2)^2r^{2n},\qquad
\sum_{n\ge N}p_n=r^{2N}(N+1-Nr^2).
\end{equation}
Diese Größen sind konforme Gradgewichte eines kohärenten Zustands. Sie sind keine automatisch stattfindenden Detektorklicks. Zwischen zwei Bewegungsständen gilt
\begin{equation}
\langle A,s|B,t\rangle=\delta_{AB}\operatorname{sech}^2\!\left(\frac{s-t}{2}\right).
\end{equation}

\subsection{Warum eine endliche Projektion Gedächtnis verliert}
Für den Projektor $P_1$ auf die anfänglichen Grad-eins-Ströme folgt
\begin{equation}
P_1U(s)P_1=\operatorname{sech}^2(s/2)P_1.
\end{equation}
Bei $s=\log3$ ist diese Amplitude $3/4$. Die ununterbrochene Bewegung bis $2s$ hat Amplitude $9/25$, während zweimaliges Projizieren $9/16$ ergibt. Die scheinbare endliche Dynamik komponiert somit nicht wie der volle Prozess. Sie hat Anteile entfernt, die später zurückwirken können.

\picturetext{Ein angestoßener Ton verteilt sich auf Obertöne. Wer zwischendurch nur den Grundton behält, beschreibt beim nächsten Schritt ein anderes Instrument. Die rationale Gedächtnisdarstellung und der genaue Rest sagen, was beim Vergröbern erhalten bleiben muss.}

Mit $g_s(z)=(z+r)/(1+rz)$ und der Cayley-Koordinate $x=i(1-z)/(1+z)$ gilt $x\mapsto e^{-s}x$. In der üblichen geometrischen Einintervall-Modularkonvention ist daher $s=2\pi t$ für $\Delta_I^{it}$ \cite{cklw}. Mehrere disjunkte Intervalle besitzen zusätzliche Struktur; die einfache Formel darf dort nicht unverändert eingesetzt werden \cite{longomodular}.

Die Bewegung wirkt auf jede orthonormierte interne Stromrichtung gleich. Sie erzeugt somit selbst keine Flavorhierarchie. Der konforme Parameter ist ferner noch nicht mit Laborzeit, kosmischer Zeit oder den historischen Compilerclocks identifiziert. Die neue Aussage ist die ausdrücklich konstruierte gemeinsame Quellenbewegung mit kontrolliertem Nachfahrengedächtnis.

\section{Der physikalische Anschluss nach diesen Rechnungen}\label{sec:physicalfollow}

\key{Die Quelle, ihre vollständigen Wurzelantworten, ihr positiver Geschichtskern und ihre geometrische Bewegung passen jetzt in einer Konstruktion zusammen. Die physische Aufgabe ist die Abbildung dieser Konstruktion auf tatsächlich realisierte Raumzeitfelder. Ein bloßer Ersatz durch eine passende Kovarianz würde bereits die berechneten höheren Antworten verändern.}

\subsection{Warum die Massenrekonstruktion noch keine Massenherkunft ist}
Die ursprüngliche \file{v258_dirac_covariance_induction.py} verwendet
\begin{equation}
H_{\rm mod}=\log((I-C_F)C_F^{-1}),\qquad
D_F^{\rm ind}=\mu_{\rm geo}\,\Pi_{\rm odd}H_{\rm mod}.
\end{equation}
Dabei ist $\mu_{\rm geo}$ der geometrische Maßstab und $\Pi_{\rm odd}$ die Projektion auf den chiraliätsungeraden Anteil. Im ausgeführten chiralen Beispiel ist $\mu_{\rm geo}=1$; die Projektion ändert den eingesetzten Operator nicht. Die Rekonstruktion setzt neun Massen in $D$ ein, erzeugt daraus $C$ und gewinnt anschließend $D$ zurück. Der physische Pfeil $C_F=P_FC_\Sigma P_F$ steht als bedingte Identifikation mit einer auf \texttt{True} gesetzten Prüfung im Programm. Die spätere gezielte Quellprovenienzprüfung vom 21. September behandelt genau diesen noch fehlenden Vorwärtspfeil. Dieser Befund betrifft die spezielle Kovarianzroute; er hebt die ursprünglichen Quellenmassenformeln und ihre Herleitungen aus Abschnitt~\ref{sec:physicalconstraints} nicht auf.

Im invarianten E8-Vakuum haben gleich gesmearte orthonormierte Ströme dagegen die Paarung $K(f,f)I$. Eine reine isometrische interne Auswahl bleibt skalar. Selbst bei einer zusätzlich postulierten CAR-Identifikation wäre ihr Logit skalar. Unterschiedliche Massen können daraus nicht durch bloßes Umbenennen der 24 Indizes entstehen.

Die neue Kompositdarstellung behebt einen möglichen Statistikfehler, aber sie wählt weder den physischen Materieprojektor $P_F$ noch feldabhängige räumliche Profile. Dafür müssen die ursprünglichen P2- und Clockmarkierungen mit einer tatsächlich begründeten räumlichen und geladenen Restriktion verbunden werden. Der echte höhere Stromring und die C-Erweiterungsantwort müssen dabei erhalten bleiben.

\subsection{Der räumliche Ausschnitt ist ein Teil der Herleitung}
Die originale Vierintervallrechnung \file{v480_multilocal_four_interval.py} verwendet vier echte Teilintervalle mit Lücken eines freien fermionischen Rings: 256 Orte, jeweils 48 gewählte Orte im Abstand 64. Vier markierte Punkte allein wählen diese Lücken nicht aus. Im stark additiven vollen Netz erzeugen nur durch Punkte getrennte anstoßende Bögen wieder die globale Algebra; das ist nicht derselbe reduzierte Zustand.

Der vorhandene Kompressionssatz
\begin{equation}
C_F-C_F^2=(QC_\Sigma P_F)^\dagger(QC_\Sigma P_F),\qquad Q=I-P_F,
\end{equation}
für einen reinen projektorwertigen Quellzustand zeigt, wo die gemischte endliche Kovarianz herkommen kann: aus der Kopplung an ausgelassene Freiheitsgrade. Er ist ein bereits vorhandenes Ergebnis. Die explizite Möbiusbewegung zeigt jetzt am gemeinsamen Stromturm, wie eine autonome endliche Ersatzzeit genau diese Information verlieren würde.

\subsection{Welche gemeinsame Abbildung jetzt wirklich geprüft werden kann}
Für eine behauptete physische Realisierung wird ein Feldtransfer $\iota$ benötigt, so dass für alle zugelassenen Quellwörter
\begin{equation}
\omega_{4D}\bigl(\iota(J_1)\cdots\iota(J_n)\bigr)
=\omega_{E_8}\bigl(J_1\cdots J_n\bigr)
\end{equation}
gilt, mit verträglicher Lokalisation, P2-Ladung, Statistik, Adjungierung und Zeit. Die neue Produktformel macht die rechte Seite für beliebig lange endliche C/X-Wörter direkt berechenbar. Der positive Kernel prüft zusätzlich den gemeinsamen Zustand; die Möbiusformel prüft seine geometrische Bewegung.

Die bisher ausgeführte 3+1-Wandprobe besitzt eine 2+1-dimensionale Wand und spinlose Jordan-Wigner-Moden. Sie enthält noch keinen nachgewiesenen Transfer der 1+1-dimensionalen C/X-Naht als Defekt innerhalb dieser Wand. Ebenso bleibt der Seamterm der vierdimensionalen Masterwirkung eine noch zu realisierende Anschlussanforderung. Diese fehlende Abbildung verhindert derzeit, die neue Wandprobe als gemeinsame Realisierung der bereits vorhandenen Massen-, Kopplungs- und Gravitationsgesetze auszugeben. Die dafür nötige Operatoridentität folgt nicht durch Einsetzen der Zielgrößen in eine Kovarianz.

Die tragende Arbeitsrichtung ist damit zusammenhängend: den ursprünglichen geladenen Quelltransfer konstruieren, seinen Level durch den einzelnen Wurzelzeugen festlegen und denselben Transfer räumlich sowie zeitlich fortsetzen. Das erhält Double Cover, P2, $\phi_0$, $\alpha$ und die Clocks als gemeinsam zu erfüllende markierte Zielantworten. Eine optionale E8-Kaskadengeschichte wird dafür nicht zur zusätzlichen Pflichtannahme erklärt.

\section{Gesamtprüfung: Ursprung, Universalraum, Raumzeit und Materie an einem Objekt}\label{sec:jointclosure}

\key{\textbf{Ergebnis dieser Fortsetzung:} Die stärkste zusammenhängende Route hat zwei noch nicht ausgefüllte Eingänge: die \emph{Auswahl der gemeinsamen Quelle} und die \emph{Auswahl ihrer physischen Lokalisation}. Dazwischen liegen echte Konstruktions- und Rekonstruktionssätze. Diese unterscheiden wir von den Eingängen, die sie voraussetzen. Dadurch wird sichtbar, welche zusätzliche Herleitung mehrere bisher getrennte Probleme gleichzeitig verbinden würde.}

\subsection{Ein wirklicher topologischer Zusammenhang: Warum der Faktor 16 erscheint}
\label{sec:bordism16}
Die originale Kummer-/K3-Konstruktion besitzt einen präzisen Anschluss an die Quantisierung einer gravitativen Antwort. Er beruht auf den Gruppen glatter geschlossener Viermannigfaltigkeiten modulo Bordismus:
\begin{equation}
\Omega_4^{\mathrm{Spin}}\cong\Z,\qquad
\Omega_4^{SO}\cong\Z.
\end{equation}
K3 erzeugt die erste Gruppe, $\mathbb{CP}^2$ die zweite, jeweils bis auf Orientierung. Die Signaturen sind $-16$ und $1$. Beim Vergessen der Spinstruktur gilt daher
\begin{equation}\label{eq:forgetspin}
f:\Omega_4^{\mathrm{Spin}}\longrightarrow\Omega_4^{SO},
\qquad [\mathrm{K3}]\longmapsto-16[\mathbb{CP}^2].
\end{equation}
Für die ganzzahligen Antwortfunktionen $q_{SO}=\sigma$ und $q_{\mathrm{Spin}}=\sigma/16$ lautet der duale Pfeil
\begin{equation}
f^*q_{SO}=16q_{\mathrm{Spin}}.
\end{equation}
Dies ist eine echte Abbildung mit erklärtem Definitions- und Zielbereich. Sie identifiziert keine sechzehn Knoten mit sechzehn Teilchen.

Für die hier betrachtete invertible rein gravitative Antwort in $2+1$ Dimensionen, mit Konvention $c_-=\nu/2$, verändert eine geschlossene Viermannigfaltigkeit $X$ als Füllungsdifferenz die Wirkung um
\begin{equation}\label{eq:gravextension}
\Delta S_{\rm grav}=2\pi\nu\,\frac{\sigma(X)}{16}.
\end{equation}
Spinfüllungen erlauben ganzzahliges $\nu$, weil ihre Signatur durch 16 teilbar ist. Soll die Antwort unabhängig von einer Spinwahl auf allen orientierten Hintergründen definiert sein, verlangt Füllungsunabhängigkeit $\nu\in16\Z$. Der primitive positive nichttriviale Generator hat dann
\begin{equation}
\boxed{\nu=16,\qquad c_-=8.}
\end{equation}
Die Normierung und die Unterscheidung von Spin- und orientierten Füllungen stehen ausdrücklich im Anhang von Kapustin et al.~\cite{kapustinbordism}. Der Übergang ist zunächst ein Satz über quantisierte Antworten. Seine Interpretation als mikroskopische Phasenäquivalenz benötigt die entsprechenden Voraussetzungen einer invertiblen Niedrigenergietheorie; diese Reichweite wird bei Freed--Hopkins unterschieden~\cite{freedhopkins}.

Hier ist $c_-=c_L-c_R$ die \emph{chirale Differenz}. Erst eine nachgewiesen rein chirale Realisierung mit $c_R=0$ macht daraus die Zentralzahl $c_L=8$. Zusätzliche gegenläufige Sektoren können dieselbe gravitative Antwort besitzen. Der Antwortpfeil allein liefert daher noch nicht die im Levelselektor benötigte obere Zentralzahlgrenze.

\textbf{Der entscheidende gemeinsame Test liegt außerhalb von K3.}
Auf K3 ist $\exp(i\Delta S_{\rm grav})=\exp(-2\pi i\nu)=1$ für jedes ganzzahlige $\nu$. K3 allein unterscheidet deshalb einen einzelnen Majorana-Index nicht von sechzehn. Auf dem nicht-Spin-Hintergrund $\mathbb{CP}^2$ gilt dagegen
\begin{equation}
\exp(i\Delta S_{\rm grav})=
\begin{cases}
e^{\pi i/8},&\nu=1,\\
1,&\nu=16.
\end{cases}
\end{equation}
Der zusätzliche Inhalt ist also die \emph{Existenz derselben Antwort ohne Spinwahl}. Ein Double Cover liefert diesen Abstieg noch nicht.

\picturetext{Wer nur auf Wegen prüft, deren Länge immer ein Vielfaches von 16 ist, kann Einerschritte und Sechzehnerschritte nicht unterscheiden. Erst ein zulässiger Weg der Länge eins trennt sie. K3 und $\mathbb{CP}^2$ übernehmen hier diese beiden Rollen; die gemeinte Größe ist die Signatur.}

Die Operatorwahl bleibt wesentlich. Am selben K3 liefern der komplexe Spin-Diracoperator den Index $\widehat A=2$, der Signaturkomplex den Index $-16$ und der Eulerkomplex den Index $24$. Der reelle/quaternionische Diracindex ist der primitive Wert $1$; seine Komplexifizierung ergibt $2$~\cite{absindex,atiyahsinger3}. Es gibt folglich keinen operatorunabhängigen Schluss von K3 auf den Wert 16. Ebenso würde die \emph{ganze} K3-Schnittform $U^3\oplus E_8(-1)^2$ als abelsche Chern--Simons-$K$-Matrix $c_-=-16$ liefern, nicht den einzelnen $E_8$-Wert $8$.

Auch $\Omega_4^{\mathrm{Pin}^+}\cong\Z_{16}$ darf nicht eingesetzt werden, als sei es derselbe ganzzahlige thermische Index. Hier ändern sich Symmetrietyp und Raumzeitdimension; die entsprechende Zeitumkehr ist antiunitär. Eine $\mu_4$-Markierung oder die Dimension eines Cliffordmoduls identifiziert diese Struktur nicht.

\textbf{Was damit am Ursprung wirklich fehlt.} P1 spezifiziert eine primitive Nahtwindung. Zu zeigen wäre, dass das \emph{tatsächliche} relative Determinanten-/Pfaffianfunktional diese Windung auf den primitiven positiven Generator der obigen bosonischen gravitativen Antwort abbildet. In den nachgelesenen Originalen ist das noch nicht geschehen: \file{v367_seam_s3_lattice.py} setzt $N_{\rm Maj}=2^{g_{\rm car}-1}=16$ vor der Berechnung von $c_-$ ein; \file{v470_alpha_inflow_level.py} verbindet den $U(1)$-Chernindex anschließend über denselben eingesetzten Komponentenfaktor mit $c_-$. \file{v260_k3_kummer_unification.py} prüft die K3-Gitterstruktur. Keine dieser Rechnungen wertet diesen fehlenden Antwortpfeil des Rohoperators aus.

\subsection{Wie dieser Auswahlpfeil an die schon vorhandene E8-Verklebung anschließt}
Der Anschluss benutzt die ursprüngliche markierte Struktur:
\begin{equation}
\mathcal A=(D_5)_1\otimes(A_3)_1,\qquad c_{\mathcal A}=5+3=8.
\end{equation}
Ihre $16$ irreduziblen invertiblen Sektoren sind $(p,q)\in\Z_4^2$. Die Twistphasen sind
\begin{equation}
\theta_{p,q}=\exp\!\left(2\pi i\,\frac{5p^2+3q^2}{8}\right).
\end{equation}
Die diagonale Untergruppe $H=\langle(1,1)\rangle$ hat vier Elemente. Ihre tatsächlichen niedrigsten Gewichte sind $0,1,1,1$; insbesondere sind alle Twistphasen eins. Wird die entsprechende \emph{lokale} einfache Stromerweiterung realisiert, gilt
\begin{equation}
[\mathcal B:\mathcal A]=4,\qquad
\mu_{\mathcal B}=\frac{16}{4^2}=1,\qquad
c_{\mathcal B}=8.
\end{equation}
Die Originalroute identifiziert damit in ihrer erklärten Klasse das holomorphe $E_{8,1}$-Netz. Über den Zwischenschritt $(D_8)_1$ sind die beiden lokalen Erweiterungsindizes jeweils zwei. Die alternative Spinorerweiterung von $(D_8)_1$ ist auf VOA-Ebene in der Primärliteratur explizit konstruiert~\cite{chuzheng}; der Netzübergang verwendet die bereits genannten Unitaritäts- und Lokalitätsvoraussetzungen~\cite{cklw}.

Damit verbindet sich die topologische Auswahl aus Abschnitt~\ref{sec:bordism16} mit der bereits ausgeführten Quelle: Sie liefert unter der erklärten Antwortidentifikation $c_-=8$; bei zusätzlich nachgewiesener reiner Chiralität ist dies $c=8$. Die markierte lokale Erweiterung liefert dann die tatsächlichen Felder und ihre Produkte. Ein Gravitationsantwortwert allein bestimmt jedoch noch kein vollständiges Randnetz; die Erweiterungs- und Chiralitätshypothesen bleiben Teil des Schlusses. Insbesondere bezeichnet der alte Kommentar in \file{v154_simple_current_theorem.py} die Identifikation des physischen RP-Seamnetzes mit $\mathcal A$ ausdrücklich als externen Eingang.

\textbf{Dies löst zugleich eine Verwechslung beim Materieanschluss.} Im vollen E8-Gitternetz gibt es wegen $E_8^*/E_8=0$ nur den Vakuumsektor~\cite{dongxu}. Die 248 Ströme sind Zustände dieses Sektors. Bei Einschränkung auf $\mathcal A$ zerfällt der E8-Vakuumraum in die vier Glue-Sektoren. Deren geladene Stromblöcke haben $64+60+64$ Komponenten. Das interne Spinorfeld in $(16,4)$ besitzt gemeinsames konformes Gewicht $5/8+3/8=1$ und bosonische Statistik.

Die vollen 16 Sektoren des Unternetzes besitzen dagegen teilweise fraktionelle Austauschphasen; etwa hat $(1,0)$ mit sich die Monodromie $i$. Sie sind somit nicht unverändert die symmetrische Kategorie punktlokalisierter Ladungen einer $3+1$-dimensionalen Theorie. Die Glue-Erweiterung rekonstruiert tatsächlich die Quelle; sie ist noch keine Herleitung von vierdimensionalen Fermionen oder der Standardmodell-Eichgruppe. Die nötige physische Auswahl betrifft die lokalen Observablen, ihre Ladungssektoren und deren Raumzeitwirkung gemeinsam~\cite{doplicherroberts}.

\subsection{Warum der Universalraum die lokale Zuordnung braucht}\label{sec:localmodular}
Der positive Geschichtskern aus Abschnitt~\ref{sec:constructive} liefert Zustandsvektoren und deren Überlappungen. Für die erklärte irreduzible E8-Vakuumdarstellung erzeugen sämtliche lokalen Felder global die schwache Algebra $B(\mathcal H)$. Daraus entsteht \emph{nicht} bereits eine globale modulare Zeit des Vektorzustands.

Das folgt unmittelbar: Sei $\|\Omega\|=1$, $v\perp\Omega$ und $v\ne0$. Für
\begin{equation}
A=|\Omega\rangle\langle v|\in B(\mathcal H)
\end{equation}
gilt $A\Omega=0$, aber $A^\dagger\Omega=v$. Die Vorschrift
$S(A\Omega)=A^\dagger\Omega$ wäre folglich nicht wohldefiniert. $\Omega$ ist global nicht separierend. Dies ist eine Aussage über den gegebenen reinen irreduziblen Vakuumzustand; sie betrifft nicht alle GNS-Darstellungen oder gemischten Zustände.

Beim wirklichen lokalen Intervallpaar $(\mathcal A_{E_8}(I),\Omega)$ ist $\Omega$ dagegen zyklisch und separierend. Dort existieren Tomitaoperator, Modulargruppe und die bereits konstruierte geometrische Bewegung. Die \emph{Wahl der lokalen Operationen} ist also keine nachträgliche Beschriftung eines fertigen Universalraums. Sie gehört zu den Daten, aus denen seine intrinsische lokale Zeit entsteht.

Beim Wechsel zum markierten Unternetz muss auch die Darstellung stimmen: Auf dem ganzen $\mathcal H_{E_8}=\bigoplus_{h\in H}\mathcal H_h$ erreicht $\mathcal A\Omega$ nur den Vakuumsektor $\mathcal H_0$. Dort wird die Vakuumdarstellung des Unternetzes aufgebaut. Der Wechsel der beobachtbaren Algebra und die Zustandsrekonstruktion dürfen nicht unabhängig voneinander vorgenommen werden.

\picturetext{Aus der Liste aller möglichen gemeinsamen Antworten folgt noch nicht, was ein örtlicher Beobachter selbst verändern kann. Hat man seine Operationen bestimmt, sieht derselbe Gesamtzustand für ihn anders aus. Erst an diesem Paar aus zugänglichen Operationen und Zustand setzt die modulare Bewegung an.}

\subsection{Was die bisherigen Symmetrien gemeinsam leisten -- und was sie nicht erzeugen}
Betrachtet man genau die bislang konstruierten geometrischen Möbiusbewegungen und die mit ihnen kommutierende kompakte interne E8-Wirkung, ist ihre zusammenhängende Liealgebra
\begin{equation}
\mathfrak{sl}_2(\R)\oplus\mathfrak e_{8,\mathrm{kompakt}}.
\end{equation}
Aus diesen Symmetrien allein lässt sich keine nichttriviale Wirkung der Lorentzalgebra $\mathfrak{so}(1,3)$ durch eine Liealgebraabbildung in dieses Produkt bilden.

\textbf{Beweis.} $\mathfrak{so}(1,3)$ ist reell einfach. Jede Projektion einer solchen Abbildung auf einen Faktor ist daher null oder injektiv. In $\mathfrak{sl}_2(\R)$ ist eine Injektion wegen $6>3$ ausgeschlossen. Eine Unteralgebra einer kompakten Liealgebra erbt eine positiv definite ad-invariante Form; eine nichtkompakte einfache Liealgebra wie $\mathfrak{so}(1,3)$ besitzt diese nicht. Auch die zweite Projektion ist null. Damit verschwindet die ganze Abbildung.

Dieser Satz betrifft ausdrücklich die genannte Symmetriegruppe, \emph{nicht} die gesamte unendlichdimensionale Stromalgebra. Endliche interne Clocks ändern ihre zusammenhängende Liealgebra nicht. Zusätzliche relative lokale Algebren und ihre modularen Wirkungen sind damit nicht ausgeschlossen. Der Satz erklärt präzise, weshalb immer neue Kombinationen derselben internen Markierungen keine zusätzlichen Lorentzboosts liefern.

\subsection{Der konstruktive gemeinsame Anschluss: Lokalisation bestimmt Bewegung}
Die vorhandene chiral-konforme Theorie hat bereits einen echten holographischen Anschluss. Rehrens Konstruktion setzt
\begin{equation}
\mathcal B(W)=\mathcal A(\alpha(W)),\qquad
\mathcal B(O)=\bigcap_{W\supset O}\mathcal B(W).
\end{equation}
Dabei übersetzt $\alpha$ AdS-Keile in Randgebiete. Für das vorhandene schwach additive chirale Netz entsteht so ein AdS$_2$-Netz; Vakuum und positive Energie werden übertragen. Eine entsprechende AdS$_4$-Konstruktion benötigt bereits einen $2+1$-dimensionalen konformen Rand~\cite{rehren}. Dies ist der genaue vorhandene Bulk-Anschluss, keine Herleitung unserer vierdimensionalen Physik.

Für den nächsten Dimensionsschritt ist der präziseste bekannte Ansatz modular: Gesucht wird \emph{aus den TFPT-Daten} eine Familie
\begin{equation}\label{eq:localizationtarget}
(\mathcal A_{E_8},\Omega,\text{P1/P2-Markierungen})
\ \longmapsto\
\{\mathcal M_W\}_{W\in\mathcal W_{1,3}}
\end{equation}
mit der richtigen Inklusions- und Kommutationsordnung. Der Satz von Buchholz--Dreyer--Florig--Summers rekonstruiert aus einer geeigneten nichtredundanten Keil-Algebraordnung, einem gemeinsamen zyklisch-separierenden Zustand, geometrischer modularer Wirkung, Transitivität und Netzkontinuität die Poincaré-Wirkung. Modulare Stabilität führt zur Spektralbedingung bis auf die Zeitrichtung. Nach deren positiver Festlegung sind die positiven Translationen unter den dortigen Voraussetzungen durch Netz und Zustand eindeutig~\cite{bdfs}.

Der wichtige positive Schluss lautet: Ein zusätzlich frei gewählter Hamiltonoperator ist auf diesem Weg nicht der erste fehlende Eingang. Die \emph{relative lokale Algebraordnung} kann Raumbewegung und Zeit zusammen bestimmen. Der genannte Satz setzt allerdings die vierdimensionale Keilstruktur voraus; er wählt sie nicht aus P1/P2 aus. Auch nichttriviale lokale Schnittalgebren und die gewünschten Ladungssektoren müssen für die konkrete Konstruktion nachgewiesen werden.

Die bislang als interne, vakuumerhaltende Automorphismen realisierten P1/P2-Markierungen erhalten die bestehenden Intervallalgebren. Geometrisches Möbius bewegt deren Intervalle. Beides liefert noch keine transversale Lokalisation. Genau diese zusätzliche Information wird auch in der Primärliteratur zur Lichtfrontholographie benötigt~\cite{schroer}. Die Lorentzform $\det X=t^2-\mathbf x^2$ auf $\mathrm{Herm}_2(\C)$ und die positive Dreierkomponente der K3-Schnittform sind hierfür geometrische Vergleichsstrukturen; sie liefern noch keine Wirkung auf denselben lokalen Feldalgebren.

Eine gelungene modulare Rekonstruktion wäre ein wesentlicher physischer Anschluss, aber noch kein automatischer Beweis der Standardmodell-Wechselwirkungen oder dynamischer Einstein-Gravitation. Diese müssen anschließend aus den Observablen und Antworten \emph{desselben} Netzes folgen. Ebenso beschreibt $c_-=8$ oben eine gravitative Antwort einer niedrigerdimensionalen invertiblen Theorie; sie ist nicht bereits die vierdimensionale Einsteinwirkung.

\subsection{Die ursprüngliche Prime Completion bleibt brauchbar, wenn die Algebren richtig gewählt sind}
Für eine normale unital injektive Einbettung $\iota$ einer \emph{ganzen von-Neumann-Quellenalgebra} und eine exakt autonome Zeitwirkung durch Automorphismen
\begin{equation}
\tau_t^{\rm bulk}\circ\iota=\iota\circ\tau_t^\Sigma
\end{equation}
gilt notwendig
\begin{equation}
\operatorname{vN}\!\left(\bigcup_t
\tau_t^{\rm bulk}(\iota(\mathcal A_\Sigma))\right)
=\iota(\mathcal A_\Sigma).
\end{equation}
Jedes Zeitbild ist dieselbe Algebra; ihre Vereinigung erzeugt keine größere. Eine schon abgeschlossene autonome Quelle wächst durch Wiederholung ihrer eigenen Zeit deshalb nicht zu einer größeren Algebra.

Der Originalvertrag \file{BULK.PRIME_COMPLETION} fordert jedoch bereits eine Familie von Unteralgebren. Seine endliche Prüfung \file{v1020_spectator_sector_obstruction.py} beginnt mit einem echten Randunterraum und erzeugt dessen Krylovraum unter einem gelieferten Bulkoperator. Die obige Identität widerlegt diesen Vertrag nicht. Sie bestimmt seine zwei möglichen konsistenten Lesarten: Entweder verlässt die physische Zeit eine ausgewählte echte lokale Randalgebra, oder der Bulk ist in der Gesamtalgebra bereits kodiert und benötigt eine andere, ausdrücklich konstruierte lokale Zuordnung. Die Auswahl dieser Wirkung beziehungsweise Zuordnung ist der gemeinsame fehlende Konstruktor.

\subsection{Warum die älteren Gesamtvariationen die Auswahl noch nicht ersetzen}
Die archivierte dreiteilige Gesamtkonstruktion wurde an ihren entscheidenden Eingangssätzen erneut geprüft. Ihr Defektvektor liegt in $\mathbb N^4$; ein lexikographischer Minimalwert ist unter den genannten Nichtleerheitsannahmen auswählbar. Daraus folgt jedoch noch keine Eindeutigkeit des Operators, der diesen Wert besitzt. Mehrere nichtäquivalente Operatoren können denselben Defektvektor haben. Benötigt wird die Klassifikation dieser Faser oder ein weiterer, begründeter Trennsatz.

Die anschließende Mastervariation arbeitet bereits über dem gewählten $B_{\min}$ und enthält eine Sperre, die für $B\not\cong B_{\min}$ unendlich ist. Eine solche Sperre setzt dessen Auswahl durch; sie beweist sie nicht. Die OS-Rekonstruktion setzt ihrerseits das vollständige geeignete Schwingerfunktional voraus. Für vierdimensionale Physik gehören dessen vierdimensionale Argumente, Lokalität, Kovarianz und Wachstumsbedingungen zum Eingang~\cite{os1975}. Die exakten chiralen C/X-Wortantworten ersetzen diese räumlichen Argumente nicht.

Diese Befunde decken sich mit den bereits datierten Verträgen
\file{UR.SOURCE.THREE_ROUTE_CLOSURE.01} und den Herkunftsprüfungen vom 22.~September. Die Fortsetzung behauptet keine bewiesene Nicht-Eindeutigkeit nach \emph{allen} TFPT-Bedingungen. Sie zeigt, an welcher Stelle die zitierten Eindeutigkeitsschlüsse die dafür nötige Auswahl noch voraussetzen.

\subsection{Die zusammenhängende Schlusskette und ihre genaue Reichweite}
\begin{center}
\begin{tikzpicture}[node distance=7mm]
\node[box,text width=13.2cm] (o) {\textbf{Originaler P1/P2-Prozess}\\Nahtoperator, relative Antwort, Markierungen};
\node[box,below=of o,text width=13.2cm] (t) {\textbf{Topologische Auswahl}\\Abstieg ohne Spinwahl und primitive gravitative Antwort\\$\nu=16$, $c_-=8$; passende lokale markierte Erweiterung};
\node[box,below=of t,text width=13.2cm] (s) {\textbf{Gemeinsame Quelle und Universalraum}\\E8-Feldprodukte, positives Vakuum, vollständige Geschichten};
\node[box,below=of s,text width=13.2cm] (l) {\textbf{Physische lokale Zuordnung}\\Welche Operationen gehören zusammen, welche sind räumlich getrennt?};
\node[box,below=of l,text width=13.2cm] (p) {\textbf{Antworten desselben lokalisierten Prozesses}\\Modulare Bewegung, Ladungssektoren, Wechselwirkungen und Spektren};
\draw[dasharr] (o)--(t);
\draw[arr] (t)--node[right,font=\scriptsize]{bedingt durch Erweiterung}(s);
\draw[dasharr] (s)--(l);
\draw[dasharr] (l)--node[right,font=\scriptsize]{jeweilige Rekonstruktionshypothesen}(p);
\end{tikzpicture}
\end{center}

Die gestrichelten Übergänge sind nicht durch das Diagramm bewiesen. Der erste verlangt die tatsächliche Antwortberechnung des ursprünglichen Operators. Der zweite verlangt eine aus den Markierungen ausgewählte lokale Algebraordnung. Danach werden Standardmodell, Kopplungen und Gravitation weiterhin an denselben Antworten geprüft; sie werden durch das Wort Rekonstruktion nicht vorweggenommen.

Zuses lokale Komposition und Wolframs kausale Graphidee passen als Forderung an diese Konstruktion: Zulässige Umordnungen derselben kausalen Komposition müssen dieselbe Antwort liefern. Die hierfür schon bearbeiteten Ward-/KZ-Kohärenzregeln bleiben erhalten. Kausale Konsistenz allein wählt aber weder die Übergangsregel und ihren Zustand noch die Dimension; die untersuchten Graphansätze führen Grenz- und Dynamikvoraussetzungen zusätzlich mit~\cite{gorardcausal}. Damit wird die Graphidee an derselben Gesamtstruktur geprüft, statt eine zweite Welt mit frei gewählten Regeln daneben zu setzen.

\key{\textbf{Was jetzt tatsächlich zusammenklickt:} Topologie kann die zulässige Antwortklasse auswählen; die lokale Erweiterung realisiert ihre geladenen Felder; der positive Kern stellt ihre Geschichten im Universalraum dar. Der folgende Abschnitt konstruiert zusätzlich den ursprünglichen Spin-Träger und eine von den Quadratmarken bestimmte lokale Algebraordnung innerhalb der Quelle. Damit wird die bislang nur allgemein formulierte Lokalisationsfrage teilweise geschlossen. Die Auswahl des tatsächlichen Nahtoperators und die vierdimensionale Fortsetzung bleiben gesondert zu beweisen.}

\section{Die neue gemeinsame Verbindung: Originalträger, vier Marken und dieselbe Antwortquelle}\label{sec:nativeclosure}

\key{\textbf{Fortsetzung vom 29. September: Drei zuvor getrennte Teile greifen konkret ineinander.}
Die gemeinsame Verklebung beseitigt die Spin-Obstruktion des ursprünglichen Trägers. Die vier Quadratmarken wählen zwei getrennte Nahtabschnitte ohne zusätzliche Abstandsparameter. Dieselbe vollständige E8-Quelle bestimmt deren Korrelation, eine nichttriviale modulare Mischung und eine exakt berechenbare Zweifach-Replica-Antwort. Die folgenden Herleitungen unterscheiden die jetzt konstruierten Verbindungen von den noch fehlenden physischen Identifikationen.}

\subsection{Der ursprüngliche gemeinsame Träger hebt wirklich nach Spin(16)}
Das ursprüngliche Rang-fünf-Bündel $E$ und die Familienverklebung liefern den reellen Träger
\begin{equation}\label{eq:nativeV16}
V_{16}=E_{\R}\oplus V_6,\qquad
V_6=\operatorname{Fix}\bigl(J_6:\Lambda^2F\longrightarrow\Lambda^2F\bigr).
\end{equation}
Hier ist $J_6$ die reelle Struktur der selbstdualen SU(4)-Darstellung. Die tatsächliche Darstellung lautet $(10,1)\oplus(1,6)$. Der komplexe Spin(10)-Halbspinor mit 16 komplexen Komponenten ist ein anderes Objekt; er darf nicht als dieser reelle Einteilchenträger eingesetzt werden.

Der ursprüngliche gemeinsame Lift ist eine ganze Familie. Mit $D=\det E$ und
\begin{equation}
m_j\in\Z,\qquad \sum_{j=1}^4m_j=-2,\qquad \beta_j=m_j+\frac12
\end{equation}
haben die lokalen Familiengewichte die Form $\beta_j c_1(D)$. Obwohl $F$ selbst nur projektiv gegeben sein kann, ist
\begin{equation}
\Lambda^2F_m=\bigoplus_{i<j}D^{m_i+m_j+1}
\end{equation}
global wohldefiniert. Komplementäre Paare tragen entgegengesetzte Exponenten. Drei Vertreter sind $a_i=m_i+m_4+1$, $i=1,2,3$, mit $\sum_i a_i=1+2m_4$. Damit gilt
\begin{equation}
w_2(V_6)=c_1(D)\bmod2=w_2(E_{\R}),\qquad
\boxed{w_2(V_{16})=0.}
\end{equation}
Die beiden einzeln ungeraden Stiefel--Whitney-Beiträge heben sich auf. Ein Quadratwurzelbündel von $D$ muss dafür nicht hinzuerfunden werden.

Auch die Gruppenwirkung ist explizit. Für den ursprünglichen diagonalen Zentrumsgenerator $h=(z,-iI_4)$ wirkt $h$ auf $V_{16}$ als $-I$, $h^2$ jedoch trivial. Der Träger lebt deshalb auf
\begin{equation}
K_f=\frac{\mathrm{Spin}(10)\times SU(4)}{\langle h^2\rangle}
\subset\mathrm{Spin}(16),\qquad K_f/\Z_2=K.
\end{equation}
Der vorhandene gemeinsame $U(5)$-Lift hebt bereits nach $K_f$. Dies ist ein tatsächlicher globaler Spin-Anschluss des Originalbündels. Es ist noch keine Aussage, dass der physische Nahtoperator auf genau diesem Bündel wirkt.

\picturetext{Die Fünferseite allein kommt nach einer Schleife mit einem Vorzeichenfehler zurück. Die Familienseite besitzt genau den passenden zweiten Vorzeichenfehler. Gemeinsam schließen sie sich zu einem konsistenten Träger. Die sechzehn reellen Richtungen entstehen aus dieser vorhandenen Verklebung, nicht aus sechzehn nachträglich ausgesuchten Teilchen.}

\subsection{Eine gemeinsame Formel für Eich- und Gravitationsantwort}\label{sec:nativeanomaly}
Die charakteristische Klasse des Trägers lässt sich für \emph{alle} ursprünglichen Lifts berechnen:
\begin{equation}\label{eq:nativelambda}
d=\frac12\sum_jm_j^2\in\Z_{\ge1},\qquad
\boxed{\lambda(V_{16})=\frac{p_1(V_{16})}{2}
=d\,c_1(E)^2-c_2(E).}
\end{equation}
Denn $(V_{16})_{\C}=E\oplus E^*\oplus\Lambda^2F$ und
$\operatorname{ch}_2(\Lambda^2F)=2\operatorname{ch}_2(F)$ bei $c_1(F)=0$.
Ganzzahligkeit folgt aus $m_j^2\equiv m_j\bmod2$; das Minimum $d=1$ erreichen genau die Permutationen von $(0,0,-1,-1)$.

\textbf{Bedingter, aber vollständiger Antwortsatz.}
Besitzt der tatsächliche Nahtoperator das gleichchirale reelle Spin-Dirac-Symbol mit Koeffizient $V_{16}$ und die entsprechende Pfaffian-Struktur, dann ist sein absolutes Anomaliepolynom
\begin{equation}\label{eq:nativeI4}
\boxed{
I_4^{\mathrm{Pf}}
=\frac12\bigl[\widehat A(T)\operatorname{ch}((V_{16})_{\C})\bigr]_4
=d\,c_1(E)^2-c_2(E)-\frac13p_1(T).
}
\end{equation}
Verwendet wird $\widehat A=1-p_1/24+\cdots$; Chiraliätsumkehr kehrt das gemeinsame Vorzeichen um. Die Normierung stammt aus dem Familienindexsatz für die Determinantenlinie, mit der Pfaffian-Hälfte für reelle Fermionen~\cite{freeddeterminant}. Die Gravitationsantwort entspricht für jeden Lift $c_-=8$. Ihre Berechnung benötigt die tatsächliche Symbolklasse, nicht bereits das ganze Spektrum oder jede Mehrpunktfunktion.

Auf dem ursprünglichen P2-Generator
\begin{equation}
Y=\diag\!\left(-\frac13,-\frac13,-\frac13,\frac12,\frac12\right),
\qquad x_Y=\frac{F_Y}{2\pi},
\end{equation}
verschwindet $c_1(E)$. Der nur über $\det E$ wirkende Familienanteil bleibt neutral. Somit folgt unabhängig von $m$
\begin{equation}
I_4^{\mathrm{Pf}}\big|_Y
=\frac5{12}x_Y^2-\frac13p_1(T),\qquad
\left|\frac{5/12}{-1/3}\right|=\frac54.
\end{equation}
Der bekannte Faktor $5/4$ erscheint damit \emph{in der ursprünglichen P2-Normierung} als Koeffizientenverhältnis derselben Eich-/Gravitationsantwort. Er ist kein normierungsunabhängiger neuer Invariant: Für $X=6Y$ lautet der Eichterm $15x_X^2$, und das Verhältnis wäre $45$. Ein Exponent der $\alpha$-Gleichung oder eine vierdimensionale Newtonkopplung folgt daraus noch nicht.

Eine zweite Restriktion trennt eine verführerische Gleichsetzung. Auf der Determinantenrichtung $E=\mathbf1^4\oplus L$ ist
\begin{equation}
I_4^{\mathrm{Pf}}=d\,x_D^2-\frac13p_1(T),\qquad
\boxed{k_D=2d=\sum_jm_j^2\in2\Z_{\ge2}.}
\end{equation}
Beim minimalen Lift ist $V_6=D_{\R}\oplus\R^4$; zusammen mit $E_{\R}$ gibt es zwei komplexe Ladung-eins-Kanäle. Ein einzelner ursprünglicher Pump-eins-Kanal besitzt dagegen $k_D=1$ und den Eichterm $x_D^2/2$. \emph{Kein Lift dieser ganzen natürlichen Familie ist bei identischer Kreis- und Ladungsnormierung dieser einzelne Kanal.} Primitive Spin-Klasse $d=1$ und primitiver U(1)-Pump $k=1$ sind verschiedene Aussagen.

\subsection{Der Ursprungspfeil wird auf einen konkreten Operator- und Feldanschluss verdichtet}
Der neue Bündelsatz liefert die rechte Seite der gesuchten Identifikation. Für einen wirklichen Diracoperator $D_\Sigma$ müsste sein Clifford-Multiplizitätsbündel
\begin{equation}
W_{D_\Sigma}:=\operatorname{Hom}_{\mathrm{Cl}}(S,\mathcal E_{D_\Sigma})
\end{equation}
mit $(V_{16})_{\C}$ identifiziert werden, einschließlich Basisgeometrie, Realität, Chiraliät und ursprünglicher Markierung. Ein lediglich elliptischer Operator erster Ordnung erfüllt nicht automatisch die Cliffordrelation seines Hauptsymbols~\cite{baerballmann}. Der Übergang vom allgemeinen elliptischen Collar-Generator des Archivs zu einem Spin-Dirac-Symbol muss deshalb bewiesen werden.

Die Quantisierung benötigt dann eine lokale CAR-Realisierung in der \emph{graduierten ursprünglichen Fermionen-Feldalgebra}:
\begin{equation}
\psi_{\mathrm{raw}}:L^2(\Sigma,(V_{16})_{\C})
\longrightarrow\mathcal F_{\mathrm{raw}}^{\mathrm{odd}},
\qquad C_{\mathrm{raw}}\cong P_C(D_\Sigma).
\end{equation}
Die Gleichheit der physischen Zweipunktkovarianz mit der Calderón-Projektion auf Cauchydaten ist ein OS-/Hardy-Anschluss. $P_C^2=P_C$ allein beweist ihn nicht.

Ist dieser Anschluss mit einer reinen Basisprojektion $P=P^*=P^2$, $JPJ=1-P$ gegeben, folgt der vollständige Zustand bereits eindeutig \emph{unter allen CAR-Zuständen}. Verschwindende Zweipunktnormen lassen die entsprechenden Vernichter den Vakuumvektor auslöschen; CAR-Normalordnung reduziert jedes höhere Wort auf die Zweipunktdaten. Quasi-Freiheit ist danach keine weitere unabhängige Annahme. Dies ist Arakis Projektionslemma~\cite{arakicar}.

Diese CAR-Felder liegen nicht als gewöhnliche lokale Felder im bosonischen E8-Netz: Ihre konformen Gewichte sind $1/2$. Die korrekte Folge lautet
\begin{equation}
\mathcal F_{\mathrm{raw}}^{\mathrm{even}}\cong(D_8)_1
\quad\longrightarrow\quad E_{8,1}
\end{equation}
durch die ursprüngliche lokale Spinorerweiterung vom Gewicht eins. Die $D_5+A_3$-Glue-Markierung bestimmt den beschrifteten Anschluss.

Derselbe tatsächliche Einteilchenoperator kann auch den Transfer liefern. Auf einem positiven Halbzylinder mit $D=\partial_\tau+B$, $B=B^*$ und passender reeller Struktur ist für $\tau>0$
\begin{equation}
G(\tau)=e^{-\tau B}P_+,\quad P_+=1_{(0,\infty)}(B),\qquad
T(\tau)=\Gamma(e^{-\tau h}),\quad h=B|_{\operatorname{Ran}P_+}>0.
\end{equation}
Hier ist $\Gamma$ die Fermion-Zweitquantisierung, und $T(\tau)\Omega=\Omega$. Nullmoden und Randbedingungen sind gesondert zu behandeln; auf einem thermischen Collar ist die Kovarianz im Allgemeinen keine Projektion.

Damit lässt sich ein Fehler der älteren Fixraumroute präzise korrigieren: Die Polarisation wählt Erzeuger und Vernichter; \emph{der Vakuumvektor} ist Fixpunkt des zweitquantisierten Transfers. Würde ein selbstdualer reeller Einteilchenkontraktor $T$ mit $TJ=JT$ die ganze $\operatorname{Ran}P$ punktweise fixieren, so fixierte er wegen $J\operatorname{Ran}P=\operatorname{Ran}(1-P)$ auch ihr Komplement und wäre $T=I$. Ein nichttrivialer Gap könnte so nicht entstehen. Die acht eingesetzten Einsen eines endlichen Spektrums belegen daher keine Calderón-Polarisation.

\textbf{Geometrische und analytische Reflexionen auseinanderhalten.}
Der archivierte Satz, der geometrische Decklift sei schlicht $2P_C-I$, trägt für einen lokalen geometrischen Lift nicht. Im üblichen Dirac-Collar hat die Calderón-Reflexion das Hauptsymbol
\begin{equation}
\sigma_0(2P_C-I)(x,\xi)=i\,c(\nu)c(\xi)/|\xi|.
\end{equation}
Es ist ungerade in der tangentialen Kovariable $\xi$. Ein Decklift, der die Nahtpunkte fixiert, wirkt durch eine von $\xi$ unabhängige Fasermatrix. Beide können nicht bei $\xi$ und $-\xi$ übereinstimmen. Die positive Ersatzstruktur ist die \emph{nichtlokale} Calderón-/Streureflexion: In der dazugehörigen konform kompakten Streukonstruktion gilt $S_D(0)=2P_C-I$~\cite{calderonscattering}. Dies identifiziert weder jeden kompakten Operator mit einer Streutheorie noch diese Reflexion mit Decklift oder Tomita-Konjugation.


\subsection{Die vier ursprünglichen Marken wählen eine echte lokale Zweiteilung}\label{sec:nativemodular}
Die erklärte Quadrat-Markierung wird jetzt konkret benutzt:
\begin{equation}
p_j=i^j,\qquad
E=(p_0,p_1)\cup(p_2,p_3),\qquad
E'=(p_1,p_2)\cup(p_3,p_0).
\end{equation}
Gemeint sind die positiv orientierten offenen Kreisbögen. Die beiden Komponenten von $E$ haben disjunkte Abschlüsse; ihre Zwischenräume sind genau die Komponenten von $E'$. Es wird keine zusätzliche Gapbreite gewählt. \emph{Die Quadratgestalt selbst bleibt eine ursprüngliche Rahmenvoraussetzung: Die Anzahl vier allein erzwingt die Kreuzratio $1/2$ nicht.}

\begin{center}
\begin{tikzpicture}[scale=1.05]
\draw[gray!60] (0,0) circle (1.5);
\draw[teal,line width=4pt] (4:1.5) arc (4:86:1.5);
\draw[teal,line width=4pt] (184:1.5) arc (184:266:1.5);
\draw[amber,line width=4pt] (94:1.5) arc (94:176:1.5);
\draw[amber,line width=4pt] (274:1.5) arc (274:356:1.5);
\foreach \a/\lab in {0/1,90/i,180/{-1},270/{-i}}{
\fill[navy] (\a:1.5) circle (1.8pt);
\node at (\a:1.86) {$\lab$};}
\node[align=left,text width=8.8cm] at (6.7,0) {
\textcolor{teal}{\textbf{Grün:}} Operationen in den beiden ausgewählten Abschnitten.\\[4pt]
\textcolor{amber}{\textbf{Orange:}} Operationen in ihrem Komplement.\\[4pt]
Eine Vierteldrehung vertauscht beide Seiten. Derselbe Vakuumzustand korreliert die beiden grünen Abschnitte; seine modulare Bewegung kann sie deshalb nicht durchgehend unabhängig bewegen.};
\end{tikzpicture}
\end{center}

Im vollen holomorphen E8-Vakuumnetz ist $\mu=|E_8^*/E_8|=1$. Die Zweintervall-Dualität liefert
\begin{equation}
M:=\mathcal A(E),\qquad M'=\mathcal A(E').
\end{equation}
$\Omega$ ist zyklisch und separierend~\cite{dongxu}. Es gibt daher einen wirklichen modularen Operator $\Delta_E$ und eine Tomita-Konjugation $J_E$ auf \emph{demselben} Quellhilbertraum.

Mit der vakuumnormierten geometrischen Vierteldrehung
\begin{equation}
\rho=e^{i\pi L_0/2},\qquad K_E=\log\Delta_E
\end{equation}
gelten exakt
\begin{equation}\label{eq:nativequarter}
\rho^4=I,\quad \rho M\rho^{-1}=M',\quad
\rho K_E\rho^{-1}=-K_E,\quad
[\rho,J_E]=0,\quad[\rho^2,K_E]=0.
\end{equation}
Dies folgt aus der Eindeutigkeit modularer Daten unter vakuumerhaltender unitärer Konjugation und $\Delta_{M'}=\Delta_M^{-1}$, $J_{M'}=J_M$.
Der antiunitäre Operator $\Theta=J_E\rho$ erfüllt
\begin{equation}
\Theta M\Theta^{-1}=M,\qquad
\Theta^2=\rho^2,\qquad \Theta^4=I,\qquad
\Theta K_E\Theta^{-1}=K_E.
\end{equation}
Wegen Antiunitarität bedeutet die letzte Gleichung
$\Theta\Delta_E^{it}\Theta^{-1}=\Delta_E^{-it}$.
Dagegen gilt $J_EK_EJ_E=-K_E$, aber $J_E\Delta_E^{it}J_E=\Delta_E^{it}$.

Dies ist keine automatische $\mathrm{Pin}^+$-Identifikation:
$\rho^2$ ist eine geometrische Halbdrehung, nicht Fermionparität.
Schon die bosonischen Ströme vom Grad eins erhalten unter ihr das Vorzeichen minus.
Ebenso ist die geometrische Vierteldrehung nicht ohne zusätzlichen Intertwiner die interne Compiler-, Cartan- oder Quartik-Clock.

Eine weitere echte Folge ist
\begin{equation}
J_E(V_1)\subseteq\widehat{\bigoplus}_{m\ge0}V_{4m+3}.
\end{equation}
Denn $\rho$ hat auf $V_1$ den Eigenwert $i$; die mit $\rho$ kommutierende antiunitäre Konjugation macht daraus $-i$. Die 248 Ströme und insbesondere die 24 C/X-Erzeuger sind unter dieser modularen Reflexion nicht abgeschlossen. Der bereits implementierte Nachfahrenturm ist für die gemeinsame Quelle notwendig.

\subsection{Die Quelle bestimmt einen neutralen Mischer der geladenen lokalen Felder}
Setze $N=\mathcal A((p_0,p_1))$ und $R=\mathcal A((p_2,p_3))$.
Die Split-Eigenschaft liefert $M=N\vee R\cong N\bar\otimes R$.
Die Vakuummodulargruppe kann $N$ nicht für alle Zeiten erhalten.

\textbf{Beweis.} Wäre $N$ invariant, gäbe es nach Takesaki eine vakuumerhaltende bedingte Erwartung $E_N:M\to N$~\cite{takesakiexpectation}. Für $y\in R$ liegt $E_N(y)$ wegen Bimodularität im Zentrum des Faktors $N$, also $E_N(y)=\omega(y)I$. Folglich wäre für $x\in N$
\begin{equation}
\omega(xy)=\omega(E_N(xy))=\omega(x)\omega(y).
\end{equation}
Die tatsächliche Quelle hat jedoch nichtverschwindende Cartanstrom-Zweipunktantworten zwischen den getrennten Abschnitten. Streng mit beschränkten Operatoren verwendet man deren lokale Weyl-Operatoren und differenziert die Antwort bei null. Der Vakuumzustand ist kein Produktzustand. Dies widerspricht der Folgerung; die entsprechende Mehrintervall-Obstruktion ist auch in der Primärliteratur behandelt~\cite{longomodular}.

Damit ist eine echte Bewegung außerhalb der bisherigen Kombination aus geometrischem Möbius und interner E8-Wirkung nachgewiesen. Bewiesen ist fehlende Invarianz unter der \emph{ganzen} Modulargruppe, nicht eine Veränderung bei ausnahmslos jedem $t\ne0$.

Die Mischung besitzt eine genaue operatorielle Beschreibung. Aus den schon festen Vakuummarginalen definiert die Split-Isomorphie die eindeutige Produktreferenz
\begin{equation}
\phi(xy)=\omega(x)\omega(y),\qquad x\in N,\ y\in R.
\end{equation}
Sie ist eine Rechenreferenz, keine neue physische Präparation. Ihre Modulargruppe ist das Produkt der bekannten Intervallbewegungen. Der Connes-Kokzyklus liefert~\cite{connesradon}
\begin{equation}\label{eq:nativecocycle}
u_t=(D\omega:D\phi)_t\in M,\qquad
\sigma_t^\omega=\operatorname{Ad}(u_t)\circ\sigma_t^\phi,\qquad
N_t=\sigma_t^\omega(N)=u_tNu_t^*.
\end{equation}
Quelle und ursprüngliche lokale Aufteilung bestimmen diesen Mischer vollständig als Operatorobjekt. Eine geschlossene Formel für seinen vollen E8-Kernel wurde noch nicht berechnet. Interne vakuumerhaltende E8-/P2-Automorphismen erhalten $\omega$ und $\phi$; daher ist $u_t$ unter ihnen invariant. Der Mischer transportiert die geladenen Felder, ohne eine neue interne Ladungsregel zu wählen.

Man erhält darüber hinaus auf derselben $\mathcal H$ die ausdrücklich bestimmte Familie
\begin{equation}
\mathcal A_t(I)=\Delta_E^{it}\mathcal A(I)\Delta_E^{-it},
\qquad U_t(g)=\Delta_E^{it}U(g)\Delta_E^{-it}.
\end{equation}
Jedes $\mathcal A_t$ ist ein konjugiertes positives E8-Netz mit demselben Vakuum. Vorhandene Intervallinklusionen und ihre halbseitigen modularen Beziehungen werden durch die Konjugation erhalten. Die neue Information liegt in der \emph{relativen Lage unterschiedlicher} $t$-Kopien.

Ein exakter Satz verhindert dabei eine voreilige Raumzeitableitung. Die Halbdrehung $\rho^2$ vertauscht $N$ und $R$, kommutiert aber mit $\Delta_E^{it}$. Zusammen mit $R=N'\cap M$ folgt
\begin{equation}
N_t\subseteq N_s\quad\Longrightarrow\quad N_t=N_s.
\end{equation}
Denn die Halbdrehung überträgt die Inklusion auf die zugehörigen $R$-Algebren, während relative Kommutanten sie umkehren; beide müssen gleich sein. Die Orbitfamilie allein ist keine strikt geschachtelte halbseitige modulare Inklusion.

Die nächsten tatsächlichen Objekte sind die gemischten Schnitte
$\mathcal A_t(I)\cap\mathcal A_s(J)$. Modulare Rekonstruktionssätze benötigen je nach Ziel zusätzliche Standardheits-, halbseitige Inklusions- und Reflexionsbedingungen. Wiesbrocks konkreter Intersektionssatz liefert unter seinen Bedingungen zunächst eine zweidimensionale Liegruppe~\cite{wiesbrockintersection}; er darf nicht als fertige vierdimensionale Rekonstruktion eingesetzt werden. Innerhalb jeder einzelnen Netzkopie sind die ursprünglichen Eigenschaften übertragen; zwischen verschiedenen Kopien folgen sie nicht aus der bloßen Definition. Eine $3+1$-dimensionale Keilordnung ist noch nicht konstruiert.


\subsection{Die exakte Replica-Antwort verbindet die Viermarkengeometrie mit der Clockzahl}\label{sec:nativereplica}
Die vier geordneten Quadratmarken besitzen die Kreuzratio $x=1/2$. Für allgemeine geordnete $0<x<1$ ist die an den vier Endpunkten verzweigte doppelte Überlagerung ein Torus mit
\begin{equation}
\tau(x)=i\,\frac{K(1-x)}{K(x)}.
\end{equation}
$K$ bezeichnet das vollständige elliptische Integral mit Parameter $x$. Das Quadrat entspricht $\tau=i$. Dieses Replica-Double-Cover wird aus der Intervallantwort konstruiert; seine Gleichheit mit einem anders definierten ursprünglichen physischen Double Cover wäre durch eine Abbildung zu zeigen.

Die E8-Vakuumcharakterfunktion ist
\begin{equation}
\chi_{E_8}(\tau)=\frac{E_4(\tau)}{\eta(\tau)^8}.
\end{equation}
Die bekannte Überlagerungsformel enthält einen wesentlichen Weyl-Faktor. Für den chiralen Anteil mit $c=8$ lautet sie~\cite{headrickreplica,castroanomaly}
\begin{equation}
F_{2,\mathrm{ch}}(x)
=\left[\frac{x(1-x)}{16}\right]^{c/12}\chi(\tau(x)),
\qquad
R_{2,\mathrm{ch}}(x)=(1-x)^{-c/8}F_{2,\mathrm{ch}}(x).
\end{equation}
$R_2$ ist die normierte Twist-Vierpunktfunktion, geteilt durch die beiden zugehörigen Zweipunktfunktionen. Der geordnete reelle Ast, die positive Torusspur und die Vakuum-OPE fixieren die Phase. Lokale Framefaktoren werden in Zähler und Nenner identisch geführt.

Die Thetaidentitäten
\begin{equation}
x=\frac{\vartheta_2^4}{\vartheta_3^4},\quad
1-x=\frac{\vartheta_4^4}{\vartheta_3^4},\quad
E_4=\vartheta_3^8(1-x+x^2),\quad
2\eta^3=\vartheta_2\vartheta_3\vartheta_4
\end{equation}
geben ohne Fit für die ganze Familie
\begin{equation}\label{eq:replicarational}
\boxed{F_{2,E_8}(x)=1-x+x^2,\qquad
R_{2,E_8}(x)=\frac{1-x+x^2}{1-x}
=1+\frac{x^2}{1-x}.}
\end{equation}
Am ursprünglichen Quadrat sind daher
\begin{equation}
\boxed{\chi_{E_8}(i)=12,\qquad F_2=\frac34,\qquad
R_2=\frac32,\qquad I_2:=\log R_2=\log\frac32.}
\end{equation}
Die alte skalare Lücke erhält die exakte Verbindung
\begin{equation}\label{eq:clockreplicaequality}
\boxed{\Delta_1=6\log(3/2)=6I_2,\qquad
e^{-\Delta_1}=R_2^{-6}=(2/3)^6.}
\end{equation}

\picturetext{Man vergleicht die gemeinsame Antwort zweier getrennter Nahtabschnitte mit ihren Einzelantworten. Für diese Rechnung werden zwei identische Kopien passend zusammengenäht. Die vier ursprünglichen Ecken machen daraus genau den quadratischen Torus. Die vollständige E8-Quelle liefert dort den Faktor $3/2$. Es wurde kein neuer Übergangswert eingestellt.}

Die Entkopplungsgrenze $R_2\to1$ für $x\to0$ und $F_2(x)=F_2(1-x)$ kontrollieren die Normierung. Der erste Koeffizient ist
\begin{equation}
R_2=1+x^2+x^3+\cdots,\qquad
1=\frac{248}{256}+\frac8{256}.
\end{equation}
Die Strom-Bilineare und der Orbifold-Stresstensor liefern beide Beiträge. Eine unabhängige gepaarte Theorie $E_8\times\overline{E_8}$ ergäbe $R_2=9/4$. Das ist eine Normierungskontrolle und fügt der chiralen Quelle keinen rechten Sektor hinzu.

Für Typ-III-Lokalalgebren existieren keine scharfen Intervall-Dichtematrizen. Die Formel
$R_2=\operatorname{Tr}\rho_{AB}^2/
(\operatorname{Tr}\rho_A^2\operatorname{Tr}\rho_B^2)$
ist nur mit passendem gemeinsamem Regulator zu verwenden. Die primäre Definition bleibt die endliche Twist-Ratio.
$I_2$ ist die übliche Rényi-Entropiekombination, keine Petz-Rényi-Relativentropie und nicht das volle Spektrum von $\Delta_E$~\cite{petzrenyi}.
Die Replica-Methode ist etabliert; auch E8-Replica-Orbifolds wurden bereits untersucht~\cite{e8replica}. Die neue konkrete Verbindung hier ist ihre Auswertung an den ursprünglichen TFPT-Marken.

\subsection{Dahinter steht tatsächlich ein Fermionfeld derselben replizierten Quelle}
Die Antwort lässt sich auch unmittelbar auf Operatoren zurückführen.
Die zwei Kopien besitzen die konforme Zerlegung
\begin{equation}
T_{E_{8,1}\otimes E_{8,1}}
=T_{E_{8,2}}+T_{\mathrm{Ising}},\qquad
16=\frac{31}{2}+\frac12.
\end{equation}
Entscheidend ist nicht die Zentralzahlgleichheit, sondern die tatsächliche Verzweigung des verdrehten Sektors~\cite{collazuolmelnikov}:
\begin{equation}
\chi_{E_8}(\tau/2)
=\chi^I_{1/2}\chi^{E_{8,2}}_{\mathbf1}
+\chi^I_0\chi^{E_{8,2}}_{\mathbf{3875}}
+\chi^I_{1/16}\chi^{E_{8,2}}_{\mathbf{248}}.
\end{equation}
Der eindeutige Replica-Twistgrundzustand $\Lambda$ ist ein affiner $E_{8,2}$-Singulett mit Ising-Gewicht $1/2$. Sein Vierpunkt ist der fermionische Pfaffian. In der festgelegten Reihenfolge $0,x,1,\infty$ folgt
\begin{equation}
G_4(x)=\frac1x-1+\frac1{1-x},\qquad R_2(x)=xG_4(x).
\end{equation}
Am Quadrat sind die drei kohärenten Kontraktionsbeiträge $(2,-1,2)$ und ihre normierten Beiträge
\begin{equation}\label{eq:replicaclockamplitudes}
\frac{(2,-1,2)}{G_4(1/2)}
=\left(\frac23,-\frac13,\frac23\right).
\end{equation}
Damit treten beide historischen Clockbeträge in \emph{einer} konkreten Antwort derselben replizierten Quelle auf. Das negative Vorzeichen ist fermionischer Austausch und darf nicht durch einen Betrag ersetzt werden.

Das Feld lebt im Replica-Twistsektor, nicht als zusätzliches lokales Majoranafeld des einzelnen bosonischen E8-Netzes.
Es erfüllt $\Lambda\times\Lambda=\mathbf1$~\cite{e8replica}; daher
\begin{equation}
\dim\operatorname{Hom}(\mathbf1,\Lambda^{\otimes4})=1.
\end{equation}
Die drei Wickterme sind also Beiträge zu \emph{einem} Fusionsblock, keine drei orthogonalen Zustände und keine drei positiven Übergangswahrscheinlichkeiten. Die Zahl drei der verschiedenen Orbifold-Charakterfunktionen liefert ebenfalls keinen Dreizustandsraum: Zwei der vier Primärsektoren besitzen denselben Charakter; die reduzierte Charaktermatrix ist nicht die unitäre S-Matrix der vollständigen Sektoren.

\subsection{Der zweite Clockwert entscheidet über die noch fehlende Prozessidentität}
Die Replica-Verbindung identifiziert eine konkrete gemeinsame Quellantwort, aber noch keinen alten Transferoperator. Das lässt sich direkt am vollständigen Clockpaar prüfen. Der ursprüngliche Einzelschritt auf den drei Cusp-Populationen ist
\begin{equation}
B=\frac1{18}
\begin{pmatrix}13&1&4\\1&13&4\\4&4&10\end{pmatrix},
\qquad \operatorname{spec}(B)=\{1,2/3,1/3\}.
\end{equation}
Die Moden sind $(1,1,1)$, $(1,-1,0)$ und $(1,1,-2)$; der registrierte Sechsschritt ist $B^6$. Die Originalverifikation \file{v486_transfer_full_rule.py} beschreibt die Richtung Spektrum zu Übergangsregel; \file{v814_k5_sixstep_transport.py} kennzeichnet die geprüfte Sechsschrittidentifikation als Clockaussage.

Beim Aufschneiden des Replica-Torus entsteht dagegen der Zylinderoperator
\begin{equation}
Q=e^{-2\pi(L_0-c/24)},\qquad Q_0=e^{-2\pi L_0}
\end{equation}
nach Vakuumnormierung. Auf $\Omega^\perp$ gilt $\|Q_0\|=e^{-2\pi}<2/3$. Es gibt deshalb keinen vakuumerhaltenden isometrischen Intertwiner von $B$ in diesen Transfer.

Auch eine bloße Zeitumkalibrierung $Q_a=e^{-aL_0}$ repariert dies nicht. Wegen der ganzzahligen E8-Vakuumgrade müssten für beide alten Moden positive ganze $m,n$ existieren mit
\begin{equation}
e^{-am}=2/3,\qquad e^{-an}=1/3.
\end{equation}
Dann wäre $(2/3)^n=(1/3)^m$, also $2^n=3^{n-m}$, unmöglich.
Der Ausschluss betrifft die gemeinsame reine $L_0$-Zylinderzeit, nicht jeden Quellenkanal oder jede modulare Dynamik. Ein partieller Readout mit Zustandsverlust kann andere Eigenwerte besitzen; seine Beobachtungsalgebra und Kompositionsregel sind dann zusätzliche zu begründende Daten.

Der nun präzise Prozessanschluss verlangt eine \emph{aus dem Originalprozess hergeleitete} Cusp-Auslese mit Effekten $E_i$ und eine komponierbare Operation $\mathcal S$, sodass in festgelegter Konvention
\begin{equation}\label{eq:nativeclockmap}
\mathcal S^*(E_i)=\sum_jB_{ij}E_j.
\end{equation}
Zusätzlich sind die sechs ursprünglichen aufeinanderfolgenden Anwendungen zu identifizieren. Weder die drei signierten Wickterme noch die Torusspur legen diese Auslese fest. Ein beliebiges Stinespring-Modell mit eingesetztem $B$ würde die Gleichung zwar realisieren, aber ihre Herkunft nicht herleiten.

\subsection{Was diese Verbindung am Gesamtbild ändert}
Der Universalraum enthält innerhalb der erklärten Quelle nun auch die Operationen in zwei ursprünglichen Nahtabschnitten, ihr Komplement, ihre modulare Mischung und ihre Replica-Antwort. Auf der Ursprungsseite ist der gemeinsame globale Spin-Träger konstruiert; seine charakteristische Antwort ist für die ganze natürliche Liftfamilie bekannt.

Die zusammenhängende Kette ist damit enger und konkreter:
\begin{equation}
\begin{gathered}
\text{ursprüngliche markierte Verklebung}
\longrightarrow V_{16}\ \text{mit Spin-Lift},\\
\text{tatsächlicher Naht-/CAR-Anschluss}
\longrightarrow(D_8)_1\longrightarrow E_{8,1},\\
(E_{8,1},\Omega,\text{Quadratmarken})
\longrightarrow
\bigl(M_E,\Delta_E,J_E,u_t,R_2=3/2\bigr).
\end{gathered}
\end{equation}
Die erste und die letzte Zeile sind unter ihren erklärten Voraussetzungen konstruiert. Der Eingang der mittleren Zeile und die Identifikation des historischen Transfers sind noch offene Operatorfragen. Die vierdimensionale Fortsetzung muss die relative lokale Ordnung derselben geladenen Quelle realisieren.

\picturetext{Die Bauteile haben jetzt gemeinsame Anschlüsse: Die Verklebung liefert das Trägermaterial; das Feldgesetz legt die möglichen Antworten fest; die vier Marken legen fest, welche Abschnitte zusammen betrachtet werden. Daraus entstehen Mischung und die konkrete Antwort $3/2$. Noch nicht angeschlossen ist die Vorrichtung, die diese kohärente Antwort als die drei historischen Cusp-Populationen ausliest und sechs Mal weiterführt. Deren tatsächliche Konstruktion entscheidet über den Clockanschluss.}

Zuses lokale Komposition und Wolframs Graphideen passen hier als Darstellungen derselben Operationen. Ein Graph muss die geladenen Felder, den Zustand, die fermionischen Vorzeichen und die vollständigen Antworten erhalten. Seine bloße Knotenanzahl oder ein passendes Übergangsspektrum wählt die fehlende Auslese und die Raumzeitordnung nicht aus.


\section{Prüfung des neuen Gesamtvorschlags: Die Clock als Auslese mit Gedächtnis}\label{sec:pairingcompletion}

\key{\textbf{Der neue Vorschlag enthält eine tragfähige zusätzliche Verbindung.}
Die drei Replica-Paarungen liefern exakt normierte positive Gewichte. Zusammen mit den bereits vorhandenen nativen Reflexionen entsteht daraus eine konkrete Operation, deren Erwartungswerttransfer die gesamte historische Matrix $B$ ergibt. Dafür muss kein Clockeigenwert eingesetzt werden. Dieselbe Rechnung zeigt, welche Information beim Weglassen des Ereignisgedächtnisses verloren geht und wie ein anderes Instrument diese Information erhalten kann. Die Herkunft dieses Instruments bleibt eine Auswahlfrage. Der zusätzliche Schluss von drei Paarungen auf drei Raumrichtungen trägt dagegen für die vorhandenen Kreisalgebren nicht.}

Geprüft wird der neue eingefügte Text mit Kennung
\file{55f6a61b-610f-41af-9626-87e399bbb9d1}.
Seine vorgeschlagene gemeinsame Struktur ist als Forschungsziel sinnvoll:
\begin{equation}
\mathfrak T=
(\Sigma,\mu_4,V_{16},D_\Sigma,P_C,\mathcal F_{\rm CAR},
\Omega,\mathcal A_{E_8},\{\mathcal M_O\}).
\end{equation}
Ein solches Tupel benennt zusammengehörige Daten. Es beweist noch nicht, dass die ersten Komponenten die späteren eindeutig bestimmen. Die folgenden Rechnungen prüfen gerade diese Übergänge, statt ihre Auswahl durch einen neuen Namen zu verdecken.

\subsection{Die normierten Paarungsgewichte folgen tatsächlich aus der Quelle}
Verwende die Reihenfolge der Paarungen
$m_1=(01)(23)$, $m_2=(03)(12)$, $m_3=(02)(13)$ und setze
\begin{equation}
a(x)=\left(\frac1x,\frac1{1-x},-1\right),\qquad
G(x)=\sum_k a_k(x),\qquad f(x)=1-x+x^2.
\end{equation}
Für $0<x<1$ gilt mit $a_1a_2=a_1+a_2$ exakt
\begin{equation}\label{eq:pairingnorm}
\boxed{G(x)^2=a_1(x)^2+a_2(x)^2+a_3(x)^2.}
\end{equation}
Damit sind die Gewichte
\begin{equation}\label{eq:pairingweights}
\boxed{\pi(x)=
\frac{\bigl((1-x)^2,\ x^2,\ x^2(1-x)^2\bigr)}{f(x)^2},
\qquad \sum_k\pi_k(x)=1}
\end{equation}
positiv und durch die gewählte reale Viermarkengeometrie festgelegt.
Am Quadrat ist $\pi=(4,4,1)/9$. Hier wird kein Nenner nachträglich auf die Clock angepasst.

Die Identität ist kein Sonderglück der Zahl $1/2$. Für beliebige gerade Punktzahlen besitzt der Cauchy-Kernel die bekannte Identität
\begin{equation}
\left[\operatorname{Pf}\left(\frac1{z_i-z_j}\right)\right]^2
=
\operatorname{Hf}\left(\frac1{(z_i-z_j)^2}\right).
\end{equation}
Die Hafnianseite summiert die quadrierten Paarungsprodukte ohne fermionische Vorzeichen. Dies ist die entsprechende Bosonisierungsidentität, ausdrücklich Gleichung~(19) bei Ardonne--Sierra~\cite{ardonneising}. Auf der reellen Geraden, beziehungsweise für gemeinsam auf einen Kreis abgebildete Punkte mit konsistenten lokalen Faktoren, folgt die Betragsquadratsumme. Es handelt sich um bekannte Mathematik; der neue Anschluss ist ihre Verwendung für die vorhandenen TFPT-Replica-Gewichte.

Für allgemeine komplexe Kreuzratio darf man das algebraische Quadrat nicht mit dem Betragsquadrat verwechseln:
\begin{equation}
|G(x)|^2-\sum_k|a_k(x)|^2
=-\frac{4(\operatorname{Im}x)^2}{|x(1-x)|^2}.
\end{equation}
Die ursprüngliche reale beziehungsweise Kreisgeometrie ist daher eine wesentliche Voraussetzung der positiven Gewichte.

\subsection{Die Halbierung erhält eine konkrete operative Bedeutung}
Der Anhang setzt zunächst $B=I-\alpha L$ an und bestimmt $\alpha$ aus
$1-6\alpha=R_2^{-1}$. Diese Gleichsetzung ist im ursprünglichen Vorschlag noch eine zusätzliche Prozessannahme. Sie lässt sich durch eine genau benannte Operation ersetzen.

Die frühere native Quellenrechnung enthält bereits die drei selbstadjungierten Reflexionen $R_{13},R_{14},R_{59}$ auf der tatsächlichen markierten $\C^2$-Ebene. Sie sind Einschränkungen der ursprünglichen Quellreflexionen. Auf den drei Trine-Effekten $E_1,E_2,E_3$ wirken sie als
\begin{equation}
R_{13}:(23),\qquad R_{14}:(13),\qquad R_{59}:(12).
\end{equation}
Dies ist genau die im Anhang vorgeschlagene komplementäre Kantenzuordnung.

Für eine dieser Involutionen liefert ihre nichtselektive Spektralmessung
\begin{equation}
\Pi_{k,\pm}=\frac{I\pm R_k}{2},\qquad
\mathcal D_k(X)=\sum_{\pm}\Pi_{k,\pm}X\Pi_{k,\pm}
=\frac{X+R_kXR_k}{2}.
\end{equation}
Wählt man die drei Messungen mit den Gewichten~\eqref{eq:pairingweights}, entsteht der vollständig positive, spurtreue und unitale Kanal
\begin{equation}\label{eq:replicaluders}
\boxed{\Phi_x(X)=\sum_k\pi_k(x)\mathcal D_k(X).}
\end{equation}
Die Krausoperatoren $\sqrt{\pi_k}\Pi_{k,\pm}$ zeigen Positivität und Normierung unmittelbar. Die Halbierung stammt jetzt aus dieser Messregel, nicht aus einem passend gewählten Zeitschritt.

Auf dem vorhandenen Erwartungswertsystem gilt exakt
\begin{equation}\label{eq:pairingtransfer}
\Phi_x(E_i)=\sum_j B_{ij}(x)E_j,\qquad
B(x)=\frac12I+\frac12\sum_k\pi_k(x)P_k
=I-\frac{L(x)}{2G(x)^2}.
\end{equation}
$P_k$ vertauscht die beiden anderen Labels; $L$ besitzt die komplementären Kantengewichte $w_{ij}=a_k^2$. Die gesamte Familie ist symmetrisch und stochastisch. Insbesondere
\begin{equation}
L(1/2)=
\begin{pmatrix}5&-1&-4\\-1&5&-4\\-4&-4&8\end{pmatrix},
\quad G(1/2)^2=9,
\quad
\boxed{B(1/2)=\frac1{18}
\begin{pmatrix}13&1&4\\1&13&4\\4&4&10\end{pmatrix}.}
\end{equation}
Die beiden Moden liefern $2/3$ und $1/3$. Ihre sechsten Potenzen ergeben die historischen Lücken, \emph{falls} sechs Anwendungen dieses Instruments dem ursprünglichen Sechsschritt entsprechen.

\textbf{Exakter Status:} Die Gewichte folgen aus der Replica-Antwort. Die native Reflexionswirkung und der daraus berechnete Matrixanschluss sind exakt. Die Auswahl dieser Reflexionsmessung, ihre Verbindung mit den Paarungslabels und das Verwerfen ihrer Ausgänge bleiben physische Instrumentdaten. Sie sind durch den Vierpunkt allein nicht bewiesen.

\picturetext{Die Quelle liefert drei fest bestimmte Gewichte. Drei bereits vorhandene Schalter können damit betätigt werden. Wählt man die beschriebene Betätigung, bewegen sich die drei Anzeigen exakt wie die alte Clock. Damit kennen wir jetzt eine konkrete Verbindung von Quelle und Anzeige. Noch zu begründen ist, warum der ursprüngliche Prozess genau diese Betätigung verwendet.}

\subsection{Eine Anzeige ist noch keine Kette gemessener Ausgänge}
Die drei nativen Effekte sind eine Trine auf $\C^2$:
$E_i=\tfrac23 P_i$ mit Rang-eins-Projektoren $P_i$.
Sie kommutieren nicht. Die Aussage des Anhangs, sie bildeten bereits eine kommutative Dreierauslesealgebra, ist auf diesem ursprünglichen Träger deshalb falsch.

Eine wirkliche wiederholte Trinemessung präpariert nach Ausgang $j$ den Zustand $P_j$. Für die nächste Messung erhält man
\begin{equation}
Q_{ij}=\tr(E_iP_j),\qquad
Q=\frac12I+\frac16\mathbf1\mathbf1^T.
\end{equation}
Mit $\Phi$ zwischen den Messungen folgt der tatsächlich gemessene Übergang
\begin{equation}
T_{ij}=\tr(E_i\Phi(P_j)),\qquad
\boxed{T=QB=BQ.}
\end{equation}
Am Quadrat sind seine Eigenwerte $1,1/3,1/6$.
Sie unterscheiden sich von $1,2/3,1/3$.
Eine zusätzliche kommutative Registeralgebra ist möglich; ihre Präparation und Kopplung wären aber weitere zu identifizierende Prozessdaten.

\subsection{Der neue Anschluss trifft einen bereits bekannten Punkt der ganzen Prozessfamilie}
Die exakte ursprüngliche $S_3$-Wortfamilie hat in der Reihenfolge
$I,c_{123},c_{132},p_{12},p_{13},p_{23}$ die Gewichte
\begin{equation}
w_t=\left(\frac12+t,t,t,\frac1{18}-t,
\frac29-t,\frac29-t\right),\qquad 0\le t\le\frac1{18}.
\end{equation}
Alle diese Prozesse besitzen dieselbe Matrix $B$ auf den Anzeigen. Auf dem vollständigen markierten Operatorsystem $(I,A,F,G=iAF)$ ist jedoch
\begin{equation}
\Phi_q(I)=I,\quad \Phi_q(A)=\frac23A,\quad
\Phi_q(F)=\frac13F,\quad \Phi_q(G)=qG,\qquad q=6t.
\end{equation}
Die neue Replica-Lüderskonstruktion wählt \emph{genau} $t=0$, also $q=0$. Sie ist damit an die vorhandene Quellenrechnung angeschlossen; eine neue unabhängige Hilfswelt wird nicht benötigt.

Die schon untersuchte GNS-Antwort mit linken Einsetzungen ist
\begin{equation}
\frac12\tr\!\left[
F\,\Phi_q\!\left(A\,\Phi_q\!\left(A\,\Phi_q(F)\right)\right)\right]
=\frac q9.
\end{equation}
Der neue reduzierte Kanal liefert null. Der ältere Wert $1/27$ würde $q=1/3$ verlangen, sofern er unabhängig als dieselbe ursprüngliche Quellenantwort identifiziert ist. Die Rechnung widerlegt daher eine automatische Gleichsetzung der vollständigen Prozesse; sie widerlegt nicht die neue Matrixidentität.

Das verlorene $G$ ist kein fremder Zusatzoperator. Es ist bereits aus neutralen Paaren der ursprünglichen Ströme rekonstruierbar. Die beiden zulässigen Eingaben
\begin{equation}
\rho_\pm=\frac{I\pm G}{2}
\end{equation}
haben vor der Operation die Antworten $\tr(G\rho_\pm)=\pm1$, danach jedoch beide $\Phi_0(\rho_\pm)=I/2$.
Beim beschriebenen Rang-eins-Lüdersinstrument gilt sogar für jeden einzelnen klassisch gespeicherten Messausgang
\begin{equation}
\Pi_{k,\pm}\rho_+\Pi_{k,\pm}
=\Pi_{k,\pm}\rho_-\Pi_{k,\pm}.
\end{equation}
Dieses Instrument kann deshalb keine verlustfreie Vergröberung für \emph{alle} ursprünglichen Zukunftstests sein.

\subsection{Die konstruktive Fortsetzung: Derselbe Anzeigenübergang mit erhaltenem Ereignisrecord}\label{sec:replicacarry}
Die letzte Aussage beendet die Verbindung nicht. Sie bestimmt, was die vollständige Fortschreibung zusätzlich behalten muss.

Derselbe reduzierte Kanal lässt sich durch ein anderes natives Instrument realisieren: Mit Wahrscheinlichkeit $1/2$ wird $I$ angewendet, mit Wahrscheinlichkeit $\pi_k/2$ die tatsächliche Reflexion $R_k$. Der angewendete Operator wird als klassisches Ereignislabel gespeichert. Seine unnormierten Ausgangszustände sind
\begin{equation}
\mathcal J_0(\rho)=\frac12\rho,\qquad
\mathcal J_k(\rho)=\frac{\pi_k}{2}R_k\rho R_k.
\end{equation}
Vergisst man das Label, folgt wieder exakt $\Phi_x$. Behält man es, sind sämtliche Eingangszustände auf dieser markierten Ebene durch kontrollierte Rückführung rekonstruierbar:
\begin{equation}\label{eq:recordrecovery}
\boxed{\sum_e U_e^*\mathcal J_e(\rho)U_e=\rho,\qquad
U_0=I,\quad U_k=R_k.}
\end{equation}
Dies ist keine Messung des Reflexionseigenwerts. Es ist eine protokollierte Anwendung der Reflexion. Deshalb widerspricht die Rückführung nicht dem Informationsverlust des vorherigen Lüdersinstruments.

Die fehlende $G$-Antwort ist dabei besonders einfach lokalisiert. Die Reflexionen kehren $G$ um. Schreibe $\varepsilon_0=+1$ und $\varepsilon_k=-1$. Dann gilt
\begin{equation}
\boxed{\sum_e\varepsilon_e\,\tr\bigl(G\,\mathcal J_e(\rho)\bigr)
=\tr(G\rho).}
\end{equation}
Die Korrelation zwischen Ereignisparität und Quelle trägt genau die Antwort, die nach dem Wegwerfen des Labels verschwindet. Für Folgen ausschließlich dieser geschlossenen nativen Ereignisse ist ihr Produkt ein Element des bereits berechneten $S_3$; das Gruppenlabel erlaubt die kontrollierte Umkehr der ganzen Folge. Dies behauptet keinen Sechszustandsspeicher für beliebige geladene C/X-Geschichten oder den vollständigen Nachfahrenturm.

\begin{center}
\begin{tikzpicture}[node distance=7mm]
\node[box,text width=13cm] (s) {Dieselbe Replica-Antwort und dieselben nativen Reflexionen};
\node[box,below=of s,text width=13cm] (p) {Protokollierte Operation: Quellzustand \textbf{plus angewendetes Ereignis}};
\node[box,below=10mm of p,xshift=-3.6cm,text width=6cm] (a) {Record mitführen:\\kontrollierte Rückführung;\\$G$ bleibt unterscheidbar};
\node[box,below=10mm of p,xshift=3.6cm,text width=6cm] (b) {Record verwerfen:\\reduzierter Kanal $\Phi_0$;\\Anzeigen folgen $B$, $G$ verschwindet};
\draw[arr] (s)--(p);\draw[arr] (p)--(a);\draw[arr] (p)--(b);
\end{tikzpicture}
\end{center}

\picturetext{Eine zufällige Drehung verwischt das Bild, wenn niemand aufschreibt, welche Drehung ausgeführt wurde. Mit dem Drehprotokoll kann man jedes Bild zurückdrehen. Eine wirkliche Messung, die vorher beide Bilder auf dasselbe Ergebnis zusammendrückt, lässt sich dagegen durch das Aufschreiben dieses Ergebnisses nicht rückgängig machen. Beide Vorgänge können dieselbe gemittelte Anzeige erzeugen. Für die Gesamtlösung zählt deshalb der Vorgang mit seinem Gedächtnis, nicht nur die Anzeige.}

Hier treffen Clock, Universalraum und die Ereignisgraphidee konkret zusammen:
Ein Graph kann die ursprünglichen Operatoren und ihre Records komponieren. Der Universalraum muss alle später unterscheidbaren Antworten dieses gewählten Operationssystems erhalten. Die Clock darf eine reduzierte Anzeige dieses reicheren Prozesses sein. \emph{Welche} dieser Instrumente der ursprüngliche Nahtprozess ausführt, ist weiterhin herzuleiten; die neue Rechnung liefert zwei explizite, anhand ihrer Zukunftsantworten verschiedene Kandidaten.

Auch identische Diagonalnormen des Vierpunkts ersetzen diese Prüfung nicht. Für zwei reelle Geometrieparameter gilt
\begin{equation}
G(x)G(y)-a(x)\cdot a(y)
=\frac{(x-y)^2}{x(1-x)y(1-y)}.
\end{equation}
Die normierten Paarungsvektoren haben bei $x=1/2,y=1/5$ das Skalarprodukt $6/7$, ihre skalaren Summen jeweils eins. Daraus folgt allein: Die Normidentität~\eqref{eq:pairingnorm} identifiziert nicht automatisch einen ganzen bilinearen Kernel. Verschiedene Geometrien werden hier nicht zu identischen physischen Zuständen erklärt; der tatsächliche GNS-Mischkernel benötigt seine eigenen Quelleinsetzungen.

\subsection{P1 als primitive Symbolklasse: sinnvolle Präzisierung, keine automatische Operatorwahl}
Der neue Text schlägt vor, primitive Windung durch eine primitive reelle K-Homologie- oder Dirac-Symbolklasse zu ersetzen. Dafür sind zunächst die Typen festzulegen. Auf dem Nahtkreis besitzt die unverdrehte Fundamentalklasse einen primitiven freien Anteil. Ein orientiertes reelles $V_{16}$ ist als Bündel auf dem Kreis trivial; die Verdrehung mit Rang 16 vervielfacht diesen freien Anteil mit 16. Eine primitive \emph{unverdrehte} Klasse und ein primitiver \emph{bereits mit $V_{16}$ verdrehter} Operator sind daher nicht dieselbe Aussage. Eine äquivariante oder Morita-normalisierte Variante müsste eigens definiert werden.

Selbst bei identischem Hauptsymbol, Koeffizientenbündel und Vakuum bleibt eine konkrete Auswahlfrage. Auf dem vorhandenen antiperiodischen Nahtkreis betrachte
\begin{equation}
B_a=-i\partial_\theta+aY_{16},\qquad |a|<1.
\end{equation}
$Y_{16}$ ist die ursprüngliche P2-Ladungswirkung: auf dem komplexifizierten $E_\R$ die fünf Ladungen und ihre Gegenteile, auf $V_6$ null. Die kanonische Realstruktur vertauscht die beiden Ladungsrichtungen und erfüllt $JB_aJ=-B_a$. Für NS-Moden $r\in\Z+1/2$ und $|q|\le1/2$ gilt
\begin{equation}
\operatorname{sign}(r+aq)=\operatorname{sign}(r).
\end{equation}
Somit sind Hauptsymbol, dessen Klasse und die positive Polarisation für die ganze Familie gleich. Der damit bestimmte reine CAR-Zustand ist derselbe. Dennoch unterscheiden sich die positiven zeitlichen Antworten:
\begin{equation}
G_{a,q}(\tau)=\sum_{n\ge0}e^{-\tau(n+1/2+aq)}
=\frac{e^{-aq\tau}}{2\sinh(\tau/2)}.
\end{equation}
Bei $a=1/2$ erhält der Ladung-$1/2$-Kanal den Faktor $e^{-\tau/4}$, der neutrale Kanal nicht. Es handelt sich also nicht bloß um einen gemeinsamen Wechsel der Zeiteinheit.

Auch auf der vollen vorhandenen E8-Quelle gibt es die entsprechende Richtung $H_a=L_0+aY_0$. Aus $\|Y\|^2=5/6$ und den ganzzahligen positiven Vakuumgraden folgt für $|a|<\sqrt{3/5}$ ein positives Spektrum außerhalb desselben eindeutigen Vakuums. Die festen geometrischen Marken kommutieren weiterhin mit dieser Bewegung. \emph{Nicht} gleich bleiben die flache P2-Holonomie oder die Behauptung, die physische Zeit sei genau der geometrische beziehungsweise modulare Fluss.

Das ist ein Gegenbeispiel zur Auswahl aus Symbolklasse plus Positivität allein, kein bereits gegen alle TFPT-Constraints geprüfter Ersatzprozess. Eine ursprüngliche feste Holonomie oder eine bewiesene Identifikation mit der lokalen Modularzeit könnte diese Freiheit beseitigen. Genau eine solche zusätzliche ursprüngliche Auswahl muss der vorgeschlagene P1-Satz noch liefern.

\subsection{Drei Paarungen liefern keine drei räumlichen Achsen}
Setze $I_j=(p_j,p_{j+1})$ für die vier zyklischen Marken. Die ersten beiden Paarungen geben die bereits vorhandenen komplementären Mengen
\begin{equation}
M=\mathcal A(I_0)\vee\mathcal A(I_2),\qquad
M'=\mathcal A(I_1)\vee\mathcal A(I_3).
\end{equation}
Ihre modularen Daten sind miteinander verknüpft; sie sind keine zwei unabhängigen räumlichen Richtungen. Insbesondere $M\cap M'=\C I$ ist auf dem nichttrivialen Vakuumraum kein zyklischer Schnitt.

Die kreuzende Paarung $(02)(13)$ besitzt keine zwei disjunkten verbindenden Kreisbögen. Wählt man die Halbkreise
$K=\mathcal A((p_0,p_2))$, $L=\mathcal A((p_1,p_3))$,
folgen aus den Intervall-Netzeigenschaften
\begin{equation}
K\cap L=\mathcal A(I_1),\qquad
K\vee L=\mathcal A((p_0,p_3)).
\end{equation}
Dies liefert echte lokale Schnitte, aber zunächst die gewöhnlichen Intervallinklusionen der bestehenden chiralen Möbiusordnung. Ihre Generatoren liegen in derselben $\mathfrak{sl}_2(\R)$; dort existieren keine drei unabhängigen paarweise kommutierenden Translationen.

Die Vierteldrehung wirkt zudem nur als
\begin{equation}
m_1\longleftrightarrow m_2,\qquad m_3\longmapsto m_3.
\end{equation}
Auf den Paarungslabels hat sie Ordnung zwei. Die geometrische Halbdrehung ist auf ihnen bereits unsichtbar. Das reine Dreierregister trägt deshalb nicht die volle ursprüngliche $\mu_4$-Wirkung; auch hier wäre die verlorene Phaseninformation gesondert mitzuführen.

Schließlich sind drei positive Generatoren nicht automatisch drei gewöhnliche Raumimpulse. Eine räumliche Rotation kann $P_j$ in $-P_j$ überführen. Wären beide positiv, wäre $P_j=0$. Die Positivitätsbedingung gehört zu zukünftigen zeitartigen oder lichtartigen Kombinationen, etwa $P_0\pm P_j$, und muss zusammen mit den Lorentzrelationen nachgewiesen werden.

Der tragfähige weitere Weg bleibt die schon aus der Quelle bestimmte Familie
\begin{equation}
\mathcal G=\langle U(\text{Möb}),\Delta_E^{it}\rangle,\qquad
\mathcal A_g(I)=g\mathcal A(I)g^*,\quad g\in\mathcal G.
\end{equation}
Die nicht rein geometrische Zweintervall-Modularbewegung eröffnet darin neue relative Positionen. Für Schnitte verschiedener Kopien müssen Standardität, halbseitige Inklusionen und gemeinsame Gruppenrelationen tatsächlich bewiesen werden. Die Anzahl der Ausgangspaarungen übernimmt diese Arbeit nicht.

\subsection{Was der gemeinsame Abschluss jetzt tatsächlich verlangen muss}
Die Gesamtidee des Anhangs wird präziser, wenn sie dieselben Operatoren und dasselbe Gedächtnis durch alle Auslesungen verfolgt:
\begin{enumerate}
\item Die ursprünglichen Daten wählen einen \textbf{konkreten} Nahtoperator samt Markierung und Verbindungsdaten. Eine bloße Symbolklasse genügt dafür nicht.
\item Dieser Operator trägt die bereits konstruierte $V_{16}$-Realstruktur und liefert die tatsächliche CAR-Polarisation; ihre gerade Algebra und markierte Erweiterung müssen mit derselben E8-Quelle identifiziert werden.
\item Der ursprüngliche Prozess wählt ein \textbf{Instrument einschließlich Ereignisrecord}. Die neue Rechnung zeigt einen positiven nativen Anschluss an $B$ und eine reversible Version mit Record. Der bisherige gemittelte Transfer entscheidet zwischen diesen Vorgängen nicht.
\item Die lokalen Algebren derselben Quelle besitzen eine relative Ordnung, welche die verlangten Raumzeitrelationen realisiert. Drei Pairings allein liefern diese Ordnung nicht.
\end{enumerate}
Erst auf einer solchen Raumzeitstruktur sind lokale Ladungssektoren, Lorentzspin, Massenspektren und Wechselwirkungen gemeinsam auszuwerten. Ein transportabler Ladungssektor ist dabei nicht schon ein einzelnes Teilchen: Er kann Mehrteilchenzustände und verschiedene Massen enthalten. Ebenso sind Fusionsregeln noch keine bestimmten Kopplungskonstanten. Die bereits markierte P2-Zentralisatorstruktur bleibt eine konkrete interne Vorgabe, deren vierdimensionale Eichrealisierung zu zeigen ist.

\key{\textbf{Der neue Fortschritt am Gesamtbild:}
Es gibt jetzt eine exakt berechnete Verbindung zwischen Replica-Geometrie, nativen Reflexionen und der ganzen alten Clockmatrix. Außerdem ist an denselben Operatoren nachgewiesen, wie die Clock als reduzierte Anzeige eines Prozesses mit erhaltener Information auftreten kann. Das bindet Universalraum, Ereignisgraph und Clock enger zusammen. Eine behauptete eindeutige Urregel muss künftig dieses Instrument und seinen Record auswählen; sie darf nicht bloß den passenden Anzeigenverlauf reproduzieren.}

\section{Anker, Viertelphase und Ladungsfasern: dieselbe Clockfamilie in zwei Darstellungen}\label{sec:anchorfibre}

\key{\textbf{Die nächste Verbindung ist eine Identität für die ganze Vierpunktfamilie.}
Die ursprünglichen Replica-Paarungen und der neue Anker-Phasenschritt erzeugen exakt dieselben Matrizen $B(x)$. Am Quadrat liefert dies den ursprünglichen Anker $1:1:2$. Der kohärente Schritt besitzt eine ausdrückliche Einbettung in die markierte C-Quelle. Die ältere Trine-Auslese besitzt wiederum eine genaue orthogonale Dreier-Erweiterung. Damit sind mehrere bisher getrennte Beschreibungen mathematisch verbunden. Ihre Instrumente und ihre ursprünglichen Träger dürfen dabei weiterhin unterschiedliche Zukunftsantworten haben.}

Dieser Abschnitt prüft den eingefügten Text
\file{14360bff-eb1f-42b4-a10f-04c973ecd011}.
Die dort verlinkten Dateien unter \file{sandbox:/mnt/data/} sind im übergebenen Material nicht enthalten. Die angegebenen 75 Prüfbedingungen werden deshalb nicht als hier ausgeführt übernommen. Die folgenden Formeln wurden aus dem Text und den lokalen Originalquellen eigenständig rekonstruiert und im vorhandenen Rechenkern geprüft.

\subsection{Die beiden Clockvorschläge stimmen für alle reellen Viermarkengeometrien überein}
Benutze dieselben normierten signierten Beiträge wie zuvor,
\begin{equation}
s_i(x)=\frac{a_i(x)}{G(x)},\qquad
\sum_i s_i=1,\qquad \sum_i s_i^2=1,\qquad 0<x<1.
\end{equation}
Definiere daraus
\begin{equation}
p_i=\frac{1-s_i}{2},\qquad
f=1-x+x^2,\qquad
p=\frac{(x^2,(1-x)^2,1)}{2f}.
\end{equation}
Dann folgen ohne Zahlenanpassung
\begin{equation}\label{eq:anchorreplicajoint}
\sum_i p_i=1,\qquad \sum_i p_i^2=\frac12,\qquad
\boxed{s_k^2=4p_ip_j\quad\text{für }\{i,j,k\}=\{1,2,3\}.}
\end{equation}
Für die letzte Gleichung benutzt man
$s_1s_2+s_1s_3+s_2s_3=0$.

Der normierte Vektor und sein Viertelphasenschritt lauten
\begin{equation}
v(x)=\frac{(x,1-x,1)^T}{\sqrt{2f}},\quad
P(x)=v(x)v(x)^*,\quad
U(x)=I+(i-1)P(x).
\end{equation}
Wegen $P^2=P$ ist $U$ unitär und $U^4=I$.
Außerhalb der Diagonale gilt $|U_{ij}|^2=2p_ip_j$.
Mit~\eqref{eq:anchorreplicajoint} ergibt dies
\begin{equation}\label{eq:wholeclockfamily}
\boxed{|U(x)|_{\rm eintragsweise}^2
=I-\frac{L(x)}{2G(x)^2}
=B(x).}
\end{equation}
Die rechte Seite ist genau der im vorangehenden Abschnitt aus den Replica-Gewichten rekonstruierte Transfer. Am Quadrat wird
\begin{equation}
v(1/2)=\frac{(1,1,2)^T}{\sqrt6},\qquad
U(1/2)=\frac16
\begin{pmatrix}
5+i&-1+i&-2+2i\\
-1+i&5+i&-2+2i\\
-2+2i&-2+2i&2+4i
\end{pmatrix}.
\end{equation}
Die Viertelphase bleibt für die ganze Familie gleich. Die geometrische Kreuzratio verändert die Ankeramplituden.

Die eingeschränkte Umkehrung des Anhangs stimmt ebenfalls. Für
$U=I+(e^{i\theta}-1)vv^*$ erzwingen die drei historischen Nebendiagonaleinträge
$p_1=p_2$, $p_3=4p_1$ und damit $p=(1,1,4)/6$ sowie
$\cos\theta=0$. Die Vorzeichenwahl $\theta=\pm\pi/2$ und die Komponentenphasen von $v$ bleiben durch $B$ unbestimmt. Dies ist eine Eindeutigkeitsaussage innerhalb der Rang-eins-Phasenklasse, keine Auswahl dieser Klasse aus P1/P2.

\picturetext{Die beiden Vorschläge beschreiben dieselbe Anzeige von zwei Seiten. Die Paarungsrechnung liefert die Gewichte der sichtbaren Übergänge. Der Anker beschreibt eine kohärente Drehung, deren Betragsquadrate genau diese Gewichte ergeben. Diese Verbindung gilt beim Verformen der vier Marken weiter; sie beruht nicht nur auf einer einzelnen passenden Zahl am Quadrat.}

\subsection{Die vorhandene markierte Quelle trägt die Viertelbewegung tatsächlich}
Die im Originalvertrag bereits bewiesenen Charaktere sind
\begin{equation}
k(q)=2\sum_{j=0}^4q_j,\quad
a(q)=q_7-q_6,\quad
t(q)=4q_5+3q_6+5q_7,
\qquad T=G^2A.
\end{equation}
Ihre Wirkung wird aus den tatsächlichen ursprünglichen C-Wurzeln ausgelesen. In der Reihenfolge $C_0,C_1,C_2,C_3$ folgt
\begin{equation}
A_C=\diag(1,1,i,-i),\qquad
T_C=-A_C.
\end{equation}
Der Drei-Unterraum $C_0,C_1,C_2$ ist unter $T_C$ invariant.
Für eine orthogonale Matrix $R$ mit $Re_3=v(x)$ und die natürliche Inklusion
$\iota:\C^3\to\C^4$ setze $V=\iota R^T$. Dann
\begin{equation}
V^*V=I,\qquad
\boxed{T_CV=-VU(x).}
\end{equation}
Eine explizite Wahl verwendet
$n=(1,1,-1)^T/\sqrt3$ und $R=(n\times v,n,v)$.
Hier ist $n\cdot v(x)=0$ für die gesamte Familie.

Dies ist ein echter algebraischer Intertwiner. Seine Konstruktion wählt aber die orthogonale Detektorbasis $R$; die ursprüngliche Cusp-Messung wurde damit noch nicht aus dem Nahtoperator hergeleitet. Die Detektorprojektoren mischen verschiedene Familiengewichte. Wären sie bloß diagonale Funktionen der Familienladungen, würden sie mit $T$ kommutieren und überhaupt keine nichttrivialen Übergänge erzeugen. Die physische Auslese muss deshalb kohärente Familienkombinationen zugänglich machen.

Die ursprüngliche Familienreflexion vertauscht die letzten beiden C-Gewichtslinien, einschließlich ihrer vorgegebenen Kokzyklusphasen. Der Dreier muss für diesen Operationssatz zum Vierer erweitert werden. Für \emph{alle} ursprünglichen Stromoperationen reicht auch der Vierer nicht: Schon
$[C_0^\dagger,C_0]$ liegt im Cartanraum außerhalb des C-Vierers. Da die 24 Felder samt ihren Adjungierten die einfache Algebra $E_8$ erzeugen, ist die unter allen ihren Nullmoden erzeugte Hülle eines nichtverschwindenden Stromzustands die gesamte 248-dimensionale adjungierte Darstellung. Mit den übrigen Strommoden kommt der unendliche Nachfahrenturm hinzu.

Das gibt dem Universalraumprinzip einen konkreten Umfang:
Die nötige Hülle wächst mit den tatsächlich zugelassenen Fortsetzungen.
Die Folge $3\to4\to248\to\mathcal H_{E_8}$ bezeichnet hier unterschiedliche Operationsabschlüsse, keine Herleitung von Raumzeitdimensionen.

Die geometrische Vierteldrehung ist ebenfalls getrennt zu halten. Auf allen Stromzuständen vom Grad eins wirkt $e^{i\pi L_0/2}=iI$. Der markierte Charakter $T_C$ besitzt dagegen verschiedene Eigenwerte. Seine konkrete Einbettung beseitigt diesen Unterschied nicht.

\subsection{Die alte Trine und die neue Dreiermessung haben einen exakten gemeinsamen Anschluss}
Die bereits vorhandenen nativen Effekte auf $\C^2$ sind genau
\begin{equation}
E_i=W^*P_iW,\qquad
W=\begin{pmatrix}
1/\sqrt2&1/\sqrt6\\
-1/\sqrt2&1/\sqrt6\\
0&-2/\sqrt6
\end{pmatrix},
\qquad P_i=|i\rangle\langle i|.
\end{equation}
Es gilt
\begin{equation}
W^*W=I_2,\qquad WW^*=I_3-\frac13\mathbf1\mathbf1^T.
\end{equation}
Damit ist die Trine eine Einschränkung der orthogonalen Dreierauslese, eine konkrete minimale Naimark-Erweiterung. Die Effekte sind aus den ursprünglichen Strahlen 13, 14 und 59 nachgerechnet, nicht durch passend gewählte neue Projektoren ersetzt.

Die Erweiterung erklärt zugleich die bereits gefundene Messgrenze. Nach einem orthogonalen Ausgang $P_i$ verbleibt bei Rückprojektion in den alten Zweierraum die Spur
\begin{equation}
\tr(W^*P_iW)=\frac23.
\end{equation}
Der zusätzliche Detektoranteil trägt $1/3$. Eine normierte Rückpräparation in den Zweierraum und erneute Trinemessung bringt genau die Matrix
$Q=\tfrac12I+\tfrac16\mathbf1\mathbf1^T$ hinein. Daher bleibt $QB$ der Übergang bei diesem Resetablauf. Lässt man die orthogonale Dreierpräparation bestehen, ist ein anderer Ablauf möglich.

Die Gleichung der Effekte ist ausdrücklich \emph{kein} Intertwiner der ganzen Quellenwirkung. Der ursprüngliche Trine-Zweierraum besteht aus Cartan-Stromzuständen; der Charakter $T$ fixiert diese. Die neue C-Dreier-Einbettung benutzt Wurzelstromzustände mit nichtkonstanten $T$-Phasen. Eine Identifikation der physischen Prozesse benötigt somit weiterhin die tatsächliche Präparation, Kopplung und Rücksetzung des Detektors.

\subsection{Die ausgeschlossenen Kohärenzen liefern den Gedächtniskern ausdrücklich}
Für die erklärte orthogonale Auslese ist das Instrument
\begin{equation}
\mathcal S=\mathcal D\circ\operatorname{Ad}_U\circ\mathcal D,\qquad
\mathcal D(X)=\sum_iP_iXP_i
\end{equation}
vollständig angegeben. Seine Krausoperatoren
$K_{ij}=U_{ij}|i\rangle\langle j|$ erfüllen
$\sum K_{ij}^*K_{ij}=I$ und
\begin{equation}
\mathcal S^*(P_i)=\sum_jB_{ij}P_j.
\end{equation}
Wird diese Dephasierung zwischen jedem Schritt tatsächlich vollzogen, folgen die Populationen $B^n$. Sechs Anwendungen ergeben die bekannten Clocklücken; ihre physische Zuordnung zu genau sechs ursprünglichen Ereignissen bleibt gesondert zu zeigen.

Ohne solche Zwischenoperationen gilt stattdessen
\begin{equation}
C_n=|U^n|_{\rm eintragsweise}^2
=B+(I-B)\cos(n\pi/2),\qquad C_4=I.
\end{equation}
Für den anfänglich ersten Detektor ist die Vier-Schritt-Rückkehr kohärent eins, mit Zwischen-Dephasierung dagegen $(B^4)_{11}=211/486$.
Eine Aufzeichnung muss hier eine tatsächliche Unterscheidung der Wege mit anschließender Reduktion bezeichnen. Bloßes Betrachten einer Formel erzeugt keine Dephasierung.

Der im Anhang angesprochene Feshbach-/Gedächtnisanschluss lässt sich für diesen konkreten Quellenoperator vollständig berechnen. Zerlege den neundimensionalen Operatorraum in die drei Diagonalen und sechs Kohärenzen. Für
$\mathsf U=\operatorname{Ad}_U$ schreibe die Blöcke
\begin{equation}
\mathsf U=
\begin{pmatrix}B&X\\Y&Z\end{pmatrix}.
\end{equation}
Elimination der Kohärenzen gibt im formalen Resolventengebiet um $z=0$
\begin{equation}
\mathcal G(z)=\sum_{n\ge0}C_nz^n
=\bigl[I-zB-z^2M(z)\bigr]^{-1},
\qquad M(z)=X(I-zZ)^{-1}Y.
\end{equation}
Aus der exakten Folge folgt gleichzeitig
\begin{equation}
\mathcal G(z)=\frac{B}{1-z}+\frac{I-B}{1+z^2}.
\end{equation}
Für jede Eigenmode $b$ von $B$ ist damit der ganze Gedächtniskern
\begin{equation}\label{eq:anchormemory}
\boxed{M_b(z)=
\frac{-(1-b)^2+(1-b^2)z}{1-(1-b)z+bz^2}.}
\end{equation}
Dies wurde unabhängig durch den Schurkomplement-Ausdruck der tatsächlichen $9\times9$-Operatorwirkung geprüft. Insbesondere ist $M(0)=-(I-B)^2$. Diese kohärente Rückwirkung ist kein zusätzlicher positiver Sprungprozess.

\picturetext{Die unsichtbaren Phasen sind nicht einfach weg. Ohne Messung wirken sie später wieder auf die Anzeigen zurück und ermöglichen die vollständige Rückkehr. Der Gedächtniskern berechnet genau diese Rückwirkung. Wird nach jedem Schritt dephasiert, entfernt dieser andere Ablauf die Rückwirkung und erzeugt den relaxierenden Transfer.}

Eine weitere Konsequenz folgt ohne Modellergänzung: Ein festes geschlossenes endliches unitäres System kann einen strikt relaxierenden $B^n$-Verlauf nicht für alle $n$ durch dieselbe feste Reduktion erzeugen. Seine unitäre Entwicklung kehrt auf einer Teilfolge beliebig nahe zur Identität zurück, während $B^n$ gegen die stationäre Projektion konvergiert. Ein autonomer Prozess mit dauerhafter solcher Relaxation benötigt daher etwa erneuerte Records, eine unendliche Umgebung oder einen entsprechenden Grenzübergang. Die vollständige E8-Quelle besitzt bereits unendlich viele Moden; ihre tatsächliche Rolle als Umgebung ist dadurch noch nicht ausgewählt.

\subsection{Die Dreier-Gitterbeschreibung ist eine exakte Ladungsfaserung}
In den ursprünglichen $D_5+D_3$-Koordinaten ist
\begin{equation}
0\longrightarrow D_5\longrightarrow E_8
\overset{\pi_F}{\longrightarrow}D_3^*\longrightarrow0,\qquad
D_3^*=A_3^*=\Z^3\cup(\Z+\tfrac12)^3.
\end{equation}
Der Kern besteht aus ganzzahligen Fünfervektoren gerader Summe.
Für jede ganzzahlige oder gleichmäßig halbganzzahlige Dreierkoordinate kann die Parität der fünf anderen Koordinaten passend ergänzt werden; die Projektion ist surjektiv. Jede Faser ist ein bestimmter $D_5$-Koset. Die Folge zerfällt als Folge freier abelscher Gruppen, aber nicht automatisch als orthogonales Produkt der markierten Ladungsgitter.

Die vollständige Quelle lässt sich deshalb verlustfrei umordnen:
\begin{equation}
\mathcal H_{E_8}=\bigoplus_{y\in D_3^*}\mathcal K_y,\qquad
\mathcal K_y=\bigoplus_{x:(x,y)\in E_8}\mathcal F_{(x,y)}.
\end{equation}
Die inneren Ladungen, Kokzyklusphasen und alle Oszillatormoden bleiben in den Fasern erhalten. Die vier tatsächlichen Familiengewichte
\begin{equation}
f_0=\tfrac12(-1,-1,-1),\quad f_1=\tfrac12(-1,1,1),\quad
f_2=\tfrac12(1,-1,1),\quad f_3=\tfrac12(1,1,-1)
\end{equation}
erfüllen
\begin{equation}
\sum_a f_a=0,\qquad f_a\cdot f_b=\delta_{ab}-\tfrac14,
\qquad\sum_a f_af_a^T=I_3.
\end{equation}
$C_a$ und alle $X_{ia}$ verschieben $y$ um dasselbe $f_a$, ihre vollständigen Ladungen unterscheiden sich aber in der Fünferfaser.

Die vier getrennten radialen C-Einsetzungen ergeben einen besonders deutlichen Zeugen:
\begin{equation}
\sum_a\beta(C_a)=(-2,-2,-2,-2,-2;0,0,0).
\end{equation}
Der projizierte Weg ist geschlossen. Der erzeugte Quellzustand besitzt jedoch eine andere Ladung als das Vakuum, ist zu ihm orthogonal und hat Mindestgrad zehn. Für die konkreten Radialpositionen
$(4/5,3/5,2/5,1/5)$ liefert der bestehende vollständige Geschichtskern eine strikt positive Norm. Eine Projektion auf die Dreierposition allein würde also zwei später unterscheidbare Zustände gleichsetzen.

Vier bloße Grad-eins-Erzeugungsmoden liefern diesen Zustand nicht: Ihr Gesamtgrad vier reicht für die Ladung vom Mindestgrad zehn nicht aus. Die tatsächlichen getrennten Vertexeinsetzungen enthalten den erforderlichen Nachfahrenturm. Damit bleibt der Raumvorschlag an dieselbe vollständige Quellenkonstruktion angeschlossen.

\subsection{Was für physische Bewegung und Lokalität noch fehlt, ist direkt prüfbar}
Die gleiche Faserbeschreibung zeigt eine bisher nur implizite Grenze:
\begin{equation}
T|_{\mathcal K_y}=i^{\,4y_0+3y_1+5y_2}I,\qquad
[T,P_y]=[L_0,P_y]=0.
\end{equation}
Die vorhandene markierte Clock und die konforme Zeit verschieben die Dreierkoordinate nicht. Die C/X-Felder können sie verschieben; daraus folgt aber noch keine ausgewählte Zeitentwicklung, die diese Einsetzungen mit bestimmten Raten ausführt.

Der angegebene Graphlaplacian mit gleichen Raten ist unter dieser Zusatzwahl korrekt:
\begin{equation}
\lambda(k)=2\sum_a[1-\cos(k\cdot f_a)]
=8\left[1-\prod_{j=0}^2\cos(k_j/2)\right]
=|k|^2+O(|k|^4).
\end{equation}
Die drei isotropen langwelligen Richtungen sind Eigenschaften dieses Ladungsgraphen mit dieser Regel. Die quadratische Isotropie setzt noch keine Lorentzrelation und keine physische Geschwindigkeit fest.

Auch eine direkte lokale Algebrazuordnung zu den Fasern scheitert an einem kleinen Strukturtest. Für $y\ne0$ gilt $P_y\Omega=0$. Die unitalisierte Faser-Algebra
$\C I+P_y\mathcal B(\mathcal H)P_y$ erzeugt aus dem Vakuum daher nur $\C\Omega$. Sie ist keine vakuumzyklische relativistische Lokalalgebra auf dem gesamten Quellraum. Die Direktsumme nach Ladungen ist somit noch keine Anordnung unabhängiger räumlicher Labore. Eine andere, tatsächlich quellenbestimmte Lokalisation müsste zusätzlich konstruiert werden.

\subsection{Der uniforme Sewing-Test wird zu einem exakt zertifizierten Auswahlsatz}
Innerhalb der unitären affinen E8-Vakuumklasse gilt
\begin{equation}
c(k)=\frac{248k}{k+30},\qquad
d_2(k)=31124\ (k\ge2),\qquad d_2(1)=4124.
\end{equation}
Die erste affine Integrabilitäts-Nullrelation tritt bei Grad $k+1$ auf; deshalb ist der Grad-zwei-Raum für $k\ge2$ noch
$248+\dim\operatorname{Sym}^2(248)=31124$.
Für die erklärte Vakuumblock-Sewinggröße
\begin{equation}
\mathcal R_\chi(k)=2^{-3c(k)/8}\chi_{k,0}(i)
=e^{c(k)\varepsilon}\sum_{n\ge0}d_n e^{-2\pi n},
\quad
\varepsilon=\frac{\pi}{12}-\frac38\log2
\end{equation}
folgt wegen $c(k)\ge31/2$ und positiver weiterer Koeffizienten die im Anhang angegebene Schranke
\begin{equation}
\mathcal R_\chi(k)\ge
e^{(31/2)\varepsilon}
\bigl(1+248e^{-2\pi}+31124e^{-4\pi}\bigr)
=1.61786701766\ldots>\frac32.
\end{equation}
Eine gröbere, vollständig rationale Zertifizierung genügt bereits:
$512<e^{2\pi}<540$ impliziert $\varepsilon>0$ und damit
\begin{equation}\label{eq:rationalsewing}
\boxed{\mathcal R_\chi(k)>
1+\frac{248}{540}+\frac{31124}{540^2}
=\frac{114161}{72900}
>\frac32,\qquad k\ge2.}
\end{equation}
Die Exponentialgrenzen werden mit den klassischen rationalen Schranken
$157/50<\pi<22/7$ und positiven Taylorreihen samt geometrisch beschränktem Rest geprüft. Der Abstand $4811/72900$ ist damit kein Rundungsbefund.

Für $k\ge2$ ist der einzelne Vakuumcharakter allerdings noch nicht die vollständige physische Replica-Antwort einer nicht-holomorphen Theorie.
Eine saubere zusätzliche Vergleichsklasse sind volle unitäre linke/rechte E8$_k$-Theorien mit nichtnegativen Sektormultiplizitäten und Vakuummultiplizität eins. Bei $\tau=i$ gilt
$Z(i)\ge\chi_{k,0}(i)^2$.
Die vollständige Zweintervallformel~\cite{headrickreplica} liefert daher
\begin{equation}
R_{2,\mathrm{voll}}(1/2)
=2^{-3c(k)/4}Z(i)
\ge\mathcal R_\chi(k)^2
>\left(\frac{114161}{72900}\right)^2
>\frac94.
\end{equation}
Eine \emph{unabhängig} hergeleitete vollständige Antwort $9/4$ würde alle höheren Level auch in dieser vollen Vergleichsklasse ausschließen. Damit wird der chiralen TFPT-Quelle kein rechter Sektor hinzugefügt. Beide Selektoren bleiben bedingt darauf, dass der Rohprozess die jeweilige Antwort und ihre Sewingbedeutung unabhängig festlegt. Ein aus E8$_1$ berechnetes $3/2$ oder $9/4$ darf E8$_1$ nicht selbst begründen.

\subsection{Der verbundene nächste Herkunftsschritt}
Die neuen Rechnungen verbinden jetzt ausdrücklich vier Ebenen:
Die gleiche Viermarkengeometrie liefert die Paarungsgewichte; diese liefern einen Viertelphasen-Anker; die markierte Quelle trägt dessen kohärente Wirkung; eine konkret definierte Auslese erzeugt die alte Clockmatrix. Derselbe vollständige Quellenraum enthält außerdem die Ladungsfasern und das auf geschlossenen projizierten Wegen verbleibende Gedächtnis.

Die Auswahlfrage ist damit enger geworden. Der ursprüngliche Prozess muss die kohärente Familien-Auslesebasis, den tatsächlichen Umgang mit Zwischenrecords und eine Bewegung beziehungsweise lokale Ordnung der vollen Fasern festlegen. Der Clockcharakter allein bewegt die Fasern nicht. Die Faserkoordinate allein liefert keine lokalen Labore. Die zusätzliche Sewing-Schranke ist ein scharfer Vergleichstest, sobald eine unabhängige Rohquellenantwort vorliegt.

\begin{center}
\begin{tikzpicture}[node distance=7mm]
\node[box,text width=13.2cm] (s) {\textbf{Eine vollständige geladene Quelle}\\Feldoperationen, Zustand, Ladungen, Phasen und Nachfahrenturm};
\node[box,below=10mm of s,xshift=-3.6cm,text width=6cm] (c) {\textbf{Clockansicht}\\Markierter Phasenschritt und gewählte Auslese ergeben $B(x)$. Kohärenzen bestimmen die spätere Rückwirkung.};
\node[box,below=10mm of s,xshift=3.6cm,text width=6cm] (f) {\textbf{Faseransicht}\\Die Dreierkoordinate ordnet Ladungen. Innere Faser und Moden unterscheiden auch geschlossene projizierte Wege.};
\node[box,below=45mm of s,text width=13.2cm] (u) {\textbf{Universalraumregel für beide Ansichten}\\Behalte alle Unterschiede, die eine zulässige spätere Operation wieder sichtbar machen kann.};
\draw[arr] (s)--(c);\draw[arr] (s)--(f);
\draw[arr] (c)--(u);\draw[arr] (f)--(u);
\end{tikzpicture}
\end{center}

\key{\textbf{Was der Gesamtlösung jetzt konkret näherkommt:}
Die vorher getrennten Clockkonstruktionen sind für alle $x$ verbunden; der Quellen-Intertwiner, die orthogonale Erweiterung der Trine und die Rückwirkung der ausgelassenen Kohärenzen sind ausdrücklich berechnet. Der volle Quellenzustand erklärt zugleich, warum ein geschlossener projizierter Weg Information behalten kann. Die noch zu beweisende Herkunft betrifft die gemeinsame Auswahl dieses Instruments und einer physischen lokalen Dynamik auf derselben Quelle. Eine vollständige Raumzeit- und Wechselwirkungsherleitung ist durch die neuen Darstellungen allein weiterhin nicht geliefert.}

\section{Gemeinsame Zeit, wirkliche Übertragung und die Auswahl des Quellprozesses}\label{sec:sharedtime}

\key{\textbf{Der neue Anhang verbindet die richtige Frage mit einer noch zu engen Ereignisklasse.}
Die gemeinsame Clockhierarchie lässt sich exakt ausführen. Ihre kontinuierliche Markovfortsetzung ist jetzt für die ganze erklärte Klasse entschieden. Gleichzeitig zeigt die tatsächliche Quellenalgebra, warum diese Hierarchie allein keine wechselwirkende Welt beschreibt und warum der vorgeschlagene räumliche Lauf nicht unverändert auf die C-Felder übertragen werden kann. Ein positiver Anschluss ist bereits in derselben Quelle vorhanden: Lokale Stromimpulse besitzen unter ihrer bestehenden konformen Zeit eine exakt berechenbare transportierte Wirkung.}

Geprüft wird der eingefügte Text \file{f548f0d4-2ae2-4067-8f27-3b19b7c82ae4}. Seine verlinkten Dateien unter \file{sandbox:/mnt/data/} wurden nicht mitgeliefert. Die genannten 37 neuen und 75 älteren Bedingungen sind deshalb keine hier übernommenen Ausführungsnachweise. Die zentralen Formeln wurden eigenständig aus den vorhandenen Quellenoperatoren rekonstruiert. Die Einzelregister-GKLS-Fortsetzung und die freie Richtung $q$ waren bereits im Rechenkern enthalten; neu geprüft werden insbesondere die gemeinsame Registerhierarchie, ihr vollständiger Einbettungsbereich und die quelltreue räumliche Übertragung.

\subsection{Der ganze Quantenkanal und die gemeinsame Ereignisregel}
Auf der tatsächlichen markierten Zweierebene gilt in der Pauli-Reihenfolge $X,Y,Z$
\begin{equation}
\Phi_q(X)=\tfrac23X,\qquad \Phi_q(Y)=qY,\qquad \Phi_q(Z)=\tfrac13Z.
\end{equation}
Die ursprüngliche Trine erfüllt $\Phi_q^*(E_i)=\sum_jB_{ij}E_j$.
Am historischen, weiterhin bedingt ausgewählten Punkt $q=1/3$ ist
\begin{equation}
\Phi_{1/3}(\rho)=\tfrac7{12}\rho+\tfrac14X\rho X
+\tfrac1{12}Y\rho Y+\tfrac1{12}Z\rho Z.
\end{equation}
Der partiell transponierte normierte Choioperator hat den Eigenwert $-1/12$. Dieser Kanal kann also Verschränkung mit einer Referenz erhalten. Die dreidimensionale vollständige Dephasierung aus Abschnitt~\ref{sec:anchorfibre} ist dagegen ein anderer Quantenprozess. Die Übereinstimmung der Matrix $B$ bleibt richtig, identifiziert aber nicht die beiden vollständigen Instrumente.

Die angegebene Einzelregister-Fortsetzung stimmt:
\begin{equation}
\mathcal L_1=\tfrac{\log6}{4}(\operatorname{Ad}_X-\Id)
+\tfrac{\log(3/2)}4(\operatorname{Ad}_Y-\Id)
+\tfrac{\log(3/2)}4(\operatorname{Ad}_Z-\Id),\qquad
 e^{\mathcal L_1}=\Phi_{1/3}.
\end{equation}
Die Operatoren liegen in der von den nativen Reflexionen erzeugten Algebra. Die Definition einer Rate setzt jedoch eine Zeiteinheit voraus und macht eine solche kohärente Kombination noch nicht zu einem aus P1/P2 ausgewählten Sprungereignis.

Die sechs bereits identifizierten nativen Wörter $u_g$ besitzen auf dieser Ebene die $S_3$-Gewichte
\begin{equation}\label{eq:sharedweights}
w(q)=\left(\tfrac12+\tfrac q6,\tfrac q6,\tfrac q6,
\tfrac1{18}-\tfrac q6,\tfrac29-\tfrac q6,\tfrac29-\tfrac q6\right),
\quad 0\le q\le\tfrac13,
\end{equation}
in der Reihenfolge $e,c,c^{-1},r_{59},r_{14},r_{13}$.
Für $N$ Register und einen gemeinsamen Record gilt
\begin{equation}\label{eq:sharedhierarchy}
\Psi_{N,q}(\rho)=\sum_gw_g(q)u_g^{\otimes N}\rho(u_g^{\otimes N})^*,
\qquad
\mathcal V_{N,q}\psi=\sum_g\sqrt{w_g(q)}\ket g\otimes u_g^{\otimes N}\psi.
\end{equation}
Direkt folgen $\mathcal V^*\mathcal V=I$ und die Verträglichkeit mit jeder partiellen Spur. Bei unabhängig erneuerten Ereignissen und ohne Zwischeneingriffe ist die gesamte ganzzahlige Entwicklung geschlossen:
\begin{equation}
\begin{split}
w_e^{*m}&=(1+h+2r+2s)/6,\\
w_c^{*m}=w_{c^{-1}}^{*m}&=(1+h-r-s)/6,\\
w_{r_{59}}^{*m}&=(1-h-2r+2s)/6,\\
w_{r_{14}}^{*m}=w_{r_{13}}^{*m}&=(1-h+r-s)/6,
\end{split}\qquad
r=(2/3)^m,\quad s=(1/3)^m,\quad h=q^m.
\end{equation}
Damit gilt $\Psi_{N,q}^m=\sum_gw_g^{*m}\operatorname{Ad}_{u_g^{\otimes N}}$ für alle $N$ und $m\ge0$; bei $m=0$ gilt $w^{*0}=\delta_e$, die Punktmasse am Identitätsereignis. Eingriffe zwischen den Schritten müssen ausdrücklich eingefügt werden. Die Faltungsformel ist kein Ersatz für solche geordneten Geschichten.

\subsection{Die gemeinsame kontinuierliche Zeit ist für die ganze Klasse entschieden}
\key{\textbf{Neuer, stärkerer Satz:}
Die Hierarchie~\eqref{eq:sharedhierarchy} besitzt genau dann eine konsistente zeitunabhängige GKLS-Fortsetzung auf denselben Registern, wenn $2/9\le q\le1/4$. Die Notwendigkeit folgt bereits aus zwei Registern und setzt den hypothetischen Generator nicht vorher als klassische oder reversible $S_3$-Mischung an.}

Zunächst bestätigt sich der einfache Endpunktzeuge des Anhangs. Für $q=1/3$ ist
\begin{equation}
w*w=(10/27,5/54,5/54,1/27,11/54,11/54).
\end{equation}
Die sechs $V_g=u_g\otimes u_g$ sind unabhängig. Ihre Gram-Matrix ist
\begin{equation}
G_{gh}=\tr(V_g^*V_h)=\diag(3I_3+\mathbf1\mathbf1^T,3I_3+\mathbf1\mathbf1^T),
\quad \det G=2916.
\end{equation}
Daher ist der Choi-Rang nach einem Schritt fünf und nach zwei Schritten sechs. Alle Tensorpotenzen $u^{\otimes N}$ mit $N\ge2$ enthalten die drei irreduziblen $S_3$-Darstellungen; der gleiche Rangsprung gilt für sämtliche dieser $N$.

Für eine endlichdimensionale zeitunabhängige GKLS-Halbgruppe ist der Krausrang für $t>0$ konstant. Eine Begründung benutzt die positiven Sprungtrajektorien: Nach Abspaltung der invertiblen driftenden Linksmultiplikation wird ihr linearer Raum von Produkten der analytisch konjugierten Sprungoperatoren erzeugt. Die Auswertung auf jedem offenen Zeit-Simplex spannt denselben analytischen Funktionenraum auf; seine Dimension hängt nicht von der positiven Intervalllänge ab. Der Rangsprung schließt somit jede solche Halbgruppe am Endpunkt aus. Die Hauptquadratwurzel liefert zusätzlich den unabhängigen negativen normierten Choi-Erwartungswert
\begin{equation}
\frac{3+\sqrt3-2\sqrt6}{20}<0.
\end{equation}

\textbf{Vollständiger Beweis im offenen Bereich.}
Für $0<q<1/3$ zerfällt der Zweiregisterraum multiplikitätsfrei als
$H=H_1\oplus H_\varepsilon\oplus H_\sigma$ mit Dimensionen $1,1,2$.
Seine orthonormale Basis kann als
\begin{equation}
\tfrac1{\sqrt2}(1,0,0,1),\quad
\tfrac1{\sqrt2}(0,1,-1,0),\quad
\tfrac1{\sqrt2}(1,0,0,-1),\quad
\tfrac1{\sqrt2}(0,1,1,0)
\end{equation}
gewählt werden. Schreibe $P_1,P_\varepsilon,P_\sigma$ für die zugehörigen Projektoren. Der Fixraum von $\Psi_{2,q}$ ist
$\C P_1\oplus\C P_\varepsilon\oplus\C P_\sigma$; seine extremalen normierten Fixzustände sind $P_1,P_\varepsilon,P_\sigma/2$.

Angenommen, $\Psi_{2,q}=e^{\mathcal K}$ mit einem beliebigen GKLS-Generator.
Die Halbgruppe kommutiert mit ihrem Zeitschritt und erhält daher dessen Fixraum. Auf dem Fixzustandssimplex ist sie positiv und spurtreu und nach Zeit eins die Identität. Sie besitzt dort positive Inversen. Solche Simplex-Automorphismen sind Permutationen; Stetigkeit erzwingt die Identität. Also fixiert sie jeden $P_i$ zu jeder Zeit. \emph{Erst daraus} folgt ihre Unitalität. Der Multiplikativitätsbereich fixer Projektoren erhält nun alle neun Räume $\operatorname{Hom}(H_j,H_i)$.

Auf diesen einzelnen Räumen sind die positiven Spektren einfach:
\begin{center}
\begin{tabular}{ll}
\toprule
Block & Spektrum\\\midrule
$11,\varepsilon\varepsilon$ & $1$\\
$1\varepsilon,\varepsilon1$ & $q$\\
$1\sigma,\sigma1,\varepsilon\sigma,\sigma\varepsilon$ & $1/3,2/3$\\
$\sigma\sigma$ & $1,2/3,q,1/3$\\\bottomrule
\end{tabular}
\end{center}
Das volle 16-dimensionale Spektrum ist dagegen degeneriert.
Blockweise muss $\mathcal K$ deshalb in derselben orthonormalen Eigenbasis diagonal sein. Seine möglichen Logarithmuszweige haben identischen reellen Anteil:
\begin{equation}
\tfrac12(\mathcal K+\mathcal K^*)=\log\Psi_{2,q}.
\end{equation}
Weil die Halbgruppe unital und spurtreu ist, sind $\mathcal K$ und sein Hilbert--Schmidt-Adjungierter beide GKLS-Generatoren. Auch ihr Mittel muss einer sein.

Setze $a=\log3$, $b=\log(3/2)$ und $c=-\log q$. Der Hauptlogarithmus hat die fünf Raten
\begin{equation}
\gamma_c=\gamma_{c^{-1}}=\frac{a+b-c}{6},\quad
\gamma_{r_{59}}=\frac{c-2a+2b}{6},\quad
\gamma_{r_{14}}=\gamma_{r_{13}}=\frac{a-b+c}{6}.
\end{equation}
Nach Projektion orthogonal zu $\operatorname{vec}I$ bleiben die anderen fünf Gruppenoperatoren linear unabhängig. Der bedingte Choioperator des Generators ist die Summe ihrer Rang-eins-Formen mit Koeffizienten $\gamma_g$. Seine Positivität ist daher äquivalent zu $\gamma_g\ge0$ für jedes $g\ne e$. Das allgemeine bedingte Choi-Kriterium ist in~\cite{wolfjoint} angegeben; die vorstehende Reduktion aller Logarithmuszweige verwendet die hier nachgerechnete Sektorstruktur.

Es folgt $2(a-b)\le c\le a+b$, also $2/9\le q\le1/4$.
Umgekehrt erzeugen genau diese nichtnegativen Raten
$\mathcal K_N=\sum_{g\ne e}\gamma_g(\operatorname{Ad}_{u_g^{\otimes N}}-\Id)$
die ganze kompatible Hierarchie. Bei $q=0$ verhindert bereits Nichtinvertierbarkeit eine endliche Generator-Exponentialdarstellung; $q=1/3$ wurde oben ausgeschlossen. Damit ist die Klasse vollständig entschieden.

Das Ergebnis geht über ein zufälliges Nullgewicht hinaus: Bei $q=3/10$ sind alle sechs Gewichte positiv, beide Choi-Ränge sechs, aber
$\gamma_{r_{59}}=\log(5/6)/6<0$. Auch dort ist jede zeitunabhängige GKLS-Einbettung auf dem Zweiregisterraum ausgeschlossen.

\picturetext{Eine Uhr kann für sich allein einen glatten Sekundenzeiger besitzen. Werden zwei Uhren durch dieselben Ereignisse angetrieben, muss auch ihre gemeinsame Reaktion zwischen den angezeigten Schritten passen. Genau diese Forderung schränkt die mögliche Fortsetzung stärker ein.}

Falls die historische Antwort $FAAF=1/27$ tatsächlich unabhängig auf diese Quelle übertragen wird, verlangt sie $q=1/3$. Die drei Anforderungen \emph{diese Antwort}, \emph{die gemeinsame Ereignisregel} und \emph{eine zeitunabhängige Markovzeit auf denselben endlichen Registern} sind dann unvereinbar. Daraus folgt keine fundamentale Diskretheit der Welt: Eine größere Quelle mit Records, zeitabhängige Ausführung und die vorhandene konforme oder modulare Quellenzeit bleiben möglich.

\subsection{Die gemeinsame Ereignisregel transportiert noch keine Wirkung zwischen Registern}
Der Bellzustand $\Omega_2=(\ket{00}+\ket{11})/\sqrt2$ besitzt unter einem gemeinsamen Ereignis Rückkehrwahrscheinlichkeit eins; unter unabhängigen Ereignissen beträgt sie $5/12$. Diese im Anhang behaupteten Werte stimmen. Die Ursache ist eine geschützte gemeinsame Symmetrie.

Eine weitergehende Gleichung entscheidet aber über die Eignung als Weltmechanik. Für jede Registerteilmenge $S$ gilt
\begin{equation}\label{eq:nosupportgrowth}
\boxed{\Psi_{N,q}^*(A_S\otimes I)
=\Psi_{|S|,q}^*(A_S)\otimes I.}
\end{equation}
Dies folgt unmittelbar aus dem Produkt jedes einzelnen Krausoperators. Jede lokale Operatoralgebra bleibt in sich; das gilt auch für jede Potenz. Eine Änderung am Eingang eines anderen Registers kann daher die lokale Ausgabe nicht beeinflussen. Separierbare Eingaben bleiben separierbar, weil ihre Ausgaben Mischungen von Produktzuständen sind. Es handelt sich um lokale Operationen mit gemeinsamem klassischem Zufall, eine bekannte Klasse~\cite{gutoskijoint}.

\key{\textbf{Gemeinsame Zufälligkeit ist eine gemeinsame Ursache, aber noch keine Wechselwirkung zwischen den Registern.}
Die Hierarchie kann vorhandene Verschränkung schützen und klassische Korrelationen erzeugen. Für sich allein erzeugt sie weder Informationsübertragung zwischen den Registern noch Verschränkung aus Produktzuständen. Sie kann deshalb kein vollständiger Ersatz für den gesuchten wechselwirkenden Feldprozess sein.}

Dies bleibt wahr, wenn man solche zustandsunabhängigen gemeinsamen Ereignisse auf Kanten oder Hyperkanten verteilt und komponiert. Eine Zeichnung mit verbindenden Kanten erzeugt noch keine transportierende Operatorwirkung. Kohärent rückgekoppelte Records, echte Mehrregisteroperatoren oder rückwirkende dynamische Vermittler wären zusätzliche Operationen; sie sind nicht bereits durch~\eqref{eq:sharedhierarchy} implementiert.

\subsection{Warum weitere simultane Symmetrieprüfungen die Ereignisgewichte nicht auswählen}
Die tatsächlichen komplexen Reflexionen sind Produkte zweier reeller E8-Weylreflexionen:
\begin{equation}
R_\alpha=s_\alpha s_{J\alpha},\qquad (\alpha,J\alpha)=0.
\end{equation}
Für die drei verwendeten Strahlen wurde dies an der ursprünglichen Wurzelwirkung geprüft. Solche Gitterisometrien besitzen passende Vakuum-erhaltende Lifts zur Gitterquelle~\cite{dongnagjoint}. Der Lift samt Kokzyklus gehört zu den Daten: Eine $S_3$-Wortgleichheit auf der markierten Zweierebene ist nicht automatisch eine kanonische $S_3$-Aktion auf allen geladenen Feldern.

Für jeden gültigen einzelnen Lift $\alpha_g$ und alle wohldefinierten Vakuumwörter gilt jedoch bereits
\begin{equation}\label{eq:vacuumweightblind}
\omega\bigl(\alpha_g(A_1)\cdots\alpha_g(A_n)\bigr)
=\omega(A_1\cdots A_n).
\end{equation}
Jede normierte Mischung dieser \emph{simultanen} Antworten ist somit von $w_g$ unabhängig. Auch beliebig viele solche Ward- und Symmetrietests können die Ereignisgewichte nicht identifizieren. Dagegen können Operationen \emph{zwischen} den Einsetzungen oder Antworten präparierter angeregter Zustände die Gewichte unterscheiden. Genau in diese zweite Klasse fällt die bedingte $q/9$-Antwort.

\picturetext{Viele genaue Fotos eines symmetrischen Zahnrads verraten nicht, wie sein Motor zwischen zwei Fotos betätigt wird. Dafür muss man eine Änderung auslösen und verfolgen, wo und wann sie wieder erscheint. Die bisherigen Strukturprüfungen bleiben richtig; ihre Unempfindlichkeit gegenüber den Gewichten erklärt, warum ihre Anzahl die dynamische Auswahl nicht automatisch schließt.}

\subsection{Der vorgeschlagene Raumlauf scheitert bei direkter Übertragung auf die Quelle}
Die neue Faserenergie ist korrekt:
\begin{equation}
E_{\min}(y)=\tfrac12|y|^2+
\begin{cases}
0,&y\in\Z^3,\ \sum y\text{ gerade},\\
1/2,&y\in\Z^3,\ \sum y\text{ ungerade},\\
5/8,&y\in(\Z+1/2)^3.
\end{cases}
\end{equation}
Sie folgt durch Minimieren der ersten fünf Ladungskoordinaten unter der E8-Parität. Ein energieerhaltender homogener Ortstransport lässt sich daher nicht einfach mit beliebig großen Familienladungen und unverändertem $L_0$ identifizieren.

Auf einem zusätzlich als Ortsgitter interpretierten $A_3^*$ ist
\begin{equation}
W(k)=e^{-ik_xX/2}e^{-ik_yY/2}e^{-ik_zZ/2}
=\sum_{\epsilon\in\{\pm1\}^3}e^{-ik\cdot\epsilon/2}A_\epsilon,
\quad A_\epsilon=\tfrac18(I+\epsilon_xX)(I+\epsilon_yY)(I+\epsilon_zZ)
\end{equation}
tatsächlich unitär und besitzt die behauptete Weylgrenze. Die beiden Nullquasienergie-Kegel bei $0$ und $(\pi,\pi,-\pi)$ besitzen entgegengesetzte Orientierungen $1/8$ und $-1/8$. Zusätzlich liegen zwei inequivalente $W=-I$-Kegel bei $(2\pi,0,0)$ und $(\pi,\pi,\pi)$. Solche freien QCA-Grenzen sind etablierte Konstruktionen~\cite{darianojo}; sie wählen weder die TFPT-Wechselwirkungen noch Massen aus.

\textbf{Der entscheidende neue Test verwendet die ursprünglichen Ladungsfaktoren.}
Für die vier C-Wurzeln $\alpha_a=((-1/2)^5;f_a)$ gilt
\begin{equation}
(\alpha_a,\alpha_b)=1+\delta_{ab},\qquad
\sum_a f_a=0,\qquad \sum_a\alpha_a=((-2)^5;0)\ne0.
\end{equation}
Die vier Wurzeln und ihre Differenzen schließen auf 20 Wurzeln vom Typ $A_4$ mit Cartanrang vier. Die Projektion hat also eine tatsächliche innere Ladungsrichtung ausgeblendet; daraus folgt keine Raumzeitdimension.

Die unitären Ladungsfaktoren auf $\ell^2(E_8)$ erfüllen
\begin{equation}
U_\alpha\ket\lambda=\varepsilon(\alpha,\lambda)\ket{\lambda+\alpha},
\qquad U_{\alpha_a}U_{\alpha_b}=-U_{\alpha_b}U_{\alpha_a}\quad(a\ne b).
\end{equation}
Das ist die originale Gitter-Kokzyklusrelation~\cite{dongnagjoint}. Der Fourierbeweis für $W$ verwendet dagegen gewöhnliche kommutierende Translationen.

Ersetzt man die acht Translationen unmittelbar durch die originalen $\pm C_a$-Ladungsshifts samt den vorhandenen Feldphasen $p_\alpha$, entsteht
\begin{equation}
\widetilde W=\sum_{\alpha=\pm\alpha_a}A_{2\pi_F(\alpha)}\otimes p_\alpha U_\alpha.
\end{equation}
Die exakte Auswertung des vorhandenen Kokzyklus ergibt in $\widetilde W^*\widetilde W$ für
$\Delta=(0,0,0,0,0;0,1,1)$ den Koeffizienten
\begin{equation}\label{eq:nativehopfailure}
\boxed{\begin{pmatrix}0&0\\1/2&0\end{pmatrix}\ne0.}
\end{equation}
Die Nullverschiebung hat weiterhin Koeffizient $I$; insgesamt bleiben zwölf nichtverschwindende zusätzliche Koeffizienten. Bereits mit vollständiger Ladung und vorübergehend unterdrücktem Kokzyklus bleiben vier. Nur die gewöhnliche projizierte Gitterrechnung erfüllt alle Unitaritätsbedingungen.

Dies widerlegt genau die unmittelbare Übertragung mit diesen Richtungskoeffizienten. Die $U_\alpha$ sind dabei lediglich die unitären Ladungsfaktoren, nicht die vollständigen Vertexfelder mit ihren Nachfahren. Andere Operatoren auf der vollständigen Faser sind durch diesen Gegenfall nicht ausgeschlossen. Ein einfaches Umbenennen der C-Felder in räumliche Hops ist jedoch keine quelltreue Lösung.

\subsection{Ein positiver gemeinsamer Prozess ist bereits in der vollständigen Quelle vorhanden}
Der stärkste nächste Ansatz verwendet die vorhandene lokale E8-Quelle selbst. Betrachte einen normierten Cartan-Strom mit
\begin{equation}
[J_n,J_m]=n\delta_{n+m,0}I,\qquad [L_0,J_n]=-nJ_n.
\end{equation}
Dies sind ihre bestehenden Level-eins- und Zeitrelationen. Für reelle glatte Kreisfunktionen seien
\begin{equation}
f_n=\frac1{2\pi}\int_0^{2\pi}f(\theta)e^{-in\theta}\,d\theta,
\quad J(f)=\sum_n f_nJ_n,
\quad W(f)=e^{iJ(f)},
\quad \sigma(f,g)=\frac1{2\pi}\int f g'.
\end{equation}
Dann gilt $[J(f),J(g)]=i\sigma(f,g)I$ und unter der bereits vorhandenen Quellenzeit
\begin{equation}
\alpha_t(W(g))=W(g_t),\qquad g_t(\theta)=g(\theta-t).
\end{equation}
Die zugrunde liegende lokale Stromnetzstruktur ist Standard~\cite{tanimotojoint}; die folgende Rechnung ist ihre direkte Anwendung auf den vorhandenen Cartan-Strom.

Präpariere lokal $\Omega_\epsilon=W(\epsilon f)\Omega$, wobei $f$ in einem echten Intervall getragen wird. Aus dem zentralen Kommutator folgen ohne Näherung
\begin{equation}\label{eq:localcausalresponse}
\boxed{\omega_\epsilon(\alpha_t(W(g)))
=e^{i\epsilon\sigma(f,g_t)}\omega(W(g)),\qquad
\omega_\epsilon(J(g_t))-\omega(J(g_t))=\epsilon\sigma(f,g_t).}
\end{equation}
Die lokale Präparation und ihre spätere Ausgabe gehören zu demselben Zustand, derselben Algebra und demselben $L_0$. Sie kostet die ebenfalls bestimmte Quellenenergie
\begin{equation}
\Delta\langle L_0\rangle
=\epsilon^2\sum_{n>0}n^2|f_n|^2
=\frac{\epsilon^2}{4\pi}\int(f')^2.
\end{equation}
Die Testfunktion ist eine gewählte Eingabe an den festgelegten Quellenprozess; hier wird kein neuer Hamiltonoperator eingeführt.

Ein strikt lokaler positiver Zeuge ist ausdrücklich konstruierbar. Wähle eine nichtkonstante reelle glatte Beule $h$ mit kleinem Träger, $f=h'$ und $g(\theta)=h(\theta+t_0)$, so dass die Träger von $f$ und $g$ anfangs getrennt sind. Dann ist $\sigma(f,g)=0$, aber
\begin{equation}
\sigma(f,g_{t_0})=\frac1{2\pi}\int(h')^2>0.
\end{equation}
Die Eingriffswirkung erreicht also tatsächlich ein anderes Intervall. Dies ist chiral-konformer Transport entlang der Naht, noch keine $3+1$-Geometrie oder bereits nachgewiesene nichtabelsche Streuung.

Die lokale Operation erhält dabei den vollen reinen Vakuumzustand nicht. Damit passt sie auch zum korrekten lokalen Ausschluss des Anhangs: Sind alle beschränkten Krausoperatoren $K_a$ in einer Algebra mit separierendem $\Omega$ und gilt $\sum_aK_a\ket\Omega\bra\Omega K_a^*=\ket\Omega\bra\Omega$, so erzwingt Reinheit $K_a\Omega=c_a\Omega$ und Separiertheit $K_a=c_aI$. Die globale Projektion auf einen angeregten Clockzustand kann deshalb nicht unverändert eine solche lokale Operation sein. Der lokale Puls liefert eine tatsächliche Alternative innerhalb der Quelle und zeigt ihren Energiepreis.

\subsection{Die fundamentale Auswahlfrage ist jetzt eine gemeinsame Antwortgleichung}
\begin{center}
\begin{tikzpicture}[node distance=7mm]
\node[box,text width=12.7cm] (source) {\textbf{Dieselbe vollständige Quelle}\quad lokale Felder, Kokzyklus, Nachfahren, Vakuum und konforme Zeit};
\node[box,below=of source,text width=12.7cm] (pulse) {\textbf{Tatsächlicher Eingriff und spätere Wirkung}\quad lokale Präparation $\rightarrow$ Quellenentwicklung $\rightarrow$ messbare Antwort};
\node[box,below=of pulse,text width=12.7cm] (map) {\textbf{Gesuchte ursprüngliche Identifikation}\quad P1/P2 wählen Quelle und Instrument; eine räumliche Fortsetzung erhält dieselben gemeinsamen Antworten};
\draw[arr] (source)--(pulse);
\draw[dasharr] (pulse)--(map);
\end{tikzpicture}
\end{center}

Die bereits explizite Antwort~\eqref{eq:localcausalresponse} verhindert, dass wir die Quelle auf eine gemeinsam rauschende Anzeige reduzieren. Eine korrekte Rekodierung muss ihre geordneten Eingriffe und späteren Wirkungen erhalten, zusätzlich zu Ladungen, Statistik, höherer Antwort und Positivität. Für den Universalraum bedeutet dies: Er bewahrt die Unterschiede von \emph{beeinflussbaren zukünftigen Experimenten}, einschließlich ihres Orts- und Zeitbezugs. Eine gemeinsame Ereignisnummer oder derselbe marginale Kanal reicht dafür nicht.

Die Arbeit am fundamentalen Abschluss wird daher auf zwei gekoppelte Herkunftsgleichungen konzentriert: Der tatsächliche ursprüngliche Nahtoperator muss die erklärte Quelle samt Zustand und Instrument liefern; die räumliche Realisierung muss deren vollständige lokale Antwortstruktur erhalten und daraus die geforderten gemeinsamen physikalischen Ausgaben gewinnen. Die lokalen Impulsantworten, die gemeinsame Clockschranke und der explizite Kokzyklusgegenfall sind jetzt konkrete Bedingungen an diese Gleichungen. Sie ersetzen sie nicht durch ein weiteres frei gewähltes Weltmodell.

\key{\textbf{Was diese Fortsetzung tatsächlich entscheidet:}
Die gemeinsame zeitunabhängige Markovklasse ist vollständig klassifiziert. Eine reine gemeinsame $S_3$-Zufallsregel ist als vollständiger wechselwirkender Prozess ausgeschlossen; die naive C-Hop-Übertragung ist ebenfalls ausgeschlossen. Zugleich ist ein echter lokaler Eingriffs- und Transportprozess in der vorhandenen Quelle ausdrücklich berechnet. Die vollständige fundamentale Weltlösung ist damit nicht bewiesen. Die Suche hat aber nun einen vorhandenen positiven Quellenprozess und genaue gemeinsame Identitätsbedingungen, an denen jede vorgeschlagene Gesamtlösung bestehen muss.}


\section{Der konstruktive Rückweg: Eine lokale Feldregel bestimmt Quelle, Zustand und Zeit}\label{sec:inverseorigin}

\key{\textbf{Die entscheidende Verdichtung der Ursprungsfrage.} Die bereits gewählte, markierte E8-Level-1-Quelle lässt sich zu genau ihrer passenden Theorie mit sechzehn reellen Fermionen zurückverfolgen. Dabei werden Vakuum, vollständige Feldantworten und Konformalzeit gemeinsam rekonstruiert. Der Herkunftsnachweis muss deshalb nicht vier unabhängig gewählte Zutaten rechtfertigen. Er muss einen bestimmten lokalen Feldkeim im ursprünglichen P1/P2-Prozess nachweisen. Der folgende Satz beschreibt diesen Keim und beweist, was daraus folgt.}

\subsection{Das vorhandene E8 enthält den Rückweg bereits in seiner Markierung}
Die ursprüngliche Verklebung enthält die Zwischenstufe
\begin{equation}
D_5\oplus D_3\subset D_8\subset E_8,
\qquad D_8=\{n\in\Z^8:\textstyle\sum_i n_i\text{ gerade}\}.
\end{equation}
Sie wird nicht nachträglich ausgesucht: Im ursprünglichen $\Z_4$-Glue sind es die beiden geraden Klassen. Dessen Charakter $G$ hat auf den Klassen $j$ den Wert $i^j$; der Fixpunkt von $G^2$ ist die $D_8$-Gitteralgebra. In den ursprünglichen verdoppelten Wurzelkoordinaten $r=2\alpha$ kann man $j=\sum_{i=1}^5r_i\bmod4$ benutzen. Die vier Wurzelklassen enthalten $52,64,60,64$ Elemente. Die geraden Klassen liefern die $112$ Wurzeln $\pm e_i\pm e_j$; mit den acht Cartanrichtungen sind das die $120$ Ströme von $\mathfrak{so}(16)$. Die beiden ungeraden Klassen liefern den ursprünglichen Halbspinor mit $128$ Komponenten. Dies ist dieselbe ursprüngliche Verzweigung aus \file{tfpt_1_architecture_e8.tex}, nicht ein neuer E8-Ansatz.

Die vier Sektoren des $D_8$-Gitters sind
\begin{equation}
0,\quad v=e_1+D_8,\quad s=\tfrac12(1,\ldots,1)+D_8,\quad c=s+v,
\qquad (h_0,h_v,h_s,h_c)=(0,\tfrac12,1,1).
\end{equation}
Hier bezeichnet $h$ das niedrigste konforme Gewicht, also die halbe minimale quadrierte Gitterlänge. Genau $v$ hat den fermionischen Drehfaktor $e^{2\pi ih}=-1$. Seine Erweiterung ist ausdrücklich
\begin{equation}\label{eq:inverseZ8}
\boxed{D_8\cup(D_8+v)=\Z^8,\qquad V_{\Z^8}\cong\mathcal F_{16}.}
\end{equation}
Die zweite Gleichung ist die Boson-Fermion-Isomorphie: Jede der acht Einheitsrichtungen liefert ein komplexes Fermion, also zwei reelle Majoranafelder. Die explizite Rang-eins-Abbildung und ihre Tensorpotenzen stehen in~\cite{barronfermion}. Die sechzehn Norm-eins-Vektoren $\pm e_i$ liefern die Felder vom Gewicht $1/2$.

Die Erweiterung ist bis zu Feldphasen und markierungserhaltender Isomorphie festgelegt. Die Standard-Kokzyklen der ungeraden Gitterbeschreibung dürfen dabei auf $D_8$ durch eine Eins-Koketten-Umphasung an den ursprünglichen Kokzyklus angeglichen werden; sie dürfen nicht still an dessen Stelle treten. Für die zurückgewonnenen C/X-Antworten wird weiterhin der originale E8-Kokzyklus verwendet.

\begin{center}
\begin{tikzpicture}[node distance=8mm]
\node[box,text width=12.9cm] (m) {\textbf{Ursprünglich markiertes $E_{8,1}$}\\$D_5+A_3$, Glue und ausgewählter Halbspinor $s$};
\node[box,below=of m,text width=12.9cm] (d) {\textbf{Gemeinsame gerade Quelle $(D_8)_1$}\\Fixpunkt der halben Glue-Wirkung; derselbe Vakuum- und Stress-Tensor};
\node[box,below=11mm of d,xshift=-3.45cm,text width=5.6cm] (f) {\textbf{Fermionische Erweiterung}\\$0\oplus v$, sechzehn Majoranas\\$\mathcal H_F=\mathcal H_0\oplus\mathcal H_v$};
\node[box,below=11mm of d,xshift=3.45cm,text width=5.6cm] (e) {\textbf{Ursprüngliche bosonische Erweiterung}\\$0\oplus s$, wieder $E_{8,1}$\\$\mathcal H_{E_8}=\mathcal H_0\oplus\mathcal H_s$};
\draw[arr] (m)--(d);\draw[arr] (d)--node[left,font=\scriptsize]{einzig fermionischer Sektor}(f);\draw[arr] (d)--node[right,font=\scriptsize]{markierter Spinorsektor}(e);
\end{tikzpicture}
\end{center}

Die beiden Erweiterungen besitzen unterschiedliche Hilberträume. Auf dem ganzen E8-Raum wirkt das feste $D_8$-Unterobjekt in der Darstellung $0\oplus s$; sein eigener zyklischer Vakuumraum ist $\mathcal H_0=\overline{V_{D_8}\Omega}$. Die Majoranafelder sind keine bisher übersehenen lokalen $h=1/2$-Felder im bosonischen E8-Vakuumraum. Der markierte Halbspinor wird beim Rückweg mitgeführt und führt über $D_8+s$ wieder zur ursprünglichen E8-Quelle~\cite{chuzheng}. Insbesondere gehören die ursprünglichen C/X-Felder zum Spinor-Erweiterungssektor; sie sind nicht die neuen elementaren Majoranafelder.

\subsection{Die zurückgewonnene Darstellung ist der ursprüngliche Träger}
Die sechzehn Gewichte zerfallen entlang der ursprünglichen ersten fünf und letzten drei Gitterrichtungen als
\begin{equation}
16\downarrow_{\mathrm{Spin}(10)\times SU(4)}=(10,1)\oplus(1,6)=V_{16}.
\end{equation}
Jede der $112$ ursprünglichen $D_8$-Wurzeln verbindet tatsächlich zwei Gewichtspaare dieser Darstellung. Auf P2 lauten ihre Ladungen
\begin{equation}
\{-\tfrac13\ (3\times),+\tfrac13\ (3\times),
-\tfrac12\ (2\times),+\tfrac12\ (2\times),0\ (6\times)\},
\qquad \tr_{V_{16}}Y^2=\frac53.
\end{equation}
Das rekonstruiert die vorher geometrisch gewonnene Darstellung mitsamt ihrer Wirkung. Es identifiziert diese Richtungen weiterhin nicht mit sechzehn vierdimensionalen Teilchensorten.

Auch die Überlagerung stimmt: Für $h=(z,-iI_4)$ ist $h^2$ der Kern der Einbettung nach Spin(16), während $h$ auf $V_{16}$ als $-I$ und auf dem gewählten Halbspinor als $+I$ wirkt. Somit ist die schon konstruierte Gruppe $K_f$ gerade die Paritätsüberlagerung der bosonischen Markierungsgruppe $K$. Eine allgemeine interne E8-Symmetrie wirkt innerhalb dieser Fermionisierung nur, wenn sie das gewählte $D_8$ erhält. Selbst dann können die Gruppenlifts die zentrale Fermionparität benötigen. Das ursprüngliche Gruppenwörterbuch bleibt maßgeblich.

\textbf{Die originale Viererstruktur lässt sich dabei ausdrücklich zurückgewinnen.} Für Cartanoperatoren $H_{0,i}$ mit Gewichten $\pm e_i$ ist der oben definierte duale Glue-Charakter
\begin{equation}\label{eq:inverseinternalglue}
G=\exp\!\left(i\pi\sum_{i=1}^5H_{0,i}\right).
\end{equation}
Auf den ersten zehn reellen Fermionen wirkt er mit $-1$, auf den letzten sechs mit $+1$. Im NS-Fockraum ist seine Wirkung deshalb $G_F=(-1)^{F_{10}}$ und $G_F^2=I$. Sein Spin-Lift besteht aber aus einer Halbdrehung in jeder der fünf ursprünglichen reellen Zweiebenen. Mit euklidischen Cliffordgeneratoren gilt
\begin{equation}
g=\prod_{j=1}^5(\gamma_{2j-1}\gamma_{2j}),\qquad g^2=(-1)^5=-1.
\end{equation}
Das zentrale $-1$ wirkt auf Vektoren trivial, auf Spinoren mit $-1$. Entsprechend tragen die ursprünglichen Spinorfelder $G$-Phasen $+i$ oder $-i$ und sehen Ordnung vier. Die $D_5$- und $D_3$-Ströme sehen $+1$, die gemischten $60$ Ströme $-1$. Genau daraus entsteht wieder die originale Klassenzählung $52+64+60+64$. Die Vorzeichenkonvention des Spin-Lifts kann $G$ und $G^{-1}$ vertauschen; die festgehaltene Orientierung bestimmt sie.

Damit ist $\mu_4$ eine tatsächliche Wirkung auf den ursprünglichen Sektoren. Der duale Glue-Charakter $G$, der gemeinsame Quotientsgenerator $h$ und die geometrische Vierteldrehung sind verschiedene, jetzt ausdrücklich beschriebene Operatoren. Ihre Unterscheidung erlaubt es, die alte Überlagerungsgeometrie an dieselbe Quelle anzuschließen, ohne ihre Wirkungen gleichzusetzen.

\subsection{Vakuum und Konformalzeit folgen zusammen}
Für die rekonstruierte Neveu--Schwarz-Darstellung sind die Moden und das Vakuum durch
\begin{equation}
\{\psi_r^a,\psi_s^b\}=\delta_{ab}\delta_{r+s,0},\qquad
r,s\in\Z+\tfrac12,\qquad \psi_r^a\Omega=0\quad(r>0)
\end{equation}
festgelegt. Der gemeinsame geerbte Stress-Tensor gibt
\begin{equation}\label{eq:inverseL0}
\boxed{L_0=\sum_{a=1}^{16}\sum_{r>0}r\,\psi_{-r}^a\psi_r^a,
\qquad L_0\Omega=0,\qquad c=8.}
\end{equation}
Die Realisierung des geraden CAR-Netzes, seiner Sektoren und der Virasorowirkung ist in~\cite{bockenhauercar} explizit konstruiert. Hier wird sie an die ursprüngliche TFPT-Markierung angeschlossen.

In Kreisradius-eins-Einheiten gilt $B_{\rm NS}=-i\partial_\theta$ mit antiperiodischer Randbedingung. Der positive Anregungsraum hat Energien $r=1/2,3/2,\ldots$ und $L_0=d\Gamma(h)$ mit $h=B_{\rm NS}|_{\mathcal K_+}$. Eine Zweitquantisierung von $|B_{\rm NS}|$ auf dem ganzen verdoppelten selbstdualen Raum würde die Moden doppelt zählen. Bei der Feldkonvention $B(e_r)=\psi_r$ ist die CAR-Zustandskovarianz die negative Fourierprojektion; der positive Anregungsraum und diese Kovarianz sind entsprechend konjugiert.

Die Antworten werden dadurch explizit. Auf der Ebene ist
\begin{equation}
\langle\psi^a(z)\psi^b(w)\rangle=\frac{\delta_{ab}}{z-w}.
\end{equation}
Für den positiven Halbzylinder, $\tau>0$, $\vartheta=\theta-\phi$ und Kreismaß $d\theta/(2\pi)$ liefert dieselbe positive Modensumme
\begin{equation}\label{eq:inversekernel}
K_{ab}(\tau,\vartheta)
=\delta_{ab}\sum_{n\ge0}e^{-(n+1/2)(\tau-i\vartheta)}
=\frac{\delta_{ab}}{2\sinh((\tau-i\vartheta)/2)}.
\end{equation}
Gerade höhere Majorana-Antworten sind Pfaffianen dieses Zweipunktkerns; ungerade verschwinden. Die alten C/X-Spinorfelder behalten ihre bereits berechneten Gittervertex-Antworten. Ihr gemeinsamer gerader Teil hat dieselben Vakuumantworten in beiden Erweiterungen.

Eine geometrische Vierteldrehung hat nun den tatsächlichen Spin-Lift
\begin{equation}
\rho_F=e^{i\pi L_0/2},\qquad \rho_F^4=(-1)^F,\qquad \rho_F^8=I.
\end{equation}
In der bosonischen Erweiterung $0\oplus s$ sind die Gewichte ganzzahlig, also $\rho_{E_8}^4=I$. Damit haben Double Cover, Vierteldrehung und Spinstatistik einen gemeinsamen Operatoranschluss. Geometrische Rotation, interner Glue-Charakter und eine Wilsonholonomie am rohen Nahtoperator sind trotzdem verschiedene Wirkungen. Eine originale Viertel-Wilsonholonomie muss als konkreter verdrehter Sektor identifiziert werden; sie wird nicht still zur NS-Randbedingung. Der festgehaltene konforme Stress-Tensor lässt innerhalb dieser Quelle keine unabhängig gewählte alternative Konformalzeit zu. Eine physische Zeiteinheit ist damit noch nicht kalibriert.

\subsection{Der lokale Rekonstruktionssatz: Was am Ursprung tatsächlich genügt}\label{sec:originkernelcriterion}
Der Rückweg liefert ein scharfes hinreichendes Ursprungskriterium. Angenommen, der ursprüngliche Prozess erzeugt eine unitäre, positiv graduierte chirale Vertexsuperalgebra mit eindeutigem Vakuumraum $V_0=\C\Omega$ und sechzehn reellen, gegenseitig graduiert lokalen Primärfeldern $\psi^a$ vom Gewicht $1/2$. Ihr positives Zweipunktprodukt sei nicht ausgeartet und in einer orthonormalen Basis normiert; ihre markierte Darstellung sei das originale $V_{16}$. Schließlich erzeugen diese Felder die betrachtete fermionische Quelle vollständig. Die zugehörige ursprüngliche Spinormarkierung bleibt erhalten.

\textbf{Dann sind die lokale Fermionenquelle, ihr Vakuum und ihr konformer Generator gemeinsam festgelegt.} Der Beweis benötigt keine angepasste Folge höherer Antworten:
\begin{enumerate}
\item Das singuläre Vertexprodukt $\psi^a_{(n)}\psi^b$ besitzt das Gewicht $1/2+1/2-n-1=-n$. Positive Energie verbietet $n\ge1$. Für $n=0$ ist wegen $V_0=\C\Omega$ nur das normierte Vakuumprodukt möglich. Deshalb ist der ganze singuläre Anteil zwingend
\begin{equation}
\psi^a(z)\psi^b(w)\sim\delta_{ab}(z-w)^{-1}.
\end{equation}
Gewicht und gewöhnliche graduierte Lokalität schließen hier weitere singuläre Felder aus.
\item Die Moden erfüllen somit CAR. Positive Energie erzwingt $\psi_r\Omega=0$ für $r>0$. Normalordnung bestimmt jedes höhere Vakuumwort; vollständige Erzeugung macht die CAR-Fockwörter dicht. Der Zustand ist nicht nachträglich frei wählbar.
\item Der freie Stress $T_F=-\frac12\sum_a:\psi^a\partial\psi^a:$ besitzt dieselbe Primärwirkung. Die Differenz seiner Moden zu einem zweiten solchen Stress kommutiert mit sämtlichen CAR-Moden. Sie bildet das Vakuum daher auf ein skalares Vakuum ab; Graduierung und Vakuumnormierung setzen diesen Skalar auf null. Durch Erzeugung verschwindet die Differenz auf der ganzen endlichen Fockdomäne. Es folgt genau Gleichung~\eqref{eq:inverseL0}. Ein unabhängiger zusätzlicher Sektor wäre mit vollständiger Erzeugung unvereinbar.
\item Der gerade Teil ist $(D_8)_1$. Die mitgeführte ursprüngliche Spinorerweiterung rekonstruiert $E_{8,1}$ mit den ursprünglichen Marken und demselben Stress-Tensor.
\end{enumerate}
Dies ist ein Rekonstruktionssatz in der angegebenen chiralen Feldklasse, keine Behauptung, P1/P2 hätten seine Voraussetzungen bereits geliefert. Seine entscheidende Leistung ist die Kompression: Ein \emph{lokaler markierter Feldkeim mit bewiesener Konformalwirkung} ersetzt die getrennte Suche nach kompletter Algebra, Zustand, höherer Korrelationshierarchie und Zeitgenerator. Eine Symbolklasse, ein bloßer Rang oder CAR zu einem einzelnen Zeitpunkt allein erfüllen dieses Kriterium noch nicht.

\subsection{Der mikroskopische Ausgangspunkt ist im Originalbestand bereits vorhanden}
Der Vertrag \file{microscopic-charged-car-limit/README.md} konstruiert am tatsächlich verwendeten QWZ-Zylinder einen lokalen geladenen CAR-Kanal mit beiden Adjungierten, seinem gefüllten Quellvakuum und kontrolliertem Zeitgrenzwert. Die tragende ursprüngliche Operatoridentität lautet für die obere Randzeile
\begin{equation}
h(p)r=-\sin(p)\,r+(1-\cos p)\,r_-,
\qquad r=\delta_{y,7}\chi_+,\quad r_-=\delta_{y,7}\chi_-.
\end{equation}
Sie liefert den chiralen linearen Anteil und einen quadratisch kleinen Mischterm. Daraus werden im Original alle festen endlichen Feldwörter samt Adjungierten, der positive Zustand und der Generatorgrenzwert auf einem gemeinsamen Energiekern gewonnen. Das ist ein tatsächlicher mikroskopischer Feldanschluss, auf dem weiter aufgebaut werden kann.

Die dortige ursprüngliche Holonomie wird ausdrücklich erhalten. Im Sektor $a=1/4$ ist der physische Quellgenerator
\begin{equation}
H_a=L_0+(a-\tfrac12)Q=L_0-\tfrac14 Q.
\end{equation}
Die im Vertrag \file{source-localized-spinor-20260919/PROOF.txt} bewiesene lokale Phasentransformation zum NS-Kernel ist nicht auf dem ganzen Kreis periodisch und entfernt diesen Term nicht. Der neue Rekonstruktionssatz bestimmt $L_0$; der vorhandene Quellgrenzwert zeigt konkret, wie die zusätzliche ursprüngliche Holonomie an die reale Zeit gekoppelt bleiben muss.

Die entscheidende noch zu liefernde Identifikation ist jetzt ebenso konkret: Der ursprüngliche markierte Nahtoperator muss den geometrisch bereits hergeleiteten $V_{16}$ tragen, einschließlich seiner gemischten $10\leftrightarrow6$-Wirkungen und der passenden Holonomie. Der vorhandene mikroskopische Vertrag liefert einen komplexen Randkanal. Sein nachfolgender Spinorvertrag setzt acht Kopien ausdrücklich voraus. Die transversale Breite acht ist keine Ableitung dieser acht Arten; der originale Kanaltest berechnet für alle transversalen Strombilineare einen führenden Gramrang eins. Dieser Befund richtet die Arbeit auf die tatsächliche Träger- und Operatorherleitung aus, statt denselben Ein-Kanal-Grenzwert erneut zu optimieren.

Auch die ursprüngliche Symmetrie markiert genau die erforderliche Verbindung: Unter $K_f$ sind $(10,1)$ und $(1,6)$ verschiedene irreduzible Summanden; allein diese Wirkung erzwingt keinen gemeinsamen Zeitgenerator beider Teile. Die gemischten $60$ $D_8$-Ströme verbinden sie tatsächlich. Wird ihre Wirkung am ursprünglichen Operator nachgewiesen, ist eine nur intern wirkende Kommutante auf $V_{16}$ skalar. Das bindet beide Blöcke an denselben skalaren Raum-/Zeitkern; es beweist noch nicht dessen konkrete Form. Genau diese Operatorwirkung muss aus dem Originalprozess kommen. Sie ist im rekonstruierten Ziel bereits explizit vorhanden.

\textbf{Dafür genügt zusammen mit der vorhandenen Untergruppenwirkung bereits ein wirklicher Mischstrom.} Der invariante Feldkern hat zunächst die Form $K=K_{10}P_{10}+K_6P_6$. Für einen vakuum- und zeitverträglichen reellen Kreuzgenerator gilt
\begin{equation}
X=\begin{pmatrix}0&A\\-A^{\mathsf T}&0\end{pmatrix},\quad A\ne0,
\qquad [K,X]=0\ \Longrightarrow\ K_{10}=K_6.
\end{equation}
Die Untergruppenbahn eines solchen Generators spannt die irreduziblen $10\otimes6$ Kreuzrichtungen. Ihre Lieklammern schließen zur ganzen $\mathfrak{so}(16)$-Wirkung. In der ursprünglichen, phasenfesten C/X-Tabelle ist ein passender komplexer Wurzelstrom schon ausdrücklich
\begin{equation}\label{eq:rawmixwitness}
\boxed{[C_0,X_{0,1}]_{\mathfrak e_8}=E_{-e_1-e_6}.}
\end{equation}
Der Koeffizient ist in der tatsächlich verwendeten Kokzykluskonvention $+1$. Zusammen mit dem adjungierten Wurzelstrom liefert er reelle Kreuzgeneratoren. Der unmittelbare Herkunftstest lautet deshalb: Lassen sich die lokalen $V_{16}$-Felder und diese markierte Stromwirkung im Originalprozess so realisieren, dass ihre Wardantwort den gemeinsamen Kernel erhält? Das ist eine konkrete verbindende Operatorfrage; sie verlangt nicht die unabhängige Anpassung sämtlicher späterer Antworten. Die Identität~\eqref{eq:rawmixwitness} ist im gewählten Quellenmodell berechnet, ihr mikroskopischer Rohfeld-Intertwiner bleibt zu konstruieren.

\subsection{Wie das den Weg zur fundamentalen Lösung verändert}
Der Rücklauf
\begin{equation}
(E_{8,1};D_8,s)\longmapsto\mathcal F_{16}\ \text{mit Markierung }s
\longmapsto(E_{8,1};D_8,s)
\end{equation}
ist innerhalb der gewählten Quellenklasse geschlossen. Die schon vorhandene rekursive Rückbestätigung erhält damit ein Feld-, Vakuum- und Zeitgegenstück. Dieser Rücklauf liefert ein präzises Ziel für die Herkunftsherleitung; er ersetzt sie logisch nicht.

Am Originalprozess ist nun genau der in Abschnitt~\ref{sec:originkernelcriterion} bezeichnete Keim nachzuweisen: seine tatsächlichen lokalen Felder auf dem bereits konstruierten $V_{16}$, ihre halbe konforme Gewichtswirkung und ihre vollständige Erzeugung. Dazu gehört die originale Holonomie-/Spinorsektorabbildung. Diese Bedingungen sind gemeinsam zu bearbeiten. Neue Oszillatoren, zusätzliche Clockraten oder eine andere Gitterwelt würden diesen Nachweis nicht liefern.

Für den Universalraum bedeutet das: Sein positiver Geschichtskern ist nach dieser Identifikation bereits durch dieselbe Feldregel bestimmt. Auf dem gemeinsamen geraden Teil stimmen die beiden Darstellungen überein; geladene Geschichten müssen ihren NS-/Spinorsektor behalten. Zuses Komposition und ein kausaler Ereignisgraph können diese Felder und Geschichten organisieren, sofern sie ihre Produkte, Sektoren und Antworten erhalten. Sie müssen die Grundregel dann nicht erneut frei wählen. Die tatsächliche Clock-Auslese und die vierdimensionale lokale Realisierung werden weiter an dieser einen Quelle angeschlossen; eine chirale Rekonstruktion allein bestimmt sie noch nicht.

\picturetext{Bislang sah es so aus, als müssten wir für ein Instrument Saiten, Grundzustand, Klanggesetz und Takt einzeln finden. Der Rückweg zeigt: Beim vorhandenen Instrument gehören diese Teile bereits zwingend zusammen. Wir kennen jetzt die kurze lokale Regel, aus der sie gemeinsam entstehen. Die fundamentale Aufgabe ist, genau diese Regel am ursprünglichen Nahtprozess nachzuweisen. Danach müssen die räumlichen Messungen als wirkliche Betätigungen desselben Instruments erkennbar werden.}

\subsection{Ausführbare Prüfung und ihr genauer Umfang}
Die bestehende Quellenimplementierung wurde um \file{reconstruct_marked_fermion_source} ergänzt. Sie benutzt die originalen E8-Wurzeln und Feldmarkierungen, berechnet den geraden Glue-Fixpunkt, die $10+6$-Wirkung und die P2-Ladungen. Die Tests prüfen die Glue-Graduierung an sämtlichen ursprünglichen Wurzelklammern und die tatsächliche Wirkung der $D_8$-Wurzeln auf den sechzehn Gewichten.

Zwei unabhängig berechnete Charaktere stimmen überein: die NS-Fermionenbesetzung $\prod_{n\ge0}(1+t^{2n+1})^{16}$ und der ungerade Gittercharakter mit acht bosonischen Oszillatoren. Für $t^{2L_0}$ beginnen sie mit $1+16t+120t^2+576t^3+2076t^4$. Der mit originaler Chiraliät berechnete Spinorsektor beginnt mit $128t^2+2048t^4$. Gerade NS-Quelle plus Spinor ergibt $1+248t^2+4124t^4+\cdots$ und wurde bis $L_0=6$ zusätzlich gegen $E_4/\prod_{n\ge1}(1-q^n)^8$ geprüft. Diese endlichen Vergleiche kontrollieren Implementierung und Markierung; der allgemeine Rekonstruktionssatz steht im Beweis darüber. Er ist hier nicht als neuer Lean-Beweis ausgegeben.

\section{Wie die Gesamtlösung konkret erreichbar würde}\label{sec:completion}

\key{Die Aufgabe bleibt die gemeinsame \textbf{Herkunft, Prozessrekursion und physische Realisierung}. Die vorhandenen physikalischen Herleitungen aus Abschnitt~\ref{sec:physicalconstraints} geben dabei konkrete Auswahlbedingungen vor. Vorrang hat ihre simultane Realisierung am ursprünglichen markierten Flavor- und Nahtoperator mit derselben Determinantenfamilie. Der neue E8-Kandidat, sein Clockinstrument und seine Raumzeitfortsetzung müssen diese Antworten gemeinsam erhalten. Eine weitere Rechnung nur im symmetrischen Quellmodell genügt dafür nicht.}

\subsection{Erste Brücke: Die Quantisierung aus dem Ursprung auswählen}
Die Zielaussage lautet: Die ursprünglichen P1/P2-Daten samt den tatsächlich schon begründeten Markierungen erzwingen oder charakterisieren die minimale positive lokale Level-1-Vakuumquantisierung des markierten Gitters innerhalb einer ausdrücklich definierten Vergleichsklasse.

Der Auswahlsatz für die unitäre affine Vergleichsklasse ist in Abschnitt~\ref{sec:levelselection} jetzt konkret geführt: Eine unabhängig begründete Zentralzahl höchstens acht oder ein geeigneter langer Wurzelstrom mit Doppelmoden-Nullrelation erzwingt Level eins. Der nächste Herkunftsnachweis betrifft deshalb die ursprüngliche Realisierung genau dieses Stroms beziehungsweise dieser Zentralzahlgrenze. Primitive U1-Windung, Rang acht und die normierte Fünferrekursion dürfen nicht einzeln als bereits hinreichende Ersatzbegründung verwendet werden.

Die rekursiven P1/P2-Rückbestätigungen gehören in diesen Satz. Sie dürfen weder ignoriert noch so verwendet werden, als seien sämtliche Voraussetzungen ihres Abschlussrahmens bereits hergeleitet. Eine bloße Suche nach beliebigen alternativen Kanälen ohne die ursprünglichen Marken wäre hier kein Fortschritt.

Die Fortsetzung in Abschnitt~\ref{sec:nativeanomaly} liefert dafür einen direkteren Eingang: Der ursprüngliche globale $V_{16}$-Träger ist konstruiert. Wird seine tatsächliche gleichchirale Dirac-/Pfaffian-Realisierung nachgewiesen, folgt $c=8$ aus dem Familienindex, ohne zuvor E8 in die Rechnung einzusetzen. Der echte CAR-/Projektionsanschluss bestimmt dann den vollständigen reinen Zustand; die markierte lokale Spinorerweiterung schließt an E8 an. Die bordismustheoretische Route aus Abschnitt~\ref{sec:bordism16} prüft ergänzend den Abstieg der Antwort ohne Spinwahl. Der K3-Diracindex oder ein eingesetzter Komponentenfaktor ersetzt keinen dieser tatsächlichen Operatoranschlüsse.

\subsection{Zweite Brücke: Ein vollständiger Prozess unter Rekursion}
Die Zielgleichungen sind
\begin{equation}
\widetilde A V=VA,\qquad \widetilde U(t)V=VU(t),
\end{equation}
für die gemeinsam zugelassenen geladenen Operationen und die tatsächlich identifizierte Zeit. Die vorhandenen $W$-Intertwiner, Hecke-Sektorprojektoren und geladenen Carry-Verbraucher sind konkrete Bausteine. Der Nachfahrenturm und die bei endlichen Abständen sichtbaren Restsektoren dürfen dabei nicht abgeschnitten werden, wenn die Behauptung vollständige Antworttreue lautet.

Für radiale Wurzelfeldgeschichten ist diese Verbindung jetzt ausdrücklich konstruiert: $Q,A,P(t)$ bestimmen den vollen Vektor, seine zulässige Felderweiterung und den positiven Kernel; die Möbiusformel erhält den vollständigen Nachfahrenturm. Eine vorgegebene historische Vergröberungsregel muss auf diesen Daten als antworterhaltende Abbildung identifiziert werden. Der nächste Schluss ist somit kein beliebiges Anbauen von Restregistern, sondern der Nachweis, dass genau die vorhandene Rekodierung denselben Zustand und dieselben Operationen erhält. Das ist noch kein allgemeiner kausal normierter Detektorprozess.

Die Replica-Konstruktion präzisiert nun die Clockseite: Beide Beträge $2/3$ und $1/3$ treten in der signierten Vierpunktantwort auf. Die positiven dreidimensionalen Populationen erfordern jedoch eine wirkliche Auslese gemäß Gleichung~\eqref{eq:nativeclockmap}. Die unzulässige Gleichsetzung dreier Wickterme mit drei Zuständen und die reine $L_0$-Zeitumkalibrierung sind ausgeschlossen. Damit ist klar, welche Art von Prozessabbildung noch sinnvoll zu suchen ist.

\subsection{Dritte Brücke: Dieselbe Quelle in der physischen Raumzeit}
Die Originalverträge benennen bereits den richtigen gemeinsamen Prüfgegenstand. In passender Notation lautet er
\begin{equation}\label{eq:physical}
\operatorname{Res}_\Sigma W_{4D}[j]=W_{E_{8,1}}[j],
\end{equation}
zusammen mit
\begin{equation}
\operatorname{Res}_\Sigma\det D_{4D}\cong\det D_{\rm seam}
\quad\text{als Linienbündel mit Verbindung}.
\end{equation}
Zu spezifizieren sind die Einsetzungs- und Restriktionsabbildung, die zulässigen Quellen $j$, die Feldnormalisierung, P2-Ladungen, Clockmarkierungen, Zustand, Zeit und analytischen Gebiete. Die Gleichheit eines ausgewählten Zahlenwertes oder bloßer Determinantenbeträge reicht nicht.

Die entsprechenden Originalkennungen sind
\file{TFPT4D.LATTICE.ACTION.01},
\file{SEAM.DETLINE.UNIFICATION.01} und
\file{FTRANSFER.GENERATING.01}.
Die vorhandene endliche $\Z_6$-Funktionalroute und die $2+1$-dimensionale QWZ/$U(1)$-Determinantenroute sind dafür Vorarbeiten. Sie enthalten noch nicht den vollständigen chiralen $3+1$-$E_8$-Anschluss.

\subsection{Der gemeinsame Prozess bestimmt noch nicht seine räumliche Zuordnung}
Der stärkere Originalvertrag \file{BULK.PRIME_COMPLETION.01} verlangt ausdrücklich eine durch die Randoperationen erzeugte Realisierung ohne unsichtbaren Zusatzsektor. Unter dieser Bedingung dürfen beliebige angehängte Zuschauer nicht als Gegenbeweis gegen Eindeutigkeit verwendet werden. Sind auf einer zyklisch erzeugenden Wortalgebra sämtliche Antworten
\begin{equation}
D(A,B)=\omega(A^*B),\qquad C_t(A,B)=\omega(A^*\alpha_t(B))
\end{equation}
gegeben, wobei $\alpha_t$ eine $*$-Automorphismengruppe mit $\omega\circ\alpha_t=\omega$ ist, bestimmt $D$ die GNS-Darstellung. Die Zustandsinvarianz macht die Wortverschiebung $[A]\mapsto[\alpha_t(A)]$ isometrisch; die Gruppeninversen machen sie unitär. $C_t$ bestimmt sämtliche Matrixelemente dieser Entwicklung auf der dichten Menge der Wortzustände. Unter starker Stetigkeit ist damit auch ihr Generator bis auf die entsprechende unitäre Identifikation festgelegt. Das ist die positive Aussage des ursprünglichen KMS-Kohärenzvertrags; innerhalb dieses vollständigen Datenvertrags bleibt kein unabhängig frei wählbarer Hamiltonoperator übrig.

Noch nicht enthalten ist dadurch die Abbildung
\begin{equation}
O\subset\R^{1,3}\longmapsto\mathcal A_{\rm bulk}(O).
\end{equation}
Sie ordnet den Operatoren ihre räumliche Zugänglichkeit zu und muss mit Symmetrie, Lokalität, Quelle und Wechselwirkung gemeinsam verträglich sein. Bei der vorhandenen holomorphen E8-Netzstruktur liefern zwei alternierende Intervalle und ihre Vierteldrehung aufgrund der Zweiintervall-Dualität lediglich $M$ und $M'$. Ihr Schnitt ist $\C I$; das Vakuum erzeugt daraus nur eine Linie. Dies ist keine standarde Vierkeilkonfiguration für eine vierdimensionale Modularrekonstruktion. Der passende Rekonstruktionssatz benötigt gerade eine bestimmte gemeinsame Stellung der Algebren als Voraussetzung~\cite{kahlerwiesbrock}. Vier Marken allein liefern sie nicht.

\subsection{Konforme Kreiszeit und physische Translation müssen getrennt werden}\label{sec:physicaltimeaudit}
Auch die Formulierung, der vorhandenen Konformalzeit fehle lediglich eine physische Einheit, wäre für den exakten Poincar\'e-Vertrag zu schwach. Auf dem E8-Vakuumraum gilt $e^{-2\pi iL_0}=I$. Jede affine Wahl $H=aL_0+bI$ mit $a\ne0$ besitzt daher eine skalare Periode. Beim fermionischen Cover verdoppelt sich die Periode; das Argument bleibt bestehen.

\textbf{Satz.} In einer global Poincar\'e-kovarianten unitären Darstellung mit Translationen $V(x)$ erzwingt eine skalare Periode $V(Te_0)=\zeta I$, $T>0$, die Trivialität sämtlicher Translationen.

\textbf{Beweis.} Gegenläufige Boosts entlang der $i$-ten Raumrichtung liefern
\begin{equation}
V(T\cosh\eta,\,\pm T\sinh\eta\,e_i)=\zeta I.
\end{equation}
Ihr Quotient ergibt $V(0,2T\sinh\eta\,e_i)=I$. Weil $\sinh\eta$ alle reellen Werte annimmt, sind sämtliche Raumtranslationen trivial. Der Quotient einer geboosteten Periode mit der ursprünglichen ergibt danach $V(T(\cosh\eta-1),0)=I$. Diese Zeitabstände durchlaufen $[0,\infty)$; mit den Inversen sind auch alle Zeittranslationen trivial. Ein Vakuumshift $b$ oder eine Kalibrierung $a$ verhindert den Schluss nicht.

Dieser Satz schließt die direkte Identifikation der kompakten Quellenrotation mit einer nichttrivialen Minkowski-Zeittranslation aus. Er widerlegt keine approximative Raumzeit in einem Grenzübergang und keine gekrümmte kosmologische Realisierung ohne globale Poincar\'e-Symmetrie. Er gilt für die exakte Identifikation im genannten Originalvertrag.

Die vorhandene Möbius-Darstellung enthält bereits einen möglichen anderen Zeittyp. In einer Cayley-Konvention ist
\begin{equation}
P_{\rm line}=L_0-\tfrac12(L_1+L_{-1})
\end{equation}
der positive Generator der parabolischen Translation auf dem aufgeschnittenen Kreis. Er stammt aus demselben Stress-Tensor und derselben Quelle, ist aber keine affine Umrechnung von $L_0$. Die Wahl des ausgezeichneten Unendlichkeitspunktes gehört zu dieser Identifikation. Die Unterscheidung von Rotations-, Translations- und Sonderkonformgenerator und die zugehörige Positivität sind Standard in der Möbius-Darstellungstheorie~\cite{koesterconformal,morinellirehren}. Damit ist eine nichtperiodische chirale Bewegung vorhanden; ihre Identifikation mit physischer Zeit und die fehlenden transversalen Raumwirkungen folgen daraus noch nicht. Der vorhandene Vertrag \file{source-local-line-20260920/LOCAL_LIMIT_REVIEW.txt} liefert außerdem einen kontrollierten lokalen Großradiusgrenzwert. Dieser verändert die Spektralstruktur und ist durch den Satz über eine feste periodische Darstellung ausdrücklich nicht ausgeschlossen.

Die in derselben Primärarbeit beschriebenen räumlichen Deformationen zeigen, dass ein fester Zeitprozess durchaus verschiedene räumliche Realisierungen tragen kann. Dort werden jedoch eine passende vierdimensionale Einteilchendarstellung und ihre räumlichen Generatoren verwendet; die Konstruktion betrifft freie Felder. Sie darf weder als Beweis eines allgemeinen Dimensionshindernisses noch als bereits vorhandene wechselwirkende TFPT-Raumzeit ausgegeben werden.

\subsection{Warum die 24 Felder die Aufgabe wesentlich verkleinern}
Die 24 ursprünglichen Felder zusammen mit ihren Adjungierten erzeugen die gesamte Stromalgebra. Wenn eine physische Realisierung ihre vollständigen Level-1-Relationen, ein gemeinsames positives zyklisches Vakuum, analytische Lokalität und lokale Erzeugung nachweist, ist der erzeugte Vakuumsektor dadurch stark festgelegt. Unter den passenden Rekonstruktionsvoraussetzungen bestimmen diese Daten dieselben höheren Randantworten und dasselbe Netz.

Das ist die entscheidende Kompression des Problems: Statt tausend Antworten unabhängig anzupassen, werden die Erzeuger, ihr vollständiges Produktgesetz, der Zustand und die Zeit identifiziert. Dieser Satz allein bestimmt noch keine eindeutige vierdimensionale Erweiterung. Er macht aber ihren Randanschluss zugleich umfassend und überprüfbar.

Die Sechserantwort $1/576$ ist dabei ein notwendiger strenger Kontrollfall für jede behauptete Realisierung \emph{dieser} Quelle. Sie ist weder allein hinreichend noch ein Widerlegungstest für Modelle, die diesen Quellenanschluss gar nicht behaupten. Ihr Nutzen besteht darin, eine bloße Paarungs- oder Registerimitation unmittelbar zu erkennen.

\subsection{Welche physikalischen Zweige danach am selben Objekt hängen}
\begin{longtable}{L{0.21\linewidth}L{0.38\linewidth}L{0.33\linewidth}}
\toprule
Zweig & Gemeinsames Objekt & Entscheidende noch benötigte Identifikation\\\midrule\endhead
Eichstruktur & Geladene Ströme, Wardidentitäten, lokale Gauss-Wirkung & Interne Symmetrie als tatsächliche lokale Eichwirkung\\
Chirale Materie & Physische Feldoperatoren und ihre Spektren & Chirale Raumzeitwirkung und Spiegelentkopplung\\
Flavor/Massen & Pole und Mischungsantworten derselben Mehrpunktfunktionen & Yukawa-Dynamik und Maßstab, nicht nur Rangzählung\\
$\alpha$ & Elektromagnetische Wardantwort und Normierung & Herkunft genau der registrierten Selbstkonsistenzgleichung\\
Gravitation & Metrische Variation desselben physischen Funktionals & Dynamische geometrische Antwort und Quantisierungsbereich\\
Kosmologie & Mehrzeitfunktional einschließlich Anfangszustand & Physische Zeit, CP-/Nichtgleichgewichtsantworten und Transfer\\\bottomrule
\end{longtable}

Für Nichtgleichgewichtsfragen gehört der Anfangszustand ausdrücklich zum Schwinger--Keldysh-Objekt. Ein konformer Vakuumkorrelator allein ist kein kosmologischer Anfangszustand. Die bekannten bedingten Vorhersageformeln behalten ihren Wert als gemeinsame Zielantworten, ihre Transferannahmen müssen jedoch im selben physischen Objekt erklärt werden.

\subsection{Was nach diesen Prüfungen tatsächlich herzuleiten bleibt}
Die Suche setzt am jetzt bewiesenen Rekonstruktionskriterium aus Abschnitt~\ref{sec:originkernelcriterion} an. Die originale Naht muss die tatsächlichen lokalen $V_{16}$-Felder, ihre halbe konforme Gewichtswirkung und vollständige Erzeugung liefern. Daraus folgen CAR-Kernel, reiner Zustand, Konformalzeit und die markierte Erweiterungsroute gemeinsam. Der Rückweg von E8 liefert bereits das genaue Zielobjekt; die Herkunftsherleitung muss es unabhängig am ursprünglichen Prozess wiederfinden. Der skalare DtN-Wert $|k|$ aus den älteren Verifikationen bestimmt dabei noch nicht das Cliffordmodul oder diese Feldwirkung.

Auf der Clockseite wird die tatsächliche Cusp-Auslese in derselben Quelle gesucht, nicht ein weiterer Kanal mit vorab eingesetzten Eigenwerten. Ihr Erfolgskriterium ist die Operatorgleichung~\eqref{eq:nativeclockmap} einschließlich der ursprünglichen sechs Kompositionen. Die schon bekannte freie $S_3$-Kanalfamilie darf diese Herkunftsfrage nicht durch erneute Parameterwahl verdecken. Die neue gemeinsame Klassifikation schließt außerdem $q=1/3$ als autonome zeitunabhängige Markoventwicklung derselben Register aus. Deshalb muss ein vorgeschlagener Quellenanschluss seine Records und Zwischenoperationen mitführen. Ein bloßer Ein-Schritt-Abgleich $\mathcal E\circ\alpha_\delta\circ\iota=\Phi_q$ reicht nicht: Die identischen Abbildungen müssen auch die tatsächlich beanspruchten geordneten Mehrzeitantworten erhalten. Der lokale Puls liefert hierfür eine bereits definierte Quellenintervention statt einer zusätzlichen frei gewählten Dynamik.

Für die Raumfrage reicht die kanonische Vierteldrehung der schon vorhandenen Netze nicht: Die vorstehende Prüfung liefert keine geeignete vierdimensionale Stellung. Gesucht ist eine aus den Ursprungsoperationen begründete zusätzliche relative Lokalisierung mit tatsächlich transversalen Wirkungen. Sie muss die geordneten Antworten erhalten und die Raumzeitsymmetrie zusammen mit dem richtigen Translationstyp tragen. Neue räumliche Argumente dürfen nicht stillschweigend in eine vorausgesetzte $4D$-Masterwirkung eingesetzt werden.

Die so verbundenen Feld-, Zustands-, Ladungs- und Zeittests müssen gemeinsam bestehen. Die neue Realisierung muss dabei die bereits vorliegenden Massen-, Yukawa-, $\alpha$- und Newtonantworten aus Abschnitt~\ref{sec:physicalconstraints} unter ihren jeweiligen Originalvoraussetzungen reproduzieren. Die Fortsetzung hat mehrere vorher offene strukturelle Anschlüsse konstruiert, jedoch diese verbleibenden Operatoridentitäten und die vollständige Raumzeitphysik noch nicht bewiesen.

\section{Neue Anwendungen: Wo die Verbindung bereits produktiv wird}

``Neu'' hat hier zwei verschiedene Bedeutungen. Eine Anwendung kann für TFPT neu sein, obwohl das mathematische Verfahren bekannt ist. Wissenschaftliche Erstmaligkeit oder ein bahnbrechender Leistungssprung verlangt zusätzlich einen Vergleich mit der einschlägigen Literatur und den stärksten vorhandenen Verfahren. Die folgenden Anwendungen ergeben sich konkret aus den Sitzungsbefunden.

\begin{longtable}{L{0.18\linewidth}L{0.38\linewidth}L{0.36\linewidth}}
\toprule
Bereich & Konkrete Anwendung & Woran ihr Nutzen erkennbar wäre\\\midrule\endhead
Geometrie & Aus derselben markierten Quelle geladene Blockabbildungen und effektive Ausbreitung berechnen & Korrekte Antworten unter geänderter Zellgeometrie, einschließlich harmonischem Korrektor\\
Topologie & Holonomie und Deckphasen als endlichen beziehungsweise rationalen fortsetzbaren Zustand speichern & Gleiche nackte Gitterwege werden nur dann identifiziert, wenn alle erlaubten späteren Antworten gleich bleiben\\
Mathematik & Ein symbolischer Rechner für gemischte Stromwörter, verbundene Antworten und zulässige Fusionswege & Neue exakte Antworten ohne neue fitbare Ereignisregeln; automatische Gegenfälle gegen falsche Faktorisierung\\
Quanten\-information & Geladene Codes und Records mit nachgewiesenem Verbraucher bauen & Intertwiner erhalten die tatsächliche Operation; Leakage und Nachfahrensektoren werden quantitativ kontrolliert\\
Physik & Kandidaten für den gemeinsamen Rand-/Raumzeitanschluss an denselben höheren Antworten prüfen & Ein Kandidat besteht Ladungs-, Zeit-, Positivitäts- und Mehrpunktprüfung zusammen\\
Experiment\-design & Anhand von Ladung und Operatorstruktur gezielt höhere gemeinsame Antworten suchen & Weniger Versuche bis zu einem trennenden Befund als bei zufälliger oder rein paarweiser Auswahl\\
KI/Gedächtnis & Zustände nach späterer Handlungswirkung komprimieren und Gegenbeispiele als nötige Erinnerung behalten & Bessere wiederverwendbare Vorhersagen und Entscheidungen bei gleichem Speicher und Interaktionsbudget\\\bottomrule
\end{longtable}

Die besonders produktive mathematische Anwendung ist die automatische Entdeckung \emph{falscher Vereinfachungen}: Ein Rechner findet das kleinste gemischte Wort, an dem zwei vermeintlich gleiche Darstellungen unterschiedliche Antworten liefern. Das verbindet Beweisarbeit, Experimentwahl und Programmoptimierung unmittelbar.

Eine spekulative physikalische Anwendung ist der gezielte Entwurf von Systemen mit unterdrückten niedrigen und sichtbaren höheren Antwortkoeffizienten. Der Sechserzeuge liefert dafür ein exaktes algebraisches Muster. Die Umsetzung in ein konkretes Material oder Experiment würde jedoch eine zusätzliche Kopplungs-, Präparations- und Messabbildung erfordern. Das Dokument behauptet keine bereits konstruierte Hardware und keine neue fundamentale Sechskörperkraft.

\section{Welche Algorithmen daraus entstehen könnten -- und was daran bahnbrechend wäre}\label{sec:algorithms}

\key{Der stärkste gemeinsame Algorithmusgedanke ist ein \textbf{Compiler für zukünftiges Verhalten}: Er gewinnt aus tatsächlichen Antworten die nötigen Zustandsunterschiede, speichert sie in einer geeigneten Struktur und übersetzt zulässige Operationen in ein kleines, nachweislich antworttreues Programm.}

Dieser Gedanke verbindet die fünf folgenden Bausteine. Sie sollten als ein Arbeitsweg entwickelt werden, nicht als fünf neue Modellwelten. TFPT liefert dafür eine ungewöhnlich konkrete Quelle exakter Antworten und Gegenfälle. Die Universalraumidee liefert das Kriterium, welche Information erhalten bleiben muss. Symmetrie liefert eine Ordnung dieser Information; Verifikation sichert die erlaubten Umformungen.

\subsection{Der Neuheitsmaßstab und ein bereits vorhandener Gegenvergleich}
Zustände durch zukünftige Tests zu definieren ist bekannt. Aktive Automatenlerner, Predictive State Representations, Observable Operator Models und gewichtete Automaten verfolgen verwandte Ziele. Auch Verbindungen zu Tensor-Netzen sind publiziert \cite{littman,thon,adhikary,lsharp}. Die bloße Zusammenführung dieser Begriffe ist daher kein neuer Universalalgorithmus.

Ein historischer eigener Vergleich vom 25. September wurde für dieses Dokument am Originalbericht nachgelesen. Er untersuchte vier endliche Aufgabenfamilien, zehn Seeds, zwölf Aktionen und ein Budget von 30000 Operationen pro Lauf. Der reparierte TFPT-inspirierte Causal Compiler rekonstruierte 37 von 40 vollständigen Minimalmodellen; $L^\#$ mit derselben Heuristik erreichte 40 von 40. Die frühere Variante lag bei 30 von 40.

\begin{center}
\begin{tabular}{lrr}
\toprule
Familie & Reparierter Compiler & $L^{\#}$ mit Heuristik\\\midrule
Shutter & 20554.5* & 7675\\
Dense & 6159.5 & 3349.5\\
Stack & 11109 & 8375.5\\
SQLite & 12016 & 9327.5\\\bottomrule
\end{tabular}
\end{center}

Die Tabelle zeigt den Median der Operationen bis zum ersten vollständigen Minimalmodell. Eine Operation ist ein Reset oder eine Aktion einschließlich Readout. Beim mit Stern markierten Shutter-Wert werden ungelöste Fälle mit dem vollen Budget 30000 gezählt; dies ist ein gedeckelter Score.

Die reine Lernschleifenzeit über alle vierzig Läufe bis zum gemeinsamen Budgetende betrug dagegen 3,368 Sekunden für den reparierten Kern und 27,558 Sekunden für $L^{\#}$ mit Heuristik. Das ist keine Zeit bis zu einer selbst erkannten vollständigen Lösung. Weniger Umweltoperationen bedeuten also nicht automatisch weniger Rechenzeit.

Die Verbesserung des eigenen Lerners beruhte auf bekannter binärer Gegenbeispielanalyse. Bei einer passwortartigen Aufgabe ohne informative Zwischenantworten löste keiner der sieben Lerner die zehn Fälle im Budget. Ein SQLite-Beispiel zeigte zwar, dass gleiche sichtbare Tabellen nach einem späteren Rollback verschieden reagieren können; gewöhnlicher Zustands- beziehungsweise Transaktionsspeicher kann diese Information jedoch ebenfalls erhalten.

Der Originalvergleich umfasste 280 Läufe und unabhängige Modell-, Zeugen- und SQLite-Kontrollen. Er ist ein datierter früherer Befund, kein in dieser Sitzung erneut ausgeführter Benchmark und kein aktueller Hylæan-Test. Er verbietet, denselben Automatenlerner nun unter geometrischem Namen als Durchbruch zu vermarkten. Eine neue Methode muss einen anderen, konkreten Zusatz leisten: etwa positive operatorwertige Prozesse kontrolliert komprimieren, neue Zukunftsfehler mit weniger Eingriffen finden oder antworttreue Programme mit überprüfbaren Fehlergrenzen erzeugen.

\subsection{Baustein 1: Ein schneller Compiler für das vollständige Quellgesetz}
Die Forschungsfortsetzung hat hierfür bereits einen konkreten Teil umgesetzt: Für sämtliche endlichen C/X-Wurzelwörter wertet \file{evaluate_lattice_word} das ursprüngliche Gesetz durch quadratisch viele exakte Paaroperationen aus. Die tatsächlichen Kokzyklusphasen und geladenen Erweiterungssektoren bleiben erhalten. Die Ward-Rekursion ist der unabhängige Vergleich. Das ist eine im TFPT-Rechner nutzbare Anwendung einer bekannten Gitterformel; weder wissenschaftliche Erstmaligkeit noch ein gemessener allgemeiner Geschwindigkeitsrekord wird behauptet. Die Einschränkung auf Wurzelfeldeinsetzungen und die wachsende rationale Bitgröße gehören zum Leistungsumfang.

Die bereits ausgeführte Ward-Regel ist der Rechenkern. Ein weiterentwickelter Compiler würde gleiche Teilwörter nur einmal auswerten, ladungsunmögliche Beiträge vorab entfernen und tatsächlich äquivalente Anfragen kanonisieren. Für den bewiesenen harmonischen $X$-Sektor darf er den schnellen Weg über $W_5\otimes W_4$ verwenden. Für gemischte $C/X$-Wörter muss er auf die vollständige Regel zurückfallen.

Die konkrete Neuheitschance liegt in einer nachgewiesen guten Auswahl der Auswertungsreihenfolge, die Ladung, Sektor, Sparsität und Kosten gemeinsam berücksichtigt. Equality Saturation und E-Graphs sind dafür bekannte Werkzeuge \cite{egg}. Die erlaubten Gleichheiten müssten hier jedoch die tatsächlichen geladenen Operatorbedingungen erhalten. Die Alternative $W_6\otimes W_4$ ist beispielsweise durch den Wert $1/2$ gegen $1$ bereits ausgeschlossen.

\textbf{Prüfung.} Vergleich gegen einen guten Ward-Auswerter mit Memoisierung und dünn besetzter Arithmetik. Neue Prüfwörter enthalten gemischte Felder, Adjungierte und höhere Ordnungen. Ablationen entfernen Ladungsfilter, Symmetriekanonisierung und gemeinsame Teilrechnung getrennt.

\textbf{Möglicher Durchbruch.} Eine Größenordnung mehr exakt beantwortbare Anfragen bei gleichem Speicher und gleicher Zeit, auch außerhalb der zur Optimierung benutzten Wortfamilien. Dieses Ziel ist noch nicht erreicht. Ein begrenzter Geschwindigkeitsgewinn wäre dennoch nützlich, aber keine allgemeine neue Berechnungstheorie. Höhere Kumulanten besitzen kombinatorisch viele Partitionen; es gibt hier keine bewiesene universelle polynomielle Abkürzung.

\subsection{Baustein 2: Ein aktionsabhängiger Gedächtniscompiler}
Für bekannte endliche lineare Dynamiken sind die Bedingungen konkret. Ein Encoder $R$ erhält sämtliche Antworten, wenn für jede erlaubte Aktion $a$
\begin{equation}\label{eq:compiler}
RT_a=\widehat T_aR,\qquad\ell=\widehat\ell R.
\end{equation}
Dann folgt durch Induktion
\begin{equation}
\ell T_{a_n}\cdots T_{a_1}x
=\widehat\ell\widehat T_{a_n}\cdots\widehat T_{a_1}Rx
\end{equation}
für jedes endliche Wort. Dies ist ein wirklicher Vertrag über die Zukunft, keine Ähnlichkeit nur des momentanen Zustands.

Bei endlichen Quantenkanälen lässt sich der nötige Beobachtungsraum rückwärts schließen:
\begin{align}
\mathcal V_0&=\Span\{I,E_1,\ldots,E_m\},\\
\mathcal V_{k+1}&=\Span\!\left(\mathcal V_k\cup
\{T_a^*(E):E\in\mathcal V_k,\ a\in\mathcal A\}\right).
\end{align}
Im $d$-dimensionalen System gibt es höchstens $d^2$ unabhängige Operatorrichtungen. Für echte lineare Minimalität müssen zusätzlich die auf allen erreichbaren Zuständen unsichtbaren Richtungen entfernt werden. Der resultierende lineare Quotient muss nicht automatisch eine kleinere positive Quantendarstellung besitzen.

Für unbekannte Systeme liefert ein neues trennendes Zukunftswort die fehlende Zustandsunterscheidung. Der TFPT-Carry ist das konkrete Muster: Was heute unsichtbar ist, kann nach einer geladenen Operation wieder sichtbar werden. Die Innovation müsste in der Auswahl weniger, hochinformativer Fortsetzungen und in kontrollierter Kompression operatorwertiger oder stochastischer Prozesse liegen. Eine bloße Gegenbeispielschleife bleibt etablierte aktive Identifikation.

Auch eine angenäherte Version kann einen klaren Vertrag tragen. Sind $T_a$ und $\widehat T_a$ in den verwendeten Normen kontraktiv, gilt $\norm{RT_a-\widehat T_aR}\le\epsilon$, $\norm{\ell-\widehat\ell R}\le\delta$ und $\norm x\le1$, dann liefert die Teleskopsumme für Wortlänge $n$
\begin{equation}
\bigl|\ell T_wx-\widehat\ell\widehat T_wRx\bigr|
\le\delta+n\epsilon\norm{\widehat\ell}.
\end{equation}
Dies ist eine elementare Fehlerfortpflanzungsgrenze unter den genannten Annahmen, kein neuer allgemeiner Satz über Lernen. Sie zeigt aber, was ein praktischer Compiler versprechen kann. Eine kleine Ein-Schritt-Abweichung garantiert ohne zusätzliche Struktur keine kleine Abweichung für unbegrenzte Zukunft.

\textbf{Prüfung.} Bei deterministischen Aufgaben gegen $L^\#$; bei stochastischen Aufgaben gegen starke PSR/OOM-Verfahren und einen geeigneten Sequenzprädiktor. Verglichen werden neue längere Fortsetzungen, neue Kombinationen bekannter Aktionen, persistenter Speicher, Umweltabfragen und Wiederverwendung nach vollständigem Neustart.

\textbf{Möglicher Durchbruch.} Gleiche oder bessere Entscheidungs- und Vorhersageleistung mit wesentlich weniger dauerhaftem Speicher oder Interaktionen über mehrere verschiedenartige Aufgabenfamilien. Verschwindet der Vorteil bei gleichen Budgets und starken Baselines, ist die Überlegenheitsbehauptung zu verwerfen. Für beliebige unbeschränkte programmierbare Welten kann ein endlicher Testbestand keine universelle Zukunftsgarantie liefern.

\subsection{Baustein 3: Symmetriekompression mit erhaltenen Multiplizitäten}
Eine Symmetriedarstellung zerlegt sich als
\begin{equation}
\mathcal H=\bigoplus_\lambda(M_\lambda\otimes V_\lambda).
\end{equation}
Die Darstellung $V_\lambda$ beschreibt die Symmetriewirkung; die Multiplizitätsräume $M_\lambda$ tragen weitere tatsächlich veränderliche Information. Geladene Operationen können zwischen Sektoren wechseln. Ein korrektes Kompressionsverfahren muss deshalb sowohl die Sektoren als auch die relevanten Multiplizitäten gegenüber der ganzen Aktionsfamilie erhalten.

Symmetrische Tensor-Netze trennen solche strukturellen und freien Tensoranteile bereits \cite{singh,singhvidal}. Die hier zu entwickelnde Zusatzleistung wäre eine aktionsabhängige Auswahl der benötigten Multiplizitäten mit Antwortgarantie. Der TFPT-Befund $5\to5+70$ ist dafür eine konkrete Warnung: Eine kleine Darstellung allein ist noch kein unter späteren Operationen geschlossener Zustandsspeicher.

\textbf{Prüfung.} Vergleich gegen symmetrische Tensorverfahren und SVD-Kompression bei gleichem Speicher; eine Ablation entfernt nur die Berücksichtigung zukünftiger Aktionen. Gemessen wird der Fehler bisher ungesehener Operationsfolgen. Wenn jede Erweiterung des Handlungssatzes fast den ganzen Raum benötigt, bietet die Aufgabe wenig verlustfreie Kompression. Das ist eine Eigenschaft der Aufgabe, kein Anlass, nötige Information wegzudefinieren.

\subsection{Baustein 4: Gezielte Suche nach höheren gemeinsamen Antworten}
Der Sechserzeuge legt eine konkrete Suchregel nahe: Suche zuerst minimal ladungsneutrale Gruppen, deren echte Teilgruppen ladungsungleichgewichtig sind. Prüfe dann ihre vollständige und verbundene Antwort. So findet man Fälle, die eine rein niedrige Antwortbeschreibung übersehen kann. Minimalität der Ladungsneutralität ist eine hinreichende Auswahlhilfe für diese Zeugenklasse, keine vollständige Liste aller verbundenen Korrelationen.

Höhere Kumulanten und irreduzible gemeinsame Abhängigkeiten sind bekannte Gegenstände \cite{domino,schneidman2003}. Neu wäre ein nachweislich wirksamer Auswahlalgorithmus, der tatsächliche Operatorstruktur benutzt, um mit weniger Versuchen neue Vorhersagefehler aufzudecken. Für ein Handlungssystem könnten die Gruppen aus Kontexten und Aktionen bestehen; deren ``Ladungen'' müssten jedoch aus einer begründeten Aufgabenstruktur stammen, nicht aus einer dekorativen Physikanalogie.

\textbf{Prüfung.} Vergleich gegen zufällige Auswahl, ein paarweises Maximum-Entropie-Modell, geeignete latente Modelle und einen starken nichtlinearen Prädiktor. Paarweise Wechselwirkungen können bereits komplexe gemeinsame Statistiken erzeugen \cite{schneidman2006}. Ein nichtverschwindender höherer Kumulant beweist weder eine fundamentale Mehrkörperkraft noch Kausalität.

\textbf{Möglicher Durchbruch.} Deutlich weniger Eingriffe bis zum Auffinden unbekannter systematischer Fehler oder bis zu einer messbar besseren Entscheidung. Eine größere Sammlung berechneter Hyperkanten ohne Vorhersagewirkung wäre kein Erfolg.

\subsection{Baustein 5: Lean als Quelle ausführbarer, überprüfbarer Regeln}
Lean besitzt bereits eine Programmiersprache und Codegenerierung \cite{lean}. Die sinnvolle automatische Übernahme hängt von der Art des Originals ab:
\begin{center}
\begin{tabularx}{\linewidth}{L{0.30\linewidth}Y}
\toprule
Original in Lean & Nutzbarer Anschluss\\\midrule
Berechenbare Definition & Direkt ausführbarer Rechenkern oder Export\\
Konstruktiver Wert mit Eigenschaft & Algorithmus samt formalem Vertrag\\
Gleichheits-/Erhaltungssatz & Zulässige Optimierungsregel mit Voraussetzungen\\
Klassischer Existenzsatz & Noch kein ausführbarer Suchalgorithmus\\
Axiom oder \texttt{sorry} & Sichtbare offene Annahme beziehungsweise Beweislücke\\\bottomrule
\end{tabularx}
\end{center}

Beweise in \texttt{Prop} werden im üblichen Laufzeitcode entfernt. Eine Existenzbehauptung ohne berechenbare Konstruktion wird dadurch nicht zu einem Generator. Verifizierte Extraktion ist selbst ein etabliertes Forschungsgebiet \cite{metacoq}.

Der konkrete Weg wäre eine kleine gemeinsame Zwischensprache für rationale Arithmetik, Gitterinklusionen, Ladungen, Projektionen, Operatorwörter und Rekursion. Lean-Sätze begründen die erlaubten Umformungen. Der Quellencompiler nutzt sie nur unter ihren Bereichsbedingungen. Dadurch können wissenschaftliche Verifikationen zu wiederverwendbaren ausführbaren Regeln werden, statt in isolierten Prüfscripten zu enden.

\textbf{Prüfung.} Neue Sektoren, neue Operatoren und höhere Wortlängen müssen sowohl den formalen Vertrag als auch unabhängige Originalfälle bestehen. Ein Beweis einer falsch übersetzten Spezifikation bleibt falsch angeschlossen. Der Nutzen wäre ein schneller Kern, dessen Optimierungen zusätzliche Aufgaben abdecken und relevante Integrationsfehler verhindern; die reine Zahl grüner Lean-Sätze misst diesen Nutzen nicht.

\subsection{Ein gemeinsamer Algorithmus statt fünf neuer Architekturen}
\begin{center}
\begin{tikzpicture}[node distance=7mm]
\node[box,text width=12.5cm] (a) {Tatsächliche Antworten der vorhandenen Quelle oder Umwelt};
\node[box,below=of a,text width=12.5cm] (b) {Trennende Fortsetzungen finden: Welche Vergangenheit wirkt später noch?};
\node[box,below=of b,text width=12.5cm] (c) {Nötigen Zustand und Carry bestimmen; Symmetrie und Multiplizitäten erhalten};
\node[box,below=of c,text width=12.5cm] (d) {Antworttreues Programm kompilieren; bewiesene Abkürzungen verwenden};
\node[box,below=of d,text width=12.5cm] (e) {Nach Neustart anwenden und an neuer tatsächlicher Wirkung prüfen};
\draw[arr] (a)--(b);\draw[arr] (b)--(c);\draw[arr] (c)--(d);\draw[arr] (d)--(e);
\end{tikzpicture}
\end{center}
Der ambitionierte, messbare Anspruch lautet: Ein System gewinnt aus den Folgen seiner Handlungen ein kleines wiederverwendbares Programm und nutzt es nachweislich besser als bestehende Verfahren. Eine solche Leistung wäre ein möglicher Durchbruch. Die vorliegende Sitzung hat seine mathematischen Bausteine und Gegenfälle geschärft; sie hat die Überlegenheit noch nicht gezeigt.

\section{Was das für Hylæan konkret bedeuten würde}\label{sec:hylaean}

Der Anwendungsteil nimmt vorläufig an, dass mit ``hallieren'' Hylæan beziehungsweise Field Intelligence gemeint ist. Sein aktueller Produktionsstand wurde für dieses Dokument nicht neu geprüft. Die folgenden Aussagen sind ein konkreter Integrations- und Forschungsentwurf, keine Behauptung über bereits gelandete Fähigkeiten.

Der stärkste Transfer besteht darin, Gedächtnis nach späterer Wirkung auszuwählen. Eine gespeicherte Unterscheidung ist relevant, wenn sie nach einem Neustart eine Vorhersage, Entscheidung oder Handlung verändert. Eine interne Phase, ein Projektionskoeffizient oder ein Statusfeld allein belegt diese Wirkung nicht.

Ein vorhandener Hylæan-Pfad müsste dafür vollständig durchlaufen werden:
\begin{equation*}
\text{Eingabe}\to\text{Zustand}\to\text{Speicherung}\to\text{Restore}
\to\text{spätere Aktion}\to\text{beobachtbares Ergebnis}.
\end{equation*}
Die TFPT-Regel ist an jeder Stelle dieselbe: Nichts zusammenlegen, was eine zugelassene Fortsetzung noch unterscheiden kann. Dafür muss Hylæan weder künstlich ein $E_8$-Gitter erhalten noch Quantenphasen simulieren. Eine klassische Zustandsvariable kann genau der nötige Carry sein.

Der erste sinnvolle Einsatz wäre an einem bereits vorhandenen wiederkehrenden Handlungstyp. Zwei Vorgeschichten mit gleicher aktueller Anzeige, aber unterschiedlicher späterer Wirkung werden als Gegenfall benutzt. Der normale Producer muss die nötige Unterscheidung in den normalen Speicher schreiben; der reguläre Consumer muss sie nach einem frischen Prozessstart tatsächlich verwenden. Eine Ablation entfernt genau diese Unterscheidung und muss den erwarteten Fehler zurückbringen.

Erst danach ist die weitergehende adaptive Kompression sinnvoll: Welche Teile des Speichers können unter dem echten Handlungssatz entfernt werden, ohne die gemessene Zukunftsantwort zu verschlechtern? Ein trennender neuer Fall erweitert den Speicher gezielt. Der Quellencompiler liefert hier das Vorbild für kontrollierte Umformungen; die reale Umwelt liefert die Autorität über Bedeutung und Erfolg.

Als radikale Verbesserung wäre denkbar, dass Hylæan viel weniger irrelevante Historie dauerhaft speichert und trotzdem mehr nutzbare Erfahrung über Aufgaben hinweg wiederverwendet. Die passende Messgröße ist eine gemeinsame Kurve aus externem Erfolg, Interaktionen, dauerhaftem Speicher und Laufzeit. Ein einzelner Quotient darf nicht bessere Genauigkeit durch versteckte Zusatzabfragen oder größere Hilfscaches vortäuschen.

Die stärksten Kontrollen sind ein guter gewöhnlicher Zustandsprädiktor, derselbe Hylæan-Pfad ohne den neuen Selektor, ein frischer Neustart, unbekannte Fortsetzungen und neue Kombinationen bekannter Operationen. Ein Vorteil nur auf dem ursprünglichen Demonstrationsfall wäre kein allgemeiner Lernfortschritt. Ein gegenüber starken Baselines stabiler Gewinn könnte dagegen einen echten, systematisch nutzbaren Beitrag aus dem Universalraumprinzip machen.

\section{Verifikation, Lean und Experimente: Der tatsächliche Sitzungsstand}\label{sec:verification}

Die historischen Ergebnisse dieses Abschnitts wurden aus den gespeicherten Originalprotokollen der Sitzung gelesen. Die Fortsetzung hat die neuen Herkunfts- und Realisierungsrechnungen gezielt frisch ausgeführt. Der gesamte tausendteilige Bestand wurde dafür nicht erneut ausgeführt; die folgenden Prüfstände bleiben voneinander unterscheidbar.

\subsection{Frisch ausgeführte Forschungsfortsetzung}
Die früheren gezielten Läufe bestanden 45 beziehungsweise 58 Tests. Nach Integration des Anker-/Faservorschlags bestand der neue gemeinsame Lauf \textbf{67 Tests}. Zusammen geprüft wurden \file{test_source_selection.py}, \file{test_source_realization.py}, \file{test_source_program.py} und \file{test_marked_source_process.py}. Die Zahlen werden nicht addiert; die früheren Testmengen sind enthalten. Die anschließende Prüfung der gemeinsamen Ereignishierarchie und des nativen Raumlifts bestand in denselben vier Testdateien \textbf{75 Tests}. Dies ist der neue lokale Lauf; die zufällig gleiche Zahl eines extern genannten Prüfpakets bezeichnet eine andere, hier nicht vorliegende Testmenge. Die im neuen Anhang behaupteten 75 Prüfbedingungen eines nicht mitgelieferten Pakets sind hiervon getrennt.

Die neuen Tests prüfen die Paarungsnormierung symbolisch für allgemeine reelle Kreuzratio, ihren Gegenfall außerhalb der realen Geometrie, die gesamte Matrix aus einer unabhängigen Permutationsmischung, die tatsächlichen nativen Krausoperatoren, $B$ gegenüber $QB$, die höhere GNS-Antwort und den Verlust von $\rho_+$ gegenüber $\rho_-$. Ein symbolischer allgemeiner $2\times2$-Eingang prüft die kontrollierte Rückgewinnung mit erhaltenem Ereignisrecord. Die Reflexionszuordnung, Instrumentgrenzen und die Hauptsymbol-Gegenfamilie wurden zusätzlich unabhängig gegengelesen.

Die jüngste Erweiterung prüft die gemeinsame Replica-/Ankeridentität symbolisch für alle $x$, konkrete native Einbettungen bei $x=1/5,1/2,4/5$, die Wirkung der Originalcharaktere auf allen 240 Wurzeln und die orthogonale Erweiterung der wirklichen Trine. Der Gedächtniskern wurde gegen das Schurkomplement der vollständigen neundimensionalen Operatorwirkung geprüft. Weitere Kontrollen betreffen den geschlossenen C-Weg samt positiver Norm und verschwindendem Vakuumüberlapp, den tetraedrischen Richtungsrahmen sowie die rationale uniforme Levelschranke. Diese Tests sind keine physische oder neue Lean-Bestätigung.

Die Replica-Prüfung berechnet bei $x=1/5,1/2,4/5$ die elliptischen Perioden und unabhängig davon $E_4$ als Divisor-$q$-Reihe sowie $\eta$ als Produkt. Bei 60 Dezimalstellen Arbeitspräzision stimmen die normierten Antworten mit der rationalen Formel bis besser als $10^{-45}$ überein. Eine feste NS--NS-Fermionquelle kontrolliert denselben Überlagerungsfaktor unabhängig. Dies sind numerische Kontrollen der analytischen Theta-Herleitung, kein numerischer Beweis für alle $x$.

Die Anomalieprüfung verwendet die tatsächlichen 16 komplexifizierten Gewichte von $E+E^*+\Lambda^2F$ und drei freie symbolische Liftparameter. Sie bestätigt die ganze Familie $\lambda=d c_1^2-c_2$, ihre P2-Restriktion und den geraden Determinantenkreislevel algebraisch. Neue Rechnungen und ihre Reichweite wurden außerdem unabhängig gegengeprüft. Die modulare Operatorfamilie ist analytisch bestimmt; ihr vollständiger E8-Kokzyklus wurde nicht numerisch implementiert und kein neuer Lean-Beweis für diese Fortsetzung erzeugt.

Die mathematischen Level- und Realisierungsformeln wurden unabhängig gegengelesen. Der Geschichtskern wurde zusätzlich an einem gemischten C/X-Zeugen mit Wert $192/35$ über die ursprüngliche Ward-Regel bestätigt; die rationale Momentdarstellung und ihr Einfügungsfaktor wurden separat geprüft. Die ausführbaren neuen Resultate liegen unter \file{tfpt_explorer/runtime/source_selection.json} und \file{tfpt_explorer/runtime/source_realization.json}; die APIs unter \file{source_selection.py} und \file{source_realization.py}.

Nach dem konstruktiven E8/CAR-Rückweg bestand der gemeinsame Lauf derselben vier Dateien \textbf{77 Tests}. Die anschließende Präzisierung der internen Glue-Wirkung bestand zusätzlich den gezielten Lauf der beiden neuen Tests; diese zwei werden nicht zu 77 addiert. Ein unabhängiges Gegenlesen bestätigte den Rekonstruktionsbeweis, seine Kovarianzkonvention und die tatsächliche Markierungswirkung. Die native LaTeX-Kompilierung wurde nach der Integration geprüft. Weder die Testzahl noch die Kompilierung beweist die noch offene P1/P2-Herkunft oder vierdimensionale Physik.

Diese Fortsetzung verändert den nachfolgend dokumentierten älteren 316er-Schnappschuss nicht rückwirkend. Nach Integration des ursprünglichen Wegtransports, des neutralen Quellenzustands, des internen Spinlifts und der tatsächlichen Paaroperatoren bestand am 29.~September ein gemeinsamer gezielter Lauf \textbf{39 Tests} in den fünf zugehörigen Testdateien einschließlich der Syntheseintegration. Der frisch gestartete Explorer lieferte \textbf{24 Stufen mit 343 bestandenen internen Prüfungen}. Die neue Quellenidentität $4/z^2$ und die neuen Quellenverweise waren in der laufenden API vorhanden. Die erweiterten Spinlift- und Paaransichten wurden anschließend in der laufenden Tour geprüft; beide blieben nach dem Wechsel auf sechs Halbwege erhalten. Die sichtbare Paaransicht zeigte auch den ursprünglichen Familientransport und die neutrale Antwort $4/z^2$, ohne beobachtete Konsolenfehler. Die Paaroperatorrechnung und die neue Dokumentsektion wurden unabhängig gegengelesen. Dies sind Ausführungs- und Regressionsergebnisse; ihre Anzahl ist kein Beweis einer vollständigen physikalischen Lösung.

\paragraph{Gemeinsames Quellenwörterbuch, vorheriger Stand vom 29.~September.}
Die Fortsetzung integriert die tatsächliche Hyperladungsantwort, den expliziten
CAR-Paaranschluss, das vollständige schwere Majorana-Tensorwörterbuch mit
anschließender bestehender Seesaw-Rechnung sowie den gesamten neutralen
Flavor-Defektkern. Nach Korrektur der linken/rechten RHP-Feldkonvention bestand
im gemeinsamen gezielten Lauf \textbf{72 Tests} in den acht betroffenen
Testdateien. Der neu gestartete Explorer liefert \textbf{24 Stufen mit
367 bestandenen internen Prüfungen}. Die API enthält den korrigierten rationalen
RR-Kern und die vollständige schwere Tensorabbildung.
Die neue Ansicht \emph{Gemeinsame Physik} verbindet diese Antworten.
Die Phasenbeispiele wurden per Tastatur durchgeschaltet: Bei unveränderten
Massen ändern sich die berechneten Interferenzwerte, und der Phasenzustand
bleibt beim Wechsel des Familienhalbwegs erhalten. Dabei wurden keine
Konsolenfehler beobachtet. Der gemeinsame Mathematikanschluss wurde unabhängig
gegengelesen; die Rückprüfung der ursprünglichen Differentialgleichung hat
insbesondere die falsche Matrixreihenfolge als Gegenkontrolle ausgeschlossen.
Diese Prüfungen sichern die erklärte Konstruktion und ihre Anzeige. Sie ersetzen
keinen Herkunftssatz oder empirischen Nachweis der vollständigen Physik.

\paragraph{Gemeinsamer Massentransport, anschließender Stand vom 29.~September.}
Abschnitt~\ref{sec:masssourceconnection} verbindet die ursprüngliche
Determinantenwindung mit dem vollständigen Paartransport und der neutralen
Antwort. Im gemeinsamen gezielten Lauf bestanden \textbf{82 Tests} in neun
Dateien, einschließlich zehn neuer Prüfungen des Massentransports. Geprüft
wurden die tatsächlichen C-Gewichte, allgemeine nichtkommutierende
Verbindungen, die Umkehrabbildung, die endliche Zusammensetzung, die duale
Massenform samt Metrik und die lokale kovariante Punkttrennung. Zwei separat
integrierte Differentialgleichungen kontrollieren den ursprünglichen Weg
in beiden Richtungen. Der neu gestartete Explorer liefert \textbf{24 Stufen
mit 372 bestandenen internen Prüfungen}. Diese Zahl wurde am laufenden
Rechenstand abgelesen. Sie ergänzt die algebraischen Beweise, zählt sie aber
nicht als unabhängige physikalische Beweise.

\paragraph{Wand-Lesart und vollständiger geladener Quellenkanal, aktueller Stand.}
Die neue Prüfung in Abschnitt~\ref{sec:wallaudit} testet alle zehn symmetrischen
Paarzustände und ihre zehn Gegenladungen an sämtlichen acht nichtdiagonalen
SM-Generatoren, den drei nichtabelschen Cartan-Gewichten und der Hyperladung.
Ein farbgeladener Strom dient als Gegenkontrolle: Er wird nicht annihiliert.
Die Quelle besitzt die behaupteten Ladung-zwei-Operatoren; eine physische
Skalarabbildung oder Kondensation wird dadurch nicht behauptet.
Die zugehörigen neun gezielt ausgeführten Testdateien bestanden
\textbf{99 Tests}; der zusätzlich ausgeführte bestehende Syntheseintegrationstest
bestätigte die Weitergabe an den tatsächlichen Tourdatenpfad.
Der frisch gestartete laufende Explorer liefert \textbf{24 Stufen mit 375 bestandenen internen Prüfungen}; die API enthält den eingegrenzten Status \texttt{EXACT\_SOURCE\_CHARGE\_TWO\_OPERATOR\_EXISTS}. Die neue Ansicht \emph{Wand \& Neutrinos} und ihre aufklappbare Auswahlregeltabelle wurden am laufenden Browser geprüft, ohne beobachtete Warnungen oder Konsolenfehler.
Die Rang-vier-Spinorbündelidentitäten wurden symbolisch ausgewertet,
die Kommutanten an den tatsächlichen 240 E8-Wurzeln geprüft und die
Neutrino- sowie Bündelschlüsse unabhängig gegengelesen. Die fremden
27 Prüfungen wurden mangels lokalem Contract nicht reproduziert.
Der vorhandene Originaltest \file{v254_inner_fluctuations.py} lief ebenfalls
mit fünf grünen Prüfungen; sein bedingter Skalaron-Eintrag ist jedoch ein
festes \texttt{True} und kein Identifikationsbeweis.

\subsection{Der Explorer ist vollständig ausgeführt, der Originalbestand nicht vollständig grün}

\textbf{Fortsetzung vom 29.~September: wirklicher Majorana-Modenanschluss.}
Abschnitt~\ref{sec:masszeromode} erweitert denselben ausführbaren Pfad um die
Paarfeldwirkung auf allen 248 Stromrichtungen, ihre direkte Einsetzung in
den ursprünglichen 96-dimensionalen endlichen Diracoperator und die
Seesaw-/Gedächtnisantwort desselben schweren Blocks. Der gemeinsame
gezielte Lauf von Paaroperatoren, Neutrinowörterbuch, Massentransport und
Pipeline-Integration bestand \textbf{40 Tests}. Er enthält komplexe
Koeffizienten, korrekte Adjungierte, den Level-zwei-Gegenfall, alle unteren
Paargrade und die nichtstatische Resolventenantwort. Ein unabhängiges
Gegenlesen bestätigte den Modenbeweis und präzisierte die Einheiten der
Determinantenidentität. Weder die frühere Wand-Prüfung noch diese neue
interne Identifikation wird als Auswahl eines physischen Kondensats
oder als endgültige Gesamtlösung gezählt.
Der abschließende Explorer-Schnappschuss vom 28. September 2026, 18:42:38 UTC, enthält 24 Stufen und 316 bestandene interne Prüfungen, keine fehlgeschlagene. Die Stufenliste steht in Anhang \ref{app:stages}. Der letzte protokollierte vollständige Python-Testlauf bestand 216 Tests. Nach späteren kleinen Korrekturen wurden die betroffenen Quellwort-, Tour- und HTTP-Tests gezielt erneut ausgeführt: 14 Quellworttests, 18 Quellwort-/Tourtests und zwei HTTP-Tests. Diese Mengen überlappen und werden hier nicht zu einer erfundenen Gesamtzahl addiert.

Die fünf zuletzt ergänzten Tourstationen wurden im Browser geprüft, ohne beobachtete Konsolenfehler. Der live geänderte Sechserwert von $1/576$ auf $1/352$ zeigt, dass die Oberfläche den echten Quellenrechner anspricht. Eine funktionierende Tour beweist jedoch nicht die physische Identifikation ihrer Quelle.

Der umfassende Originalmodullauf endete am selben Tag um 15:12:50 UTC mit folgendem Status:
\begin{center}
\begin{tabular}{lr}
\toprule
Abgeschlossene Module & 1028\\
Bestandene Module & 1025\\
Fehlgeschlagene Module & 3\\
Bestandene Einzelprüfungen & 16564\\
Fehlgeschlagene Einzelprüfungen & 3\\\bottomrule
\end{tabular}
\end{center}
Die drei Fehlschläge haben unterschiedliche Bedeutungen:
\begin{longtable}{L{0.25\linewidth}L{0.67\linewidth}}
\toprule
Modul & Tatsächlicher Befund\\\midrule\endhead
\file{v879_envelope_partition_closure} & 27 Prüfungen bestanden, eine Verdict-Prüfung fehlgeschlagen. Der ausgegebene Status lautet \texttt{ENVELOPE-EMPIRICAL-ONLY + PARTITION-CLASS-CLOSED}. Eine empirische Envelope-Aussage wurde nicht zum verlangten Abschluss aufgewertet.\\
\file{v907_halfgap_registered_target} & 23 bestanden, eine fehlgeschlagen. Der blinde Holdout lieferte \texttt{WARD-BROKEN}; sechs statt erwarteter vierzehn Musterprüfungen, Fehler W5. Die eingefrorene Halbkonstante darf nicht nachträglich passend verändert werden.\\
\file{v916_epstein_weil_violation} & 19 bestanden, eine fehlgeschlagen. C1 meldet einen abweichenden Hash des Zertifikatsvektors. Die separate negative Quadratikprüfung C2 bestand; der Pinfehler bleibt dennoch ein echter Reproduktionsfehler. Das geprüfte Epstein-Objekt ist ausdrücklich kein RH-Beweis oder RH-Gegenbeweis.\\\bottomrule
\end{longtable}

Ein Modulfehler bedeutet hier nicht automatisch, dass eine fundamentale TFPT-Aussage falsch ist. Er darf aber auch nicht hinter den grünen Explorerprüfungen verschwinden. Gerade registrierte Gegenfälle und gescheiterte Abschlussbehauptungen sind Teil des Forschungsbestands.

\subsection{Der umfassendere Build-Audit}
Der Build-Audit von 08:12:54 bis 10:42:14 UTC endete ebenfalls mit Status \texttt{failed}. Synchronisierung, RH-Katalog und Theoriegraph bestanden; die RH-Arbeitsstufe scheiterte. Das Protokoll enthält 264 bestandene und eine fehlgeschlagene RH-Prüfung sowie 264 Probeeinträge, davon 263 bestanden und einer fehlgeschlagen. Diese zwei Zählungen bezeichnen unterschiedliche Ebenen.

Der fehlgeschlagene Probeeintrag ist \file{block_completion_probe.py}, Runde r484, mit \texttt{BLOCK\_COMPLETION FAILED}. Ein solcher endlicher Probeausgang entscheidet nicht über die Riemannsche Vermutung. Der Status des Audits bleibt gleichwohl fehlgeschlagen.

\subsection{Was tatsächlich aus Lean ausgeführt wurde}
Das Lean-Replay von 08:14:09 bis 08:15:00 UTC ist genauer als sein oberstes Zusammenfassungsfeld zu lesen:
\begin{center}
\begin{tabular}{ll}
\toprule
Phase & Status\\\midrule
Carrier-Kern & bestanden\\
RH-Bibliothek & bestanden\\
Native Brücke & bestanden\\
Automatischer Carrier-Export & bestanden\\
Automatischer RH-Export & fehlgeschlagen\\\bottomrule
\end{tabular}
\end{center}
Von 64 automatisch ausgewählten Kandidaten wurden 62 ausgeführt. Die beiden verbleibenden Definitionen \file{RH.PrimeWindow.cap} und \file{RH.VonMangoldtWindow.cap} benötigen konkrete Fensterobjekte als Argumente. Der Generator versuchte stattdessen, die Funktionen unmittelbar als Zahlen auszugeben; Lean meldete fehlende \texttt{ToString}-Instanzen für Funktionstypen. Das ist ein konkreter Adapterbedarf, kein Beweisfehler der zugrunde liegenden Sätze.

Das oberste Replayfeld lautet trotzdem \texttt{passed}. Eine vollständige Darstellung muss deshalb die Phasenlage berichten und darf dieses Feld nicht als ``alle Exporte erfolgreich'' lesen. Vier native Gegenvergleiche bestanden, darunter die Ankerpotenzen $[3,4,6,10,18,34,66,130,258]$ für Eingaben $0$ bis $8$ sowie Wurzelzahl 240, Rang acht und Dimension 248.

Der frühe browserseitige Inventarstand der Sitzung zählte 178 Lean-Dateien, 4571 Deklarationen, 39 Axiome und fünf \texttt{sorry}-Stellen. Das sind Inventarzahlen dieses damaligen Standes, keine Behauptung, jede Deklaration sei voraussetzungslos bewiesen. Sehr große Wall-Ladder-Zertifikate wurden beim Replay nicht erneut ausgeführt. Ein erfolgreicher RH-Bibliotheksbuild ist kein Beweis der RH.

\subsection{Nicht reproduzierte Belege und Quellenalter}
Die Konsolidierungs-PDF nennt zusätzlich 177 plus 19 Prüfungen. Die beiden dafür genannten Originalprogramme konnten im untersuchten Bestand nicht gefunden und deshalb nicht erneut ausgeführt werden. Ihre Behauptungen bleiben als Dokumentaussagen erkennbar; sie werden nicht zu den tatsächlich ausgeführten Prüfungen addiert.

Der verwendete Theoriegraph hatte bei seiner Prüfung 5969 Knoten und 64164 Kanten, stammte aber aus einem Stand vom 22. September. Die neuen Konsolidierungen vom 27./28. September waren darin nicht vollständig erfasst. Die Sitzung hat deshalb die neuen Originaldateien, den ausführbaren Explorer und die Protokolle direkt herangezogen. Ein Graphpfad ist eine Suchhilfe und keine neue Beweisquelle.

\section{Die jetzt tragende Gesamtaussage}

Die Sitzung hat mehrere lose Teile an derselben konkreten Quelle verbunden: Die ursprünglichen 24 Felder samt ihren Adjungierten erzeugen ganz $E_8$; der $X$-Block erklärt $A_8$; die alten und neuen Zerlegungen teilen denselben Stress-Tensor; die Pfadrekursion folgt aus einem echten Vierstromtensor; geladene Hecke-Records besitzen einen tatsächlichen Verbraucher; die vollständige Ward-Regel erzeugt höhere verbundene Antworten.

Die Fortsetzung schließt weitere mathematische Verbindungen an genau dieser Quelle: Die Stromfelder besitzen eine explizite Komposit-/Gittervertexdarstellung, alle endlichen C/X-Wörter eine geschlossene Produktformel, ihre radialen Geschichten einen positiven Kernel und ein rational darstellbares vollständiges Modengedächtnis. Die geometrische Quellenzeit bewegt diese Moden mit bekanntem Rest. Die Levelauswahl ist auf einen tatsächlichen Herkunftszeugen komprimiert; die normierte Fünferrekursion allein kann diesen Zeugen nicht ersetzen.

Das ist der gemeinsame mathematische Kern. Sein einfachstes Bild lautet: \emph{Ein Zustand ist das, was von einer Geschichte für alle erlaubten späteren Wirkungen übrig bleiben muss. Die Quelle legt diese Wirkungen fest. Rekursion und Graphen sind korrekte Darstellungen, wenn sie genau diese Wirkungen erhalten.}

Die Fortsetzung vom 29.~September liefert zusätzliche konkrete Anschlüsse: Die originale gemeinsame Verklebung besitzt einen globalen reellen Spin-Träger $V_{16}$. Seine gleichchirale Dirac-Realisierung hätte die vollständig berechnete gemeinsame Antwort $d c_1^2-c_2-p_1/3$. Innerhalb der erklärten E8-Vakuumquelle wählen die vier Quadratmarken eine echte Zweintervallalgebra samt modularem Mischer. Ihre Zweifach-Replica besitzt einen Ising-Fermion-Twist; dessen normierte Antwort ist exakt $3/2$ und dessen kohärente Paarungsbeiträge enthalten $2/3,-1/3,2/3$.

Das verbindet Herkunftsgeometrie, lokale Quelle, Universalraum und Clockzahlen wesentlich enger. Es identifiziert noch nicht den ursprünglichen Nahtoperator mit dem $V_{16}$-Diracfeld, die kohärente Replica-Auswertung mit dem historischen Populationstransfer oder die modularen Netzkopien mit einer $3+1$-Raumzeit. Die Suche ist deshalb jetzt auf diese konkreten gemeinsamen Operatoranschlüsse gerichtet. Die mathematischen Teilkonstruktionen sind geliefert; eine allumfassende physische Lösung wird weiterhin nicht als bewiesen ausgegeben.

Die anschließende Prüfung des neuen Gesamtvorschlags liefert über die Clockzahlen hinaus die ganze historische Matrix aus normierten Replica-Gewichten und einer nativen Operation. Sie zeigt zugleich: Derselbe gemittelte Kanal kann aus einer irreversiblen Messung oder aus protokollierten, einzeln umkehrbaren Quellereignissen entstehen. Der Ereignisrecord erhält die Zukunftsantwort, die in der einfachen Clockanzeige unsichtbar ist. Das ist die jetzt konkret konstruierte Verbindung von Clock, Universalraum und Ereignisgraph; ihre Auswahl aus dem ursprünglichen Nahtprozess ist der entscheidende nächste Herkunftsnachweis. Die vorgeschlagene Dimensionszählung aus drei Paarungen ist durch die tatsächlichen Kreisalgebren ausgeschlossen, und eine primitive Symbolklasse allein lässt sogar bei identischem Vakuum verschiedene Zeitantworten zu.

Der jüngste Ankeranschluss verstärkt dies: Die ganze Vierpunktfamilie besitzt denselben Viertelphasen-Transfer wie die Replica-Konstruktion und einen expliziten Intertwiner zum vorhandenen markierten C-Operator. Sein Gedächtnis ist direkt aus den ausgelassenen Kohärenzen berechnet. Die orthogonale Erweiterung der Trine erklärt die Detektorgrenze; die exakte Ladungsfaserung erklärt die innere Information geschlossener projizierter Wege. Zugleich zeigen Originalcharakter und Vakuumtest, warum diese Faserung noch keine räumliche Dynamik ist. Die vollständig zertifizierte Sewing-Schranke trennt höhere affine Level, sobald die passende Rohquellenantwort unabhängig vorliegt.

Der gemeinsame Zeit- und Lokalitätstest in Abschnitt~\ref{sec:sharedtime} liefert einen weiteren entscheidenden Abschluss: Die vollständige gemeinsame $S_3$-Markovhierarchie ist genau für $2/9\le q\le1/4$ möglich; dies schließt ihren historischen $q=1/3$-Endpunkt aus, obwohl dessen Einzelregister eine kontinuierliche Fortsetzung besitzt. Die geteilten Ereignisse allein übertragen keine Wirkung zwischen Registern. Die unmittelbare C-Hop-Identifikation verletzt zudem die originale Unitaritätsbedingung. Dem steht eine positive Rechnung an derselben Quelle gegenüber: Ein lokaler Stromimpuls erzeugt mit genau bekanntem Energieaufwand eine unter $L_0$ transportierte Antwort. Darauf und auf seiner ursprünglichen Auswahl beruht die nächste gemeinsame Herleitung; weitere simultane Vakuumsymmetrieprüfungen können die Ereignisgewichte prinzipiell nicht auswählen.

Der konstruktive Rückweg in Abschnitt~\ref{sec:inverseorigin} vereinigt schließlich die zuvor getrennten Herkunftsaufgaben: Die markierte E8-Quelle besitzt über ihr originales $D_8$ eine eindeutig fermionische Erweiterung mit genau dem ursprünglichen $V_{16}$-Träger. Vakuum, vollständiger Kern und Konformalzeit werden dabei gemeinsam zurückgewonnen. Der lokale Rekonstruktionssatz beweist, welche ursprüngliche Feldregel genügt, um diesen ganzen Anschluss zu schließen. Die noch zu liefernde Herkunft liegt präzise in dieser Feldregel und ihrer ursprünglichen lokalen und konformen Wirkung; sie wird nicht mehr durch vier getrennte Wahlprobleme ersetzt. Dieser bedingte Abschluss des chiralen Quellenproblems ist weiterhin von der physischen Clockauslese und der vierdimensionalen Realisierung zu unterscheiden.

Für Algorithmen ergibt sich parallel ein ernsthafter Forschungsweg: Die nötige Zustandsinformation aus Zukunftsantworten gewinnen, strukturiert komprimieren und als korrektes Programm wiederverwenden. Ob dies bahnbrechend wird, entscheidet sein Vorteil gegenüber den genannten starken Verfahren und am tatsächlichen Anwendungsweg. Die exakten Quellantworten und Gegenfälle der Sitzung machen diese Frage prüfbar.

\appendix
\section{Die 24 ausführbaren Stufen des letzten Explorerstands}\label{app:stages}
Die folgende Tabelle dokumentiert den tatsächlichen Lauf. ``Bestanden'' bedeutet, dass die hinterlegte interne Prüfung bestand. Die physischen Herkunfts- und Transferannahmen der Stufen bleiben davon getrennt. Die vollständigen Einzelwerte sind im gespeicherten \file{TFPT_Ablauf.json} enthalten.

\begin{longtable}{r L{0.22\linewidth}L{0.53\linewidth}r}
\toprule
Nr. & Kennung & Inhalt & Prüfungen\\\midrule\endhead
1 & \file{origin} & Ursprungsanker & 6\\
2 & \file{seam} & Seam und erhaltener Seed & 4\\
3 & \file{carrier} & Fünferträger & 3\\
4 & \file{e8} & $D_5+A_3$-Verklebung und $E_8$ & 37\\
5 & \file{clocks} & Coxeter- und Galois-Clocks & 4\\
6 & \file{flavor} & Flavor-Residuen und Auslesungen & 4\\
7 & \file{alpha} & Elektromagnetische Fixpunktgleichung & 2\\
8 & \file{transfer} & Transfer und rekursive Antwort & 4\\
9 & \file{predictions} & Eingefrorenes Register: 16 neu berechnete Werte & 2\\
10 & \file{gravity} & Bedingte gravitative Skalenbrücke & 2\\
11 & \file{cosmology} & Bedingter kosmologischer Zweig & 2\\
12 & \file{hamming} & Binärcode & 2\\
13 & \file{rays} & Sechzig Richtungen & 5\\
14 & \file{code} & Spektral ausgewählter Fünferraum & 13\\
15 & \file{observables} & Fünfzehn Antwortflächen & 3\\
16 & \file{quartic} & Nichtlineare Präparationsfläche & 2\\
17 & \file{binding} & Bindung durch gemeinsame Ereignisse & 3\\
18 & \file{recursion} & Ältere Dreierrekursion & 9\\
19 & \file{space} & Markierter räumlicher Kandidat & 12\\
20 & \file{sourcechannel} & Gemeinsame endliche Quellenfamilie & 34\\
21 & \file{matterbridge} & Typisierter physischer Anschluss & 8\\
22 & \file{observability} & Auslese und Kontrollabschluss & 6\\
23 & \file{normalization} & Vollständige Normierungsreparatur & 7\\
24 & \file{assembly} & Pfad, Raum, Hecke, geladene Quelle und Quellgesetz & 142\\\midrule
 & & \textbf{Gesamt, alle bestanden} & \textbf{316}\\\bottomrule
\end{longtable}

\section{Befundregister: Ergebnis, Reichweite und offene Verbindung}\label{app:claims}
\begin{longtable}{L{0.21\linewidth}L{0.38\linewidth}L{0.31\linewidth}}
\toprule
Gegenstand & Ergebnis und Status & Offene Verbindung\\\midrule\endhead
P1/P2-Bootstrap & \E\ Rückkopplung der markierten diskreten Struktur in ihrer Vergleichsklasse & Herkunft der vollständigen Quantisierungswahl\\
Markiertes $E_8$ & \E\ Gerade unimodulare Verklebung; wirkliche Ladungsdarstellung & Physische lokale Realisierung\\
Markierter E8/CAR-Rückweg & \E\ $D_8\to\Z^8$ rekonstruiert $V_{16}$, NS-Vakuum, Konformalzeit und ursprüngliche Spinorerweiterung & Herkunft des lokalen $h=1/2$-Feldkeims aus P1/P2\\
Quartikfünfer & \E\ Erster positiver Pauli-Invariantengrad unter Minimalgradregel & Herkunft dieser Auswahlregel\\
Alter Dreiecksencoder & \E\ Code und Ereigniskovarianz & Kein automatischer Anschluss an dreieckfreies Netz\\
Ladungspfad $W$ & \E\ Echter P2-Intertwiner, später aus Quellenantwort gewonnen & Vollständige endliche-Abstands-Rekursion\\
Neun-Port-Regel & \E\ Invarianz des Codes auf gegebenem endlichen Blockgraphen & Auswahl der zusätzlichen Kompositionsregel\\
Markierter Raum & \N\ Dreidimensionale Blochantwort mit Zellkorrektor & Auswahl der Perioden, physische Zeit und Lorentzstruktur\\
Clock-Fixalgebren & \Xclaim\ Falscher voller markierter $M_4$-Anschluss ausgeschlossen & Passender physischer Zeitanschluss der nativen Ebene\\
Trine-Erwartungen & \Xclaim\ Keine Folge gleicher Trinemessungen & Richtige experimentelle Auslese\\
Diskrete Ereignisfaser & \E\ Native Wortrealisierung und präzise RP-Historien & Herkunft der ausgewählten höheren Antwort\\
Kontinuierliche $S_3$-Raten & \Xclaim\ $q=1/3$ liegt außerhalb der nichtnegativen reversiblen Klasse & Andere physische Dynamikklassen separat prüfen\\
Geladener Hecke-Carry & \E\ Phasenerweiterung, Projektoren, tatsächlicher Verbraucher & Gesamte rekursive Raum-/Zeitabbildung\\
Cartan-/Quartiklift & \Xclaim\ Mehrere falsche Operatorgleichsetzungen ausgeschlossen & Physische Rolle der jeweils wirklichen Lifts\\
$C_4+X_{20}$ & \E\ Erzeugt alle 248 Ströme & Herkunft der gewählten Level-1-Quelle\\
Gemeinsamer Stress-Tensor & \E\ Gleichheit durch Projektoren, nicht nur Zentralzahlen & Laborzeit und ursprüngliche Clockantworten\\
$W_5\otimes W_4$ & \E\ Voller harmonischer $X$-Tensor & Kein Ersatz für gemischte Felder und Nachfahren\\
Sechser-Hyperkante & \E\ Exakt verbunden, $1/576$, keine echte Teilantwort & Physische Einsetzung derselben Felder\\
GNS/Universalraum & \E/\B\ Rekonstruktion bei positiver Quelle und festem Testsystem & Auswahl von Quelle und zulässigen physischen Eingriffen\\
Gittervakuumquelle & \B\ Vollständiger mathematischer Quellkandidat & Ursprungshypothese und $3+1$-Realisierung\\
Nativer $V_{16}$-Träger & \E\ Globaler Spin-Lift der ursprünglichen gemeinsamen Verklebung & Identifikation mit dem tatsächlichen Dirac-Multiplizitätsbündel\\
Gemeinsame Anomalie & \E/\B\ Klasse $\lambda=d c_1^2-c_2$; Pfaffianantwort mit $-p_1/3$ bei wirklichem gleichchiralen Symbol & Tatsächliche Symbol-, Realitäts- und CAR-Realisierung\\
Viermarken-Lokalisation & \E/\B\ Zweintervallalgebra, modulare Mischung und konjugierte Netzkopien in der erklärten E8-Quelle & Gemischte lokale Schnitte und physische Raumzeitordnung\\
Replica-/Clockanschluss & \E\ $R_2(1/2)=3/2$ aus wirklichem Ising-Twist; beide Clockbeträge in signierter Antwort & Cusp-Auslese, komponierbare Operation und sechs Schritte\\
Allumfassende Physik & \Opend\ Gemeinsame Funktional- und Determinantenidentität fehlt & Drei Brücken aus Abschnitt \ref{sec:completion}\\
Bahnbrechender Algorithmus & \Opend\ Konkrete Kandidaten mit Vorläufern und Vergleichsplan & Vorteil an unbekannten Aufgaben und echter Wiederverwendung\\\bottomrule
\end{longtable}

\section{Die wichtigsten falschen Gleichsetzungen -- und ihre Reparatur}
\begin{longtable}{L{0.42\linewidth}L{0.50\linewidth}}
\toprule
Nicht zulässiger Kurzschluss & Richtige Verbindung\\\midrule\endhead
Gleiche Dimension bedeutet gleiche Physik & Gruppenwirkung, Ladung, Graduierung und Operatorwirkung abgleichen\\
Alle fünfzehner Mengen sind dasselbe & Konkrete äquivariante Abbildungen und ihre Fasern bestimmen\\
Familien-$\sigma_F$ und Decktausch $D$ sind dieselbe Clock & Ordnung drei und Ordnung zwei samt Wirkung getrennt halten\\
Gleiche Pauli-Paarung bedeutet gleichen Reflexionslift & Eigenräume, Spuren und volle geladene Wirkung vergleichen\\
Gleiche Gitterwirkung bedeutet gleiche Geschichte & Deckcharakter beziehungsweise relative Phase mitführen\\
Ein gespeicherter Record genügt & Geladene Operationen müssen den Record tatsächlich fortschreiben\\
Untergitterindex ist Raumwachstum & Arithmetische Nebenklassen von räumlichen Perioden unterscheiden\\
Dreieckscode ist der räumliche Pfadcode & Tatsächliche Bindung, Geometrie und Intertwiner benutzen\\
Kleiner $W$-Code ist die volle Quelle & Gemischte Felder, Nachfahren und endliche-Abstands-Leakage erhalten\\
Gleiche Zentralzahl bedeutet gleiche Zeit & Hier ist der stärkere Projektor-/Virasorobeweis vorhanden; physische Zeit separat identifizieren\\
GNS-Eindeutigkeit wählt das Weltgesetz & GNS rekonstruiert die gewählte positive Quelle\\
Hankelrang ist GNS-Rang & Lineare Wortantwort und kohärenten Kreuzkern getrennt rekonstruieren\\
Verlustfreie Vergröberung löscht alle Unterschiede & Rekodierung von dissipativem Vergessen unterscheiden\\
Graphverzweigungen bestimmen Wahrscheinlichkeiten & Operatoren, Zustand und Normierung explizit herleiten\\
Höherer Kumulant beweist fundamentale Mehrkörperkraft & Antwortkoeffizient von mikroskopischem Hamiltonterm unterscheiden\\
Grüne Checks bedeuten fertige Gesamttheorie & Voraussetzungen, Auswahl, physische Abbildung und Beobachtung zusammen prüfen\\\bottomrule
\end{longtable}

\section{Quellenabdeckung und nachvollziehbare Fundstellen}\label{app:sources}

\subsection{Ursprüngliche Paper und neue Konsolidierungen}
Der lokale TFPT-Stammordner ist
\file{/Users/stefanhamann/Projekte/tfpt-theoryv4/}.
Die folgenden vom Nutzer genannten Paper bilden den ursprünglichen Quellenrahmen; soweit verfügbar wurden auch ihre durchsuchbaren TeX-Originale für die präzisen Aussagen herangezogen:
\begin{itemize}
\item \file{introduction.pdf}: Gesamtarchitektur und Lesekette.
\item \file{origin_theory.pdf}: P1/P2, Abschlussrahmen und Bootstrap; TeX insbesondere um Zeilen 87 und 955.
\item \file{tfpt_1_architecture_e8.pdf}: Träger, Gitterverklebung und $E_8$-Architektur.
\item \file{tfpt_2_standard_model.pdf}: Standardmodell- und Flavoranschlüsse.
\item \file{tfpt_3_e8_audit_bootstrap.pdf}: Audit und Rückkopplung der diskreten Struktur.
\item \file{tfpt_4_frontier.pdf}: offene physische Fronten und bedingte Fortsetzungen.
\item \file{tfpt_horizon_readouts.pdf}: Rand-/Horizontauslesungen und zugehörige Transferannahmen.
\item \file{tfpt_research_contracts.pdf}: präzise Abschlussverträge. Im TeX: direkte Quellenroute um Zeile 1271, $4D$-Masterroute ab 13235, lokale Aktion ab 13530, Determinantenanschluss 13616--13661, erzeugendes Funktional 13663--13698.
\end{itemize}

Die sechs neuen Markdown-Dokumente wurden als eigenständige datierte Quellen im Unterordner \file{_newest2/} behandelt:
\begin{enumerate}
\item \file{TFPT_Gesamtdokumentation_Code_Quartik_Rekursion_2026-09-26.md}.
\item \file{TFPT_Gesamtdokumentation_Ergebnisse_und_Herleitungen_2026-09-27.md}.
\item \file{TFPT_Gesamtdokumentation2_20260927.md}.
\item \file{TFPT_Universalraum_Gesamtdokumentation_20260927.md}.
\item \file{TFPT_Gesamtdokumentation_20260927.md}.
\item \file{TFPT_Universalraum_Gesamtdokumentation_2026-09-27.md}.
\end{enumerate}
Insbesondere enthält die letzte Fassung um die Zeilen 4348/4423 die ursprünglichen Ereignis-/Codeanschlüsse. Ähnlich benannte Dateien werden nicht stillschweigend für identisch erklärt.

Die neueste Konsolidierung ist
\file{/Users/stefanhamann/Documents/TFPT_Konsolidierung_20260928.pdf}.
Die Seiten 42--44 und 63--64 waren unter anderem für Pfad-, Raum- und Bindungsanschlüsse relevant. Eine Arbeitskopie liegt unter \file{tfpt_explorer/sources/}; die extrahierte Textfassung unter \file{tfpt_explorer/runtime/consolidation_text.txt}. Die nicht auffindbaren Programme zu den dort genannten 177+19 Prüfungen bleiben ausdrücklich unverifiziert.

\subsection{Die eingefügten Synthesevorschläge}
Die folgenden Texte liegen jeweils als \texttt{Eingef\"ugter Text.txt} unter
\file{/Users/stefanhamann/.codex/attachments/} mit den Kennungen:
\begin{description}[style=nextline,leftmargin=0pt,labelindent=0pt,labelwidth=\linewidth]
\item[400eb7a5-d4e0-4258-82b9-9cc2c14eeb9e] Ausgangspunkt bei der zuletzt berechenbaren Konstruktion.
\item[814806cc-6aaa-471c-b995-49152ae80280] Zuse und die Frage nach einem rechnenden Raum.
\item[c5fe92a6-cc32-44a2-9482-5dc0f7336b53] Zuse, neue TFPT-Ergebnisse und universeller Prozess.
\item[e3100bc3-7890-4847-87ed-c7c34300e06b] Starker Abschlussvorschlag über positive Geschichte, GNS und primitiven Fixpunkt.
\item[55f6a61b-610f-41af-9626-87e399bbb9d1] Gemeinsamer markierter Fermionprozess; neuer Pairing-Laplacian, verstärkte P1-Symbolklasse und vorgeschlagene Raumrichtungen. Geprüft in Abschnitt~\ref{sec:pairingcompletion}.
\item[14360bff-eb1f-42b4-a10f-04c973ecd011] Anker-Viertelphase, nativer C-Intertwiner, Aufzeichnungsinstrument, Ladungsfasern und uniforme Sewing-Schranke. Geprüft und weitergeführt in Abschnitt~\ref{sec:anchorfibre}.
\item[f548f0d4-2ae2-4067-8f27-3b19b7c82ae4] Gemeinsame Ereigniszeit, native Trine, Choi-Rang, Vakuumlokalität und Weyllauf. Eigenständig geprüft und zur vollständigen gemeinsamen Zeitklassifikation, Quellenliftprüfung und lokalen Impulsantwort weitergeführt in Abschnitt~\ref{sec:sharedtime}.
\item[9df45383-826d-4f53-8687-bf615454c26c] Orientierte und mehrteilige Antworten, Ladungskokzyklus und Grenzen einer bloßen Carry-Erweiterung. Geprüft in Abschnitt~\ref{sec:orientationaudit}.
\item[30ea245e-2aae-4de0-8aad-8e3655f3e12b] Spin-$\Z_4$-Anomalieauswahl, A3-Bandstruktur und Vorschläge zur physischen Prozesskette, samt vollständig gelesenem verlinktem Claude-Artefakt. Geprüft und über den tatsächlichen Spinlift, die nativen Paaroperatoren und ihre gemeinsame neutrale Antwort weitergeführt in Abschnitt~\ref{sec:spinliftsynthesis}. Die dort behaupteten externen Ausführungszahlen bleiben von den eigenen Läufen getrennt.
\end{description}
Diese Texte wurden als zu prüfende Aussagen behandelt. Ihre tragenden Ideen sind in den Abschnitten zu Quelle, Universalraum und Algorithmen verarbeitet; die zu starken Gleichsetzungen zu GNS, Primitivität, Metrik und physischem Abschluss sind im Haupttext ausdrücklich korrigiert. Anweisungen innerhalb solcher Dokumente wurden nicht mit dem Nutzerauftrag verwechselt.

\subsection{Originale Herleitungs- und Ausschlussorte}
Die folgenden Stellen tragen wichtige Detailaussagen und verhindern die Wiederholung bereits erkannter Sackgassen:
\begin{itemize}
\item \file{verification/v459_seam_lattice_voa_route.py}: direkte ursprüngliche Gitter-/VOA-Route.
\item \file{verification/v487_transfer_clock_rungs.py}: wirklicher Decktausch und Survivaloperator.
\item \file{verification/v3_em_alpha.py}: angegebene selbstkonsistente $\alpha$-Gleichung.
\item \file{verification/predictions_frozen.json}: eingefrorenes Vergleichsregister; kein Ersatz für Herleitung oder Messdatenprüfung.
\item \file{experiments/theory-contracts/matter-geometry-round15/PROOF.md}, um Zeile 126: ursprünglicher direkter Materie-/Gitteranschluss.
\item \file{experiments/theory-contracts/compiler-current-product-20260919/PROOF.txt}: geladener Stromproduktanschluss.
\item \file{experiments/theory-contracts/compiler-quartic-carry-20260919/PROOF.txt}: Carry-/Kompositionsbefund, Vertrag \file{UR.COMPILER.QUARTIC_LIFT.09}, Verdict \texttt{PARTIAL}.
\item \file{experiments/theory-contracts/compiler-quartic-holonomy-20260919/PROOF.txt}: wirklicher Quartiklift und Holonomie.
\item \file{experiments/theory-contracts/source-probe-carry-20260920/PROOF.txt}: \file{UR.SOURCE.PROBE_CARRY.01}, bedingter Anschluss, Verdict \texttt{PARTIAL}.
\item \file{experiments/theory-contracts/source-three-route-closure-20260922/SELECTION.md}, um Zeile 24: Auswahlgrenze der Quellenrouten.
\item \file{articles/2026-08-30/tfpt_thermodynamic_dynamics_en.tex}, ab Zeile 515: bereits vorhandener thermodynamischer Existenzrahmen.
\end{itemize}
Die Verdicts beziehen sich auf die jeweiligen Originalverträge. Neuere Explorerrechnungen schließen diese Verträge nicht automatisch administrativ oder logisch vollständig ab.

\subsection{Ausführbare Ergebnisse der Sitzung}
Für die Fortsetzung vom 29.~September tragen außerdem diese Originalstellen die native Anschlussfrage:
\begin{itemize}
\item \file{experiments/theory-contracts/raw-carrier-origin-gate-20260920/PROOF.txt}, ab Zeile 136: ursprünglicher reeller Träger $E_{\R}\oplus V_6$.
\item \file{experiments/theory-contracts/source-joint-spin-family-lift-20260922/PROOF.md}, ab Zeile 38: gemeinsame Liftfamilie und ihre Zentrumsglieder.
\item \file{experiments/theory-contracts/compiler-source-channel-gate-20260918/PROOF.txt}, ab Zeilen 133 und 167: Unterschied zwischen geometrischem Deckoperator und Glue-Clock sowie Erweiterungsindizes.
\item \file{_archive/tfpt-45/source_extracts/01_boundary_kernel_source.tex}, Zeilen 65--170 sowie 493 und 637: eingesetzte Dirac-Collarform, geometrische/Calderón-Reflexion und Übergang von elliptischem zu Dirac-Typ.
\item \file{verification/v155_quasifree_boundary.py}, ab Zeile 126; \file{v156_seam_net_construction.py}, ab Zeile 63; \file{v157_rigid_fixed_point.py}, ab Zeile 66; \file{v160_seam_gaussianity_from_pf.py}, ab Zeile 257; \file{v161_one_particle_reduction.py}, ab Zeile 170: freier Eingang, skalarer DtN-Wert, volle Schwinger-Primitivität und noch offene Operatoridentifikation. Die verkürzten Dateinamen liegen alle unter \file{verification/}.
\item \file{verification/v486_transfer_full_rule.py} und \file{verification/v814_k5_sixstep_transport.py}: ursprüngliche Clockregel, ihr Dreieroperator und die Reichweite der Sechsschrittidentifikation.
\end{itemize}
Die neuen ausführbaren Formeln stehen in \file{tfpt_explorer/source_selection.py::native_spin_anomaly} und \file{tfpt_explorer/source_realization.py::two_interval_replica_response}. Die bestehenden Builder geben sie gemeinsam mit den bisherigen Quellrechnungen aus. Die ältere Browser-Tour übernimmt diese Fortsetzung nicht automatisch.

Der neue Anhang ist zusätzlich in \file{source_realization.py::replica_pairing_readout},
\file{marked_source_process.py::replica_luders_action} und
\file{marked_source_process.py::replica_recorded_events} umgesetzt.
Die Funktionen verwenden die bestehenden Wurzeln und markierten Reflexionen. Die P1-Präzisierung wurde außerdem gegen \file{verification/v258_dirac_covariance_induction.py} und die Originalverträge \file{source-variation-origin-gate-20260922} sowie \file{source-boundary-lift-selection-20260922} geprüft. Dort stehen eine bedingte Kovarianzrekonstruktion beziehungsweise die offen bleibende Auswahl von Verbindung und Operator, kein schon bewiesener primitiver Dirac-Selektor.

Die jüngste Fortsetzung ergänzt \file{source_realization.py::anchor_clock_readout},
\file{anchor_clock_channel} und \file{family_charge_fibre}, außerdem
\file{marked_source_process.py::native_trine_dilation} und
\file{source_selection.py::affine_sewing_exclusion}.
Der markierte Intertwiner benutzt die Charakterformeln aus Abschnitt 4 des Originalvertrags \file{compiler-source-channel-gate-20260918/PROOF.txt}.
\file{verification/v221_seam_qecc.py} setzt die historischen sechsten Potenzen und die Kontrastrichtungen ein; \file{v486_transfer_full_rule.py} hält die Richtung Spektrum zu Raten ausdrücklich offen. Ihre Ausgaben wurden nicht als bereits bewiesene Auswahl der neuen Detektorbasis umgedeutet.

Die erneute Gesamtprüfung arbeitet analytisch an den tragenden Verbindungen; sie fügt keinen weiteren endlichen Quellen-Demozweig hinzu. Dafür wurden insbesondere die folgenden Originaleingänge gegeneinander gelesen:
\begin{itemize}
\item \file{_archive/tfpt-45/source_extracts/01_boundary_kernel_source.tex}, um Zeilen 470, 544, 629 und 937: Eingang der operativen Kategorie, Spektralfluss, Defektminimum und relativer APS-Rahmen.
\item \file{_archive/tfpt-45/source_extracts/02_carrier_source.tex}, um Zeile 3005: Mastervariation über bereits gewähltem diskretem Datum.
\item \file{_archive/tfpt-45/source_extracts/04_qft_source.tex}, um Zeilen 1750, 2150 und 2740: PhysAdm, effektive Wirkung und bedingte OS-/Streurekonstruktion.
\item \file{verification/v154_simple_current_theorem.py}, um Zeilen 53 und 95: exakte Erweiterungsarithmetik und externer physischer Netz-Eingang.
\item \file{verification/v260_k3_kummer_unification.py}, um Zeile 109; \file{verification/v367_seam_s3_lattice.py}, um Zeile 118; \file{verification/v470_alpha_inflow_level.py}, um Zeilen 125 und 161: K3, eingesetzte Majoranazahl und Unterschied zwischen Ladungs- und thermischem Index.
\item \file{tfpt_research_contracts.tex}, um Zeilen 1394, 12698, 13235, 13616 und 13663: ursprüngliche Grenzen von K3, Prime Completion, Masterwirkung, Determinantenlinie und gemeinsamem erzeugendem Funktional.
\item \file{experiments/theory-contracts/source-variation-origin-gate-20260922/PROOF.md} und \file{source-three-route-closure-20260922/SELECTION.md} im selben Vertragsverzeichnis: bereits dokumentierte Herkunfts- und Eindeutigkeitsgrenzen.
\end{itemize}
Die neuen analytischen Anschlüsse wurden unabhängig auf Gruppen, Grade, Statistik und Zustandsvoraussetzungen gegengeprüft. Sie sind hier mit Beweisen beziehungsweise Primärsätzen angegeben; eine neue Lean-Formalisierung oder eine vollständige erneute Ausführung aller Repo-Verifikationen wird damit nicht behauptet.

Die Herkunfts- und Transferprüfung der Fortsetzung verwendet zusätzlich folgende konkrete Originalstellen:
\begin{itemize}
\item \file{experiments/theory-contracts/microscopic-charged-car-limit/README.md}, Zeilen 167--202: ursprüngliche Viertelholonomie-Polarisation und geladene CAR-Moden.
\item \file{experiments/theory-contracts/compiler-source-channel-gate-20260918/PROOF.txt}, Zeilen 27--83: tatsächliche Kanalzahl und Rang eins der führenden Strom-Grammatrix.
\item \file{verification/v498_celestial_wp5b_singular_vector.py}, Abschnitte S4.10 und S4.13: Verhältnis von Viertel-Slot- und Periodenmoden sowie der bereits gitterbedingte Nullvektor.
\item \file{verification/v258_dirac_covariance_induction.py}, Präambel und Schlussprüfung: $\mu_{\rm geo}\Pi_{\rm odd}H_{\rm mod}$ und bedingter physischer Kompressionspfeil.
\item \file{experiments/theory-contracts/source-yukawa-forward-gate-20260921/yukawa_source_provenance_20260921.md}, ab Zeile 143: bereits geprüfte Vorwärtsschranken der Kovarianzroute.
\item \file{tfpt_research_contracts.tex}, Kennungen \file{TFPT4D.LATTICE.ACTION.01}, \file{SEAM.DETLINE.UNIFICATION.01} und \file{FTRANSFER.GENERATING.01}: gemeinsamer physischer Abschlussvertrag.
\end{itemize}

Alle folgenden Dateien liegen im Repository unter \file{tfpt_explorer/}. Sie sind der unmittelbare Ursprung der hier beschriebenen neuen endlichen Rechnungen; ihre Existenz ist kein zusätzlicher Beweis über ihren ausgewiesenen Geltungsbereich hinaus.
\begin{longtable}{L{0.37\linewidth}L{0.55\linewidth}}
\toprule
Datei & Inhalt\\\midrule\endhead
\file{origin_closure.py}, \file{core.py} & Originalanker, Rückkopplungen, ältere TFPT-Auslesungen\\
\file{process.py} & Ereignisse, Quartikcode, Bindung, Dreiecksrekursion\\
\file{quartic_selection.py}, \file{consolidation.py} & Molien-Auswahl, Präparationsfläche, Projektoren und Kontrollen\\
\file{composition.py} & Pfadblöcke, Neun-Port-Komposition und Leakage\\
\file{joint_constraints.py}, \file{spatial_response.py} & Gemeinsame Bindungsform und volle räumliche Antwort\\
\file{origin_transfer.py} & Ursprünglicher Transfer, Klasse und Fortsetzung\\
\file{cartan_source.py}, \file{marked_source_process.py} & Markierungsprüfung, Fixalgebren, Trine und native Ereignisfaser\\
\file{history_kernel.py} & Positiver kohärenter Geschichtskern und Rangvergleich\\
\file{hecke_source.py} & Untergitter, Charakterfortsetzung, Sektorprojektoren, geladener Carry\\
\file{twisted_source.py}, \file{torus_lift.py}, \file{source_process.py} & Tatsächliche Lifts, Phasen und ausgeschlossene Gleichsetzungen\\
\file{current_block_geometry.py} & $C_4+X_{20}$, Chevalley-Klammern, gemeinsamer Quellabschluss\\
\file{rewrite_coherence.py} & KZ-/Ward-Kompositionskonsistenz\\
\file{source_unification.py} & $A_8$-Gitter, $\Z_3$-Erweiterung und gemeinsamer Virasorovektor\\
\file{source_program.py} & Ganze Stromwörter, verbundene Antworten, harmonischer $W$-Anschluss\\
\file{source_selection.py} & Allgemeiner affiner Vierpunkttensor, Levelselektoren und explizite Level-zwei-Kontrolle\\
\file{source_realization.py} & Originale SU9-Feldabbildung, unabhängiger X-Wickweg, voller C/X-Gitterauswerter, positiver Geschichtskernel, rationale Momente und Möbius-Nachfahrenturm\\
\file{README.md} & Bedienung, Befunde und ausdrückliche Grenzen des Explorers\\\bottomrule
\end{longtable}

Die wichtigen Laufprotokolle liegen unter \file{tfpt_explorer/runtime/}:
\file{snapshot.json},
\file{verification_run_summary.json},
\file{verification_full.log},
\file{build_audit_summary.json},
\file{lean_replay.json},
\file{browser_check_summary.json},
\file{phase_selection_review.json}.
Der Stand der einzelnen Dateien unterscheidet sich zeitlich. Beispielsweise ist die frühe Browserzusammenfassung kein Nachweis aller später hinzugefügten Tourstationen.

Der exportierte vollständige Ablauf liegt unter
\file{/Users/stefanhamann/Documents/Codex/2026-09-28/unter/outputs/TFPT_Ablauf.json}.
Die frühere HTML-Gesamtprüfung,
\file{TFPT_Gesamtpruefung_2026-09-28.html},
und \file{Quellenfamilie_Pruefergebnis.json} im selben Ausgabeordner dokumentieren einen älteren Zwischenstand. Dieses LaTeX-Dokument konsolidiert die danach hinzugekommenen Verbindungen und Korrekturen.

\subsection{Historischer Algorithmusvergleich}
Der im Anwendungsteil verwendete Vergleich ist am datierten Original nachgelesen:
\file{/Users/stefanhamann/Documents/Codex/2026-09-25/und-verifiziere-und-beweise-und-teste/BEURTEILUNG.md}.
Er umfasst den reparierten endlichen Causal-Compiler, AALpy-Vergleiche, unabhängige Modellprüfungen, SQLite-Traces und Gegenfälle. Er belegt weder eine Hylæan-Integration noch eine Überlegenheit gegenüber modernen neuronalen Weltmodellen. Er ist hier aufgenommen, damit dieselbe bereits geprüfte Idee nicht erneut als ungetestete Neuheit erscheint.

\section{Begriffe in einfacher Sprache}
\begin{longtable}{L{0.22\linewidth}L{0.70\linewidth}}
\toprule
Begriff & Gemeinte Bedeutung\\\midrule\endhead
Quelle & Algebra von Operationen samt Zustand, aus denen gemeinsame Antworten berechnet werden\\
Ladung & Kennzeichnung, wie ein Feld unter einer bestimmten Symmetrie transformiert\\
Kokzyklus & Konsistente Phasenregel, die beim Zusammensetzen geladener Operationen benötigt wird\\
Carry & Mitgeführte Information, die bei späteren Operationen wieder beobachtbar werden kann\\
Intertwiner & Abbildung, bei der ``erst operieren, dann übersetzen'' gleich ``erst übersetzen, dann operieren'' ist\\
Ward-Regel & Rekursion, die höhere Stromantworten durch Paarung und Lieklammer bestimmt\\
Verbunden & Der Teil einer Antwort, der nach Abzug der Produkte echter Teilantworten übrig bleibt\\
Hyperkante & Gemeinsame Antwort mehrerer Einsetzungen; im Dokument zunächst ein Rechen- und Antwortobjekt\\
GNS & Aufbau eines Hilbertraums aus einer bereits gegebenen positiven Quelle\\
Hankelmatrix & Tabelle aller Antworten von Vergangenheit gefolgt von Zukunft; ihr Rang misst lineare Vorhersagedimension\\
Vergröberung & Wechsel der Auflösung; je nach Definition entweder antworttreue Rekodierung oder echtes Vergessen\\
Fixpunkt & Objekt, das unter einer festgelegten Abbildung unverändert bleibt; Gesetz und Zustand getrennt betrachten\\
Virasorovektor & Das algebraische Objekt, das den gemeinsamen konformen Stress-Tensor und seine Zeit-/Skalenwirkung bestimmt\\
KZ-Verbindung & Regel für konsistenten Transport konformer Blöcke beim Verschieben ihrer Einsetzungen\\
Nachfahren & Weitere Zustände/Felder derselben Quelle, erzeugt durch ihre Strommoden\\
Markierung & Physisch oder geometrisch ausgezeichnete Struktur, die bei einer angeblich gleichen Darstellung erhalten werden muss\\
Gesamtlösung & Gemeinsame Herkunft, Dynamik, Raumzeitabbildung und überprüfbare Antworten derselben Konstruktion\\\bottomrule
\end{longtable}

\begin{thebibliography}{99}
\addcontentsline{toc}{section}{Wissenschaftliche Primärliteratur}
\bibitem{razamattong} S. S. Razamat, D. Tong: \emph{Gapped Chiral Fermions}. Abschnitt 3, Gleichung (7), zur Spin-$\Z_4$-Anomalie. \url{https://arxiv.org/html/2009.05037v3}.
\bibitem{garciamontero} I. García-Etxebarria, M. Montero: \emph{Dai-Freed anomalies in particle physics}. Abschnitt 4.3 zur Spin-$\Z_4$-Struktur, dem rechtshändigen Neutrino und Majoranamassen. \url{https://arxiv.org/html/1808.00009v3}.
\bibitem{boylesmith} P. Boyle Smith, Y.-H. Lin, Y. Tachikawa, Y. Zheng: \emph{Classification of chiral fermionic CFTs of central charge $\leq16$}. Abschnitt 5.1, Gleichungen (5.3)--(5.6). \url{https://arxiv.org/html/2303.16917v2}.
\bibitem{beemchiral} C. Beem, M. Lemos, P. Liendo, W. Peelaers, L. Rastelli, B. C. van Rees: \emph{Infinite Chiral Symmetry in Four Dimensions}. Abschnitt 3 und Gleichung (57). \url{https://arxiv.org/html/1312.5344v3}.
\bibitem{costellopaquette} K. Costello, N. M. Paquette: \emph{Associativity of One-Loop Corrections to the Celestial Operator Product Expansion}. Insbesondere die Strom-OPE und die für die Anomaliekompensation zugelassenen Gruppen. \url{https://arxiv.org/html/2204.05301v1}.
\bibitem{arakidomain} S. Araki, H. Fukaya, T. Onogi, S. Yamaguchi: \emph{Symmetric Mass Generation for Domain-Wall Fermions}. Zweidimensionale euklidische Gitterrechnung mit acht Majorana-Flavors. \url{https://arxiv.org/html/2608.29963v1}.
\bibitem{hasenfratzguide} A. Hasenfratz, C. Xu: \emph{A Guide to Symmetric Mass Generation in Lattice-QCD}. Fassung 3, zur Bedeutung der konkreten UV-Wechselwirkung. \url{https://arxiv.org/html/2604.02424v3}.
\bibitem{gavrylenkomarshakov} P. Gavrylenko, A. Marshakov: \emph{Free fermions, W-algebras and isomonodromic deformations}. Abschnitte 4.1--5 zu Monodromiedefekten, Fermionenkernel und Determinanten. \url{https://arxiv.org/abs/1605.04554}.
\bibitem{bfkoriginal} D. Burghelea, L. Friedlander, T. Kappeler: \emph{Meyer--Vietoris type formula for determinants of elliptic differential operators}. Journal of Functional Analysis 107 (1992), 34--65. \url{https://doi.org/10.1016/0022-1236(92)90099-5}.
\bibitem{leebfk} Y. Lee: \emph{Burghelea--Friedlander--Kappeler's gluing formula for the zeta-determinant and its applications}. Zu lokalen Verklebungsfaktoren unter Produktstruktur am Rand. \url{https://arxiv.org/abs/math/0304250}.
\bibitem{yangmodular} T. H. Yang: \emph{Modular commutator as a robust topological invariant and approximate Markovianity}. 8.~September 2026, Theoreme 2.9 und 3.11 sowie Abschnitt 5.2.1. \url{https://arxiv.org/abs/2609.09019}.
\bibitem{gasslevinreplica} J. Gass, M. Levin: \emph{R\'enyi-like entanglement probe of the chiral central charge}. Physical Review B, 22.~Juni 2026; Gleichungen (5), (10)--(14) und ausdrücklich benannte Geltungsgrenzen. \url{https://arxiv.org/abs/2512.20608}.
\bibitem{shefferdefects} Y. Sheffer, R. Fan, A. Stern, E. Berg, S. Ryu: \emph{Probing chiral topological states with permutation defects}. Physical Review B, 15.~Mai 2026; Abschnitte II--V und Anhang I. \url{https://arxiv.org/abs/2512.04649}.
\bibitem{castroentropy} A. Castro, S. Detournay, N. Iqbal, E. Perlmutter: \emph{Holographic entanglement entropy and gravitational anomalies}. Gleichungen (25), (28) zur Rand-CFT mit festgehaltenem Regulatorrahmen. \url{https://arxiv.org/abs/1405.2792}.
\bibitem{thorngrendisentanglers} R. Thorngren, J. Preskill, L. Fidkowski: \emph{Chiral lattice gauge theories from symmetry disentanglers}. JHEP 07 (2026), 271, Abschnitte 1.1 und 5. \url{https://arxiv.org/abs/2601.04304}.
\bibitem{kapustinbordism} A. Kapustin, R. Thorngren, A. Turzillo, Z. Wang: \emph{Fermionic Symmetry Protected Topological Phases and Cobordisms}. Anhang, Gleichung (6), zur Quantisierung der gravitativen Chern--Simons-Antwort auf Spin- und orientierten Hintergründen. \url{https://arxiv.org/abs/1406.7329}.
\bibitem{freedhopkins} D. S. Freed, M. J. Hopkins: \emph{Reflection positivity and invertible topological phases}. Insbesondere Abschnitt 5.4, Conjecture 8.37, Remark 8.39 und Theorem 9.79 mit Folgerungen zur Unterscheidung topologischer Klassifikation, physischer Interpretation und Symmetrietypen. \url{https://arxiv.org/abs/1604.06527}.
\bibitem{absindex} M. F. Atiyah, R. Bott, A. Shapiro: \emph{Clifford Modules}. Topology 3, Supplement 1 (1964), 3--38. Abschnitte 6 und 11 zu reellen/komplexen Indexnormalisierungen. \url{https://www.math.ens.psl.eu/~benoist/refs/ABS.pdf}.
\bibitem{atiyahsinger3} M. F. Atiyah, I. M. Singer: \emph{The Index of Elliptic Operators: III}. Annals of Mathematics 87 (1968), 546--604. \url{https://annals.math.princeton.edu/1968/87-3/p07}.
\bibitem{chuzheng} Y.-J. Chu, Z.-J. Zheng: \emph{Vertex Operator Algebra Analogue of Embedding $D_8$ into $E_8$}. \url{https://arxiv.org/abs/0808.1458}.
\bibitem{barronfermion} K. Barron, N. Vander Werf: \emph{On permutation-twisted free fermions and two conjectures}. Abschnitt 6, Gleichungen (6.1)--(6.2) und anschließende Tensorpotenz-Isomorphie $V_{\Z^d}\cong\mathcal F^{\otimes2d}$. \url{https://arxiv.org/abs/1310.1958}.
\bibitem{bockenhauercar} J. B\"ockenhau-er: \emph{An Algebraic Formulation of Level One Wess--Zumino--Witten Models}. Abschnitte 2--4 zu den $\mathfrak{so}(N)_1$-Sektoren, NS-Moden und Virasorogeneratoren. \url{https://arxiv.org/abs/hep-th/9507047}.
\bibitem{dongxu} C. Dong, F. Xu: \emph{Conformal Nets Associated with Lattices and Their Orbifolds}. Proposition 3.15, Corollary 3.19 und Lemma 4.1 zu Gitternetzsektoren und internen Automorphismen. \url{https://arxiv.org/abs/math/0411499}.
\bibitem{doplicherroberts} S. Doplicher, J. E. Roberts: \emph{Endomorphisms of $C^*$-algebras, cross products and duality for compact groups}. Annals of Mathematics 130 (1989), 75--119. \url{https://annals.math.princeton.edu/1989/130-1/p03}.
\bibitem{rehren} K.-H. Rehren: \emph{Algebraic Holography}. Corollaries 1--2 und Proposition zum chiralen Rand und AdS$_2$. \url{https://arxiv.org/abs/hep-th/9905179}.
\bibitem{bdfs} D. Buchholz, O. Dreyer, M. Florig, S. J. Summers: \emph{Geometric Modular Action and Spacetime Symmetry Groups}. Theoreme 4.3.9 und 5.1.2 sowie Proposition 5.1.3; die Keilordnung gehört zum Eingang. \url{https://arxiv.org/abs/math-ph/9805026}.
\bibitem{schroer} B. Schroer: \emph{Lightfront Formalism versus Holography and Chiral Scanning}. Abschnitt 4 und Anhang zur zusätzlichen transversalen Lokalisation. \url{https://arxiv.org/abs/hep-th/0108203}.
\bibitem{os1975} K. Osterwalder, R. Schrader: \emph{Axioms for Euclidean Green's Functions II}. Communications in Mathematical Physics 42 (1975), 281--305. \url{https://doi.org/10.1007/BF01608978}.
\bibitem{gorardcausal} J. Gorard: \emph{Some Relativistic and Gravitational Properties of the Wolfram Model}. \url{https://arxiv.org/abs/2004.14810}; sowie \emph{Algorithmic Causal Sets and the Wolfram Model}. \url{https://arxiv.org/abs/2011.12174}. Verwendet zur Trennung kausaler Konsistenz von zusätzlichen Dimensions- und Dynamikannahmen.
\bibitem{cklw} S. Carpi, Y. Kawahigashi, R. Longo, M. Weiner: \emph{From vertex operator algebras to conformal nets and back}. \url{https://arxiv.org/abs/1503.01260}. Verwendet für den bedingten Übergang unitärer stark lokaler VOAs zu konformen Netzen.
\bibitem{donglin} C. Dong, X. Lin: \emph{Unitary vertex operator algebras}. Insbesondere Sätze 4.7--4.8 zur unitären affinen Unteralgebra. \url{https://arxiv.org/abs/1308.2361}.
\bibitem{frenkelkac} I. B. Frenkel, V. G. Kac: \emph{Basic representations of affine Lie algebras and dual resonance models}. Inventiones mathematicae 62 (1980), 23--66. \url{https://doi.org/10.1007/BF01391662}. Grundlegende Gittervertex-Realisierung; die Formeln der Fortsetzung werden mit den ursprünglichen TFPT-Wurzeln und Kokzyklen konkret ausgewertet.
\bibitem{laszlosorger} Y. Laszlo, C. Sorger: \emph{The line bundles on the stack of parabolic G-bundles over curves and their sections}. Abschnitt 2.3--2.6 zum Dynkinindex und Determinantenkokzyklus. \url{https://arxiv.org/abs/alg-geom/9507002}.
\bibitem{perelomovstates} D. Ojeda-Guillén, R. D. Mota, V. D. Granados: \emph{SU(1,1) and SU(2) Perelomov number coherent states: algebraic approach for general systems}. \url{https://arxiv.org/abs/1206.1555}. Die hier verwendete $h=1$-Reihe wird zusätzlich unmittelbar aus den Strommoden hergeleitet.
\bibitem{longomodular} R. Longo, P. Martinetti, K.-H. Rehren: \emph{Geometric modular action for disjoint intervals and boundary conformal field theory}. \url{https://arxiv.org/abs/0912.1106}. Mehrintervall-Modularstruktur und ihre Abgrenzung von der einfachen Möbiusbewegung eines Intervalls.
\bibitem{distler} J. Distler, E. Sharpe: \emph{Heterotic compactifications with principal bundles for general groups and general levels}. Insbesondere Abschnitt 6.1 zur $SU(9)/\Z_3$-Erweiterung. \url{https://arxiv.org/abs/hep-th/0701244}.
\bibitem{gui} B. Gui: \emph{Convergence of Sewing Conformal Blocks}. \url{https://arxiv.org/abs/2011.07450}. Analytische Voraussetzungen für Sewing.
\bibitem{guizhang} B. Gui, H. Zhang: \emph{Analytic Conformal Blocks of $C_2$-cofinite Vertex Operator Algebras III: The Sewing-Factorization Theorems}. Proc. Lond. Math. Soc. 132 (2026), e70130. \url{https://arxiv.org/abs/2503.23995}. Kein Ersatz für die Auswahl der physischen TFPT-Quelle.
\bibitem{carlyle} J. W. Carlyle, A. Paz: \emph{Realizations by stochastic finite automata}. Journal of Computer and System Sciences 5 (1971), 26--40. \url{https://www.sciencedirect.com/science/article/pii/S0022000071800053}.
\bibitem{segal} I. E. Segal: \emph{Irreducible representations of operator algebras}. Bulletin of the AMS 53 (1947). \url{https://www.ams.org/bull/1947-53-02/S0002-9904-1947-08742-5/S0002-9904-1947-08742-5.pdf}.
\bibitem{combs} G. Chiribella, G. M. D'Ariano, P. Perinotti: \emph{Theoretical framework for quantum networks}. Physical Review A 80, 022339 (2009). \url{https://arxiv.org/abs/0904.4483}.
\bibitem{zuse} K. Zuse: \emph{Rechnender Raum} (1969). Originalscan: \url{https://sferics.idsia.ch/pub/juergen/zuserechnenderraum.pdf}.
\bibitem{wolfram} S. Wolfram: \emph{A Class of Models with the Potential to Represent Fundamental Physics}, Abschnitt 5.19, \emph{Weighted Multiway Graphs}. \url{https://www.wolframphysics.org/technical-introduction/the-updating-process-for-string-substitution-systems/weighted-multiway-graphs/}.
\bibitem{littman} M. L. Littman, R. S. Sutton, S. Singh: \emph{Predictive Representations of State} (2001). \url{https://proceedings.neurips.cc/paper/2001/file/1e4d36177d71bbb3558e43af9577d70e-Paper.pdf}.
\bibitem{thon} M. Thon, H. Jaeger: \emph{Links between multiplicity automata, observable operator models and predictive state representations: a unified learning framework}. JMLR 16 (2015). \url{https://www.jmlr.org/papers/v16/thon15a.html}.
\bibitem{adhikary} S. Adhikary et al.: \emph{Quantum Tensor Networks, Stochastic Processes, and Weighted Automata}. AISTATS 2021. \url{https://proceedings.mlr.press/v130/adhikary21a.html}. Verwendet als Neuheitsvergleich der Modellklassen.
\bibitem{lsharp} F. Vaandrager, B. Garhewal, J. Rot, T. Wißmann: \emph{A New Approach for Active Automata Learning Based on Apartness}. \url{https://arxiv.org/abs/2107.05419}.
\bibitem{egg} M. Willsey et al.: \emph{egg: Fast and Extensible Equality Saturation}. \url{https://arxiv.org/abs/2004.03082}.
\bibitem{singh} S. Singh, R. N. C. Pfeifer, G. Vidal: \emph{Tensor network decompositions in the presence of a global symmetry}. \url{https://arxiv.org/abs/0907.2994}.
\bibitem{singhvidal} S. Singh, G. Vidal: \emph{Tensor network states and algorithms in the presence of a global SU(2) symmetry}. \url{https://arxiv.org/abs/1208.3919}.
\bibitem{domino} K. Domino, P. Gawron, Ł. Pawela: \emph{Efficient computation of higher order cumulant tensors}. \url{https://arxiv.org/abs/1701.05420}.
\bibitem{schneidman2003} E. Schneidman et al.: \emph{Network information and connected correlations}. \url{https://arxiv.org/abs/physics/0307072}.
\bibitem{schneidman2006} E. Schneidman et al.: \emph{Weak pairwise correlations imply strongly correlated network states in a neural population}. Nature 440 (2006). \url{https://www.nature.com/articles/nature04701}.
\bibitem{lean} L. de Moura, S. Ullrich: \emph{The Lean 4 Theorem Prover and Programming Language}. \url{https://lean-lang.org/papers/lean4.pdf}.
\bibitem{metacoq} M. Sozeau: \emph{Touring the MetaCoq Project}. EPTCS 337 (2021), 13--29. \url{https://arxiv.org/abs/2107.07670}.
\bibitem{freeddeterminant} D. S. Freed: \emph{Determinant Line Bundles Revisited}. Theorem 24 zur Krümmung der Determinantenlinie einer Diracfamilie. \url{https://arxiv.org/abs/dg-ga/9505002}.
\bibitem{baerballmann} C. Bär, W. Ballmann: \emph{Guide to Boundary Value Problems for Dirac-Type Operators}. Abschnitt 2 zur zusätzlichen Cliffordstruktur des Dirac-Symbols. \url{https://arxiv.org/abs/1307.3021}.
\bibitem{arakicar} H. Araki: \emph{On Quasifree States of CAR and Bogoliubov Automorphisms}. Publ. RIMS Kyoto Univ. 6 (1970/71), 385--442; Lemma 4.3, S. 390, zur Eindeutigkeit des Zustands bei gegebener Basisprojektion. \url{https://ems.press/content/serial-article-files/41673}.
\bibitem{calderonscattering} C. Guillarmou, S. Moroianu, J. Park: \emph{Bergman and Calderón projectors for Dirac operators}. Theorem 2 und Proposition 12 zur Calderón-Projektion und der zugehörigen konform kompakten Streukonstruktion. \url{https://www.imar.ro/~sergium/fisiere/calderon.pdf}.
\bibitem{takesakiexpectation} M. Takesaki: \emph{Conditional expectations in von Neumann algebras}. Journal of Functional Analysis 9 (1972), 306--321. \url{https://doi.org/10.1016/0022-1236(72)90004-3}.
\bibitem{connesradon} A. Connes: \emph{Noncommutative Geometry} (1994), Kapitel V, Abschnitt 5$\alpha$, Theorem 1(a), gedruckte S. 481: Radon--Nikodym-Kokzyklus und Zusammenhang modularer Gruppen. \url{https://alainconnes.org/wp-content/uploads/book94bigpdf.pdf}.
\bibitem{wiesbrockintersection} H.-W. Wiesbrock: \emph{Symmetries and Modular Intersections of Von Neumann Algebras}. Letters in Mathematical Physics 39 (1997), 203--212. Modulare Inklusionen und zusätzliche Reflexionsbedingungen; kein voraussetzungsloser $4D$-Rekonstruktionssatz. \url{https://doi.org/10.1023/A:1007361114049}.
\bibitem{kahlerwiesbrock} R. Kähler, H.-W. Wiesbrock: \emph{Modular theory and the reconstruction of four-dimensional quantum field theories}. Journal of Mathematical Physics 42 (2001). Die relative modulare Stellung der Algebren ist Rekonstruktionsvoraussetzung. \url{https://doi.org/10.1063/1.1327597}.
\bibitem{morinellirehren} V. Morinelli, K.-H. Rehren: \emph{Spacelike deformations: higher-helicity fields from scalar fields}. Letters in Mathematical Physics 110 (2020), 2019--2038. Abschnitte 2.1--2.4 zu Möbius-Zeit, räumlichen Generatoren und den benötigten Einteilchendarstellungen. \url{https://doi.org/10.1007/s11005-020-01294-w}.
\bibitem{koesterconformal} S. Köster: \emph{Conformal Transformations as Observables}. Letters in Mathematical Physics 61 (2002), 187--198. Abschnitt 2 und Proposition 1 zu konformer Zeit, Translation und positiver Energie. \url{https://arxiv.org/abs/math-ph/0201016}.
\bibitem{headrickreplica} M. Headrick: \emph{Entanglement Rényi entropies in holographic theories}. Gleichungen (A.1) und (4.24) zur Überlagerungsnormierung und Zweintervall-Replica-Antwort. \url{https://arxiv.org/abs/1006.0047}.
\bibitem{castroanomaly} A. Castro, S. Detournay, N. Iqbal, E. Perlmutter: \emph{Holographic entanglement entropy and gravitational anomalies}. Abschnitt II.B, insbesondere (2.16)--(2.17), zu getrennten chiralen Twistgewichten und Frameabhängigkeit. \url{https://arxiv.org/abs/1405.2792}.
\bibitem{petzrenyi} J. Kudler-Flam: \emph{Rényi Mutual Information in Quantum Field Theory}. Zur Unterscheidung der üblichen Rényi-Entropiekombination von einer relativen Rényi-Größe. \url{https://arxiv.org/abs/2211.01392}.
\bibitem{e8replica} L. A. León Andonayre, R. Poddar: \emph{Rényi entropy of single-character CFTs on the torus}. Physical Review D 112 (2025), 025011. Insbesondere Abschnitte 3.1--3.2 zur E8-Permutationsorbifold-Fusion und Charakterreduktion. Die dortige Torus-Einintervallgeometrie wird nicht mit unserer Sphären-Zweiintervallgeometrie gleichgesetzt. \url{https://arxiv.org/abs/2412.00192}.
\bibitem{collazuolmelnikov} V. Collazuol, I. V. Melnikov: \emph{A twist at infinite distance in the CHL string}. Abschnitt 4.1, Gleichungen (4.15)--(4.18), zum tatsächlichen $E_{8,2}$-Singulett-Twist mit Ising-Gewicht $1/2$. \url{https://arxiv.org/abs/2402.01606}.
\bibitem{ardonneising} E. Ardonne, G. Sierra:
\emph{Chiral correlators of the Ising conformal field theory}.
Journal of Physics A 43 (2010), 505402, Gleichung (19):
Cauchy-Pfaffian/Hafnian-Identität für beliebige gerade Punktzahlen.
\url{https://arxiv.org/abs/1008.2863}.

\bibitem{wolfjoint} M. M. Wolf, J. Eisert, T. S. Cubitt, J. I. Cirac: \emph{Assessing non-Markovian dynamics}. Physical Review Letters 101 (2008), 150402. Gleichung (2) und Anhang zum bedingten Choi-Kriterium. \url{https://arxiv.org/abs/0711.3172}.
\bibitem{gutoskijoint} G. Gutoski: \emph{Properties of Local Quantum Operations with Shared Entanglement}. Quantum Information and Computation 9 (2009), 739--764. \url{https://arxiv.org/abs/0805.2209}.
\bibitem{dongnagjoint} C. Dong, K. Nagatomo: \emph{Automorphism Groups and Twisted Modules for Lattice Vertex Operator Algebras}. Abschnitte 2.1--2.4 zu Kokzyklus, Isometrielift und Vakuumautomorphismen. \url{https://arxiv.org/abs/math/9808088}.
\bibitem{darianojo} G. M. D'Ariano, P. Perinotti: \emph{Derivation of the Dirac Equation from Principles of Information Processing}. Physical Review A 90 (2014), 062106. \url{https://arxiv.org/abs/1306.1934}.
\bibitem{tanimotojoint} Y. Tanimoto: \emph{Ground state representations of some non-rational conformal nets}. Symmetry 10 (2018), 415. Abschnitt 3.1 zum lokalen Stromnetz. \url{https://arxiv.org/abs/1807.11723}.
\bibitem{chamseddineconnes} A. H. Chamseddine, A. Connes: \emph{Noncommutative Geometry as a Framework for Unification of all Fundamental Interactions including Gravity. Part I}. Gleichungen (5.24), (5.47), (5.49) zum vollständigen gravitativen Spektralaktionskoeffizienten. \url{https://arxiv.org/abs/1004.0464}.
\bibitem{casasibarra} J. A. Casas, A. Ibarra: \emph{Oscillating neutrinos and $\mu\to e,\gamma$}. Gleichung (13) zur allgemeinen Seesaw-Faktorisierung. \url{https://arxiv.org/abs/hep-ph/0103065}.
\bibitem{davidsonibarra} S. Davidson, A. Ibarra: \emph{A lower bound on the right-handed neutrino mass from leptogenesis}. Schranke der Zerfallsasymmetrie, keine automatische Gleichsetzung mit einer PMNS-Phase. \url{https://arxiv.org/abs/hep-ph/0202239}.
\bibitem{covirouletvissani} L. Covi, E. Roulet, F. Vissani: \emph{CP violating decays in leptogenesis scenarios}. Gleichungen (4)--(8), mit Übersetzung zur Heavy-row-Konvention. \url{https://arxiv.org/abs/hep-ph/9605319}.
\bibitem{horavawittenorigin} P. Hořava, E. Witten: \emph{Heterotic and Type I String Dynamics from Eleven Dimensions}. Abschnitt 2 zur lokalen Wandanomalie. \url{https://arxiv.org/abs/hep-th/9510209}.
\bibitem{horavawittenwall} P. Hořava, E. Witten: \emph{Eleven-Dimensional Supergravity on a Manifold with Boundary}. Abschnitte 2--3 zu Wandmultiplet und Anomaliefaktorisierung. \url{https://arxiv.org/abs/hep-th/9603142}.
\bibitem{braunheterotic} V. Braun, Y.-H. He, B. A. Ovrut, T. Pantev: \emph{A Heterotic Standard Model}. Gleichungen (1)--(10): Eichbündelindex und Calabi--Yau-Topologie. \url{https://arxiv.org/abs/hep-th/0501070}.
\bibitem{braundeformation} V. Braun, P. Candelas, R. Davies, R. Donagi: \emph{The MSSM Spectrum from (0,2)-Deformations of the Heterotic Standard Embedding}. Einleitung, Abschnitt 3.1 und Anhang B. \url{https://arxiv.org/abs/1112.1097}.
\bibitem{heeckquadruple} J. Heeck, W. Rodejohann: \emph{Neutrinoless Quadruple Beta Decay}. Abschnitt II, Gleichungen (2)--(5): Ladung-zwei-Feld und zustandsabhängiger Dirac-Schutz. \url{https://arxiv.org/abs/1306.0580}.
\bibitem{liyautorsion} J. Li, S.-T. Yau: \emph{The Existence of Supersymmetric String Theory with Torsion}. Sätze 2.1 und 4.3, mit gegebener holomorpher Deformationsfamilie. \url{https://arxiv.org/abs/hep-th/0411136}.
\bibitem{dongradical} C. Dong, H. Li, G. Mason, P. Montague:
\emph{The radical of a vertex operator algebra}. Definition der gradbewahrenden Mode $o(v)$.
\url{https://arxiv.org/abs/q-alg/9608022}.
\bibitem{barronhigherzhu} K. Barron, N. Vander Werf, J. Yang:
\emph{Higher level Zhu algebras and modules for vertex operator algebras}.
\url{https://arxiv.org/abs/1710.04767}.
\end{thebibliography}

\end{document}
